\documentclass[12pt,a4paper,oneside]{book}

% Минимальные настройки для первой сборки.
% Позже основное оформление перенесём в preamble.tex.
\usepackage{fontspec}
\usepackage[english,russian]{babel}

\setmainfont{DejaVu Serif}
\setsansfont{DejaVu Sans}
\setmonofont{DejaVu Sans Mono}

\usepackage[
    a4paper,
    top=22mm,
    bottom=24mm,
    headheight=26pt,
    left=27mm,
    right=22mm
]{geometry}

\usepackage{graphicx}
\usepackage{eso-pic}
\usepackage{xcolor}
\usepackage{hyperref}

\hypersetup{
    unicode=true,
    colorlinks=false,
    hidelinks,
    pdftitle={Распределённая оптимизация и машинное обучение},
    pdfauthor={Хомченко Илья, Шелегеда Наталья}
}

\urlstyle{same}

% Здесь позднее будут находиться основные настройки оформления.
% ============================================================
% PREAMBLE.TEX
% Распределённая оптимизация и машинное обучение
% ============================================================

% ------------------------------------------------------------
% МАТЕМАТИЧЕСКИЕ ПАКЕТЫ
% ------------------------------------------------------------

\usepackage{amsmath}
\usepackage{amssymb}
\usepackage{amsthm}
\usepackage{mathtools}
\usepackage{bm}
\usepackage{mathrsfs}

% Дополнительные математические символы и интегралы
\usepackage{esint}

% Удобная работа с условиями и командами
\usepackage{etoolbox}

% ------------------------------------------------------------
% ТИПОГРАФИКА
% ------------------------------------------------------------

\usepackage{microtype}
\usepackage{ragged2e}
\usepackage{enumitem}

% Красивое оформление заголовков
\usepackage{titlesec}
\usepackage{tocloft}

% Колонтитулы
\usepackage{fancyhdr}

% Улучшенные подписи
\usepackage{caption}
\usepackage{subcaption}

% ------------------------------------------------------------
% ТАБЛИЦЫ
% ------------------------------------------------------------

\usepackage{booktabs}
\usepackage{array}
\usepackage{tabularx}
\usepackage{longtable}
\usepackage{multirow}

% ------------------------------------------------------------
% РИСУНКИ И ГРАФЫ
% ------------------------------------------------------------

\usepackage{tikz}

\usetikzlibrary{
    arrows.meta,
    calc,
    positioning,
    fit,
    shapes.geometric,
    shapes.misc,
    decorations.pathreplacing,
    decorations.markings,
    patterns,
    backgrounds,
    graphs,
    quotes,
    matrix
}

\graphicspath{
    {figures/}
    {images/}
    {./}
}

% ------------------------------------------------------------
% АЛГОРИТМЫ И ПСЕВДОКОД
% ------------------------------------------------------------

\usepackage{float}
\usepackage{algorithm}
\usepackage{algpseudocode}

% ------------------------------------------------------------
% ЛИСТИНГИ ПРОГРАММНОГО КОДА
% ------------------------------------------------------------

\usepackage{listings}

% ------------------------------------------------------------
% ССЫЛКИ
% hyperref уже загружен в main.tex
% ------------------------------------------------------------

\usepackage{xurl}
\usepackage{bookmark}
\usepackage[nameinlink,noabbrev]{cleveref}

% ============================================================
% ЦВЕТА
% ============================================================

\definecolor{MaulRed}{HTML}{8A1010}
\definecolor{MaulDarkRed}{HTML}{5A0808}
\definecolor{TextBlack}{HTML}{1E1E1E}
\definecolor{DarkGray}{HTML}{444444}
\definecolor{MiddleGray}{HTML}{777777}
\definecolor{LightGray}{HTML}{F2F2F2}
\definecolor{CodeGray}{HTML}{F6F6F6}

% ============================================================
% ОБЩАЯ ТИПОГРАФИКА
% ============================================================

\setlength{\parindent}{1.25cm}
\setlength{\parskip}{0pt}

\setlength{\emergencystretch}{2em}

\clubpenalty=10000
\widowpenalty=10000
\displaywidowpenalty=10000

\setcounter{secnumdepth}{3}
\setcounter{tocdepth}{3}

% Ширина полей нумерации в оглавлении.
\setlength{\cftchapnumwidth}{2.8em}
\setlength{\cftsecnumwidth}{3.8em}
\setlength{\cftsubsecnumwidth}{4.8em}
\setlength{\cftsubsubsecnumwidth}{5.8em}

\allowdisplaybreaks[2]

% Немного увеличиваем расстояние между строками в формулах
\setlength{\jot}{5pt}

% Расстояния между плавающими объектами и текстом
\setlength{\textfloatsep}{18pt plus 2pt minus 4pt}
\setlength{\floatsep}{14pt plus 2pt minus 3pt}
\setlength{\intextsep}{14pt plus 2pt minus 3pt}

% ============================================================
% РУССКИЕ НАЗВАНИЯ
% ============================================================

\addto\captionsrussian{%
    \renewcommand{\chaptername}{Лекция}%
    \renewcommand{\contentsname}{Содержание}%
    \renewcommand{\figurename}{Рисунок}%
    \renewcommand{\tablename}{Таблица}%
    \renewcommand{\appendixname}{Приложение}%
    \renewcommand{\bibname}{Литература}%
}

\renewcommand{\proofname}{Доказательство}
\renewcommand{\qedsymbol}{\(\blacksquare\)}

% ============================================================
% ОФОРМЛЕНИЕ ЗАГОЛОВКОВ
% ============================================================

% Заголовок лекции
\titleformat{\chapter}[display]
    {\normalfont\bfseries}
    {\filleft
        \color{MaulRed}
        \fontsize{15}{18}\selectfont
        \MakeUppercase{\chaptername}\ \thechapter
    }
    {0.5cm}
    {\color{TextBlack}\Huge}
    [%
        \vspace{0.35cm}
        \color{MaulRed}
        \titlerule
    ]

\titleformat{name=\chapter,numberless}[display]
    {\normalfont\bfseries}
    {}
    {0pt}
    {\color{TextBlack}\Huge}
    [%
        \vspace{0.35cm}
        \color{MaulRed}
        \titlerule
    ]

\titlespacing*{\chapter}
    {0pt}
    {-15pt}
    {32pt}

% Разделы
\titleformat{\section}
    {\normalfont\Large\bfseries\color{TextBlack}}
    {\textcolor{MaulRed}{\thesection}}
    {0.75em}
    {}

\titlespacing*{\section}
    {0pt}
    {3.2ex plus 1ex minus 0.2ex}
    {1.3ex plus 0.2ex}

% Подразделы
\titleformat{\subsection}
    {\normalfont\large\bfseries\color{TextBlack}}
    {\textcolor{MaulRed}{\thesubsection}}
    {0.75em}
    {}

\titlespacing*{\subsection}
    {0pt}
    {2.6ex plus 0.8ex minus 0.2ex}
    {1ex plus 0.2ex}

% Пункты
\titleformat{\subsubsection}
    {\normalfont\normalsize\bfseries\color{DarkGray}}
    {\textcolor{MaulRed}{\thesubsubsection}}
    {0.7em}
    {}

% Абзацные заголовки
\titleformat{\paragraph}[runin]
    {\normalfont\normalsize\bfseries}
    {}
    {0pt}
    {}
    [.]

\titlespacing*{\paragraph}
    {0pt}
    {1.5ex}
    {0.6em}

% ============================================================
% КОЛОНТИТУЛЫ
% ============================================================

\setlength{\headheight}{26pt}

\pagestyle{fancy}
\fancyhf{}

\fancyhead[L]{%
    \small
    \nouppercase{\leftmark}
}

\fancyhead[R]{%
    \small
    \thepage
}

\renewcommand{\headrulewidth}{0.4pt}
\renewcommand{\footrulewidth}{0pt}

\renewcommand{\chaptermark}[1]{%
    \markboth{\chaptername\ \thechapter.\ #1}{}
}

\fancypagestyle{plain}{%
    \fancyhf{}
    \fancyfoot[C]{\thepage}
    \renewcommand{\headrulewidth}{0pt}
}

% ============================================================
% СПИСКИ
% ============================================================

\setlist[itemize,1]{
    label=\textcolor{MaulRed}{\textbullet},
    leftmargin=2em,
    itemsep=0.25em,
    topsep=0.5em
}

\setlist[itemize,2]{
    label=\textcolor{DarkGray}{\(\circ\)},
    leftmargin=2em,
    itemsep=0.15em,
    topsep=0.3em
}

\setlist[enumerate,1]{
    label=\textcolor{MaulRed}{\arabic*.},
    leftmargin=2.2em,
    itemsep=0.25em,
    topsep=0.5em
}

\setlist[enumerate,2]{
    label=\textcolor{DarkGray}{\alph*)},
    leftmargin=2em,
    itemsep=0.15em,
    topsep=0.3em
}

\setlist[description]{
    font=\normalfont\bfseries\color{MaulRed},
    style=nextline,
    leftmargin=2em,
    itemsep=0.4em
}

% ============================================================
% ПОДПИСИ К РИСУНКАМ И ТАБЛИЦАМ
% ============================================================

\captionsetup{
    font=small,
    labelfont={bf,color=MaulRed},
    labelsep=period,
    justification=centering,
    singlelinecheck=true
}

\captionsetup[subfigure]{
    font=small,
    labelfont=bf
}

% ============================================================
% НУМЕРАЦИЯ ФОРМУЛ И АЛГОРИТМОВ
% ============================================================

\numberwithin{equation}{chapter}
\numberwithin{algorithm}{chapter}

% ============================================================
% ТЕОРЕМНЫЕ СТИЛИ
% ============================================================

\newtheoremstyle{resultstyle}
    {9pt}
    {9pt}
    {\itshape}
    {}
    {\bfseries\color{MaulRed}}
    {.}
    {0.6em}
    {}

\newtheoremstyle{definitionstyle}
    {9pt}
    {9pt}
    {\normalfont}
    {}
    {\bfseries\color{MaulRed}}
    {.}
    {0.6em}
    {}

\newtheoremstyle{remarkstyle}
    {7pt}
    {7pt}
    {\normalfont}
    {}
    {\itshape\color{DarkGray}}
    {.}
    {0.6em}
    {}

% ------------------------------------------------------------
% Теоремы, леммы и следствия
% Используется единая сквозная нумерация внутри каждой лекции.
% ------------------------------------------------------------

\theoremstyle{resultstyle}

\newtheorem{theorem}{Теорема}[chapter]
\newtheorem{lemma}[theorem]{Лемма}
\newtheorem{proposition}[theorem]{Предложение}
\newtheorem{corollary}[theorem]{Следствие}
\newtheorem{claim}[theorem]{Утверждение}
\newtheorem{criterion}[theorem]{Критерий}

% ------------------------------------------------------------
% Определения, аксиомы и предположения
% ------------------------------------------------------------

\theoremstyle{definitionstyle}

\newtheorem{definition}[theorem]{Определение}
\newtheorem{axiom}[theorem]{Аксиома}
\newtheorem{assumption}[theorem]{Предположение}
\newtheorem{hypothesis}[theorem]{Гипотеза}
\newtheorem{condition}[theorem]{Условие}
\newtheorem{notation}[theorem]{Обозначение}
\newtheorem{convention}[theorem]{Соглашение}
\newtheorem{problem}[theorem]{Задача}
\newtheorem{example}[theorem]{Пример}
\newtheorem{counterexample}[theorem]{Контрпример}
\newtheorem{exercise}[theorem]{Упражнение}

% ------------------------------------------------------------
% Замечания и пояснения
% ------------------------------------------------------------

\theoremstyle{remarkstyle}

\newtheorem{remark}[theorem]{Замечание}
\newtheorem{intuition}[theorem]{Интуиция}
\newtheorem{warning}[theorem]{Предупреждение}
\newtheorem{observation}[theorem]{Наблюдение}
\newtheorem{discussion}[theorem]{Обсуждение}

% ------------------------------------------------------------
% Ненумерованные варианты
% ------------------------------------------------------------

\theoremstyle{resultstyle}

\newtheorem*{theorem*}{Теорема}
\newtheorem*{lemma*}{Лемма}
\newtheorem*{proposition*}{Предложение}
\newtheorem*{corollary*}{Следствие}

\theoremstyle{definitionstyle}

\newtheorem*{definition*}{Определение}
\newtheorem*{axiom*}{Аксиома}
\newtheorem*{assumption*}{Предположение}
\newtheorem*{example*}{Пример}

\theoremstyle{remarkstyle}

\newtheorem*{remark*}{Замечание}
\newtheorem*{intuition*}{Интуиция}
\newtheorem*{warning*}{Предупреждение}

% ============================================================
% ДОКАЗАТЕЛЬСТВА И РЕШЕНИЯ
% ============================================================

\newenvironment{proofidea}
    {\begin{proof}[Идея доказательства]}
    {\end{proof}}

\newenvironment{proofsketch}
    {\begin{proof}[Набросок доказательства]}
    {\end{proof}}

\newenvironment{solution}
    {\begin{proof}[Решение]}
    {\end{proof}}

\newenvironment{verification}
    {\begin{proof}[Проверка]}
    {\end{proof}}

% ============================================================
% СТРУКТУРА ОТДЕЛЬНОЙ ЛЕКЦИИ
% ============================================================

% Использование:
% \lecture{01}{Название лекции}
%
% или:
%
% \lecture{01}{Название лекции}[Дополнительная информация]

\NewDocumentCommand{\lecture}{m m O{}}{%
    \chapter{#2}%
    \label{lec:#1}%
    \ifstrempty{#3}{}{%
        \begin{flushright}
            \small\itshape\color{MiddleGray}
            #3
        \end{flushright}
        \vspace{0.5em}
    }%
}

% Цели лекции
\newenvironment{lecturegoals}
{%
    \par
    \medskip
    \noindent
    {\bfseries\color{MaulRed}Цели лекции}
    \begin{itemize}
}
{%
    \end{itemize}
    \medskip
}

% План лекции
\newenvironment{lectureplan}
{%
    \par
    \medskip
    \noindent
    {\bfseries\color{MaulRed}План лекции}
    \begin{enumerate}
}
{%
    \end{enumerate}
    \medskip
}

% Ключевые понятия
\newenvironment{keyconcepts}
{%
    \par
    \medskip
    \noindent
    {\bfseries\color{MaulRed}Ключевые понятия}
    \begin{itemize}
}
{%
    \end{itemize}
    \medskip
}

% Итоги лекции
\newenvironment{lecturesummary}
{%
    \begin{quote}
    \noindent
    {\bfseries\color{MaulRed}Итоги лекции.}\quad
}
{%
    \end{quote}
}

% Список обозначений конкретной лекции
\newenvironment{notationlist}
{%
    \par
    \medskip
    \noindent
    {\bfseries\color{MaulRed}Обозначения}
    \begin{description}
}
{%
    \end{description}
    \medskip
}

% Ссылка на исходную видеолекцию
\newcommand{\sourcevideo}[2]{%
    \par
    \medskip
    \noindent
    \textbf{\textcolor{MaulRed}{Видеолекция: }}
    \href{#1}{#2}
    \par
    \medskip
}

% Ссылка на дополнительный материал
\newcommand{\additionalsource}[2]{%
    \par
    \smallskip
    \noindent
    \textbf{Дополнительный источник: }
    \href{#1}{#2}
    \par
    \smallskip
}

% Заметка, которую нужно позднее доработать
\newcommand{\TODO}[1]{%
    \textcolor{MaulRed}{%
        \textbf{TODO: }#1%
    }%
}

% ============================================================
% АЛГОРИТМЫ И ПСЕВДОКОД НА РУССКОМ
% ============================================================

\floatname{algorithm}{Алгоритм}
\renewcommand{\listalgorithmname}{Список алгоритмов}

% ------------------------------------------------------------
% Базовые ключевые слова algpseudocode
% ------------------------------------------------------------

\algrenewcommand{\algorithmicrequire}{\textbf{Вход:}}
\algrenewcommand{\algorithmicensure}{\textbf{Выход:}}

\algrenewcommand{\algorithmicif}{\textbf{если}}
\algrenewcommand{\algorithmicthen}{\textbf{то}}
\algrenewcommand{\algorithmicelse}{\textbf{иначе}}

\algrenewcommand{\algorithmicfor}{\textbf{для}}
\algrenewcommand{\algorithmicforall}{\textbf{для всех}}
\algrenewcommand{\algorithmicdo}{\textbf{выполнять}}

\algrenewcommand{\algorithmicwhile}{\textbf{пока}}
\algrenewcommand{\algorithmicrepeat}{\textbf{повторять}}
\algrenewcommand{\algorithmicuntil}{\textbf{до выполнения условия}}

\algrenewcommand{\algorithmicreturn}{\textbf{вернуть}}
\algrenewcommand{\algorithmicprocedure}{\textbf{процедура}}
\algrenewcommand{\algorithmicfunction}{\textbf{функция}}
\algrenewcommand{\algorithmicend}{\textbf{конец}}

% Команда \ElsIf автоматически печатает:
% \algorithmicelse + \algorithmicif.
%
% Команды \EndIf, \EndFor и \EndWhile автоматически печатают:
% \algorithmicend + соответствующее ключевое слово.
%
% Поэтому отдельных команд algorithmicelsif,
% algorithmicendif, algorithmicendfor и
% algorithmicendwhile в algpseudocode нет.

\algrenewcommand{\algorithmiccomment}[1]{%
    \hfill
    \(\triangleright\)
    \textit{#1}%
}

% ------------------------------------------------------------
% Дополнительные команды для наших алгоритмов
% ------------------------------------------------------------

\algnewcommand{\Parameters}[1]{%
    \Statex
    \textbf{Параметры:} #1
}

\algnewcommand{\Initialize}[1]{%
    \Statex
    \textbf{Инициализация:} #1
}

\algnewcommand{\Output}[1]{%
    \Statex
    \textbf{Результат:} #1
}

\algnewcommand{\True}{%
    \textbf{истина}
}

\algnewcommand{\False}{%
    \textbf{ложь}
}

\algnewcommand{\Break}{%
    \State
    \textbf{выйти из цикла}
}

\algnewcommand{\Continue}{%
    \State
    \textbf{перейти к следующей итерации}
}

\algrenewcommand{\algorithmicindent}{1.5em}


% ============================================================
% ОФОРМЛЕНИЕ ПРОГРАММНОГО КОДА
% ============================================================

\lstdefinestyle{lecturecode}{
    basicstyle=\ttfamily\small,
    backgroundcolor=\color{CodeGray},
    frame=single,
    rulecolor=\color{MiddleGray},
    numbers=left,
    numberstyle=\scriptsize\color{MiddleGray},
    stepnumber=1,
    numbersep=8pt,
    showstringspaces=false,
    breaklines=true,
    breakatwhitespace=true,
    tabsize=4,
    columns=fullflexible,
    keepspaces=true,
    captionpos=b
}

\lstset{
    style=lecturecode
}

% ============================================================
% УМНЫЕ ССЫЛКИ
% ============================================================

\crefname{chapter}{лекция}{лекции}
\Crefname{chapter}{Лекция}{Лекции}

\crefname{section}{раздел}{разделы}
\Crefname{section}{Раздел}{Разделы}

\crefname{subsection}{подраздел}{подразделы}
\Crefname{subsection}{Подраздел}{Подразделы}

\crefname{equation}{формула}{формулы}
\Crefname{equation}{Формула}{Формулы}

\crefname{figure}{рисунок}{рисунки}
\Crefname{figure}{Рисунок}{Рисунки}

\crefname{table}{таблица}{таблицы}
\Crefname{table}{Таблица}{Таблицы}

\crefname{algorithm}{алгоритм}{алгоритмы}
\Crefname{algorithm}{Алгоритм}{Алгоритмы}

\crefname{theorem}{теорема}{теоремы}
\Crefname{theorem}{Теорема}{Теоремы}

\crefname{lemma}{лемма}{леммы}
\Crefname{lemma}{Лемма}{Леммы}

\crefname{proposition}{предложение}{предложения}
\Crefname{proposition}{Предложение}{Предложения}

\crefname{corollary}{следствие}{следствия}
\Crefname{corollary}{Следствие}{Следствия}

\crefname{claim}{утверждение}{утверждения}
\Crefname{claim}{Утверждение}{Утверждения}

\crefname{criterion}{критерий}{критерии}
\Crefname{criterion}{Критерий}{Критерии}

\crefname{definition}{определение}{определения}
\Crefname{definition}{Определение}{Определения}

\crefname{axiom}{аксиома}{аксиомы}
\Crefname{axiom}{Аксиома}{Аксиомы}

\crefname{assumption}{предположение}{предположения}
\Crefname{assumption}{Предположение}{Предположения}

\crefname{hypothesis}{гипотеза}{гипотезы}
\Crefname{hypothesis}{Гипотеза}{Гипотезы}

\crefname{condition}{условие}{условия}
\Crefname{condition}{Условие}{Условия}

\crefname{problem}{задача}{задачи}
\Crefname{problem}{Задача}{Задачи}

\crefname{example}{пример}{примеры}
\Crefname{example}{Пример}{Примеры}

\crefname{remark}{замечание}{замечания}
\Crefname{remark}{Замечание}{Замечания}

% ============================================================
% БАЗОВЫЕ ЧИСЛОВЫЕ МНОЖЕСТВА
% ============================================================

\newcommand{\NN}{\mathbb{N}}
\newcommand{\NNzero}{\mathbb{N}_0}
\newcommand{\ZZ}{\mathbb{Z}}
\newcommand{\QQ}{\mathbb{Q}}
\newcommand{\RR}{\mathbb{R}}
\newcommand{\CC}{\mathbb{C}}

\newcommand{\RRplus}{\mathbb{R}_{+}}
\newcommand{\RRplusplus}{\mathbb{R}_{++}}
\newcommand{\RRbar}{\overline{\mathbb{R}}}

% ============================================================
% СКОБКИ, НОРМЫ И МНОЖЕСТВА
% ============================================================

\DeclarePairedDelimiter{\abs}{\lvert}{\rvert}
\DeclarePairedDelimiter{\norm}{\lVert}{\rVert}
\DeclarePairedDelimiter{\ceil}{\lceil}{\rceil}
\DeclarePairedDelimiter{\floor}{\lfloor}{\rfloor}
\DeclarePairedDelimiter{\paren}{(}{)}
\DeclarePairedDelimiter{\bracks}{[}{]}
\DeclarePairedDelimiter{\angles}{\langle}{\rangle}

\newcommand{\set}[1]{%
    \left\{#1\right\}
}

\newcommand{\setgiven}[2]{%
    \left\{
        #1
        \,\middle|\,
        #2
    \right\}
}

\newcommand{\suchthat}{%
    \;\middle|\;
}

\newcommand{\card}[1]{%
    \abs*{#1}
}

\newcommand{\inner}[2]{%
    \left\langle
        #1,#2
    \right\rangle
}

\newcommand{\dualpair}[2]{%
    \left\langle
        #1,#2
    \right\rangle
}

\newcommand{\weightednorm}[2]{%
    \norm{#1}_{#2}
}

\newcommand{\dualnorm}[1]{%
    \norm{#1}_{*}
}

\newcommand{\seminorm}[1]{%
    \left|#1\right|
}

\newcommand{\evalat}[2]{%
    \left.
        #1
    \right|_{#2}
}

\newcommand{\restr}[2]{%
    #1\!\restriction_{#2}
}

% ============================================================
% ВЕКТОРЫ И МАТРИЦЫ
% ============================================================

\newcommand{\vct}[1]{\bm{#1}}
\newcommand{\mat}[1]{\bm{#1}}

\newcommand{\zero}{\mathbf{0}}
\newcommand{\one}{\mathbf{1}}
\newcommand{\Id}{\mathbf{I}}

\newcommand{\trans}{^{\mathsf{T}}}
\newcommand{\herm}{^{\mathsf{H}}}
\newcommand{\inv}{^{-1}}
\newcommand{\pinv}{^{\dagger}}

\newcommand{\kron}{\otimes}

\newcommand{\vectwo}[2]{%
    \begin{pmatrix}
        #1\\
        #2
    \end{pmatrix}
}

\newcommand{\vecthree}[3]{%
    \begin{pmatrix}
        #1\\
        #2\\
        #3
    \end{pmatrix}
}

% ============================================================
% ЛИНЕЙНАЯ АЛГЕБРА
% ============================================================

\DeclareMathOperator{\trace}{tr}
\DeclareMathOperator{\rank}{rank}
\DeclareMathOperator{\diag}{diag}
\DeclareMathOperator{\blkdiag}{blkdiag}
\DeclareMathOperator{\Ker}{Ker}
\DeclareMathOperator{\Image}{Im}
\DeclareMathOperator{\Range}{range}
\DeclareMathOperator{\Null}{null}
\DeclareMathOperator{\Lin}{Lin}
\DeclareMathOperator{\End}{End}
\DeclareMathOperator{\linspan}{span}
\DeclareMathOperator{\spec}{spec}
\DeclareMathOperator{\spectralradius}{\rho}
\DeclareMathOperator{\cond}{cond}
\DeclareMathOperator{\vecop}{vec}

\newcommand{\lambdaMin}[1]{%
    \lambda_{\min}\!\left(#1\right)
}

\newcommand{\lambdaMax}[1]{%
    \lambda_{\max}\!\left(#1\right)
}

\newcommand{\sigmaMin}[1]{%
    \sigma_{\min}\!\left(#1\right)
}

\newcommand{\sigmaMax}[1]{%
    \sigma_{\max}\!\left(#1\right)
}

\newcommand{\conditionnumber}[1]{%
    \kappa\!\left(#1\right)
}

\newcommand{\Sym}[1]{%
    \mathbb{S}^{#1}
}

\newcommand{\PSD}[1]{%
    \mathbb{S}^{#1}_{+}
}

\newcommand{\PD}[1]{%
    \mathbb{S}^{#1}_{++}
}

% ============================================================
% ПРОИЗВОДНЫЕ И ДИФФЕРЕНЦИАЛЫ
% ============================================================

\newcommand{\dd}{%
    \mathop{}\!\mathrm{d}
}

\newcommand{\od}[2]{%
    \frac{\dd #1}{\dd #2}
}

\newcommand{\odn}[3]{%
    \frac{\dd^{#3} #1}{\dd #2^{#3}}
}

\newcommand{\pd}[2]{%
    \frac{\partial #1}{\partial #2}
}

\newcommand{\pdn}[3]{%
    \frac{\partial^{#3} #1}{\partial #2^{#3}}
}

\newcommand{\grad}{\nabla}
\newcommand{\hess}{\nabla^{2}}
\newcommand{\jacobian}{\mathrm{D}}

\newcommand{\gradat}[2]{%
    \nabla #1\!\left(#2\right)
}

\newcommand{\hessat}[2]{%
    \nabla^{2} #1\!\left(#2\right)
}

% ============================================================
% ПРЕДЕЛЫ И АСИМПТОТИКА
% ============================================================

\newcommand{\BigO}{\mathcal{O}}
\newcommand{\littleo}{o}
\newcommand{\BigTheta}{\Theta}
\newcommand{\BigOmega}{\Omega}
\newcommand{\softO}{\widetilde{\mathcal{O}}}

\newcommand{\asympeq}{\sim}

% ============================================================
% БАЗОВЫЕ МАТЕМАТИЧЕСКИЕ ОПЕРАТОРЫ
% ============================================================

\DeclareMathOperator*{\argmin}{arg\,min}
\DeclareMathOperator*{\argmax}{arg\,max}

\DeclareMathOperator*{\esssup}{ess\,sup}
\DeclareMathOperator*{\essinf}{ess\,inf}

\DeclareMathOperator{\sign}{sign}
\DeclareMathOperator{\supp}{supp}
\DeclareMathOperator{\diam}{diam}
\DeclareMathOperator{\dist}{dist}

\DeclareMathOperator{\interior}{int}
\DeclareMathOperator{\closure}{cl}
\DeclareMathOperator{\boundary}{bd}
\DeclareMathOperator{\relint}{ri}

\DeclareMathOperator{\dom}{dom}
\DeclareMathOperator{\epi}{epi}
\DeclareMathOperator{\hypo}{hypo}
\DeclareMathOperator{\graphop}{gph}

\DeclareMathOperator{\conv}{conv}
\DeclareMathOperator{\cone}{cone}
\DeclareMathOperator{\aff}{aff}

\DeclareMathOperator{\Fix}{Fix}
\DeclareMathOperator{\zer}{zer}

\DeclareMathOperator{\proxop}{prox}
\DeclareMathOperator{\projop}{proj}

\DeclareMathOperator{\softmax}{softmax}
\DeclareMathOperator{\ReLU}{ReLU}

% ============================================================
% ЭКСПОНЕНТА, МНИМАЯ ЕДИНИЦА И ОПРЕДЕЛЯЮЩЕЕ РАВЕНСТВО
% ============================================================

\newcommand{\e}{\mathrm{e}}
\newcommand{\ii}{\mathrm{i}}

\newcommand{\defeq}{\coloneqq}
\newcommand{\eqdef}{\eqqcolon}

% ============================================================
% ОПТИМИЗАЦИОННЫЕ ЗАДАЧИ
% ============================================================

\newcommand{\minimize}{%
    \operatorname*{minimize}
}

\newcommand{\maximize}{%
    \operatorname*{maximize}
}

\newcommand{\subjectto}{%
    \text{при условиях}
}

\newcommand{\find}{%
    \operatorname*{find}
}

\newcommand{\optval}{%
    f^{\star}
}

\newcommand{\xopt}{%
    x^{\star}
}

\newcommand{\Xopt}{%
    X^{\star}
}

\newcommand{\Lagrangian}{%
    \mathcal{L}
}

\newcommand{\KKT}{%
    \mathrm{KKT}
}

% ============================================================
% ВЫПУКЛЫЙ АНАЛИЗ
% ============================================================

\newcommand{\subdiff}[2][]{%
    \partial_{#1} #2
}

\newcommand{\prox}[2]{%
    \operatorname{prox}_{#1}\!\left(#2\right)
}

\newcommand{\projection}[2]{%
    \operatorname{proj}_{#1}\!\left(#2\right)
}

\newcommand{\distto}[2]{%
    \operatorname{dist}\!\left(#1,#2\right)
}

\newcommand{\normalcone}[2]{%
    N_{#1}\!\left(#2\right)
}

\newcommand{\tangentcone}[2]{%
    T_{#1}\!\left(#2\right)
}

\newcommand{\indicatorfunction}[1]{%
    \delta_{#1}
}

\newcommand{\indevent}[1]{%
    \mathbf{1}_{\set{#1}}
}

\newcommand{\fenchelconjugate}[1]{%
    #1^{*}
}

\newcommand{\biconjugate}[1]{%
    #1^{**}
}

\newcommand{\bregman}[3]{%
    D_{#1}\!\left(#2,#3\right)
}

\newcommand{\simplex}[1]{%
    \Delta^{#1}
}

\newcommand{\ball}[2]{%
    \mathbb{B}_{#1}\!\left(#2\right)
}

\newcommand{\unitball}{%
    \mathbb{B}
}

\newcommand{\positivepart}[1]{%
    \left[#1\right]_{+}
}

\newcommand{\negativepart}[1]{%
    \left[#1\right]_{-}
}

% ============================================================
% ВЕРОЯТНОСТЬ И СТАТИСТИКА
% ============================================================

\providecommand{\Prob}{}
\renewcommand{\Prob}{\mathbb{P}}
\providecommand{\Expect}{}
\renewcommand{\Expect}{\mathbb{E}}

\NewDocumentCommand{\E}{o m}{%
    \mathbb{E}%
    \IfValueT{#1}{_{#1}}%
    \!\left[#2\right]
}

\NewDocumentCommand{\Prb}{o m}{%
    \mathbb{P}%
    \IfValueT{#1}{_{#1}}%
    \!\left(#2\right)
}

\DeclareMathOperator{\Var}{Var}
\DeclareMathOperator{\Cov}{Cov}
\DeclareMathOperator{\Corr}{Corr}
\DeclareMathOperator{\Law}{Law}

\newcommand{\given}{%
    \,\middle|\,
}

\newcommand{\VarOf}[1]{%
    \operatorname{Var}\!\left(#1\right)
}

\newcommand{\CovOf}[2]{%
    \operatorname{Cov}\!\left(#1,#2\right)
}

\newcommand{\NormalDist}{\mathcal{N}}
\newcommand{\UniformDist}{\mathcal{U}}
\newcommand{\BernoulliDist}{\operatorname{Bern}}
\newcommand{\CategoricalDist}{\operatorname{Cat}}

\newcommand{\iid}{%
    \mathrel{\overset{\mathrm{i.i.d.}}{\sim}}
}

\newcommand{\independent}{%
    \mathrel{\perp\!\!\!\perp}
}

\newcommand{\almostsure}{%
    \text{п.н.}
}

\newcommand{\asconv}{%
    \xrightarrow{\text{п.н.}}
}

\newcommand{\pconv}{%
    \xrightarrow{\Prob}
}

\newcommand{\dconv}{%
    \xrightarrow{d}
}

\newcommand{\loneconv}{%
    \xrightarrow{L^{1}}
}

\newcommand{\ltwoconv}{%
    \xrightarrow{L^{2}}
}

% ============================================================
% ИНФОРМАЦИОННЫЕ ВЕЛИЧИНЫ
% ============================================================

\newcommand{\Entropy}[1]{%
    H\!\left(#1\right)
}

\newcommand{\ConditionalEntropy}[2]{%
    H\!\left(#1\mid #2\right)
}

\newcommand{\MutualInformation}[2]{%
    I\!\left(#1;#2\right)
}

\newcommand{\KL}[2]{%
    D_{\mathrm{KL}}\!\left(
        #1\,\middle\|\,#2
    \right)
}

% ============================================================
% МАШИННОЕ ОБУЧЕНИЕ
% ============================================================

\newcommand{\Dataset}{\mathcal{D}}
\newcommand{\Samples}{\mathcal{S}}
\newcommand{\HypothesisClass}{\mathcal{H}}

\newcommand{\Loss}{\ell}
\newcommand{\Risk}{\mathcal{R}}
\newcommand{\EmpiricalRisk}{\widehat{\mathcal{R}}}

\newcommand{\parameters}{\theta}
\newcommand{\parameterSpace}{\Theta}

\newcommand{\prediction}{\widehat{y}}
\newcommand{\features}{x}
\newcommand{\labelvalue}{y}

\newcommand{\model}{f_{\theta}}

\newcommand{\regularizer}{r}

\newcommand{\batch}{\mathcal{B}}
\newcommand{\batchsize}{B}

\newcommand{\learningrate}{\eta}

\newcommand{\stochasticgradient}{%
    g
}

\newcommand{\gradientestimate}{%
    \widehat{\nabla f}
}

% ============================================================
% РАСПРЕДЕЛЁННАЯ ОПТИМИЗАЦИЯ
% ============================================================

\newcommand{\Graph}{\mathcal{G}}
\newcommand{\Vertices}{\mathcal{V}}
\newcommand{\Edges}{\mathcal{E}}

\newcommand{\numagents}{n}

\newcommand{\neighbors}[1]{%
    \mathcal{N}_{#1}
}

\newcommand{\degreeof}[1]{%
    d_{#1}
}

\newcommand{\Adjacency}{\mathbf{A}}
\newcommand{\DegreeMatrix}{\mathbf{D}}
\newcommand{\Laplacian}{\mathbf{L}}
\newcommand{\MixingMatrix}{\mathbf{W}}

\newcommand{\IncidenceMatrix}{\mathbf{B}}

\newcommand{\localobjective}[1]{%
    f_{#1}
}

\newcommand{\globalobjective}{%
    F
}

\newcommand{\localvariable}[1]{%
    x_{#1}
}

\newcommand{\stackedvariable}{%
    \mathbf{x}
}

\newcommand{\networkaverage}[1]{%
    \overline{#1}
}

\newcommand{\ConsensusProjector}[1]{%
    \frac{1}{#1}\one\one\trans
}

\newcommand{\DisagreementProjector}[1]{%
    \Id-\frac{1}{#1}\one\one\trans
}

\newcommand{\ConsensusSpace}{%
    \linspan\set{\one}
}

\newcommand{\consensuserror}[1]{%
    \norm{
        #1-
        \ConsensusProjector{\numagents}#1
    }
}

\newcommand{\spectralgap}[1]{%
    1-\lambda_{2}\!\left(#1\right)
}

\newcommand{\networkcondition}[1]{%
    \frac{
        \lambda_{\max}\!\left(#1\right)
    }{
        \lambda_{\min}^{+}\!\left(#1\right)
    }
}

% ============================================================
% ПОСЛЕДОВАТЕЛЬНОСТИ И ИТЕРАЦИИ
% ============================================================

\newcommand{\sequence}[2]{%
    \set{#1}_{#2}
}

\newcommand{\iter}[2]{%
    #1^{#2}
}

\newcommand{\nextiter}[1]{%
    #1^{k+1}
}

\newcommand{\currentiter}[1]{%
    #1^{k}
}

\newcommand{\previousiter}[1]{%
    #1^{k-1}
}

\newcommand{\average}[1]{%
    \overline{#1}
}

\newcommand{\runningaverage}[2]{%
    \frac{1}{#1}
    \sum_{k=1}^{#1}
    #2
}

% ============================================================
% ГЛАДКОСТЬ, ВЫПУКЛОСТЬ И ЧИСЛО ОБУСЛОВЛЕННОСТИ
% ============================================================

\newcommand{\smoothnessconstant}{L}
\newcommand{\strongconvexityconstant}{\mu}

\newcommand{\conditionratio}{%
    \kappa
    \defeq
    \frac{L}{\mu}
}

\newcommand{\LipschitzConstant}{L}
\newcommand{\StrongConvexityConstant}{\mu}

% ============================================================
% ТИПОВЫЕ ОПТИМИЗАЦИОННЫЕ ОТОБРАЖЕНИЯ
% ============================================================

\newcommand{\gradientstep}[2]{%
    #1-\learningrate\nabla #2
}

\newcommand{\resolvent}[2]{%
    \left(
        \Id+#1#2
    \right)^{-1}
}

\newcommand{\fixedpointset}[1]{%
    \operatorname{Fix}\!\left(#1\right)
}

\newcommand{\zeroset}[1]{%
    \operatorname{zer}\!\left(#1\right)
}

% ============================================================
% СПЕЦИАЛЬНЫЕ ОБОЗНАЧЕНИЯ
% ============================================================

\newcommand{\eps}{\varepsilon}

\newcommand{\weakto}{%
    \rightharpoonup
}

\newcommand{\weakstarto}{%
    \overset{*}{\rightharpoonup}
}

\newcommand{\compose}{\circ}

\newcommand{\directsum}{\oplus}

\newcommand{\bigdirectsum}{\bigoplus}

% ============================================================
% ВЫДЕЛЕНИЕ ВАЖНЫХ ФОРМУЛ
% ============================================================

\newcommand{\importantformula}[1]{%
    \boxed{\displaystyle #1}
}

% ============================================================
% TIKZ: СТИЛИ ДЛЯ ГРАФОВ И АЛГОРИТМОВ
% ============================================================

\tikzset{
    vertex/.style={
        circle,
        draw=MaulRed,
        thick,
        minimum size=8mm,
        inner sep=1pt,
        fill=white
    },
    active vertex/.style={
        vertex,
        fill=MaulRed,
        text=white
    },
    edge/.style={
        thick,
        draw=DarkGray
    },
    directed edge/.style={
        edge,
        -{Latex[length=2mm]}
    },
    communication edge/.style={
        thick,
        dashed,
        draw=MiddleGray
    },
    block/.style={
        rectangle,
        draw=MaulRed,
        thick,
        rounded corners=2pt,
        minimum width=2.8cm,
        minimum height=0.9cm,
        align=center
    },
    decision/.style={
        diamond,
        draw=MaulRed,
        thick,
        aspect=2,
        align=center
    },
    flow arrow/.style={
        thick,
        -{Latex[length=2mm]}
    }
}

% ============================================================
% КОНЕЦ PREAMBLE.TEX
% ============================================================

% Полностраничное изображение без полей.
\newcommand{\FullPageImage}[1]{%
    \AddToShipoutPictureBG*{%
        \AtPageLowerLeft{%
            \includegraphics[
                width=\paperwidth,
                height=\paperheight
            ]{#1}%
        }%
    }%
    \null
    \thispagestyle{empty}
    \clearpage
}

\begin{document}

% ============================================================
% Первый лист: полностраничная обложка
% ============================================================

\FullPageImage{maul.png}

% ============================================================
% Второй лист: титульная страница
% ============================================================

\begin{titlepage}
    \thispagestyle{empty}
    \centering

    \vspace*{1.7cm}

    {\small\bfseries ЛЕКЦИОННЫЙ КУРС}

    \vspace{0.65cm}

    \rule{0.76\textwidth}{0.8pt}

    \vspace{1.15cm}

    {\fontsize{25}{31}\selectfont
        \bfseries
        Распределённая оптимизация\\[0.25cm]
        и машинное обучение
    }

    \vspace{0.9cm}

    {\Large\itshape Конспект лекций}

    \vspace{1.15cm}

    \rule{0.76\textwidth}{0.8pt}

    \vfill

    {\small\bfseries АВТОРЫ КОНСПЕКТА}

    \vspace{0.55cm}

    {\Large Хомченко Илья}

    \vspace{0.25cm}

    {\Large Шелегеда Наталья}

    \vfill

    \begin{minipage}{0.76\textwidth}
        \centering
        \small
        Конспект подготовлен по материалам открытого видеокурса.\\[0.35cm]
        \href{
            https://www.youtube.com/playlist?list=PLOzRYVm0a65fhKDS8cYIC7YCuVGJ4FOJx
        }{
            \textbf{YouTube-плейлист курса}
        }
    \end{minipage}

    \vspace{1.5cm}

    {\large Москва}

    \vspace{0.15cm}

    {\large 2026}
\end{titlepage}

% ============================================================
% Оглавление
% ============================================================

\frontmatter

\tableofcontents

\mainmatter

% ============================================================
% Лекции
% ============================================================

\lecture{01}{Введение в распределённую оптимизацию}

\sourcevideo
    {https://www.youtube.com/playlist?list=PLOzRYVm0a65fhKDS8cYIC7YCuVGJ4FOJx}
    {Week 1: Lecture 1: Introduction to Optimization}

Распределённая оптимизация изучает задачи, в которых исходные данные,
вычисления и каналы связи распределены между несколькими агентами.
Каждый агент располагает лишь локальной информацией и обменивается
сообщениями с ограниченным числом соседей, но вся система должна
решить общую оптимизационную задачу.

\section{Архитектуры распределённого обучения}

Пусть имеется \(n\) агентов. Агент \(i\) хранит локальные данные
\(\Dataset_i\), локальную функцию \(f_i\) и переменную \(x_i\).
В централизованной постановке все функции доступны одному вычислителю,
который решает задачу
\[
    \min_{x\in\RR^d} F(x),
    \qquad
    F(x)\coloneqq\sum_{i=1}^{n}f_i(x).
\]

В распределённой постановке каждый агент работает только со своей
функцией. Эквивалентная задача записывается в виде
\[
    \begin{aligned}
        \min_{x_1,\ldots,x_n\in\RR^d}
        \quad&
        \sum_{i=1}^{n}f_i(x_i),
        \\
        \text{при условии}
        \quad&
        x_1=x_2=\cdots=x_n.
    \end{aligned}
\]
Равенство локальных переменных называется условием консенсуса.

\begin{definition}[Агент]
Агентом называется отдельный вычислительный узел, который хранит
локальные данные, выполняет локальные вычисления и обменивается
информацией с разрешёнными узлами сети.
\end{definition}

\begin{assumption}[Локальность информации]
Агент \(i\) непосредственно использует только собственную информацию
и сообщения агентов из множества соседей \(\neighbors{i}\).
\end{assumption}

Структура допустимых коммуникаций задаётся графом
\[
    \Graph=(\Vertices,\Edges),
    \qquad
    \Vertices=\{1,\ldots,n\}.
\]
В неориентированной модели наличие ребра
\((i,j)\in\Edges\) означает, что агенты \(i\) и \(j\)
могут обмениваться сообщениями.

Если вычисления выполняются локально, но все сообщения передаются
через единый сервер, архитектура является распределённой по
вычислениям, но централизованной по коммуникациям. Если и вычисления,
и обмен сообщениями происходят локально, архитектуру называют
децентрализованной.



Типичным примером распределённых вычислений служит обучение модели
предсказания текста на пользовательских устройствах. Устройство
хранит локальные данные и вычисляет градиент функции потерь, не
передавая сами данные центральному серверу.

Пусть \(w^k\in\RR^d\) --- параметры общей модели на итерации \(k\),
а \(L_i(w)\) --- функция потерь на данных устройства \(i\).
Сервер выбирает множество участвующих устройств \(S_k\), рассылает
им \(w^k\), получает локальные градиенты и выполняет обновление
\[
    w^{k+1}
    =
    w^k
    -
    \eta_k
    \frac{1}{\card{S_k}}
    \sum_{i\in S_k}
    \nabla L_i(w^k).
\]

Локальные вычисления распределены между устройствами, однако
агрегация остаётся централизованной. Частичное участие клиентов
позволяет не ожидать ответа от всех устройств на каждой итерации.
В простейшей модели предполагается, что локальные выборки имеют
близкие распределения.

Такая архитектура снижает необходимость передавать исходные данные,
но сохраняет центральный сервер как единую точку коммуникации и
потенциальную точку отказа.

\section{Экономическое распределение и необходимость координации}

Рассмотрим сеть из \(n\) электрогенераторов. Генератор \(i\)
производит мощность \(P_i\), а стоимость производства задаётся
квадратичной функцией
\[
    C_i(P_i)
    =
    a_iP_i^2+b_iP_i+c_i.
\]

\begin{problem}[Экономическое распределение мощности]
Требуется определить мощности генераторов из решения задачи
\[
    \begin{aligned}
        \min_{P_1,\ldots,P_n}
        \quad&
        \sum_{i=1}^{n}C_i(P_i),
        \\
        \text{при условиях}
        \quad&
        \sum_{i=1}^{n}P_i=P_{\mathrm{load}},
        \\
        &
        P_i^{\min}\le P_i\le P_i^{\max},
        \qquad i=1,\ldots,n.
    \end{aligned}
\]
Здесь \(P_{\mathrm{load}}\) --- суммарная требуемая мощность.
\end{problem}

Балансовое ограничение требует, чтобы суммарная генерация совпадала
с текущей нагрузкой. Ограничения
\(P_i^{\min}\le P_i\le P_i^{\max}\) описывают физические возможности
генераторов.

В централизованной постановке все коэффициенты \(a_i,b_i,c_i\)
известны одному оператору. В распределённой постановке генератор
знает только собственную функцию стоимости и может не раскрывать её
параметры другим участникам. Глобальное решение должно быть найдено
посредством локальных вычислений и обмена ограниченной информацией
между соседними генераторами.

Децентрализованная архитектура не требует центрального сервера.
Однако её отказоустойчивость зависит от топологии: отказ критической
вершины может нарушить связность сети. Коммуникационная нагрузка
отдельного агента определяется числом его соседей, а не общим числом
узлов сети.



Рассмотрим четыре локальные функции
\[
    f_i(x_i)=(x_i-i)^2,
    \qquad i=1,2,3,4.
\]
Без коммуникации каждый агент получает собственный минимум:
\[
    x_1^\star=1,
    \qquad
    x_2^\star=2,
    \qquad
    x_3^\star=3,
    \qquad
    x_4^\star=4.
\]

Если же агенты должны выбрать одно общее значение, возникает задача
\[
    \begin{aligned}
        \min_{x_1,x_2,x_3,x_4}
        \quad&
        \sum_{i=1}^{4}(x_i-i)^2,
        \\
        \text{при условии}
        \quad&
        x_1=x_2=x_3=x_4.
    \end{aligned}
\]
После подстановки общего значения \(x\) она принимает вид
\[
    \min_{x\in\RR}
    \sum_{i=1}^{4}(x-i)^2.
\]

\begin{lemma}[Минимум суммы квадратов]
Для чисел \(c_1,\ldots,c_n\in\RR\) функция
\[
    \Phi(x)=\sum_{i=1}^{n}(x-c_i)^2
\]
имеет единственный минимум в точке
\[
    x^\star=\frac{1}{n}\sum_{i=1}^{n}c_i.
\]
\end{lemma}

\begin{proof}
Имеем
\[
    \Phi'(x)
    =
    2\sum_{i=1}^{n}(x-c_i)
    =
    2n x-2\sum_{i=1}^{n}c_i.
\]
Условие \(\Phi'(x)=0\) даёт
\[
    x^\star
    =
    \frac{1}{n}\sum_{i=1}^{n}c_i.
\]
Поскольку \(\Phi''(x)=2n>0\), эта точка является единственным
минимумом.
\end{proof}

Для рассматриваемой задачи
\[
    x^\star
    =
    \frac{1+2+3+4}{4}
    =
    2.5.
\]
Следовательно, совокупность локальных минимумов не определяет
решение глобальной задачи. Для получения общего оптимума агентам
необходим обмен информацией.

\section{Консенсус и локальное усреднение}

Пусть каждый агент хранит скалярную оценку \(x_i^k\) на итерации
\(k\).

\begin{definition}[Консенсус]
Последовательность локальных оценок достигает консенсуса, если
существует \(c\in\RR\), для которого
\[
    \lim_{k\to\infty}x_i^k=c,
    \qquad
    i=1,\ldots,n.
\]
\end{definition}

\begin{definition}[Средний консенсус]
Достигается средний консенсус, если
\[
    \lim_{k\to\infty}x_i^k
    =
    \frac{1}{n}\sum_{j=1}^{n}x_j^0
\]
для каждого \(i\).
\end{definition}

Обычный консенсус требует лишь равенства предельных значений.
Средний консенсус дополнительно требует, чтобы общий предел совпадал
со средним начальных данных.



Пусть агент \(i\) получает оценки соседей и вычисляет
\[
    x_i^{k+1}
    =
    \sum_{j=1}^{n}a_{ij}x_j^k,
\]
где
\[
    a_{ij}=0
    \qquad
    \text{при }
    j\notin\neighbors{i}\cup\{i\}.
\]
Коэффициенты \(a_{ij}\) задают вес информации агента \(j\) при
обновлении агента \(i\).

Объединяя локальные оценки в вектор
\[
    x^k
    =
    \begin{pmatrix}
        x_1^k&\cdots&x_n^k
    \end{pmatrix}^{\mathsf T},
\]
получаем линейную итерацию
\[
    x^{k+1}=Ax^k,
    \qquad
    A=[a_{ij}].
\]

\begin{definition}[Матрица смешивания]
Матрицей смешивания называется неотрицательная матрица
\(A=[a_{ij}]\), определяющая обновление \(x^{k+1}=Ax^k\) и
удовлетворяющая условиям
\[
    a_{ij}=0
    \quad
    \text{при }
    j\notin\neighbors{i}\cup\{i\},
    \qquad
    \sum_{j=1}^{n}a_{ij}=1.
\]
\end{definition}

В лекции эта матрица называется матрицей смежности. Для
алгоритмического анализа точнее говорить о матрице смешивания,
поскольку её ненулевые элементы задают не только наличие связи,
но и вес передаваемой информации.

\begin{algorithm}[H]
\caption{Локальный алгоритм консенсуса}
\label{alg:lecture01-consensus}
\begin{algorithmic}[1]

\Require
граф \(\Graph=(\Vertices,\Edges)\),
матрица смешивания \(A=[a_{ij}]\),
начальные значения \(x_i^0\)

\For{\(k=0,1,2,\ldots\)}

    \ForAll{агентов \(i\in\Vertices\) параллельно}

        \State
        получить значения \(x_j^k\) от соседей
        \(j\in\neighbors{i}\)

        \State
        \(x_i^{k+1}
        \gets
        \displaystyle
        \sum_{j=1}^{n}a_{ij}x_j^k\)

    \EndFor

\EndFor

\Ensure
локальные последовательности \(\{x_i^k\}_{k\ge0}\)

\end{algorithmic}
\end{algorithm}

\section{Пример распределённого вычисления среднего}

Пусть четыре датчика имеют начальные измерения
\[
    x^0
    =
    \begin{pmatrix}
        25\\
        20\\
        24\\
        27
    \end{pmatrix}.
\]
Их среднее равно
\[
    \overline{x}^0
    =
    \frac{25+20+24+27}{4}
    =
    24.
\]

Рассмотрим матрицу смешивания
\[
    A_1
    =
    \begin{pmatrix}
        \frac12 & \frac12 & 0 & 0
        \\
        \frac14 & \frac14 & \frac14 & \frac14
        \\
        0 & \frac13 & \frac13 & \frac13
        \\
        0 & \frac13 & \frac13 & \frac13
    \end{pmatrix}.
\]
Каждая строка \(A_1\) суммируется в единицу, поэтому каждый агент
вычисляет взвешенное среднее доступных ему значений. Итерация
\[
    x^{k+1}=A_1x^k
\]
использует только локальные коммуникации.

В приведённом в лекции численном эксперименте после нескольких
итераций все компоненты сходятся приблизительно к одному значению:
\[
    x_i^k\longrightarrow 23.58.
\]
Следовательно, консенсус достигается, но его значение не равно
исходному среднему \(24\).

\begin{remark}
Нормировка строк матрицы смешивания обеспечивает корректность
локального усреднения, но сама по себе не гарантирует сохранение
глобального среднего.
\end{remark}



Определим среднее на итерации \(k\):
\[
    \overline{x}^k
    =
    \frac{1}{n}\one\trans x^k.
\]

\begin{lemma}[Инвариантность среднего]
Если матрица смешивания удовлетворяет условию
\[
    \one\trans A=\one\trans,
\]
то
\[
    \overline{x}^{k+1}=\overline{x}^k
\]
для всех \(k\).
\end{lemma}

\begin{proof}
Из равенства \(x^{k+1}=Ax^k\) следует
\[
    \overline{x}^{k+1}
    =
    \frac{1}{n}\one\trans x^{k+1}
    =
    \frac{1}{n}\one\trans Ax^k
    =
    \frac{1}{n}\one\trans x^k
    =
    \overline{x}^k.
\]
\end{proof}

\begin{theorem}[Консенсус при сохранении среднего]
Пусть среднее сохраняется и существует предел
\[
    x^k\longrightarrow c\one.
\]
Тогда
\[
    c=\overline{x}^0.
\]
Иными словами, достигнутый консенсус является средним консенсусом.
\end{theorem}

\begin{proof}
По инвариантности среднего
\[
    \overline{x}^k=\overline{x}^0.
\]
Переходя к пределу, получаем
\[
    \overline{x}^0
    =
    \lim_{k\to\infty}\frac{1}{n}\one\trans x^k
    =
    \frac{1}{n}\one\trans(c\one)
    =
    c.
\]
\end{proof}

В лекции для того же графа после изменения весов на достаточно
большой итерации численно получено
\[
    x_i^k\approx23.9989\approx24.
\]
Если новая матрица сохраняет среднее и обеспечивает консенсус, её
предельное значение равно точно \(24\). Следовательно, один и тот же
граф может приводить как к произвольному, так и к среднему консенсусу:
результат определяется не только топологией, но и весами обмена.

\section{Роль ускорения}

В прикладных задачах целевая точка может изменяться со временем.
Например, в задаче экономического распределения мощности величина
\(P_{\mathrm{load}}\) зависит от текущего потребления. Если алгоритм
сходится слишком медленно, найденное решение может соответствовать
уже устаревшему значению нагрузки.

Поэтому существенна не только сходимость метода, но и скорость
сходимости. Первая часть курса посвящена ускорению централизованных
методов оптимизации. Во второй части вводятся основы теории графов,
алгоритмы консенсуса и методы распределённого решения задач,
включая экономическое распределение мощности.


Распределённая оптимизация объединяет локальные вычисления и
ограниченный обмен информацией ради решения общей задачи.
Коммуникации необходимы, поскольку локальные минимумы в общем случае
не совпадают с глобальным решением. В алгоритмах консенсуса структура
графа определяет допустимые обмены, а матрица смешивания определяет
предельное значение. Обычный консенсус не гарантирует средний
консенсус; для последнего необходимо одновременно добиться
сходимости оценок и сохранить среднее начальных данных.
\lecture{02}{Анализ алгоритмов оптимизации в непрерывном времени}

\sourcevideo
    {https://www.youtube.com/playlist?list=PLOzRYVm0a65fhKDS8cYIC7YCuVGJ4FOJx}
    {Week 1: Lecture 2: Analyzing Optimization Algorithms
    in Continuous Time Domain}

Вычислительные методы оптимизации обычно реализуются в дискретном
времени, однако их непрерывные аналоги часто проще анализировать
средствами теории динамических систем. Такой подход позволяет
исследовать устойчивость оптимального равновесия, оценивать скорость
сходимости и конструировать новые дискретные методы.

\section{От градиентного спуска к градиентному потоку}

Для минимизации дифференцируемой функции
\[
    f\colon\RR^d\to\RR
\]
градиентный спуск имеет вид
\[
    x^{k+1}
    =
    x^k-\eta\nabla f(x^k),
\]
где \(\eta>0\) --- шаг метода. Перепишем обновление как
\[
    \frac{x^{k+1}-x^k}{\eta}
    =
    -\nabla f(x^k).
\]

Полагая \(t_k=k\eta\) и переходя к формальному пределу
\(\eta\to0\), получаем дифференциальное уравнение
\[
    \dot{x}(t)
    =
    -\nabla f(x(t)).
\]

\begin{definition}[Градиентный поток]
Градиентным потоком функции \(f\) называется динамическая система
\[
    \dot{x}=-\nabla f(x).
\]
\end{definition}

Дискретный градиентный спуск можно рассматривать как явную схему
Эйлера для градиентного потока:
\[
    x^{k+1}
    =
    x^k+\eta\dot{x}^k
    =
    x^k-\eta\nabla f(x^k).
\]

\begin{remark}
Непрерывная система используется прежде всего для анализа и
проектирования метода. Численная реализация по-прежнему требует
дискретизации.
\end{remark}

\begin{definition}[Точка равновесия]
Точка \(x^\star\) называется точкой равновесия системы
\(\dot{x}=F(x)\), если
\[
    F(x^\star)=0.
\]
Для градиентного потока условие равновесия имеет вид
\[
    \nabla f(x^\star)=0.
\]
\end{definition}

Если \(x^\star\) является внутренней точкой локального минимума
дифференцируемой функции, то она является равновесием градиентного
потока. Обратное утверждение без дополнительных предположений
неверно: критическая точка может быть максимумом или седловой
точкой.

\section{Квадратичная функция}

Рассмотрим
\[
    f(x)=\frac12x^2.
\]
Её градиент равен
\[
    \nabla f(x)=x,
\]
поэтому градиентный поток задаётся уравнением
\[
    \dot{x}=-x.
\]

Единственная точка равновесия
\[
    x^\star=0
\]
одновременно является единственным глобальным минимумом функции.

\begin{lemma}
Решение задачи
\[
    \dot{x}=-x,
    \qquad
    x(0)=x_0,
\]
имеет вид
\[
    x(t)=x_0\mathrm{e}^{-t}.
\]
\end{lemma}

\begin{proof}
Если \(x_0=0\), то \(x(t)\equiv0\). При \(x_0\ne0\) решение сохраняет
знак, и разделение переменных даёт
\[
    \frac{\mathrm{d}x}{x}=-\mathrm{d}t.
\]
После интегрирования и учёта начального условия получаем
\[
    \ln\abs{x(t)}
    =
    -t+\ln\abs{x_0},
\]
откуда
\[
    x(t)=x_0\mathrm{e}^{-t}.
\]
\end{proof}

\begin{definition}[Глобальная экспоненциальная устойчивость]
Равновесие \(x^\star\) называется глобально экспоненциально
устойчивым, если существуют \(C\ge1\) и \(\lambda>0\) такие, что для
любого начального состояния и всех \(t\ge0\)
\[
    \norm{x(t)-x^\star}
    \le
    C\mathrm{e}^{-\lambda t}
    \norm{x(0)-x^\star}.
\]
\end{definition}

Для рассматриваемой системы можно взять \(C=1\) и \(\lambda=1\).
Следовательно, минимум \(x^\star=0\) является экспоненциально
устойчивым равновесием.

Значение целевой функции убывает как
\[
    f(x(t))
    =
    \frac12x_0^2\mathrm{e}^{-2t}.
\]
Таким образом, устойчивость равновесия непосредственно задаёт
скорость сходимости алгоритма к оптимальному решению.

\section{Нормированный градиентный поток}

Обычный градиентный поток замедляется вблизи критической точки,
поскольку
\[
    \norm{\nabla f(x)}
    \longrightarrow0
    \qquad
    \text{при }
    x\to x^\star.
\]
Для изменения скорости движения можно нормировать градиент.

\begin{definition}[Нормированный градиентный поток]
Нормированным градиентным потоком называется система
\[
    \dot{x}
    =
    -
    \frac{\nabla f(x)}
         {\norm{\nabla f(x)}},
    \qquad
    \nabla f(x)\ne0.
\]
В критических точках полагается \(\dot{x}=0\).
\end{definition}

Векторное поле может быть разрывным в критических точках. Ниже
рассматриваются абсолютно непрерывные траектории, удовлетворяющие
уравнению почти всюду. Нормирование сохраняет направление
наискорейшего убывания в евклидовой геометрии и делает скорость равной
единице вне критических точек.

Для функции \(f(x)=x^2/2\) получаем
\[
    \dot{x}
    =
    -\frac{x}{\abs{x}}
    =
    -\operatorname{sign}(x),
    \qquad x\ne0.
\]

\begin{lemma}[Конечно-временная сходимость]
Пусть \(\operatorname{sign}(0)=0\). Решение системы
\[
    \dot{x}=-\operatorname{sign}(x),
    \qquad
    x(0)=x_0,
\]
имеет вид
\[
    x(t)
    =
    \operatorname{sign}(x_0)\bigl(\abs{x_0}-t\bigr)_+,
    \qquad
    (s)_+\coloneqq\max\{s,0\}.
\]
Оно достигает равновесия за время
\[
    T=\abs{x_0}.
\]
\end{lemma}

\begin{proof}
При \(x_0>0\) до достижения нуля выполняется \(\dot{x}=-1\), поэтому
\(x(t)=x_0-t\). При \(x_0<0\) выполняется \(\dot{x}=1\), поэтому
\(x(t)=x_0+t\). После достижения нуля траектория остаётся в точке
равновесия.
\end{proof}

Обычный градиентный поток приближается к нулю асимптотически, тогда
как нормированный поток достигает нуля за конечное время. При этом
далеко от оптимума нормирование может уменьшить естественно большой
градиент, поэтому выбор поля должен учитывать геометрию функции и
область движения.

Дискретизация нормированного потока приводит к обновлению
\[
    x^{k+1}
    =
    x^k
    -
    \eta_k
    \frac{\nabla f(x^k)}
         {\norm{\nabla f(x^k)}}.
\]

\begin{algorithm}[H]
\caption{Нормированный градиентный метод}
\label{alg:normalized-gradient}
\begin{algorithmic}[1]

\Require
дифференцируемая функция \(f\),
начальная точка \(x^0\),
последовательность шагов \(\{\eta_k\}_{k\ge0}\)

\For{\(k=0,1,2,\ldots\)}

    \State
    \(g^k\gets\nabla f(x^k)\)

    \If{\(\norm{g^k}=0\)}

        \State
        \Return \(x^k\)

    \EndIf

    \State
    \(x^{k+1}
    \gets
    x^k-\eta_k g^k/\norm{g^k}\)

\EndFor

\end{algorithmic}
\end{algorithm}

Идея изменения масштаба градиента появляется и в практических
методах обучения нейронных сетей. Однако конкретные алгоритмы
AdaGrad и Adam используют более сложное покоординатное масштабирование
и накопленные статистики градиентов.

\section{Квадратичная и четвертичная функции}

Сравним функции
\[
    f_1(x)=\frac12x^2,
    \qquad
    f_2(x)=\frac14x^4.
\]
Обе имеют единственный минимум в точке \(x^\star=0\), но их градиенты
различаются:
\[
    \nabla f_1(x)=x,
    \qquad
    \nabla f_2(x)=x^3.
\]

При \(0<\abs{x}<1\) выполняется
\[
    \abs{x^3}<\abs{x}.
\]
Следовательно, вблизи минимума четвертичная функция создаёт меньший
градиент и более медленное движение.

Градиентные потоки имеют вид
\[
    \dot{x}=-x
    \qquad\text{и}\qquad
    \dot{x}=-x^3.
\]

\begin{lemma}
Решение системы
\[
    \dot{x}=-x^3,
    \qquad
    x(0)=x_0,
\]
равно
\[
    x(t)
    =
    \frac{x_0}
         {\sqrt{1+2x_0^2t}}.
\]
\end{lemma}

\begin{proof}
Для \(x_0\ne0\)
\[
    \frac{\mathrm{d}x}{x^3}=-\mathrm{d}t.
\]
Отсюда
\[
    \frac{1}{2x(t)^2}
    =
    t+\frac{1}{2x_0^2},
\]
что даёт требуемую формулу. При \(x_0=0\) утверждение очевидно.
\end{proof}

Для квадратичной функции расстояние до минимума убывает
экспоненциально:
\[
    \abs{x(t)}=\abs{x_0}\mathrm{e}^{-t}.
\]
Для четвертичной функции оно убывает лишь полиномиально:
\[
    \abs{x(t)}
    =
    \frac{\abs{x_0}}
         {\sqrt{1+2x_0^2t}}
    =
    O(t^{-1/2}).
\]

Следовательно, одинаковое положение минимума не означает одинаковую
скорость оптимизации. Существенна геометрия функции вблизи
оптимального решения.

\section{Сильная выпуклость}

\begin{definition}[Сильная выпуклость]
Дифференцируемая функция \(f\colon\RR^d\to\RR\) называется
\(\mu\)-сильно выпуклой, если существует \(\mu>0\), такое что
\[
    f(y)
    \ge
    f(x)
    +
    \inner{\nabla f(x)}{y-x}
    +
    \frac{\mu}{2}\norm{y-x}^2
\]
для всех \(x,y\in\RR^d\).
\end{definition}

Для дважды дифференцируемой функции это условие эквивалентно
матричному неравенству
\[
    \nabla^2f(x)\succeq\mu I_d
    \qquad
    \text{для всех }x\in\RR^d.
\]

Функция
\[
    f_1(x)=\frac12x^2
\]
является \(1\)-сильно выпуклой, поскольку
\[
    f_1''(x)=1.
\]

Для функции
\[
    f_2(x)=\frac14x^4
\]
имеем
\[
    f_2''(x)=3x^2.
\]
Так как
\[
    \inf_{x\in\RR}f_2''(x)=0,
\]
она выпукла, но не является сильно выпуклой на всей прямой.

Сильная выпуклость препятствует чрезмерному уплощению функции около
минимума. Поэтому для сильно выпуклых задач возможно получать более
сильные оценки скорости сходимости.

\section{Почему одного градиента может быть недостаточно}

Рассмотрим функцию
\[
    f(x)=\abs{x}.
\]
При \(x\ne0\) модуль её производной равен единице:
\[
    \abs{f'(x)}=1.
\]
Следовательно, по величине градиента невозможно определить расстояние
до минимума: точки \(x=1\) и \(x=100\) дают одинаковую величину
градиента.

Аналогичная трудность сохраняется для гладких приближений функции
\(\abs{x}\): без дополнительной связи между градиентом, значением
функции и расстоянием до множества минимумов невозможно в общем
случае гарантировать ускоренную сходимость только на основании
градиентной информации.

\section{Условие Поляка--Лоясиевича}

Сильная выпуклость является достаточным, но не единственным условием
быстрой сходимости.

\begin{definition}[Неравенство Поляка--Лоясиевича]
Функция \(f\) удовлетворяет неравенству Поляка--Лоясиевича с
константой \(\mu>0\), если
\[
    \frac12\norm{\nabla f(x)}^2
    \ge
    \mu\bigl(f(x)-f^\star\bigr)
\]
для всех \(x\), где
\[
    f^\star=\inf_x f(x).
\]
\end{definition}

Это условие связывает величину градиента с текущей ошибкой по
значению функции. В отличие от сильной выпуклости, оно не требует
выпуклости функции.

Рассмотрим невыпуклый пример
\[
    f(x)=x^2+3\sin^2x.
\]
Эта функция не выпукла, поскольку её вторая производная
\[
    f''(x)=2+6\cos(2x)
\]
может быть отрицательной. Вместе с тем
\[
    f(x)\ge0,
    \qquad
    f(x)=0
    \Longleftrightarrow
    x=0,
\]
поэтому она имеет единственный глобальный минимум в нуле. Однако
невыпуклость и единственность минимума сами по себе не доказывают
PL-неравенство: для этого требуется отдельная глобальная оценка
градиента через \(f(x)-f^\star\).

Неравенство Поляка--Лоясиевича и более широкие классы функций,
допускающие ускоренный анализ, рассматриваются в следующих лекциях.

\section{Проектирование методов через динамические системы}

Непрерывная модель позволяет разделить разработку алгоритма на три
этапа. Сначала выбирается векторное поле
\[
    \dot{x}=F(x),
\]
равновесия которого совпадают с решениями оптимизационной задачи.
Затем средствами теории устойчивости исследуется сходимость
траекторий и её скорость. После этого система дискретизируется для
численной реализации.

\begin{algorithm}[H]
\caption{Проектирование оптимизационного метода через непрерывную модель}
\label{alg:continuous-design}
\begin{algorithmic}[1]

\Require
задача \(\min_x f(x)\)

\State
выбрать поле \(F(x)\), для которого
\[
    F(x)=0
    \quad\Longleftrightarrow\quad
    x\in\operatorname*{arg\,min}_{z}f(z)
\]

\State
исследовать устойчивость равновесий системы
\(\dot{x}=F(x)\)

\State
получить оценку скорости сходимости траекторий

\State
выбрать дискретизацию
\[
    x^{k+1}=x^k+\eta_kF(x^k)
\]

\State
исследовать ограничения на шаг \(\eta_k\)

\Ensure
дискретный оптимизационный алгоритм

\end{algorithmic}
\end{algorithm}


Градиентный спуск является дискретизацией градиентного потока
\(\dot{x}=-\nabla f(x)\). Оптимальные критические точки выступают
равновесиями соответствующей динамической системы, а скорость
оптимизации связана с типом их устойчивости. Для квадратичной функции
градиентный поток сходится экспоненциально, тогда как для четвертичной
функции сходимость имеет полиномиальный характер. Нормирование
градиента позволяет изменить скорость движения и в простейшем
скалярном случае обеспечивает конечно-временную сходимость. Возможность
ускорения определяется геометрическими свойствами функции, в
частности сильной выпуклостью или более общим неравенством
Поляка--Лоясиевича.

\lecture{03}{Структура курса}

\sourcevideo
    {https://www.youtube.com/playlist?list=PLOzRYVm0a65fhKDS8cYIC7YCuVGJ4FOJx}
    {Week 1: Lecture 3: Course Outline}

Курс объединяет выпуклую оптимизацию, анализ динамических систем,
ускоренные методы и распределённые алгоритмы. Основная часть
результатов формулируется для выпуклых задач. Отдельно будут
рассматриваться некоторые невыпуклые функции, удовлетворяющие
неравенству Поляка--Лоясиевича.

\section{Почему основное внимание уделяется выпуклым задачам}

Выпуклая функция не обязана иметь единственный минимум. Например,
постоянная функция является выпуклой, и каждая точка её области
определения является глобальным минимумом. Аналогично минимум может
достигаться на целом отрезке или более общем выпуклом множестве.

Ключевое отличие выпуклой задачи от общей невыпуклой задачи состоит
не в единственности решения, а в отсутствии неглобальных локальных
минимумов.

\begin{theorem}[Локальная и глобальная оптимальность]
Пусть функция \(f\) выпукла на выпуклом множестве \(X\). Тогда всякая
точка локального минимума функции \(f\) на \(X\) является точкой её
глобального минимума на \(X\).
\end{theorem}

Таким образом, если алгоритм нашёл локальный минимум выпуклой задачи,
дополнительная проверка его глобальности не требуется. Различные
точки минимума могут существовать, но значение целевой функции во
всех них одинаково.

В невыпуклой задаче функция может иметь несколько локальных минимумов
с различными значениями:
\[
    f(x_{\mathrm{loc}})
    >
    f(x^\star).
\]
Поэтому локальная сходимость алгоритма в общем случае не гарантирует
нахождение глобального решения.

\begin{remark}
Формальные определения выпуклых множеств, выпуклых функций и строгой
выпуклости вводятся в последующих лекциях. Здесь используется только
главное следствие выпуклости: локальный минимум является глобальным.
\end{remark}

\section{Выпуклая оптимизация}

Первая содержательная часть курса посвящена централизованным задачам
оптимизации. Будут рассмотрены как задачи без ограничений,
\[
    \min_{x\in\RR^d} f(x),
\]
так и задачи с ограничениями,
\[
    \begin{aligned}
        \min_{x\in\RR^d}
        \quad& f(x),
        \\
        \text{при условии}
        \quad& x\in\mathcal X.
    \end{aligned}
\]

Основной класс составляют выпуклые функции и выпуклые допустимые
множества. Помимо них рассматриваются функции, удовлетворяющие
неравенству Поляка--Лоясиевича. Такие функции могут быть невыпуклыми,
но сохранять свойства, достаточные для анализа быстрой сходимости.

В качестве мотивационного примера используется функция
\[
    f(x)=x^2+3\sin^2x.
\]
Она не является выпуклой на всей прямой, но имеет единственный
глобальный минимум в точке \(x^\star=0\). Этот пример показывает, что
невыпуклость не исключает простой глобальной структуры задачи, однако
сама по себе не гарантирует выполнения PL-неравенства.

\section{Теория устойчивости}

После основ выпуклой оптимизации курс переходит к теории устойчивости
по Ляпунову для непрерывных динамических систем
\[
    \dot{x}=F(x).
\]

Оптимизационный метод сначала представляется как динамическая
система, множество равновесий которой должно совпадать с множеством
решений исходной задачи. Затем устойчивость равновесий используется
для доказательства сходимости, а её тип задаёт скорость приближения к
решению.

Особое внимание уделяется фиксированно-временной сходимости.

\begin{definition}[Фиксированно-временная сходимость]
Равновесие \(x^\star\) обладает фиксированно-временной сходимостью,
если существует число \(T_{\max}>0\), не зависящее от начального
состояния, такое что всякая траектория достигает равновесия не позднее
момента \(T_{\max}\):
\[
    x(t)=x^\star
    \qquad
    \text{при всех }t\ge T(x_0),
    \qquad
    T(x_0)\le T_{\max}.
\]
\end{definition}

Главная особенность этой оценки состоит в независимости верхней
границы времени сходимости от начальной точки. Это сильнее обычной
асимптотической или экспоненциальной сходимости, при которой решение
приближается к равновесию, но в общем случае не достигает его за
конечное время.

\section{Проектирование ускоренных методов}

Непрерывные динамические системы используются не только для анализа
известных алгоритмов, но и для построения новых методов. Общая схема
имеет вид
\[
    \begin{gathered}
        \text{оптимизационная задача}
        \longrightarrow
        \text{непрерывная динамика},
        \\
        \text{непрерывная динамика}
        \longrightarrow
        \text{анализ устойчивости},
        \\
        \text{анализ устойчивости}
        \longrightarrow
        \text{дискретный алгоритм}.
    \end{gathered}
\]

Сначала выбирается система
\[
    \dot{x}=F(x),
\]
множество равновесий которой совпадает с множеством решений, а скорость
сходимости удовлетворяет требуемой оценке. После этого система
дискретизируется:
\[
    x^{k+1}
    =
    x^k+\eta_kF(x^k).
\]

Непрерывная система может иметь фиксированно-временную или иную
усиленную сходимость, однако после дискретизации соответствующие
гарантии необходимо исследовать заново. Поведение дискретного метода
зависит от выбранной схемы и величины шага \(\eta_k\).

\section{Переход к распределённой оптимизации}

После централизованной оптимизации и теории устойчивости начинается
распределённая часть курса. Для её построения вводятся основные
понятия теории графов:
\[
    \Graph=(\Vertices,\Edges).
\]

Вершины графа соответствуют вычислительным агентам, а рёбра задают
допустимые каналы коммуникации. Будут рассмотрены связность графа и
матрица Лапласа \(L=D-A\), где \(D\) --- матрица степеней, а \(A\) ---
матрица смежности.

\begin{definition}[Значение Фидлера]
Для неориентированного графа значением Фидлера называется второе
наименьшее собственное значение его матрицы Лапласа:
\[
    \lambda_2(L).
\]
\end{definition}

Значение Фидлера называют также алгебраической связностью. Для
неориентированного графа условие \(\lambda_2(L)>0\) эквивалентно его
связности.

После введения графовой модели рассматриваются алгоритмы, в которых
каждый агент выполняет локальные вычисления и обменивается информацией
только с соседями. Их цель состоит в решении глобальной задачи без
полного доступа к данным и функциям остальных участников.

Таким образом, содержательная последовательность курса имеет вид
\[
    \begin{gathered}
        \text{выпуклая оптимизация}
        \longrightarrow
        \text{теория устойчивости},
        \\
        \text{теория устойчивости}
        \longrightarrow
        \text{ускоренные методы},
        \\
        \text{теория графов}
        \longrightarrow
        \text{консенсус},
        \\
        \text{консенсус}
        \longrightarrow
        \text{распределённая оптимизация}.
    \end{gathered}
\]

Заключительная часть курса посвящена крупномасштабному и
федеративному обучению.

\section{Источники курса}

Курс не следует одному учебнику. Материал основывается на конспектах
лекций и исследовательских статьях, сопровождающих соответствующие
темы.

Значительная часть рассматриваемых результатов относится к
современным направлениям ускоренной, крупномасштабной и
распределённой оптимизации. Эти методы образуют теоретическую основу
распределённого машинного обучения.


Основу курса составляют выпуклая оптимизация и функции, допускающие
усиленные оценки сходимости. Затем вводится теория устойчивости
непрерывных динамических систем, включая сходимость за фиксированное время, и на её основе проектируются ускоренные методы. Далее
изучаются теория графов, консенсус, распределённая оптимизация,
крупномасштабное и федеративное обучение.
\lecture{04}{Основные постановки задач оптимизации}

\sourcevideo
    {https://www.youtube.com/playlist?list=PLOzRYVm0a65fhKDS8cYIC7YCuVGJ4FOJx}
    {Week 1: Lecture 4: Basics of Optimization Problems}

Перед построением распределённых алгоритмов необходимо определить,
какая задача решается, какие величины являются неизвестными и какие
ограничения должны выполняться. В этой лекции вводится общая
математическая постановка оптимизации и рассматриваются две
прикладные модели: разреженное восстановление сигнала и метод
опорных векторов.

\section{Общая задача оптимизации и ограничения}

Пусть
\[
    f\colon\RR^n\to\RR.
\]
Задача без ограничений имеет вид
\[
    \min_{x\in\RR^n} f(x).
\]

\begin{definition}[Целевая функция]
Функция \(f\), значение которой требуется минимизировать или
максимизировать, называется целевой функцией.
\end{definition}

\begin{definition}[Переменная решения]
Вектор
\[
    x=
    \begin{pmatrix}
        x_1&\cdots&x_n
    \end{pmatrix}^{\mathsf T}
    \in\RR^n
\]
называется переменной решения, оптимизационной переменной или
вектором параметров задачи.
\end{definition}

\begin{definition}[Задача без ограничений]
Задача
\[
    \min_{x\in\RR^n}f(x)
\]
называется задачей без ограничений, если переменная \(x\) может
принимать любое значение из области определения \(f\).
\end{definition}

Максимизацию можно свести к минимизации:
\[
    \max_x f(x)
    \quad\Longleftrightarrow\quad
    \min_x\bigl(-f(x)\bigr).
\]
Поэтому далее достаточно рассматривать задачи минимизации.



Общая задача с ограничениями записывается как
\[
    \begin{aligned}
        \min_{x\in\RR^n}
        \quad&
        f(x),
        \\
        \text{при условиях}
        \quad&
        h_i(x)\le0,
        \qquad i=1,\ldots,m,
        \\
        &
        \ell_j(x)=0,
        \qquad j=1,\ldots,r.
    \end{aligned}
\]

Функции \(h_i\) задают ограничения-неравенства, а функции
\(\ell_j\) --- ограничения-равенства.

\begin{definition}[Допустимое множество]
Множество
\[
    \mathcal X
    =
    \left\{
        x\in\RR^n:
        \begin{array}{l}
            h_i(x)\le0,\quad i=1,\ldots,m,\\
            \ell_j(x)=0,\quad j=1,\ldots,r
        \end{array}
    \right\}
\]
называется допустимым множеством задачи.
\end{definition}

\begin{definition}[Допустимая точка]
Точка \(x\in\mathcal X\), удовлетворяющая всем ограничениям,
называется допустимой.
\end{definition}

\begin{definition}[Оптимальное решение]
Допустимая точка \(x^\star\in\mathcal X\) называется глобальным
решением задачи, если
\[
    f(x^\star)\le f(x)
    \qquad
    \text{для всех }x\in\mathcal X.
\]
Соответствующее значение
\[
    f^\star=f(x^\star)
\]
называется оптимальным значением.
\end{definition}

Таким образом, задача оптимизации состоит не в поиске произвольной
допустимой точки, а в выборе среди допустимых точек той, на которой
целевая функция принимает наименьшее значение.

\section{Выпуклая задача и методы оптимизации}

Задача
\[
    \begin{aligned}
        \min_x
        \quad& f(x),
        \\
        \text{при условиях}
        \quad& h_i(x)\le0,
        \\
        & \ell_j(x)=0
    \end{aligned}
\]
является выпуклой, если \(f\) и все функции \(h_i\) выпуклы, а
ограничения-равенства аффинны:
\[
    \ell_j(x)=a_j^{\mathsf T}x-b_j.
\]

В этом случае допустимое множество выпукло, а всякий локальный
минимум является глобальным. Формальные определения выпуклых
множеств и функций вводятся в следующей лекции.



Оптимизационные методы различаются по информации о целевой функции,
которую они используют.

\begin{definition}[Метод первого порядка]
Методом первого порядка называется алгоритм, использующий значения
функции и её градиента
\[
    \nabla f(x).
\]
\end{definition}

Типичным примером является градиентный спуск:
\[
    x^{k+1}=x^k-\eta_k\nabla f(x^k).
\]

\begin{definition}[Метод второго порядка]
Методом второго порядка называется алгоритм, дополнительно
использующий матрицу Гессе
\[
    \nabla^2f(x).
\]
\end{definition}

Если матрица Гессе в точке \(x^k\) обратима, классический шаг
метода Ньютона имеет вид
\[
    x^{k+1}
    =
    x^k
    -
    \bigl(\nabla^2f(x^k)\bigr)^{-1}
    \nabla f(x^k).
\]

Методы второго порядка учитывают локальную кривизну функции, но
вычисление, хранение и обращение матрицы Гессе может быть дорого при
большой размерности. К методам, использующим приближение информации
второго порядка, относятся квазиньютоновские методы, в частности
L-BFGS. В лекции также упоминается AdaHessian.

\section{Разреженное восстановление сигнала}

Пусть неизвестный сигнал \(x\in\RR^n\) наблюдается через линейную
модель
\[
    y=Mx+v,
\]
где
\[
    y\in\RR^s,\qquad
    M\in\RR^{s\times n},\qquad
    v\in\RR^s.
\]
Здесь \(y\) --- измеренный сигнал, \(M\) --- известная матрица
измерений, а \(v\) --- шум.

Если \(s\ll n\), линейная система недоопределена: число измерений
меньше числа неизвестных. Для восстановления \(x\) требуется
дополнительное структурное предположение. В задачах разреженного
восстановления предполагается, что большинство компонент \(x\)
равно нулю.

\begin{definition}[Разреженный вектор]
Носителем вектора \(x\in\RR^n\) называется множество
\[
    \operatorname{supp}(x)
    \coloneqq
    \{i\in\{1,\ldots,n\}:x_i\ne0\}.
\]
Вектор называется разреженным, если мощность его носителя мала по
сравнению с \(n\). Число ненулевых компонент обозначают через
\[
    \lVert x\rVert_0
    \coloneqq
    \bigl\lvert\operatorname{supp}(x)\bigr\rvert,
\]
хотя \(\lVert\cdot\rVert_0\) не является нормой.
\end{definition}

Оптимизация с непосредственным штрафом \(\lVert x\rVert_0\) в общем
случае невыпукла и комбинаторно сложна. Стандартная выпуклая замена ---
\(\ell_1\)-регуляризация.

\begin{problem}[\(\ell_1\)-регуляризованный метод наименьших квадратов]
Требуется решить задачу
\[
    \min_{x\in\RR^n}
    \left\{
        \frac12\lVert Mx-y\rVert_2^2
        +
        \lambda\lVert x\rVert_1
    \right\},
    \qquad
    \lambda>0.
\]
\end{problem}

Первое слагаемое измеряет согласованность восстановленного сигнала с
наблюдениями. Второе слагаемое
\[
    \lVert x\rVert_1
    =
    \sum_{j=1}^{n}\lvert x_j\rvert
\]
штрафует суммарную величину компонент и тем самым способствует
разреженности решения. Параметр \(\lambda\) задаёт компромисс между
точностью восстановления и силой регуляризации.

Задача не содержит явных ограничений и потому относится к задачам
без ограничений. При этом слагаемое \(\lVert x\rVert_1\) негладко в
точках с нулевыми компонентами, поэтому применяются, например,
субградиентные или проксимальные методы. Такие модели используются в
сжатых измерениях, обработке изображений и разреживании параметров
машинного обучения.



Пусть один и тот же неизвестный сигнал \(x\) наблюдается \(N\)
агентами:
\[
    y_i=M_ix+v_i,
    \qquad
    i=1,\ldots,N.
\]
Агент \(i\) знает только \(M_i\) и \(y_i\). Централизованная задача
имеет вид
\[
    \min_{x\in\RR^n}
    \left\{
        \frac12
        \sum_{i=1}^{N}
        \lVert M_ix-y_i\rVert_2^2
        +
        \lambda\lVert x\rVert_1
    \right\}.
\]

Для распределённой записи каждому агенту сопоставляется локальная
копия \(x_i\):
\[
    \begin{aligned}
        \min_{x_1,\ldots,x_N}
        \quad&
        \sum_{i=1}^{N}
        \left(
            \frac12\lVert M_ix_i-y_i\rVert_2^2
            +
            \frac{\lambda}{N}\lVert x_i\rVert_1
        \right),
        \\
        \text{при условиях}
        \quad&
        x_1=x_2=\cdots=x_N.
    \end{aligned}
\]

\begin{remark}
Множитель \(\lambda/N\) распределяет исходный регуляризатор между
агентами. При выполнении условия консенсуса \(x_i=x\) сумма локальных
регуляризаторов равна \(\lambda\lVert x\rVert_1\).
\end{remark}

Исходная централизованная задача без ограничений после введения
локальных копий становится задачей с ограничениями. Эти ограничения
необходимы, поскольку все агенты восстанавливают один и тот же
физический сигнал.

Если каждый агент решает только собственную задачу
\[
    \min_{x_i}
    \left\{
        \frac12\lVert M_ix_i-y_i\rVert_2^2
        +
        \frac{\lambda}{N}\lVert x_i\rVert_1
    \right\},
\]
то локальные оценки \(x_i^\star\) в общем случае различаются.
Коммуникации позволяют учитывать измерения других узлов без передачи
всего локального набора данных.

Для связного неориентированного графа условие глобального консенсуса
эквивалентно локальным равенствам
\[
    x_i=x_j,
    \qquad
    (i,j)\in\Edges.
\]
Связность графа тогда влечёт
\[
    x_1=\cdots=x_N.
\]

Распределённый алгоритм должен одновременно уменьшать локальные
невязки
\[
    \lVert M_ix_i-y_i\rVert_2
\]
и согласовывать оценки соседних агентов. Поэтому распределённая
оптимизация объединяет два процесса:
\[
    \text{оптимизацию локальных функций}
    \quad\text{и}\quad
    \text{достижение консенсуса}.
\]

\section{Линейная классификация и геометрия гиперплоскости}

Рассмотрим обучающую выборку
\[
    \{(z_i,y_i)\}_{i=1}^{m},
    \qquad
    z_i\in\RR^n,
    \qquad
    y_i\in\{-1,+1\}.
\]
Здесь \(z_i\) обозначает вектор признаков, а \(y_i\) --- метку класса.

Линейный классификатор задаётся гиперплоскостью
\[
    w^{\mathsf T}z+b=0,
\]
где
\[
    w\in\RR^n,\qquad b\in\RR.
\]
Вектор \(w\) является нормалью к гиперплоскости, а \(b\) задаёт её
смещение относительно начала координат.

Предсказанная метка определяется знаком:
\[
    \widehat y(z)
    =
    \operatorname{sign}\bigl(w^{\mathsf T}z+b\bigr).
\]
Для правильной классификации обучающей точки требуется
\[
    y_i\bigl(w^{\mathsf T}z_i+b\bigr)>0.
\]

Через одни и те же линейно разделимые данные можно провести несколько
разделяющих гиперплоскостей. Метод опорных векторов выбирает ту, для
которой минимальное расстояние от обучающих точек до гиперплоскости
максимально.



При \(w\ne0\) расстояние от точки \(z_i\) до гиперплоскости
\[
    w^{\mathsf T}z+b=0
\]
равно
\[
    d_i
    =
    \frac{
        \lvert w^{\mathsf T}z_i+b\rvert
    }{
        \lVert w\rVert_2
    }.
\]

Для правильно классифицированной точки
\[
    \lvert w^{\mathsf T}z_i+b\rvert
    =
    y_i\bigl(w^{\mathsf T}z_i+b\bigr),
\]
поэтому минимальный геометрический отступ равен
\[
    \gamma
    =
    \min_{1\le i\le m}
    \frac{
        y_i(w^{\mathsf T}z_i+b)
    }{
        \lVert w\rVert_2
    }.
\]

Параметры гиперплоскости определены неоднозначно: для любого
\(\alpha>0\) пары
\[
    (w,b)
    \qquad\text{и}\qquad
    (\alpha w,\alpha b)
\]
задают одну и ту же гиперплоскость. Если выборка линейно разделима,
положительным масштабированием параметры нормируют так, чтобы
\[
    \min_{1\le i\le m}
    y_i(w^{\mathsf T}z_i+b)
    =
    1.
\]
Тогда ограничения принимают вид
\[
    y_i(w^{\mathsf T}z_i+b)\ge1,
    \qquad i=1,\ldots,m,
\]
а геометрический отступ равен
\[
    \gamma=\frac{1}{\lVert w\rVert_2}.
\]

Максимизация \(\gamma\) эквивалентна минимизации нормы \(w\):
\[
    \max_{w,b}
    \frac{1}{\lVert w\rVert_2}
    \quad\Longleftrightarrow\quad
    \min_{w,b}
    \frac12\lVert w\rVert_2^2.
\]

\section{Метод опорных векторов и распределённая классификация}

\begin{problem}[SVM с жёстким зазором (hard margin)]
Для линейно разделимой выборки требуется решить задачу
\[
    \begin{aligned}
        \min_{w\in\RR^n,\ b\in\RR}
        \quad&
        \frac12\lVert w\rVert_2^2,
        \\
        \text{при условиях}
        \quad&
        y_i(w^{\mathsf T}z_i+b)\ge1,
        \qquad i=1,\ldots,m.
    \end{aligned}
\]
\end{problem}

Минимизация целевой функции при заданной нормировке эквивалентна
максимизации геометрического отступа, а ограничения требуют правильной
классификации всех обучающих точек.

\begin{definition}[Опорный вектор]
Обучающая точка \(z_i\) называется опорным вектором оптимального
классификатора, если соответствующее ограничение активно:
\[
    y_i\bigl(({w^\star})^{\mathsf T}z_i+b^\star\bigr)=1.
\]
\end{definition}

Опорные векторы являются ближайшими к разделяющей гиперплоскости
точками. Именно они определяют положение гиперплоскости максимального
отступа.

\begin{remark}
Функция
\[
    \frac12\lVert w\rVert_2^2
\]
квадратична и сильно выпукла по \(w\), но не зависит от \(b\) и
потому не является сильно выпуклой по паре \((w,b)\). Такая форма
удобнее для анализа и вычислений, чем непосредственная максимизация
\(1/\lVert w\rVert_2\).
\end{remark}



Пусть обучающие данные распределены между \(N\) агентами:
\[
    \Dataset_i
    =
    \{(z_{ij},y_{ij})\}_{j=1}^{m_i}.
\]
Агент \(i\) видит только собственную выборку \(\Dataset_i\).

Если каждый агент независимо строит локальный классификатор, он
решает
\[
    \begin{aligned}
        \min_{w_i,b_i}
        \quad&
        \frac12\lVert w_i\rVert_2^2,
        \\
        \text{при условиях}
        \quad&
        y_{ij}(w_i^{\mathsf T}z_{ij}+b_i)\ge1,
        \qquad j=1,\ldots,m_i.
    \end{aligned}
\]
Полученные гиперплоскости могут существенно различаться, поскольку
локальные выборки содержат разные участки общего распределения
данных.

Централизованный классификатор должен учитывать объединённую выборку
\[
    \Dataset
    =
    \bigcup_{i=1}^{N}\Dataset_i.
\]
Если объединённая выборка линейно разделима, её распределённая
постановка имеет вид
\[
    \begin{aligned}
        \min_{\substack{w_1,\ldots,w_N\\b_1,\ldots,b_N}}
        \quad&
        \frac{1}{2N}
        \sum_{i=1}^{N}\lVert w_i\rVert_2^2,
        \\
        \text{при условиях}
        \quad&
        y_{ij}(w_i^{\mathsf T}z_{ij}+b_i)\ge1,
        \qquad
        j=1,\ldots,m_i,
        \\
        &
        w_i=w_\ell,\qquad
        b_i=b_\ell,
        \qquad
        (i,\ell)\in\Edges.
    \end{aligned}
\]

При связном неориентированном графе локальные равенства приводят к
общим параметрам:
\[
    w_1=\cdots=w_N,
    \qquad
    b_1=\cdots=b_N.
\]
После подстановки консенсуса коэффициент \(1/N\) сокращает сумму
одинаковых квадратичных слагаемых, поэтому распределённая задача
эквивалентна централизованной. При этом каждый агент использует только
собственные данные и сообщения соседей.

Практическая сложность состоит в том, что локально оптимальная
гиперплоскость может плохо классифицировать данные, отсутствующие у
данного агента. Поэтому обмен информацией необходим не только для
согласования параметров, но и для приближения решения, построенного
по совокупной выборке.

\section{Неразделимые данные}

Жёстко-разделяющая постановка требует существования \(w,b\), для
которых
\[
    y_i(w^{\mathsf T}z_i+b)\ge1
\]
выполняется для всех точек. Если классы линейно неразделимы или
некоторые метки ошибочны, допустимое множество такой задачи может
оказаться пустым.

В лекции отмечаются два стандартных обобщения. SVM с мягким зазором
допускает нарушения ограничений и штрафует их величину. Ядерный SVM
неявно работает в пространстве признаков, где разделяющая поверхность
может быть линейной. Подробный вывод этих моделей в лекции не
проводится.


Общая задача оптимизации задаётся целевой функцией, переменной решения
и системой ограничений. Допустимое множество содержит точки,
удовлетворяющие всем ограничениям. Методы первого порядка используют
градиенты, а методы второго порядка дополнительно используют
кривизну функции. В разреженном восстановлении \(\ell_1\)-регуляризация
способствует получению решений с небольшим числом ненулевых компонент. При
распределении измерений между агентами вводятся локальные копии
неизвестного сигнала и ограничения консенсуса. В методе опорных
векторов максимизация минимального геометрического отступа приводит
к выпуклой квадратичной задаче
\[
    \begin{aligned}
        \min_{w,b}\quad&\frac12\lVert w\rVert_2^2,\\
        \text{при условиях}\quad&
        y_i(w^{\mathsf T}z_i+b)\ge1,
        \qquad i=1,\ldots,m.
    \end{aligned}
\]
При распределённых данных локальные классификаторы могут существенно
отличаться от решения по объединённой выборке, поэтому оптимизация
должна сочетаться с коммуникацией и согласованием параметров.

\lecture{05}{Выпуклые множества и выпуклые функции}

\sourcevideo
    {https://www.youtube.com/playlist?list=PLOzRYVm0a65fhKDS8cYIC7YCuVGJ4FOJx}
    {Week 1: Lecture 5: Convex Sets and Convex Functions}

Выпуклость обеспечивает глобальные гарантии в задачах оптимизации.
На выпуклой области всякий локальный минимум выпуклой функции является
глобальным. Для дифференцируемых функций выпуклость можно проверять по
градиенту, а для дважды дифференцируемых --- по матрице Гессе.

\section{Выпуклые комбинации и множества}

Пусть \(x,y\in\RR^n\).

\begin{definition}[Линейная комбинация]
Точка
\[
    z=\alpha x+\beta y,
    \qquad
    \alpha,\beta\in\RR,
\]
называется линейной комбинацией \(x\) и \(y\).
\end{definition}

\begin{definition}[Аффинная комбинация]
Линейная комбинация называется аффинной, если сумма коэффициентов
равна единице:
\[
    z=\alpha x+\beta y,
    \qquad
    \alpha+\beta=1.
\]
Для двух точек её можно записать как
\[
    z=\theta x+(1-\theta)y,
    \qquad
    \theta\in\RR.
\]
\end{definition}

\begin{definition}[Выпуклая комбинация]
Аффинная комбинация называется выпуклой, если её коэффициенты
неотрицательны:
\[
    z=\theta x+(1-\theta)y,
    \qquad
    \theta\in[0,1].
\]
\end{definition}

При изменении \(\theta\) от \(0\) до \(1\) точка \(z\) пробегает
отрезок, соединяющий \(x\) и \(y\).

Для точек \(x_1,\ldots,x_m\in\RR^n\) их выпуклая комбинация имеет вид
\[
    z=\sum_{i=1}^{m}\theta_i x_i,
    \qquad
    \theta_i\ge0,
    \qquad
    \sum_{i=1}^{m}\theta_i=1.
\]



\begin{definition}[Выпуклое множество]
Множество \(C\subseteq\RR^n\) называется выпуклым, если
\[
    x,y\in C,\quad \theta\in[0,1]
    \quad\Longrightarrow\quad
    \theta x+(1-\theta)y\in C.
\]
\end{definition}

Иными словами, вместе с любыми двумя точками выпуклое множество
содержит весь соединяющий их отрезок.

Интервалы, аффинные подпространства и евклидовы шары являются
выпуклыми. Отсутствие разрывов или отверстий недостаточно: множество
выпукло именно тогда, когда содержит отрезок между каждой парой своих
точек. Например,
\[
    C=\{2,5\}
\]
не выпукло, поскольку
\[
    \frac12\cdot2+\frac12\cdot5=3.5\notin C.
\]

\begin{definition}[Выпуклая оболочка]
Выпуклой оболочкой множества \(S\subseteq\RR^n\) называется множество
всех конечных выпуклых комбинаций его точек:
\[
    \operatorname{conv}(S)
    =
    \left\{
        \sum_{i=1}^{m}\theta_i x_i
        \,\middle|\,
        m\in\mathbb N,\quad
        x_i\in S,\quad
        \theta_i\ge0,\quad
        \sum_{i=1}^{m}\theta_i=1
    \right\}.
\]
\end{definition}

\section{Выпуклые функции и минимумы}

Пусть \(C\subseteq\RR^n\) --- выпуклое множество.

\begin{definition}[Выпуклая функция]
Функция \(f\colon C\to\RR\) называется выпуклой, если для любых
\(x,y\in C\) и \(\theta\in[0,1]\)
\[
    f\bigl(\theta x+(1-\theta)y\bigr)
    \le
    \theta f(x)+(1-\theta)f(y).
\]
\end{definition}

Правая часть задаёт значение хорды между точками
\[
    (x,f(x))
    \quad\text{и}\quad
    (y,f(y)).
\]
Следовательно, граф выпуклой функции лежит не выше любой своей хорды.

Требование выпуклости области определения существенно: точка
\[
    \theta x+(1-\theta)y
\]
должна принадлежать области определения функции.

\begin{definition}[Строго выпуклая функция]
Функция \(f\colon C\to\RR\) называется строго выпуклой, если для
различных \(x,y\in C\) и каждого \(\theta\in(0,1)\)
\[
    f\bigl(\theta x+(1-\theta)y\bigr)
    <
    \theta f(x)+(1-\theta)f(y).
\]
\end{definition}

Постоянная функция выпукла, но не строго выпукла. Выпуклая функция
может иметь несколько минимумов, тогда как строго выпуклая функция
имеет не более одного глобального минимума.



\begin{theorem}
Пусть \(C\) выпукло, а \(f\colon C\to\RR\) выпукла. Тогда всякий
локальный минимум \(f\) является глобальным.
\end{theorem}

\begin{proof}
Пусть \(x^\star\) --- локальный минимум относительно множества
\(C\), но существует \(y\in C\),
для которого
\[
    f(y)<f(x^\star).
\]
Для \(\theta\in(0,1)\) положим
\[
    x_\theta=(1-\theta)x^\star+\theta y.
\]
При достаточно малом \(\theta\) точка \(x_\theta\) сколь угодно
близка к \(x^\star\). По выпуклости
\[
    f(x_\theta)
    \le
    (1-\theta)f(x^\star)+\theta f(y)
    <
    f(x^\star),
\]
что противоречит локальной минимальности \(x^\star\).
\end{proof}

\begin{corollary}
Если строго выпуклая функция достигает минимума, то этот минимум
единственен.
\end{corollary}

Для общей невыпуклой функции локальные минимумы могут иметь различные
значения. Поэтому локальная сходимость алгоритма не обеспечивает
глобальной оптимальности.



Функции
\[
    f(x)=\frac12x^2,
    \qquad
    f(x)=\frac14x^4,
    \qquad
    f(x)=\mathrm{e}^x
\]
выпуклы на \(\RR\).

Функция
\[
    f(x)=c
\]
одновременно выпукла и вогнута.

Функция
\[
    f(x)=\log x,
    \qquad x>0,
\]
не выпукла, но вогнута. Поэтому
\[
    -\log x
\]
выпукла на \(\RR_{++}\).

Функция может быть одновременно невыпуклой и невогнутой. Поэтому
невыпуклость не означает вогнутость.

\begin{definition}[Вогнутая функция]
Функция \(f\) называется вогнутой, если функция \(-f\) выпукла.
Эквивалентно,
\[
    f\bigl(\theta x+(1-\theta)y\bigr)
    \ge
    \theta f(x)+(1-\theta)f(y).
\]
\end{definition}

После смены знака минимизация вогнутой функции превращается в
максимизацию выпуклой функции, которая в общем случае не является
задачей выпуклой оптимизации. Стандартные выпуклые постановки ---
минимизация выпуклой или максимизация вогнутой функции на выпуклом
множестве.

\section{Операции, сохраняющие выпуклость}



\begin{theorem}
Пусть \(f_1,\ldots,f_m\) выпуклы на общем выпуклом множестве \(C\), а
\[
    \alpha_i\ge0.
\]
Тогда функция
\[
    f(x)=\sum_{i=1}^{m}\alpha_i f_i(x)
\]
выпукла на \(C\).
\end{theorem}

\begin{proof}
Для \(x,y\in C\) и \(\theta\in[0,1]\)
\[
\begin{aligned}
    f\bigl(\theta x+(1-\theta)y\bigr)
    &=
    \sum_{i=1}^{m}
    \alpha_i
    f_i\bigl(\theta x+(1-\theta)y\bigr)
    \\
    &\le
    \sum_{i=1}^{m}
    \alpha_i
    \bigl(
        \theta f_i(x)+(1-\theta)f_i(y)
    \bigr)
    \\
    &=
    \theta f(x)+(1-\theta)f(y).
\end{aligned}
\]
\end{proof}

Поскольку функция \(t\mapsto\lvert t\rvert\) выпукла,
\[
    \lVert x\rVert_1
    =
    \sum_{i=1}^{n}\lvert x_i\rvert
\]
также выпукла.

Аналогично функция
\[
    f(x)
    =
    \frac12x^2+\frac14x^4
\]
выпукла как сумма выпуклых функций.



\begin{theorem}
Пусть \(f_1,\ldots,f_m\) выпуклы на \(C\). Тогда функция
\[
    f(x)=\max_{1\le i\le m}f_i(x)
\]
выпукла на \(C\).
\end{theorem}

\begin{proof}
Для каждого \(i\)
\[
\begin{aligned}
    f_i\bigl(\theta x+(1-\theta)y\bigr)
    &\le
    \theta f_i(x)+(1-\theta)f_i(y)
    \\
    &\le
    \theta f(x)+(1-\theta)f(y).
\end{aligned}
\]
Переход к максимуму по \(i\) даёт требуемое неравенство.
\end{proof}

Поточечные максимумы возникают в негладких целевых функциях,
минимаксных постановках и задачах робастной оптимизации.



\begin{theorem}
Пусть \(D\subseteq\RR^m\) выпукло, \(f\colon D\to\RR\) выпукла,
\(A\in\RR^{m\times n}\) и \(b\in\RR^m\). Тогда множество
\[
    C=\{x\in\RR^n:Ax+b\in D\}
\]
выпукло, а функция
\[
    g(x)=f(Ax+b)
\]
выпукла на \(C\).
\end{theorem}

\begin{proof}
Для \(x,y\in C\) и \(\theta\in[0,1]\)
\[
    A\bigl(\theta x+(1-\theta)y\bigr)+b
    =
    \theta(Ax+b)+(1-\theta)(Ay+b)\in D,
\]
поэтому \(C\) выпукло. Кроме того,
\[
\begin{aligned}
    g\bigl(\theta x+(1-\theta)y\bigr)
    &=
    f\bigl(
        \theta(Ax+b)+(1-\theta)(Ay+b)
    \bigr)
    \\
    &\le
    \theta f(Ax+b)+(1-\theta)f(Ay+b).
\end{aligned}
\]
\end{proof}

В частности, поскольку квадрат евклидовой нормы выпуклый,
\[
    f(x)=\frac12\lVert Ax-b\rVert_2^2
\]
является выпуклой функцией.

\section{Логарифмический барьер}

Для строгого линейного ограничения
\[
    a^{\mathsf T}x<b
\]
можно использовать барьер
\[
    \phi(x)
    =
    -\log\bigl(b-a^{\mathsf T}x\bigr).
\]
Функция определена только внутри допустимой области. При приближении
к её границе
\[
    b-a^{\mathsf T}x\to0^+
\]
выполняется
\[
    \phi(x)\to+\infty.
\]

Для ограничений \(a_i^{\mathsf T}x\le b_i\) барьерный метод
рассматривает семейство задач
\[
    \min_{x:\,a_i^{\mathsf T}x<b_i}
    \left\{
        f(x)
        -
        \mu
        \sum_{i=1}^{m}
        \log\bigl(b_i-a_i^{\mathsf T}x\bigr)
    \right\},
    \qquad
    \mu>0.
\]
При уменьшении \(\mu\) решения барьерных задач приближаются к решению
исходной задачи при стандартных условиях регулярности. Поскольку
\(-\log t\) выпукла на \(t>0\), а аффинная композиция сохраняет
выпуклость, каждый барьер выпуклый на своей области определения.

\section{Условия выпуклости первого и второго порядка}

\begin{theorem}[Критерий первого порядка]
Пусть \(C\subseteq\RR^n\) открыто и выпукло, а
\(f\colon C\to\RR\) непрерывно дифференцируема. Тогда \(f\) выпукла
на \(C\) тогда и только тогда, когда
\[
    f(y)
    \ge
    f(x)
    +
    \langle\nabla f(x),y-x\rangle
\]
для всех \(x,y\in C\).
\end{theorem}

Геометрически это означает, что граф выпуклой дифференцируемой
функции лежит не ниже любой касательной гиперплоскости.

\begin{proof}
Пусть сначала \(f\) выпукла. Для \(\lambda\in(0,1]\)
\[
\begin{aligned}
    f\bigl(x+\lambda(y-x)\bigr)
    &=
    f\bigl((1-\lambda)x+\lambda y\bigr)
    \\
    &\le
    (1-\lambda)f(x)+\lambda f(y).
\end{aligned}
\]
Следовательно,
\[
    \frac{
        f\bigl(x+\lambda(y-x)\bigr)-f(x)
    }{\lambda}
    \le
    f(y)-f(x).
\]
Переход к пределу при \(\lambda\to0^+\) даёт
\[
    \langle\nabla f(x),y-x\rangle
    \le
    f(y)-f(x).
\]

Обратно, пусть неравенство первого порядка выполнено. Для
\(z=\theta x+(1-\theta)y\) применим его к парам \((z,x)\) и
\((z,y)\):
\[
\begin{aligned}
    f(x)&\ge f(z)+\langle\nabla f(z),x-z\rangle,\\
    f(y)&\ge f(z)+\langle\nabla f(z),y-z\rangle.
\end{aligned}
\]
Умножая неравенства на \(\theta\) и \(1-\theta\) и складывая,
получаем
\[
    \theta f(x)+(1-\theta)f(y)\ge f(z),
\]
поскольку
\(\theta(x-z)+(1-\theta)(y-z)=0\). Следовательно, \(f\) выпукла.
\end{proof}

\begin{corollary}[Условие глобального минимума]
Пусть \(C\) открыто и выпукло, \(f\colon C\to\RR\) выпукла и
дифференцируема. Если \(x^\star\in C\) и
\[
    \nabla f(x^\star)=0,
\]
то \(x^\star\) является глобальным минимумом \(f\) на \(C\).
\end{corollary}

\begin{proof}
По критерию первого порядка для любого \(y\in C\)
\[
    f(y)
    \ge
    f(x^\star)
    +
    \langle\nabla f(x^\star),y-x^\star\rangle
    =
    f(x^\star).
\]
\end{proof}



\begin{theorem}[Критерий второго порядка]
Пусть \(C\subseteq\RR^n\) открыто и выпукло, а
\(f\colon C\to\RR\) дважды непрерывно дифференцируема. Тогда \(f\)
выпукла на \(C\) тогда и только тогда, когда
\[
    \nabla^2f(x)\succeq0
    \qquad
    \text{для всех }x\in C.
\]
\end{theorem}

\begin{proof}
Если \(f\) выпукла, то для любых \(x\in C\) и \(v\in\RR^n\)
одномерная функция \(g(t)=f(x+tv)\) выпукла в окрестности нуля.
Следовательно,
\[
    v^{\mathsf T}\nabla^2f(x)v=g''(0)\ge0.
\]

Обратно, пусть \(\nabla^2f(x)\succeq0\) на \(C\). Для любых
\(x,y\in C\) функция
\[
    g(t)=f\bigl(x+t(y-x)\bigr),
    \qquad t\in[0,1],
\]
удовлетворяет
\[
    g''(t)
    =
    (y-x)^{\mathsf T}
    \nabla^2f\bigl(x+t(y-x)\bigr)
    (y-x)
    \ge0.
\]
Поэтому \(g\) выпукла на \([0,1]\), что равносильно выпуклости
\(f\) на каждом отрезке в \(C\).
\end{proof}

Положительная полуопределённость означает
\[
    v^{\mathsf T}\nabla^2f(x)v\ge0
    \qquad
    \text{для всех }v\in\RR^n.
\]

В одномерном случае критерий принимает вид
\[
    f''(x)\ge0.
\]

Для функции
\[
    f(x)=\frac12x^2
\]
имеем
\[
    f''(x)=1>0,
\]
поэтому она выпукла.

Для функции
\[
    f(x)=\frac14x^4
\]
получаем
\[
    f''(x)=3x^2\ge0.
\]
Следовательно, она также выпукла, хотя
\[
    f''(0)=0.
\]

\section{Схема проверки выпуклости}

\begin{algorithm}[H]
\caption{Проверка выпуклости функции}
\label{alg:convexity-check}
\begin{algorithmic}[1]

\Require
функция \(f\colon C\to\RR\)

\State
проверить выпуклость области \(C\)

\If{\(f\) дважды непрерывно дифференцируема}

    \State
    вычислить \(\nabla^2f(x)\)

    \State
    проверить
    \(\nabla^2f(x)\succeq0\) для всех \(x\in C\)

\ElsIf{\(f\) непрерывно дифференцируема}

    \State
    проверить
    \[
        f(y)
        \ge
        f(x)+
        \langle\nabla f(x),y-x\rangle
    \]
    для всех \(x,y\in C\)

\Else

    \State
    проверить определение выпуклости непосредственно

\EndIf

\Ensure
заключение о выпуклости \(f\)

\end{algorithmic}
\end{algorithm}


Множество выпукло, если оно содержит отрезок между любыми двумя
своими точками. На выпуклой области всякий локальный минимум выпуклой
функции является глобальным. Неотрицательные суммы, поточечные
максимумы и аффинные композиции сохраняют выпуклость. Для непрерывно
дифференцируемой функции выпуклость эквивалентна условию
\[
    f(y)
    \ge
    f(x)+\langle\nabla f(x),y-x\rangle,
\]
а для дважды непрерывно дифференцируемой функции --- условию
\[
    \nabla^2f(x)\succeq0.
\]
\lecture{06}{Строгая и сильная выпуклость}

\sourcevideo
    {https://www.youtube.com/playlist?list=PLOzRYVm0a65fhKDS8cYIC7YCuVGJ4FOJx}
    {Week 2: Lecture 6: Strictly and Strongly Convex Functions}

Выпуклость гарантирует глобальность локальных минимумов, но не
единственность решения и не задаёт количественной оценки кривизны.
Строгая выпуклость обеспечивает единственность минимума, если он
существует, а сильная выпуклость вводит равномерную квадратичную
нижнюю оценку. Она позволяет получать явные оценки скорости
сходимости.

\section{Строгая выпуклость и условие второго порядка}

Пусть \(C\subseteq\RR^n\) --- выпуклое множество.

\begin{definition}[Строго выпуклая функция]
Функция \(f\colon C\to\RR\) называется строго выпуклой, если для
любых различных \(x,y\in C\) и каждого \(\theta\in(0,1)\)
\[
    f\bigl(\theta x+(1-\theta)y\bigr)
    <
    \theta f(x)+(1-\theta)f(y).
\]
\end{definition}

Для обычной выпуклости допускается равенство. Поэтому выпуклая
функция может иметь плоский участок и несколько точек минимума.
Строгая выпуклость исключает такой участок между различными точками.

\begin{theorem}[Единственность минимума]
Если строго выпуклая функция достигает глобального минимума, то её
точка минимума единственна.
\end{theorem}

\begin{proof}
Предположим, что различные точки \(x^\star\) и \(y^\star\) являются
минимумами:
\[
    f(x^\star)=f(y^\star)=f^\star.
\]
Для \(\theta\in(0,1)\) строгая выпуклость даёт
\[
\begin{aligned}
    f\bigl(\theta x^\star+(1-\theta)y^\star\bigr)
    &<
    \theta f(x^\star)+(1-\theta)f(y^\star)
    \\
    &=f^\star,
\end{aligned}
\]
что противоречит минимальности \(f^\star\).
\end{proof}

Условие существования существенно: строгая выпуклость сама по себе
не гарантирует, что минимум достигается. Например,
\[
    f(x)=\mathrm{e}^x
\]
строго выпукла на \(\RR\), но не имеет точки глобального минимума.



\begin{theorem}[Достаточное условие строгой выпуклости]
Пусть \(C\subseteq\RR^n\) открыто и выпукло, а
\(f\in C^2(C)\). Если
\[
    \nabla^2 f(x)\succ0
    \qquad
    \text{для всех }x\in C,
\]
то \(f\) строго выпукла на \(C\).
\end{theorem}

Положительная определённость означает
\[
    v^{\mathsf T}\nabla^2f(x)v>0
    \qquad
    \text{для всех }v\ne0.
\]

Это условие достаточно, но не необходимо. Рассмотрим
\[
    f(x)=x^4.
\]
Функция строго выпукла и имеет единственный минимум в нуле, однако
\[
    f''(x)=12x^2,
    \qquad
    f''(0)=0.
\]
Следовательно, матрица Гессе не обязана быть положительно определена
в каждой точке строго выпуклой функции.

Для функции
\[
    f(x)=x^2
\]
имеем
\[
    f''(x)=2>0,
\]
поэтому её строгая выпуклость непосредственно следует из критерия
второго порядка.

Таким образом,
\[
    \nabla^2f(x)\succ0
    \quad\Longrightarrow\quad
    f\text{ строго выпукла},
\]
но обратная импликация в общем случае неверна.

\section{Сильная выпуклость и её характеризации}

Строгая выпуклость утверждает, что значение на внутренней точке
отрезка строго меньше значения соответствующей хорды, но не оценивает
величину разности.

\begin{definition}[Сильная выпуклость]
Функция \(f\colon C\to\RR\) называется \(\mu\)-сильно выпуклой,
где \(\mu>0\), если для любых \(x,y\in C\) и
\(\theta\in[0,1]\)
\[
\begin{aligned}
    f\bigl(\theta x+(1-\theta)y\bigr)
    \le{}&
    \theta f(x)+(1-\theta)f(y)
    \\
    &-
    \frac{\mu}{2}
    \theta(1-\theta)\norm{x-y}^2.
\end{aligned}
\]
Число \(\mu\) называется модулем сильной выпуклости.
\end{definition}

Дополнительное квадратичное слагаемое задаёт равномерную нижнюю
оценку кривизны. Чем больше допустимое значение \(\mu\), тем сильнее
функция отклоняется вниз от своих хорд.

Справедлива цепочка импликаций
\[
    \text{сильная выпуклость}
    \Longrightarrow
    \text{строгая выпуклость}
    \Longrightarrow
    \text{выпуклость}.
\]
Обратные импликации неверны.

Функция \(x^2\) сильно выпукла, тогда как \(x^4\) строго выпукла,
но не сильно выпукла на всей прямой. Действительно,
\[
    \inf_{x\in\RR}12x^2=0,
\]
поэтому её кривизну нельзя равномерно отделить от нуля.



\begin{theorem}[Критерий первого порядка]
Пусть \(C\subseteq\RR^n\) открыто и выпукло, а
\(f\in C^1(C)\). Функция \(f\) является \(\mu\)-сильно выпуклой
тогда и только тогда, когда
\[
    f(y)
    \ge
    f(x)
    +
    \inner{\nabla f(x)}{y-x}
    +
    \frac{\mu}{2}\norm{y-x}^2
\]
для всех \(x,y\in C\).
\end{theorem}

Для обычной выпуклой функции имеется только опорная гиперплоскость:
\[
    f(y)
    \ge
    f(x)+\inner{\nabla f(x)}{y-x}.
\]
При сильной выпуклости функция лежит выше этой гиперплоскости как
минимум на величину
\[
    \frac{\mu}{2}\norm{y-x}^2.
\]

\begin{corollary}[Квадратичный рост около минимума]
Пусть \(C\subseteq\RR^n\) открыто и выпукло, а
\(f\in C^1(C)\) является \(\mu\)-сильно выпуклой. Если
\(x^\star\in C\) --- точка минимума, то для всех \(x\in C\)
\[
    f(x)-f^\star
    \ge
    \frac{\mu}{2}\norm{x-x^\star}^2.
\]
\end{corollary}

\begin{proof}
Поскольку \(x^\star\) --- внутренняя точка минимума,
\(\nabla f(x^\star)=0\). Подстановка \(x=x^\star\) в критерий
первого порядка даёт требуемое неравенство.
\end{proof}

Это неравенство связывает ошибку по значению функции с расстоянием
до оптимального решения.



\begin{theorem}[Критерий второго порядка]
Пусть \(C\subseteq\RR^n\) открыто и выпукло, а
\(f\in C^2(C)\). Тогда \(f\) является \(\mu\)-сильно выпуклой на
\(C\) тогда и только тогда, когда
\[
    \nabla^2f(x)\succeq\mu I_n
    \qquad
    \text{для всех }x\in C.
\]
\end{theorem}

Эквивалентно,
\[
    v^{\mathsf T}\nabla^2f(x)v
    \ge
    \mu\norm{v}^2
\]
для всех \(x\in C\) и \(v\in\RR^n\).

В одномерном случае условие принимает вид
\[
    f''(x)\ge\mu.
\]

Для функции
\[
    f(x)=x^2
\]
имеем
\[
    f''(x)=2,
\]
поэтому она \(2\)-сильно выпукла. Она также является
\(\mu\)-сильно выпуклой для любого \(0<\mu\le2\), но максимальный
модуль равен \(2\).

Для квадратичной функции
\[
    f(x)=\frac12x^{\mathsf T}Qx+b^{\mathsf T}x+c,
    \qquad Q=Q^{\mathsf T},
\]
матрица Гессе равна
\[
    \nabla^2f(x)=Q.
\]
Следовательно,
\[
    f\text{ выпукла}
    \Longleftrightarrow
    Q\succeq0,
\]
а
\[
    f\text{ \(\mu\)-сильно выпукла}
    \Longleftrightarrow
    Q\succeq\mu I_n.
\]
Если \(Q\succ0\), наибольший допустимый модуль сильной выпуклости
равен
\[
    \mu_{\max}=\lambda_{\min}(Q)>0.
\]
При \(\lambda_{\min}(Q)=0\) функция выпукла, но не сильно выпукла на
всём пространстве.



Из критерия первого порядка следует важное свойство градиента.

\begin{theorem}
Если \(f\in C^1(C)\) является \(\mu\)-сильно выпуклой, то
\[
    \inner{\nabla f(x)-\nabla f(y)}{x-y}
    \ge
    \mu\norm{x-y}^2
\]
для всех \(x,y\in C\).
\end{theorem}

\begin{proof}
Сильная выпуклость даёт
\[
    f(y)
    \ge
    f(x)
    +
    \inner{\nabla f(x)}{y-x}
    +
    \frac{\mu}{2}\norm{y-x}^2
\]
и после перестановки \(x,y\)
\[
    f(x)
    \ge
    f(y)
    +
    \inner{\nabla f(y)}{x-y}
    +
    \frac{\mu}{2}\norm{x-y}^2.
\]
Складывая неравенства, получаем
\[
    \inner{\nabla f(x)-\nabla f(y)}{x-y}
    \ge
    \mu\norm{x-y}^2.
\]
\end{proof}

При \(y=x^\star\) и \(\nabla f(x^\star)=0\)
\[
    \inner{\nabla f(x)}{x-x^\star}
    \ge
    \mu\norm{x-x^\star}^2.
\]
По неравенству Коши--Буняковского отсюда следует
\[
    \norm{\nabla f(x)}
    \ge
    \mu\norm{x-x^\star}.
\]

Следовательно, градиент сильно выпуклой функции количественно
контролирует расстояние до минимума. Строгая выпуклость сама по себе
такой оценки не даёт: для \(f(x)=x^4\) градиент \(4x^3\) быстро
исчезает около нуля.

\section{Негладкие выпуклые функции}

Для негладкой функции градиент может не существовать. В этом случае
используется понятие субградиента.

\begin{definition}[Субградиент]
Пусть \(f\colon C\to\RR\) выпукла. Вектор \(g\in\RR^n\) называется
субградиентом \(f\) в точке \(x\in C\), если
\[
    f(y)
    \ge
    f(x)+\inner{g}{y-x}
\]
для всех \(y\in C\).
Множество всех субградиентов обозначается через
\[
    \partial f(x).
\]
\end{definition}

Если \(f\) дифференцируема в \(x\), то
\[
    \partial f(x)=\set{\nabla f(x)}.
\]

Для функции
\[
    f(x)=\abs{x}
\]
имеем
\[
    \partial f(x)
    =
    \begin{cases}
        \set{1}, & x>0,\\
        [-1,1], & x=0,\\
        \set{-1}, & x<0.
    \end{cases}
\]

В точке \(x=0\) обычная производная не существует, но каждая прямая
с наклоном \(g\in[-1,1]\) лежит не выше графика функции.

\begin{definition}[Сильная выпуклость через субградиенты]
Выпуклая функция \(f\) является \(\mu\)-сильно выпуклой, если для
любых \(x,y\in C\) и каждого \(g\in\partial f(x)\)
\[
    f(y)
    \ge
    f(x)+\inner{g}{y-x}
    +
    \frac{\mu}{2}\norm{y-x}^2.
\]
\end{definition}

Это определение позволяет анализировать сильно выпуклые функции,
содержащие негладкие слагаемые.

\section{Неравенство Поляка--Лоясиевича}

\begin{definition}[Неравенство Поляка--Лоясиевича]
Дифференцируемая функция \(f\colon\RR^n\to\RR\) удовлетворяет
неравенству Поляка--Лоясиевича с константой \(\mu>0\), если
\[
    \frac12\norm{\nabla f(x)}^2
    \ge
    \mu\bigl(f(x)-f^\star\bigr)
\]
для всех \(x\), где
\[
    f^\star=\inf_x f(x).
\]
\end{definition}

Это условие также называют условием градиентного доминирования.
Оно утверждает, что большой разрыв по значению функции сопровождается
достаточно большим градиентом.

\begin{theorem}
Всякая дифференцируемая \(\mu\)-сильно выпуклая функция
удовлетворяет неравенству Поляка--Лоясиевича с той же константой
\(\mu\).
\end{theorem}

\begin{proof}
Из сильной выпуклости для любого \(y\)
\[
    f(y)
    \ge
    f(x)
    +
    \inner{\nabla f(x)}{y-x}
    +
    \frac{\mu}{2}\norm{y-x}^2.
\]
Беря инфимум по \(y\) в обеих частях, получаем
\[
    f^\star
    \ge
    \inf_y
    \left\{
        f(x)
        +
        \inner{\nabla f(x)}{y-x}
        +
        \frac{\mu}{2}\norm{y-x}^2
    \right\}.
\]
Минимум выражения в фигурных скобках достигается при
\[
    y=x-\frac1\mu\nabla f(x)
\]
и равен
\[
    f(x)-\frac{1}{2\mu}\norm{\nabla f(x)}^2.
\]
После перестановки слагаемых получаем PL-неравенство.
\end{proof}

Обратное утверждение неверно: функция может удовлетворять
неравенству Поляка--Лоясиевича, не будучи сильно выпуклой и даже
не будучи выпуклой.

В лекции в качестве такого примера приводится функция
\[
    f(x)=x^2+3\sin^2x.
\]
Она имеет единственный глобальный минимум в нуле и не является
выпуклой на всей прямой, поскольку
\[
    f''(x)=2+6\cos(2x)
\]
принимает отрицательные значения. Для включения функции в класс PL
требуется отдельная глобальная оценка
\[
    \frac12\lvert f'(x)\rvert^2
    \ge
    \mu f(x).
\]
В лекции такая оценка не доказывается, поэтому пример фиксируется как
заявленный, но не доказанный.

Корректное соотношение классов имеет вид
\[
    \{\text{сильно выпуклые функции}\}
    \subset
    \{\text{PL-функции}\},
\]
а также
\[
    \{\text{сильно выпуклые}\}
    \subset
    \{\text{строго выпуклые}\}
    \subset
    \{\text{выпуклые}\}.
\]
При этом класс PL-функций пересекается с классом выпуклых функций,
но не содержится в нём и не совпадает с ним.

\begin{remark}
Класс PL-функций не образует одну цепочку вложений с классами
выпуклых, строго выпуклых и сильно выпуклых функций. Существуют как
выпуклые, но не сильно выпуклые PL-функции, так и невыпуклые
PL-функции.
\end{remark}

\section{Аффинная композиция}

Пусть \(f\colon\RR^m\to\RR\) является \(\mu\)-сильно выпуклой и
\[
    g(x)=f(Ax+b),
    \qquad
    A\in\RR^{m\times n}.
\]
Для любых \(x,y\in\RR^n\) и \(\theta\in[0,1]\)
\[
\begin{aligned}
    g\bigl(\theta x+(1-\theta)y\bigr)
    \le{}&
    \theta g(x)+(1-\theta)g(y)
    \\
    &-
    \frac{\mu}{2}\theta(1-\theta)
    \norm{A(x-y)}^2.
\end{aligned}
\]
Если \(A\) имеет полный столбцовый ранг, то
\[
    \norm{Av}
    \ge
    \sigma_{\min}(A)\norm{v},
\]
и потому \(g\) является
\(\mu\sigma_{\min}^2(A)\)-сильно выпуклой.

Если \(\ker A\ne\{0\}\), то для некоторого \(v\ne0\)
\[
    g(x+tv)=g(x)
    \qquad
    \text{для всех }t\in\RR.
\]
Следовательно, композиция на всём пространстве не может быть строго,
а значит, и сильно выпуклой.

\begin{remark}[Уточнение рангового условия]
Для композиции \(x\mapsto f(Ax+b)\), определённой на \(\RR^n\),
необходима инъективность отображения \(x\mapsto Ax\), то есть полный
столбцовый ранг \(A\) или, эквивалентно, \(\ker A=\{0\}\).
\end{remark}

Вырожденная квадратичная функция всё же может удовлетворять
PL-неравенству. Пусть
\[
    g(x)=\frac12\norm{Ax-b}^2
\]
и \(x^\star\) --- любое решение задачи наименьших квадратов. Тогда
\[
    g(x)-g^\star
    =
    \frac12\norm{A(x-x^\star)}^2,
\]
а
\[
    \nabla g(x)=A^{\mathsf T}A(x-x^\star).
\]
Если \(A\ne0\), то
\[
    \frac12\norm{\nabla g(x)}^2
    \ge
    \sigma_{\min,+}^2(A)\bigl(g(x)-g^\star\bigr),
\]
где \(\sigma_{\min,+}(A)\) --- наименьшее положительное сингулярное
число \(A\). При нетривиальном ядре минимум не единственен, но
PL-неравенство сохраняется.

\section{Практическое значение}

Строгая выпуклость полезна прежде всего для доказательства
единственности решения. Она не задаёт равномерной скорости роста
функции около минимума.

Сильная выпуклость даёт количественные оценки:
\[
    f(x)-f^\star
    \ge
    \frac{\mu}{2}\norm{x-x^\star}^2,
\]
\[
    \norm{\nabla f(x)}
    \ge
    \mu\norm{x-x^\star},
\]
\[
    \frac12\norm{\nabla f(x)}^2
    \ge
    \mu\bigl(f(x)-f^\star\bigr).
\]
Эти соотношения используются при доказательстве линейной и
экспоненциальной сходимости оптимизационных методов.

PL-неравенство сохраняет последнюю оценку без требования
выпуклости. Поэтому доказательства сходимости, использующие только
PL-условие и общие предположения о гладкости, применимы ко всем
сильно выпуклым задачам и одновременно охватывают некоторые
невыпуклые модели.

Единственность глобального минимума сама по себе не характеризует инвексность, поэтому этот класс функций далее не используется.


Строгая выпуклость гарантирует единственность минимума, если он
существует. Положительная определённость матрицы Гессе достаточна, но
не необходима для строгой выпуклости. Сильная выпуклость задаёт
равномерную квадратичную кривизну:
\[
    \nabla^2f(x)\succeq\mu I_n.
\]
Она влечёт строгую выпуклость, квадратичный рост, сильную
монотонность градиента и PL-неравенство. PL-условие может выполняться
и для невыпуклых функций. Для негладких функций аналогом градиента
служит субградиент.
\lecture{07}{Следствия сильной выпуклости}

\sourcevideo
    {https://www.youtube.com/playlist?list=PLOzRYVm0a65fhKDS8cYIC7YCuVGJ4FOJx}
    {Week 2: Lecture 7: Implications of Strong Convexity}

Сильная выпуклость допускает несколько эквивалентных
характеризаций. Одни из них удобны для проверки свойства функции,
другие используются при анализе сходимости оптимизационных
алгоритмов.

\section{Эквивалентные характеризации}

Пусть \(C\subseteq\RR^n\) открыто и выпукло, \(f\in C^1(C)\), а
\(\mu>0\).

\begin{theorem}[Характеризации сильной выпуклости]
Следующие условия эквивалентны.

Функция \(f\) является \(\mu\)-сильно выпуклой, то есть
\[
    f(y)
    \ge
    f(x)
    +
    \inner{\nabla f(x)}{y-x}
    +
    \frac{\mu}{2}\norm{y-x}^2
\]
для всех \(x,y\in C\).

Функция
\[
    g(x)
    =
    f(x)-\frac{\mu}{2}\norm{x}^2
\]
выпукла.

Градиент \(f\) является \(\mu\)-сильно монотонным:
\[
    \inner{
        \nabla f(x)-\nabla f(y)
    }{
        x-y
    }
    \ge
    \mu\norm{x-y}^2.
\]

Для любых \(x,y\in C\) и \(\lambda\in[0,1]\)
\[
\begin{aligned}
    f\bigl(\lambda x+(1-\lambda)y\bigr)
    \le{}&
    \lambda f(x)+(1-\lambda)f(y)
    \\
    &-
    \frac{\mu}{2}
    \lambda(1-\lambda)\norm{x-y}^2.
\end{aligned}
\]
\end{theorem}

\begin{proof}
Поскольку
\[
    \nabla g(x)=\nabla f(x)-\mu x,
\]
критерий выпуклости первого порядка для \(g\) имеет вид
\[
    g(y)
    \ge
    g(x)
    +
    \inner{\nabla f(x)-\mu x}{y-x}.
\]
Подставляя определение \(g\) и собирая квадратичные слагаемые,
получаем
\[
    f(y)
    \ge
    f(x)
    +
    \inner{\nabla f(x)}{y-x}
    +
    \frac{\mu}{2}\norm{y-x}^2.
\]
Следовательно, первые два условия эквивалентны.

Если \(g\) выпукла, то её градиент монотонен:
\[
    \inner{\nabla g(x)-\nabla g(y)}{x-y}\ge0.
\]
После подстановки
\[
    \nabla g(x)-\nabla g(y)
    =
    \nabla f(x)-\nabla f(y)-\mu(x-y)
\]
получаем третье условие.

Обратно, третье условие означает монотонность \(\nabla g\). Для
\(d=y-x\) имеем
\[
\begin{aligned}
    g(y)-g(x)-\inner{\nabla g(x)}{d}
    &=
    \int_0^1
    \inner{
        \nabla g(x+td)-\nabla g(x)
    }{d}
    \,\mathrm dt
    \\
    &\ge0,
\end{aligned}
\]
поскольку при \(t>0\) монотонность градиента даёт
\[
    \inner{
        \nabla g(x+td)-\nabla g(x)
    }{td}
    \ge0.
\]
Следовательно, \(g\) выпукла.

Наконец, выпуклость \(g\) равносильна неравенству
\[
    g\bigl(\lambda x+(1-\lambda)y\bigr)
    \le
    \lambda g(x)+(1-\lambda)g(y).
\]
Используя тождество
\[
\begin{aligned}
    \norm{\lambda x+(1-\lambda)y}^2
    ={}&
    \lambda\norm{x}^2+(1-\lambda)\norm{y}^2
    \\
    &-
    \lambda(1-\lambda)\norm{x-y}^2,
\end{aligned}
\]
получаем четвёртую характеризацию.
\end{proof}

Таким образом, сильную выпуклость можно понимать как обычную
выпуклость после вычитания квадратичной функции
\[
    \frac{\mu}{2}\norm{x}^2.
\]

\section{Нижняя оценка изменения градиента}

Из сильной монотонности и неравенства Коши--Буняковского следует
\[
\begin{aligned}
    \mu\norm{x-y}^2
    &\le
    \inner{
        \nabla f(x)-\nabla f(y)
    }{
        x-y
    }
    \\
    &\le
    \norm{\nabla f(x)-\nabla f(y)}
    \norm{x-y}.
\end{aligned}
\]
Поэтому при \(x\ne y\)
\[
    \norm{\nabla f(x)-\nabla f(y)}
    \ge
    \mu\norm{x-y}.
\]

В частности, если \(x^\star\) является минимумом и
\(\nabla f(x^\star)=0\), то
\[
    \norm{\nabla f(x)}
    \ge
    \mu\norm{x-x^\star}.
\]
Следовательно, градиент сильно выпуклой функции количественно
контролирует расстояние до оптимального решения.

\section{Липшицевость и гладкость}

Необходимо различать липшицевость самой функции и липшицевость её
градиента.

\begin{definition}[Липшицева функция]
Функция \(f\colon C\to\RR\) называется \(L\)-липшицевой, если
\[
    \abs{f(x)-f(y)}
    \le
    L\norm{x-y}
\]
для всех \(x,y\in C\).
\end{definition}

\begin{definition}[Гладкая функция]
Дифференцируемая функция \(f\) называется \(L\)-гладкой, если её
градиент \(L\)-липшицев:
\[
    \norm{\nabla f(x)-\nabla f(y)}
    \le
    L\norm{x-y}.
\]
\end{definition}

Если функция одновременно \(\mu\)-сильно выпукла и \(L\)-гладка, то
\[
    \mu\norm{x-y}
    \le
    \norm{\nabla f(x)-\nabla f(y)}
    \le
    L\norm{x-y}.
\]
Следовательно,
\[
    0<\mu\le L.
\]

Для функции
\[
    f(x)=x^2
\]
имеем
\[
    \nabla f(x)=2x,
\]
поэтому
\[
    \abs{\nabla f(x)-\nabla f(y)}
    =
    2\abs{x-y}.
\]
Она \(2\)-сильно выпукла и \(2\)-гладка.

Для дважды дифференцируемой функции сочетание сильной выпуклости и
гладкости выражается через Гессиан:
\[
    \mu I
    \preceq
    \nabla^2f(x)
    \preceq
    LI.
\]

\section{Сумма выпуклых функций}

\begin{theorem}
Пусть \(f\) является \(\mu\)-сильно выпуклой, а \(g\) выпукла.
Тогда функция
\[
    h=f+g
\]
также является \(\mu\)-сильно выпуклой.
\end{theorem}

\begin{proof}
Для \(\lambda\in[0,1]\) сильная выпуклость \(f\) даёт
\[
\begin{aligned}
    f\bigl(\lambda x+(1-\lambda)y\bigr)
    \le{}&
    \lambda f(x)+(1-\lambda)f(y)
    \\
    &-
    \frac{\mu}{2}
    \lambda(1-\lambda)\norm{x-y}^2.
\end{aligned}
\]
По выпуклости \(g\)
\[
    g\bigl(\lambda x+(1-\lambda)y\bigr)
    \le
    \lambda g(x)+(1-\lambda)g(y).
\]
Складывая эти неравенства, получаем
\[
\begin{aligned}
    h\bigl(\lambda x+(1-\lambda)y\bigr)
    \le{}&
    \lambda h(x)+(1-\lambda)h(y)
    \\
    &-
    \frac{\mu}{2}
    \lambda(1-\lambda)\norm{x-y}^2.
\end{aligned}
\]
\end{proof}

Следовательно, к выпуклой целевой функции можно добавить сильно
выпуклый регуляризатор и получить сильно выпуклую задачу. Например,
если \(g\) выпукла, то
\[
    h(x)=g(x)+\frac{\mu}{2}\norm{x}^2
\]
является \(\mu\)-сильно выпуклой.

При этом добавление регуляризатора в общем случае изменяет точку
минимума. Сохранение исходного решения требует отдельного
обоснования.

\section{Сильная выпуклость и PL-неравенство}

\begin{theorem}
Пусть \(f\in C^1(\RR^n)\) является \(\mu\)-сильно выпуклой и
\[
    f^\star=\min_x f(x).
\]
Тогда
\[
    \frac12\norm{\nabla f(x)}^2
    \ge
    \mu\bigl(f(x)-f^\star\bigr)
\]
для всех \(x\in\RR^n\).
\end{theorem}

\begin{proof}
Из сильной выпуклости для любого \(y\) следует
\[
    f(y)
    \ge
    f(x)
    +
    \inner{\nabla f(x)}{y-x}
    +
    \frac{\mu}{2}\norm{y-x}^2.
\]
Правая часть как функция \(y\) минимальна при
\[
    y
    =
    x-\frac{1}{\mu}\nabla f(x).
\]
Её минимальное значение равно
\[
    f(x)
    -
    \frac{1}{2\mu}\norm{\nabla f(x)}^2.
\]
Поскольку неравенство справедливо для всех \(y\),
\[
    f^\star
    \ge
    f(x)
    -
    \frac{1}{2\mu}\norm{\nabla f(x)}^2.
\]
После перестановки слагаемых получаем
\[
    \frac12\norm{\nabla f(x)}^2
    \ge
    \mu\bigl(f(x)-f^\star\bigr).
\]
\end{proof}

Таким образом,
\[
    \text{сильная выпуклость}
    \quad\Longrightarrow\quad
    \text{PL-неравенство}.
\]
Обратное утверждение неверно. PL-неравенство может выполняться для
функций, не являющихся сильно выпуклыми и даже выпуклыми.

В лекции вновь упоминается пример
\[
    f(x)=x^2+3\sin^2x.
\]
Он используется для иллюстрации невыпуклой функции, допускающей
PL-анализ. Подробная проверка соответствующего неравенства будет
рассматриваться позднее.

\section{Условие оптимальности на выпуклом множестве}

Рассмотрим задачу
\[
    \minimize_{x\in\mathcal X}f(x),
\]
где \(\mathcal X\subseteq\RR^n\) выпукло, а \(f\) выпукла и
дифференцируема.

\begin{theorem}[Условие первого порядка]
Точка \(x^\star\in\mathcal X\) является глобальным минимумом тогда и
только тогда, когда
\[
    \inner{\nabla f(x^\star)}{y-x^\star}
    \ge0
\]
для всех \(y\in\mathcal X\).
\end{theorem}

\begin{proof}
Пусть \(x^\star\) является минимумом. Для любого \(y\in\mathcal X\)
и \(t\in[0,1]\) точка
\[
    x^\star+t(y-x^\star)
\]
принадлежит \(\mathcal X\). Функция
\[
    \varphi(t)
    =
    f\bigl(x^\star+t(y-x^\star)\bigr)
\]
имеет минимум при \(t=0\), поэтому
\[
    \varphi'(0)
    =
    \inner{\nabla f(x^\star)}{y-x^\star}
    \ge0.
\]

Обратно, по выпуклости
\[
    f(y)
    \ge
    f(x^\star)
    +
    \inner{\nabla f(x^\star)}{y-x^\star}.
\]
Если скалярное произведение неотрицательно, то
\[
    f(y)\ge f(x^\star)
\]
для всех \(y\in\mathcal X\).
\end{proof}

Градиент \(\nabla f(x^\star)\) указывает направление возрастания
функции. Условие оптимальности требует, чтобы угол между градиентом и каждым
допустимым направлением \(y-x^\star\) не превышал \(\pi/2\).

Если минимум лежит на границе допустимого множества, градиент может
быть ненулевым. Он должен быть направлен так, чтобы движение внутрь
множества не уменьшало функцию.

\section{Задача без ограничений}

Если
\[
    \mathcal X=\RR^n,
\]
то условие оптимальности принимает вид
\[
    \inner{\nabla f(x^\star)}{y-x^\star}\ge0
\]
для всех \(y\in\RR^n\).

Можно выбрать
\[
    y=x^\star-\nabla f(x^\star).
\]
Тогда
\[
    -\norm{\nabla f(x^\star)}^2\ge0,
\]
откуда
\[
    \nabla f(x^\star)=0.
\]

Следовательно, для дифференцируемой выпуклой функции
\[
    x^\star\in\argmin_x f(x)
    \quad\Longleftrightarrow\quad
    \nabla f(x^\star)=0.
\]

\section{Квадратичная задача с равенствами}

Рассмотрим задачу
\[
    \begin{aligned}
        \minimize_{x\in\RR^n}
        \quad&
        \frac12x^{\mathsf T}Qx+c^{\mathsf T}x,
        \\
        \subjectto
        \quad&
        Ax=b,
    \end{aligned}
\]
где
\[
    Q=Q^{\mathsf T}\succeq0.
\]
Целевая функция выпукла, а допустимое множество
\[
    \mathcal X=\setgiven{x\in\RR^n}{Ax=b}
\]
аффинно и потому выпукло.

Градиент целевой функции равен
\[
    \nabla f(x)=Qx+c.
\]

По условию первого порядка допустимая точка \(x^\star\) оптимальна
тогда и только тогда, когда
\[
    \inner{Qx^\star+c}{y-x^\star}\ge0
\]
для всех \(y\), удовлетворяющих \(Ay=b\).

Поскольку
\[
    Ay=Ax^\star=b,
\]
имеем
\[
    A(y-x^\star)=0.
\]
Следовательно,
\[
    y-x^\star\in\ker A.
\]

Если \(z\in\ker A\), то допустимы оба направления \(z\) и \(-z\).
Поэтому условие оптимальности эквивалентно равенству
\[
    \inner{Qx^\star+c}{z}=0
    \qquad
    \text{для всех }z\in\ker A.
\]
Иными словами,
\[
    Qx^\star+c
    \in
    (\ker A)^\perp
    =
    \operatorname{range}(A^{\mathsf T}).
\]

В задаче без ограничений это условие сводится к
\[
    Qx^\star+c=0.
\]
Если \(Q\succ0\), то решение единственно и равно
\[
    x^\star=-Q^{-1}c.
\]

При наличии равенства \(Ax=b\) необходимо одновременно выполнить
условие допустимости и условие ортогональности:
\[
    Ax^\star=b,
    \qquad
    Qx^\star+c\perp\ker A.
\]
В следующей лекции эта система будет получена посредством функции
Лагранжа и двойственной задачи.

\begin{lecturesummary}
Сильная выпуклость эквивалентна выпуклости функции
\[
    f(x)-\frac{\mu}{2}\norm{x}^2,
\]
сильной монотонности градиента и усиленному неравенству для выпуклых
комбинаций. Если функция одновременно \(\mu\)-сильно выпукла и
\(L\)-гладка, то
\[
    \mu
    \le
    \frac{
        \norm{\nabla f(x)-\nabla f(y)}
    }{
        \norm{x-y}
    }
    \le
    L.
\]
Сумма \(\mu\)-сильно выпуклой и выпуклой функций остаётся
\(\mu\)-сильно выпуклой. Сильная выпуклость также влечёт
PL-неравенство. Для минимизации выпуклой функции на выпуклом
множестве точка \(x^\star\) оптимальна тогда и только тогда, когда
\[
    \inner{\nabla f(x^\star)}{y-x^\star}\ge0
\]
для всех допустимых \(y\). В задаче без ограничений это условие
переходит в равенство \(\nabla f(x^\star)=0\).
\end{lecturesummary}
\input{lectures/lecture_08.tex}
\lecture{09}{Условие Слейтера и сильная двойственность}

\sourcevideo
    {https://www.youtube.com/playlist?list=PLOzRYVm0a65fhKDS8cYIC7YCuVGJ4FOJx}
    {Week 3: Lecture 9: Slater's Condition}

Для прямой задачи с оптимальным значением \(p^\star\) и её
лагранжевой двойственной задачи с оптимальным значением \(d^\star\)
слабая двойственность даёт
\[
    d^\star\le p^\star.
\]
Разность между этими значениями может быть положительной. Условие
Слейтера является основным достаточным условием, при котором
двойственный разрыв исчезает.

\section{Двойственный разрыв}

Рассмотрим прямую задачу
\[
    \begin{aligned}
        \minimize_x
        \quad&
        f_0(x),
        \\
        \subjectto
        \quad&
        f_i(x)\le0,
        \qquad i=1,\ldots,m,
        \\
        &
        h_j(x)=0,
        \qquad j=1,\ldots,r.
    \end{aligned}
\]
Её функция Лагранжа и двойственная функция равны
\[
    \mathcal L(x,\lambda,\nu)
    =
    f_0(x)
    +
    \sum_{i=1}^{m}\lambda_i f_i(x)
    +
    \sum_{j=1}^{r}\nu_jh_j(x),
\]
\[
    g(\lambda,\nu)
    =
    \inf_x\mathcal L(x,\lambda,\nu).
\]

Двойственная задача имеет вид
\[
    d^\star
    =
    \sup_{\lambda\ge0,\nu}g(\lambda,\nu).
\]

\begin{definition}[Двойственный разрыв]
Если оптимальные значения \(p^\star\) и \(d^\star\) конечны, то
величина
\[
    \Delta_{\mathrm{dual}}
    =
    p^\star-d^\star
    \ge0
\]
называется двойственным разрывом.
\end{definition}

Если
\[
    p^\star=d^\star,
\]
выполняется сильная двойственность. В противном случае двойственная
задача даёт лишь строгую нижнюю оценку прямого оптимального значения.

\section{Пример положительного двойственного разрыва}

Рассмотрим задачу
\[
    \begin{aligned}
        \minimize_{x=(x_1,x_2)\in\RR^2}
        \quad&
        \exp(x_2),
        \\
        \subjectto
        \quad&
        \norm{x}_2\le x_1.
    \end{aligned}
    \tag{\ensuremath{P}}
\]

Ограничение эквивалентно
\[
    \sqrt{x_1^2+x_2^2}\le x_1.
\]
Обе части неотрицательны, поэтому после возведения в квадрат
получаем
\[
    x_1^2+x_2^2\le x_1^2,
\]
откуда
\[
    x_2=0,
    \qquad
    x_1\ge0.
\]
Следовательно, на всём допустимом множестве
\[
    \exp(x_2)=1,
\]
и потому
\[
    p^\star=1.
\]

Запишем ограничение как
\[
    h(x)
    =
    \norm{x}_2-x_1
    \le0.
\]
Функция Лагранжа имеет вид
\[
    \mathcal L(x,\lambda)
    =
    \exp(x_2)
    +
    \lambda\bigl(\norm{x}_2-x_1\bigr),
    \qquad
    \lambda\ge0,
\]
а двойственная функция равна
\[
    g(\lambda)
    =
    \inf_{x\in\RR^2}
    \left[
        \exp(x_2)
        +
        \lambda\bigl(\norm{x}_2-x_1\bigr)
    \right].
\]

Поскольку
\[
    \exp(x_2)>0,
    \qquad
    \norm{x}_2-x_1\ge0
\]
для любого \(x\), имеем
\[
    g(\lambda)\ge0.
\]
Следовательно,
\[
    d^\star\ge0.
\]

Теперь ограничим поиск последовательностью
\[
    x_1=x_2^4,
    \qquad
    x_2\to-\infty.
\]
Тогда
\[
\begin{aligned}
    \norm{x}_2-x_1
    &=
    \sqrt{x_2^8+x_2^2}-x_2^4
    \\
    &=
    \frac{x_2^2}{
        \sqrt{x_2^8+x_2^2}+x_2^4
    }
    \longrightarrow0.
\end{aligned}
\]
Одновременно
\[
    \exp(x_2)\longrightarrow0.
\]
Поэтому для каждого \(\lambda\ge0\)
\[
    g(\lambda)\le0.
\]
Вместе с ранее полученной оценкой это даёт
\[
    g(\lambda)=0,
    \qquad
    d^\star=0.
\]

Таким образом,
\[
    p^\star=1,
    \qquad
    d^\star=0,
    \qquad
    p^\star-d^\star=1.
\]
В этой задаче сильная двойственность не выполняется.

\section{Строгая допустимость}

Пусть прямая задача выпукла: функция \(f_0\) и ограничения
\(f_i\) выпуклы, а равенства имеют аффинный вид
\[
    h_j(x)=a_j^{\mathsf T}x-b_j.
\]

\begin{definition}[Строго допустимая точка]
Точка \(\bar x\) называется строго допустимой, если
\[
    f_i(\bar x)<0,
    \qquad i=1,\ldots,m,
\]
и одновременно
\[
    h_j(\bar x)=0,
    \qquad j=1,\ldots,r.
\]
\end{definition}

Строгость требуется только для неравенств:
\[
    f_i(\bar x)<0.
\]
Аффинные равенства должны выполняться точно.

\section{Условие Слейтера}

\begin{theorem}[Условие Слейтера]
Рассмотрим выпуклую задачу
\[
    \begin{aligned}
        \minimize_x
        \quad&
        f_0(x),
        \\
        \subjectto
        \quad&
        f_i(x)\le0,
        \qquad i=1,\ldots,m,
        \\
        &
        Ax=b.
    \end{aligned}
\]
Пусть \(p^\star\) конечно и существует точка
\[
    \bar x\in
    \operatorname{ri}
    \bigl(
        \operatorname{dom}f_0
        \cap
        \textstyle\bigcap_i\operatorname{dom}f_i
    \bigr),
\]
для которой
\[
    f_i(\bar x)<0,
    \qquad
    A\bar x=b.
\]
Тогда выполняется сильная двойственность:
\[
    p^\star=d^\star.
\]
При этих предположениях существует и оптимальная двойственная точка
\((\lambda^\star,\nu^\star)\).
\end{theorem}

Здесь \(\operatorname{ri}\) обозначает относительную внутренность.
Она используется потому, что допустимое множество может лежать в
аффинном подпространстве меньшей размерности.

В рассмотренном выше примере допустимое множество имеет вид
\[
    \setgiven{(x_1,x_2)}{x_1\ge0,\ x_2=0}.
\]
Для каждой его точки
\[
    \norm{x}_2-x_1=0,
\]
поэтому строгого неравенства получить нельзя. Условие Слейтера не
выполняется, что согласуется с положительным двойственным разрывом.

\begin{remark}[Логика условия]
Условие Слейтера является достаточным, но не необходимым. Из его
невыполнения нельзя заключить, что сильная двойственность обязательно
нарушается. Однако если сильная двойственность нарушена, условия
теоремы Слейтера выполнены быть не могут.
\end{remark}

\begin{remark}[Аффинные ограничения]
В уточнённой форме условия Слейтера строгие неравенства требуются
только для неаффинных выпуклых ограничений. Аффинные неравенства
достаточно выполнять нестрого. Поэтому, если все ограничения аффинны,
задача допустима и её оптимальное значение конечно, сильная
двойственность выполняется в стандартной конечномерной постановке.
\end{remark}

\section{Следствия сильной двойственности}

Если
\[
    p^\star=d^\star,
\]
прямую задачу можно решать через двойственную, не теряя оптимального
значения. Это особенно полезно, когда размерность двойственной
переменной меньше размерности прямой переменной.

Кроме того, при выполнении подходящего условия регулярности условия
Каруша--Куна--Таккера становятся не только достаточными, но и
необходимыми для оптимальности выпуклой задачи. Эти условия будут
введены отдельно.

\section{Чувствительность оптимального значения}

Рассмотрим параметрическую задачу
\[
    p(b)
    =
    \minimize_x
    \left\{
        f(x)
        \,\middle|\,
        Bx=b
    \right\}.
\]
Её функция Лагранжа равна
\[
    \mathcal L(x,\nu;b)
    =
    f(x)+\nu^{\mathsf T}(Bx-b).
\]

Пусть решение \(x^\star(b)\) и оптимальный множитель
\(\nu^\star(b)\) определены однозначно и достаточно гладко зависят
от \(b\). Тогда
\[
    p(b)
    =
    \mathcal L
    \bigl(
        x^\star(b),
        \nu^\star(b);
        b
    \bigr).
\]

Дифференцируя и используя условия стационарности
\[
    \nabla_x\mathcal L=0,
    \qquad
    \nabla_\nu\mathcal L=0,
\]
получаем
\[
    \nabla p(b)
    =
    \nabla_b\mathcal L
    =
    -\nu^\star(b).
\]

\begin{remark}[Знак множителя]
При соглашении
\[
    \mathcal L(x,\nu;b)
    =
    f(x)+\nu^{\mathsf T}(Bx-b)
\]
справедливо
\[
    \nabla p(b)=-\nu^\star.
\]
При записи ограничения в форме \(b-Bx=0\) знак меняется. Поэтому
экономическая интерпретация множителя зависит от выбранного
соглашения о знаках.
\end{remark}

Двойственная переменная характеризует чувствительность оптимального
значения к изменению правой части ограничения. Поэтому множители
Лагранжа также называют теневыми ценами или предельными стоимостями.

\section{Двойственная форма квадратичной задачи}

Рассмотрим задачу
\[
    \begin{aligned}
        \minimize_{x\in\RR^n}
        \quad&
        \frac12x^{\mathsf T}Qx+c^{\mathsf T}x,
        \\
        \subjectto
        \quad&
        Ax\le b,
    \end{aligned}
    \tag{\ensuremath{P_Q}}
\]
где
\[
    Q=Q^{\mathsf T}\succ0,
    \qquad
    A\in\RR^{m\times n}.
\]

Функция Лагранжа равна
\[
\begin{aligned}
    \mathcal L(x,\lambda)
    &=
    \frac12x^{\mathsf T}Qx
    +
    c^{\mathsf T}x
    +
    \lambda^{\mathsf T}(Ax-b)
    \\
    &=
    \frac12x^{\mathsf T}Qx
    +
    (c+A^{\mathsf T}\lambda)^{\mathsf T}x
    -
    b^{\mathsf T}\lambda,
\end{aligned}
\]
где
\[
    \lambda\ge0.
\]

Для вычисления двойственной функции минимизируем
\(\mathcal L\) по \(x\). Условие стационарности имеет вид
\[
    Qx+c+A^{\mathsf T}\lambda=0.
\]
Поскольку \(Q\succ0\),
\[
    x(\lambda)
    =
    -Q^{-1}
    \bigl(
        c+A^{\mathsf T}\lambda
    \bigr).
\]

Подставляя это выражение в функцию Лагранжа, получаем
\[
    g(\lambda)
    =
    -
    \frac12
    \bigl(
        c+A^{\mathsf T}\lambda
    \bigr)^{\mathsf T}
    Q^{-1}
    \bigl(
        c+A^{\mathsf T}\lambda
    \bigr)
    -
    b^{\mathsf T}\lambda.
\]

Следовательно, двойственная задача равна
\[
    \begin{aligned}
        \maximize_{\lambda\in\RR^m}
        \quad&
        -
        \frac12
        \bigl(
            c+A^{\mathsf T}\lambda
        \bigr)^{\mathsf T}
        Q^{-1}
        \bigl(
            c+A^{\mathsf T}\lambda
        \bigr)
        -
        b^{\mathsf T}\lambda,
        \\
        \subjectto
        \quad&
        \lambda\ge0.
    \end{aligned}
    \tag{\ensuremath{D_Q}}
\]

После раскрытия скобок
\[
\begin{aligned}
    g(\lambda)
    ={}&
    -\frac12
    \lambda^{\mathsf T}
    AQ^{-1}A^{\mathsf T}
    \lambda
    \\
    &-
    \bigl(
        AQ^{-1}c+b
    \bigr)^{\mathsf T}\lambda
    -
    \frac12c^{\mathsf T}Q^{-1}c.
\end{aligned}
\]

Матрица
\[
    AQ^{-1}A^{\mathsf T}
\]
положительно полуопределена, поэтому \(g\) вогнута. Последнее
постоянное слагаемое не влияет на положение двойственного
максимума.

Если двойственное решение \(\lambda^\star\) существует и выполняется
сильная двойственность, прямое решение восстанавливается по формуле
\[
    x^\star
    =
    -Q^{-1}
    \bigl(
        c+A^{\mathsf T}\lambda^\star
    \bigr).
\]

Если
\[
    m\ll n,
\]
двойственная задача имеет существенно меньшую размерность. Вместо
обмена векторами \(x\in\RR^n\) распределённые агенты могут
согласовывать множители
\[
    \lambda\in\RR^m.
\]

\begin{remark}
Для применения сильной двойственности к задаче \((P_Q)\) достаточно,
например, существования строго допустимой точки
\[
    \bar x
    \quad\text{такой, что}\quad
    A\bar x<b.
\]
Положительная определённость \(Q\) обеспечивает строгую выпуклость
целевой функции и единственность прямого решения, но сама по себе не
заменяет условие регулярности ограничений.
\end{remark}

\section{Связь с машинным обучением}

В методе опорных векторов прямая задача содержит вектор параметров
гиперплоскости и по одному ограничению для каждого объекта обучающей
выборки. Переход к двойственной задаче выражает решение через
множители этих ограничений.

В двойственной форме объекты входят только через скалярные
произведения. Это позволяет заменить скалярное произведение
ядровой функцией и получить нелинейный классификатор без явного
построения пространства признаков. Двойственная форма SVM и
ядровый переход рассматриваются далее.

\begin{lecturesummary}
Двойственный разрыв равен
\[
    p^\star-d^\star\ge0
\]
и может быть строго положительным даже для выпуклой задачи. Условие
Слейтера требует существования строго допустимой точки для выпуклых
неравенств при точном выполнении аффинных равенств. При этом условии
выполняется сильная двойственность:
\[
    p^\star=d^\star.
\]
Оптимальный множитель Лагранжа характеризует чувствительность
оптимального значения к изменению ограничения. Для квадратичной
задачи с \(Q\succ0\) прямая переменная исключается явно:
\[
    x(\lambda)
    =
    -Q^{-1}(c+A^{\mathsf T}\lambda),
\]
что приводит к вогнутой квадратичной двойственной задаче по
неотрицательному вектору \(\lambda\).
\end{lecturesummary}
\lecture{10}{Градиентный спуск: двойственность}

\sourcevideo
    {https://www.youtube.com/playlist?list=PLOzRYVm0a65fhKDS8cYIC7YCuVGJ4FOJx}
    {Week 3: Lecture 10: Analysis of Gradient Descent Algorithm}

В начале лекции обсуждается влияние шага градиентного спуска на
сходимость. Основная часть посвящена построению двойственных задач
для линейного программирования и метода опорных векторов. Двойственная
форма SVM приводит к ядровому методу, позволяющему строить нелинейные
границы классификации.

\section{Гладкость, сильная выпуклость и градиентный спуск}

Пусть \(f\colon\RR^n\to\RR\) дифференцируема.

\begin{definition}[\(L\)-гладкость]
Функция \(f\) называется \(L\)-гладкой, если её градиент
\(L\)-липшицев:
\[
    \norm{\nabla f(x)-\nabla f(y)}
    \le
    L\norm{x-y}
\]
для любых \(x,y\in\RR^n\).
\end{definition}

Если \(f\in C^2(\RR^n)\), то \(L\)-гладкость эквивалентна оценке
\[
    \lVert\nabla^2 f(x)\rVert_{\mathrm{op}}
    \le L
    \qquad
    \text{для всех }x\in\RR^n.
\]
В частности, для выпуклой функции это условие имеет вид
\[
    0\preceq\nabla^2 f(x)\preceq LI.
\]

Для \(\mu\)-сильно выпуклой дважды дифференцируемой функции
\[
    \nabla^2f(x)\succeq\mu I.
\]
Следовательно, если функция одновременно \(L\)-гладка и
\(\mu\)-сильно выпукла, то
\[
    \mu I
    \preceq
    \nabla^2f(x)
    \preceq
    LI,
\]
откуда
\[
    0<\mu\le L.
\]

Отношение \(\mu/L\in(0,1]\) характеризует относительную нижнюю
границу кривизны. Величину
\[
    \kappa=\frac{L}{\mu}\ge1
\]
называют числом обусловленности функции на рассматриваемом классе:
чем больше \(\kappa\), тем сильнее различаются масштабы кривизны и тем
медленнее обычно сходятся методы первого порядка.



Для задачи
\[
    \minimize_{x\in\RR^n}f(x)
\]
градиентный спуск задаётся правилом
\[
    x^{k+1}
    =
    x^k-\eta\nabla f(x^k),
\]
где \(\eta>0\) --- шаг метода.

\begin{algorithm}[H]
\caption{Градиентный спуск с постоянным шагом}
\label{alg:lecture10-gradient-descent}
\begin{algorithmic}[1]

\Require
функция \(f\), начальная точка \(x^0\), шаг \(\eta>0\)

\For{\(k=0,1,2,\ldots\)}

    \State
    \(g^k\gets\nabla f(x^k)\)

    \State
    \(x^{k+1}\gets x^k-\eta g^k\)

\EndFor

\Ensure
последовательность \(\{x^k\}_{k\ge0}\)

\end{algorithmic}
\end{algorithm}

Даже для простой квадратичной функции выбор шага определяет поведение
итераций.



Рассмотрим
\[
    f(x)=\frac12x^2.
\]
Тогда
\[
    \nabla f(x)=x,
    \qquad
    L=1,
\]
и градиентный спуск принимает вид
\[
    x^{k+1}
    =
    (1-\eta)x^k.
\]
Следовательно,
\[
    x^k
    =
    (1-\eta)^k x^0.
\]

Сходимость к \(x^\star=0\) происходит тогда и только тогда, когда
\[
    \abs{1-\eta}<1,
\]
то есть
\[
    0<\eta<2.
\]

Поведение метода различается в зависимости от шага:
\[
\begin{array}{c|c}
    \text{шаг} & \text{поведение}
    \\
    \hline
    0<\eta<1
    &
    \text{монотонное уменьшение } \abs{x^k}
    \\
    \eta=1
    &
    x^1=0
    \\
    1<\eta<2
    &
    \text{сходимость с чередованием знака}
    \\
    \eta=2
    &
    x^{k+1}=-x^k
    \\
    \eta>2
    &
    \abs{x^k}\to\infty
\end{array}
\]

Таким образом, слишком большой шаг приводит к колебаниям или
расходимости. Для рассматриваемой функции выбор
\[
    \eta=\frac1L=1
\]
даёт решение за одну итерацию.

\begin{remark}
Для произвольной выпуклой \(L\)-гладкой функции выбор
\(\eta=1/L\) является стандартным безопасным шагом. Полный вывод
оценки сходимости в данной лекции не проводится и будет дан при
последующем анализе градиентного метода.
\end{remark}

\section{Двойственная задача линейного программирования}

Рассмотрим прямую задачу
\[
    \begin{aligned}
        \minimize_{x\in\RR^n}
        \quad&
        c^{\mathsf T}x,
        \\
        \subjectto
        \quad&
        Ax\le b,
    \end{aligned}
    \tag{\ensuremath{P_{\mathrm{LP}}}}
\]
где
\[
    A\in\RR^{m\times n},
    \qquad
    b\in\RR^m,
    \qquad
    c\in\RR^n.
\]

Неравенство \(Ax\le b\) понимается покомпонентно. Введём двойственную
переменную
\[
    \lambda\in\RR^m_+.
\]

Функция Лагранжа равна
\[
\begin{aligned}
    \mathcal L(x,\lambda)
    &=
    c^{\mathsf T}x
    +
    \lambda^{\mathsf T}(Ax-b)
    \\
    &=
    \bigl(c+A^{\mathsf T}\lambda\bigr)^{\mathsf T}x
    -
    b^{\mathsf T}\lambda.
\end{aligned}
\]

Двойственная функция определяется как
\[
    g(\lambda)
    =
    \inf_{x\in\RR^n}
    \mathcal L(x,\lambda).
\]

Если
\[
    c+A^{\mathsf T}\lambda\ne0,
\]
то линейную по \(x\) функцию можно сделать сколь угодно малой,
выбирая направление \(x\). Поэтому
\[
    g(\lambda)=-\infty.
\]

Если же
\[
    A^{\mathsf T}\lambda+c=0,
\]
зависимость от \(x\) исчезает и
\[
    g(\lambda)=-b^{\mathsf T}\lambda.
\]

Таким образом,
\[
    g(\lambda)
    =
    \begin{cases}
        -b^{\mathsf T}\lambda,
        &
        A^{\mathsf T}\lambda+c=0,
        \\[1mm]
        -\infty,
        &
        A^{\mathsf T}\lambda+c\ne0.
    \end{cases}
\]

\begin{problem}[Двойственная задача LP]
Двойственная задача к \((P_{\mathrm{LP}})\) имеет вид
\[
    \begin{aligned}
        \maximize_{\lambda\in\RR^m}
        \quad&
        -b^{\mathsf T}\lambda,
        \\
        \subjectto
        \quad&
        A^{\mathsf T}\lambda+c=0,
        \\
        &
        \lambda\ge0.
    \end{aligned}
    \tag{\ensuremath{D_{\mathrm{LP}}}}
\]
\end{problem}

Двойственная задача также является линейной программой. В ней
\(m\) переменных и \(n\) равенств, тогда как в прямой задаче ---
\(n\) переменных и \(m\) неравенств. Поэтому условие \(m\ll n\)
уменьшает размер двойственного вектора, но само по себе ещё не
гарантирует меньшей вычислительной сложности.

\begin{remark}[Уточнение формулы из транскрипта]
Для конечности двойственной функции требуется именно равенство
\[
    A^{\mathsf T}\lambda+c=0.
\]
Здесь требуется именно равенство; знак неравенства не обеспечивает конечности двойственной функции.
\end{remark}



В распределённой задаче работа с прямой переменной может требовать
хранения и передачи вектора
\[
    x\in\RR^n.
\]
При переходе к двойственной постановке агенты работают с
\[
    \lambda\in\RR^m.
\]
Если число ограничений значительно меньше числа прямых переменных,
это уменьшает объём памяти и коммуникаций.

Однако решение двойственной задачи непосредственно даёт
\(\lambda^\star\), а не \(x^\star\). Восстановление прямого решения
требует дополнительных условий оптимальности. Для общей линейной
программы одного равенства
\[
    A^{\mathsf T}\lambda^\star+c=0
\]
недостаточно для однозначного определения \(x^\star\); необходимо
использовать прямую допустимость и условия дополняющей нежёсткости,
которые будут введены вместе с условиями Каруша--Куна--Таккера.

\begin{remark}[Сильная двойственность для LP]
Для линейной программы условие Слейтера со строгими неравенствами не
требуется. Если прямая задача допустима и её оптимальное значение
конечно, то двойственная задача также имеет оптимальное решение и
\[
    p^\star=d^\star.
\]
Аналогичное утверждение верно при перестановке прямой и двойственной
задач. При недопустимости или неограниченности одной из задач
необходимо отдельно разбирать соответствующий альтернативный случай.
\end{remark}

\section{Прямая задача SVM и стационарность}

Пусть дана строго линейно разделимая выборка, в которой представлены
оба класса:
\[
    \{(x_i,y_i)\}_{i=1}^{m},
    \qquad
    x_i\in\RR^n,
    \qquad
    y_i\in\{-1,+1\}.
\]

Жёстко-разделяющий метод опорных векторов решает задачу
\[
    \begin{aligned}
        \minimize_{w\in\RR^n,\ b\in\RR}
        \quad&
        \frac12\norm{w}^2,
        \\
        \subjectto
        \quad&
        y_i\bigl(w^{\mathsf T}x_i+b\bigr)\ge1,
        \qquad
        i=1,\ldots,m.
    \end{aligned}
    \tag{\ensuremath{P_{\mathrm{SVM}}}}
\]

Запишем ограничения в стандартной форме:
\[
    1-y_i\bigl(w^{\mathsf T}x_i+b\bigr)\le0.
\]
Введём множители
\[
    \lambda_i\ge0,
    \qquad
    i=1,\ldots,m.
\]

Функция Лагранжа имеет вид
\[
\begin{aligned}
    \mathcal L(w,b,\lambda)
    =
    \frac12\norm{w}^2
    +
    \sum_{i=1}^{m}
    \lambda_i
    \left[
        1-y_i\bigl(w^{\mathsf T}x_i+b\bigr)
    \right].
\end{aligned}
\]

Для построения двойственной функции необходимо минимизировать
\(\mathcal L\) по \(w\) и \(b\).



Минимизация по \(w\) даёт единственную точку
\[
    w(\lambda)
    =
    \sum_{i=1}^{m}\lambda_i y_i x_i,
\]
поскольку
\[
    \nabla_w\mathcal L
    =
    w-
    \sum_{i=1}^{m}\lambda_i y_i x_i.
\]
По переменной \(b\) функция Лагранжа линейна. Поэтому
\[
    \inf_{b\in\RR}\mathcal L(w,b,\lambda)>-\infty
\]
тогда и только тогда, когда
\[
    \sum_{i=1}^{m}\lambda_i y_i=0.
\]
При выполнении этого равенства зависимость от \(b\) исчезает.

Первая формула показывает, что нормаль к оптимальной разделяющей
гиперплоскости является линейной комбинацией объектов обучающей
выборки.

\section{Двойственная задача SVM и восстановление классификатора}

Подставим
\[
    w=\sum_{j=1}^{m}\lambda_jy_jx_j
\]
в функцию Лагранжа. Имеем
\[
\begin{aligned}
    \frac12\norm{w}^2
    &=
    \frac12
    \sum_{i=1}^{m}
    \sum_{j=1}^{m}
    \lambda_i\lambda_j
    y_iy_j
    x_i^{\mathsf T}x_j,
    \\
    \sum_{i=1}^{m}
    \lambda_i y_i w^{\mathsf T}x_i
    &=
    \sum_{i=1}^{m}
    \sum_{j=1}^{m}
    \lambda_i\lambda_j
    y_iy_j
    x_i^{\mathsf T}x_j.
\end{aligned}
\]
Слагаемое, содержащее \(b\), исчезает вследствие условия
\[
    \sum_{i=1}^{m}\lambda_i y_i=0.
\]

Поэтому двойственная функция равна
\[
    g(\lambda)
    =
    \sum_{i=1}^{m}\lambda_i
    -
    \frac12
    \sum_{i=1}^{m}
    \sum_{j=1}^{m}
    \lambda_i\lambda_j
    y_iy_j
    x_i^{\mathsf T}x_j.
\]

\begin{problem}[Двойственная задача SVM]
Двойственная задача жёстко-разделяющего SVM имеет вид
\[
    \begin{aligned}
        \maximize_{\lambda\in\RR^m}
        \quad&
        \sum_{i=1}^{m}\lambda_i
        -
        \frac12
        \sum_{i=1}^{m}
        \sum_{j=1}^{m}
        \lambda_i\lambda_j
        y_iy_j
        x_i^{\mathsf T}x_j,
        \\
        \subjectto
        \quad&
        \lambda_i\ge0,
        \qquad i=1,\ldots,m,
        \\
        &
        \sum_{i=1}^{m}\lambda_i y_i=0.
    \end{aligned}
    \tag{\ensuremath{D_{\mathrm{SVM}}}}
\]
\end{problem}

В матричной форме положим
\[
    K_{ij}=x_i^{\mathsf T}x_j,
    \qquad
    Y=\diag(y_1,\ldots,y_m).
\]
Тогда
\[
    g(\lambda)
    =
    \one^{\mathsf T}\lambda
    -
    \frac12
    \lambda^{\mathsf T}YKY\lambda.
\]

Матрица \(K\) является матрицей Грама и положительно
полуопределена. Следовательно, \(YKY\succeq0\), а функция
\(g\) вогнута.



Поскольку выборка строго линейно разделима, ограничения можно
выполнить строго после подходящего масштабирования \((w,b)\).
Следовательно, условие Слейтера выполнено, сильная двойственность
имеет место, а KKT-условия необходимы и достаточны.

После решения двойственной задачи нормаль гиперплоскости находится
по формуле
\[
    w^\star
    =
    \sum_{i=1}^{m}
    \lambda_i^\star y_i x_i.
\]

Значение \(b^\star\) определяется с помощью опорного вектора.
Для индекса \(i\), для которого
\[
    \lambda_i^\star>0,
\]
условие дополняющей нежёсткости даёт
\[
    y_i
    \bigl(
        w^{\star\mathsf T}x_i+b^\star
    \bigr)
    =
    1.
\]
Следовательно,
\[
    b^\star
    =
    y_i-w^{\star\mathsf T}x_i.
\]

В самой лекции это восстановление откладывается до введения условий
Каруша--Куна--Таккера.

Решающее правило имеет вид
\[
    \widehat y(x)
    =
    \sign
    \left(
        \sum_{i=1}^{m}
        \lambda_i^\star y_i
        x_i^{\mathsf T}x
        +
        b^\star
    \right).
\]

Точки с
\[
    \lambda_i^\star=0
\]
не входят в итоговую сумму. Ненулевые множители соответствуют
опорным векторам.

\section{Ядровый переход и отображения признаков}

Двойственная задача зависит от объектов \(x_i\) только через
скалярные произведения
\[
    x_i^{\mathsf T}x_j.
\]
Пусть
\[
    \phi\colon\RR^n\to\mathcal H
\]
отображает исходные признаки в другое пространство, где данные
могут стать линейно разделимыми. Тогда в двойственной задаче
возникают только величины
\[
    \inner{\phi(x_i)}{\phi(x_j)}_{\mathcal H}.
\]

\begin{definition}[Ядровая функция]
Функция
\[
    K(x,z)
    =
    \inner{\phi(x)}{\phi(z)}_{\mathcal H}
\]
называется ядром, соответствующим отображению признаков \(\phi\).
\end{definition}

В двойственной задаче достаточно выполнить замену
\[
    x_i^{\mathsf T}x_j
    \longmapsto
    K(x_i,x_j).
\]
Получаем
\[
    \begin{aligned}
        \maximize_{\lambda\in\RR^m}
        \quad&
        \sum_{i=1}^{m}\lambda_i
        -
        \frac12
        \sum_{i,j=1}^{m}
        \lambda_i\lambda_j
        y_iy_j
        K(x_i,x_j),
        \\
        \subjectto
        \quad&
        \lambda_i\ge0,
        \\
        &
        \sum_{i=1}^{m}\lambda_i y_i=0.
    \end{aligned}
\]

Классификатор принимает форму
\[
    \widehat y(x)
    =
    \sign
    \left(
        \sum_{i=1}^{m}
        \lambda_i^\star y_i
        K(x_i,x)
        +
        b^\star
    \right).
\]

Пространство признаков \(\mathcal H\) при этом не требуется строить
явно. Этот приём называется ядровым трюком.



Пусть в \(\RR^2\) классы расположены около окружностей
\[
    x_1^2+x_2^2=1
    \qquad\text{и}\qquad
    x_1^2+x_2^2=4.
\]
В исходных координатах такие классы нельзя разделить одной прямой.

Рассмотрим новый признак
\[
    z=x_1^2+x_2^2.
\]
Тогда внутренний класс располагается около значения \(z=1\), а
внешний --- около \(z=4\). В новом пространстве их можно разделить
линейным порогом
\[
    z=\beta,
    \qquad
    1<\beta<4.
\]

В исходных координатах этому порогу соответствует нелинейная граница
\[
    x_1^2+x_2^2=\beta.
\]
Следовательно, линейный классификатор в пространстве признаков задаёт
нелинейную границу в исходном пространстве.

Распространённые примеры ядер:
\[
    K_{\mathrm{poly}}(x,z)
    =
    \bigl(x^{\mathsf T}z+c\bigr)^p,
\]
\[
    K_{\mathrm{RBF}}(x,z)
    =
    \exp
    \left(
        -\frac{\norm{x-z}^2}{2\sigma^2}
    \right).
\]

\begin{remark}
Не всякая функция двух аргументов является допустимым ядром. Для
симметричной функции \(K\) необходимо и достаточно, чтобы для любой
конечной системы точек \(x_1,\ldots,x_m\) матрица Грама
\[
    \bigl[K(x_i,x_j)\bigr]_{i,j=1}^{m}
\]
была положительно полуопределённой.
\end{remark}

\section{Практическая роль двойственной формы}

Для линейной программы двойственный переход меняет размерность
оптимизируемого вектора с \(n\) на \(m\). Это может уменьшить объём
передаваемого состояния при \(m\ll n\), однако двойственная задача
одновременно содержит \(n\) равенств, поэтому вычислительное
преимущество зависит от структуры матриц и используемого алгоритма.

Для SVM основное преимущество двойственной формы состоит не только
в размерности. В ней данные входят исключительно через скалярные
произведения, что позволяет применить ядровый переход и получить
нелинейный классификатор.

При этом двойственная постановка также имеет вычислительные
ограничения. В SVM число двойственных переменных равно числу объектов
обучающей выборки, а матрица Грама имеет размер
\[
    m\times m.
\]
Поэтому при очень большом числе объектов двойственная задача может
оказаться дороже прямой.

\begin{lecturesummary}
Для градиентного спуска
\[
    x^{k+1}=x^k-\eta\nabla f(x^k)
\]
размер шага принципиально влияет на сходимость. В квадратичном
примере \(f(x)=x^2/2\) метод сходится при \(0<\eta<2\), а шаг
\(\eta=1/L\) приводит к минимуму за одну итерацию.

Для линейной программы
\[
    \minimize_x c^{\mathsf T}x
    \quad
    \subjectto
    \quad
    Ax\le b
\]
двойственная задача равна
\[
    \maximize_{\lambda\ge0}
    -b^{\mathsf T}\lambda
    \quad
    \subjectto
    \quad
    A^{\mathsf T}\lambda+c=0.
\]

Двойственная задача жёстко-разделяющего SVM имеет вид
\[
    \maximize_{\lambda\ge0}
    \left\{
        \sum_i\lambda_i
        -
        \frac12
        \sum_{i,j}
        \lambda_i\lambda_jy_iy_j
        x_i^{\mathsf T}x_j
    \right\},
    \qquad
    \sum_i\lambda_i y_i=0.
\]
Зависимость только от скалярных произведений позволяет заменить их
ядрами \(K(x_i,x_j)\) и получить нелинейную разделяющую поверхность
без явного построения пространства признаков.
\end{lecturesummary}
\lecture{11}{KKT-условия и градиентный спуск}

\sourcevideo
    {https://www.youtube.com/playlist?list=PLOzRYVm0a65fhKDS8cYIC7YCuVGJ4FOJx}
    {Week 3: Lecture 11: KKT Conditions}

Условия Каруша--Куна--Таккера объединяют допустимость, стационарность
и согласование прямых и двойственных переменных. Для выпуклых задач
они дают проверяемый сертификат оптимальности. Во второй части лекции
начинается анализ градиентного спуска для выпуклых \(L\)-гладких
функций и выводится стандартный шаг \(1/L\).

\section{KKT-условия для задач с ограничениями}

Рассмотрим задачу
\[
    \begin{aligned}
        \minimize_{x\in\RR^n}
        \quad&
        f(x),
        \\
        \subjectto
        \quad&
        h_i(x)\le0,
        \qquad i=1,\ldots,m,
        \\
        &
        \ell_j(x)=0,
        \qquad j=1,\ldots,r.
    \end{aligned}
    \tag{\ensuremath{P}}
\]

Предполагается, что функции \(f\), \(h_i\) и \(\ell_j\)
дифференцируемы. Функция Лагранжа равна
\[
    \mathcal L(x,\lambda,\nu)
    =
    f(x)
    +
    \sum_{i=1}^{m}\lambda_i h_i(x)
    +
    \sum_{j=1}^{r}\nu_j\ell_j(x),
\]
где
\[
    \lambda\in\RR^m,
    \qquad
    \nu\in\RR^r.
\]

Переменные \(\lambda_i\) соответствуют ограничениям-неравенствам,
а \(\nu_j\) --- ограничениям-равенствам.



\begin{definition}[Условия Каруша--Куна--Таккера]
Тройка
\[
    (x^\star,\lambda^\star,\nu^\star)
\]
удовлетворяет условиям Каруша--Куна--Таккера, если выполнены четыре
группы условий.

Стационарность:
\[
    \nabla_x
    \mathcal L
    (x^\star,\lambda^\star,\nu^\star)
    =
    0.
\]

Прямая допустимость:
\[
    h_i(x^\star)\le0,
    \qquad
    \ell_j(x^\star)=0.
\]

Двойственная допустимость:
\[
    \lambda_i^\star\ge0.
\]

Дополняющая нежёсткость:
\[
    \lambda_i^\star h_i(x^\star)=0,
    \qquad i=1,\ldots,m.
\]
\end{definition}

Условие дополняющей нежёсткости утверждает, что для каждого
ограничения выполняется хотя бы одно из равенств
\[
    \lambda_i^\star=0
    \qquad\text{или}\qquad
    h_i(x^\star)=0.
\]

\begin{definition}[Активное ограничение]
Ограничение \(h_i(x)\le0\) называется активным в точке \(x^\star\),
если
\[
    h_i(x^\star)=0.
\]
Если
\[
    h_i(x^\star)<0,
\]
оно называется неактивным.
\end{definition}

Для неактивного ограничения дополняющая нежёсткость немедленно даёт
\[
    \lambda_i^\star=0.
\]
Следовательно, в условии стационарности участвуют только активные
ограничения с потенциально ненулевыми множителями.



В выпуклой постановке функции \(f\) и \(h_i\) выпуклы и
дифференцируемы, а ограничения-равенства аффинны:
\[
    \ell_j(x)=a_j^{\mathsf T}x-b_j.
\]

\begin{theorem}[KKT как критерий оптимальности]
Пусть задача \((P)\) выпукла и все входящие в неё функции
дифференцируемы. Любая тройка
\[
    (x^\star,\lambda^\star,\nu^\star),
\]
удовлетворяющая условиям KKT, задаёт глобальное решение прямой задачи
и оптимальные двойственные множители.

Обратно, пусть \(x^\star\) является решением прямой задачи и
выполнено условие Слейтера. Тогда существуют множители
\((\lambda^\star,\nu^\star)\), с которыми тройка
\((x^\star,\lambda^\star,\nu^\star)\) удовлетворяет условиям KKT.
\end{theorem}

\begin{proof}[Достаточность]
Поскольку \(\lambda^\star\ge0\), функция
\[
    x\longmapsto
    \mathcal L(x,\lambda^\star,\nu^\star)
\]
выпукла. Стационарность поэтому означает, что \(x^\star\) является
её глобальным минимумом. Следовательно,
\[
    g(\lambda^\star,\nu^\star)
    =
    \mathcal L
    (x^\star,\lambda^\star,\nu^\star).
\]

Из прямой допустимости и дополняющей нежёсткости
\[
\begin{aligned}
    \mathcal L
    (x^\star,\lambda^\star,\nu^\star)
    &=
    f(x^\star)
    +
    \sum_i
    \lambda_i^\star h_i(x^\star)
    +
    \sum_j
    \nu_j^\star\ell_j(x^\star)
    \\
    &=
    f(x^\star).
\end{aligned}
\]
Следовательно,
\[
    g(\lambda^\star,\nu^\star)
    =
    f(x^\star).
\]
Теперь используем слабую двойственность и допустимость
\(x^\star\):
\[
    g(\lambda^\star,\nu^\star)
    \le
    d^\star
    \le
    p^\star
    \le
    f(x^\star).
\]
Крайние величины равны, поэтому равенства выполняются во всей цепочке:
\[
    g(\lambda^\star,\nu^\star)
    =d^\star=p^\star=f(x^\star).
\]
\end{proof}

Точные KKT-условия дают сертификат оптимальности. Малые KKT-невязки
часто используются как критерий остановки, но сами по себе не
гарантируют близости к решению: для количественной оценки нужны
дополнительные условия регулярности или оценки ошибки.

На практике проверяются невязки:
\[
    \norm{
        \nabla_x
        \mathcal L(x,\lambda,\nu)
    },
\]
\[
    \positivepart{h_i(x)},
    \qquad
    \abs{\ell_j(x)},
    \qquad
    \positivepart{-\lambda_i},
\]
\[
    \abs{\lambda_i h_i(x)}.
\]

\section{Примеры KKT: квадратичная задача и SVM}

Рассмотрим
\[
    \begin{aligned}
        \minimize_{x\in\RR^n}
        \quad&
        \frac12x^{\mathsf T}Qx+c^{\mathsf T}x,
        \\
        \subjectto
        \quad&
        Ax=b,
    \end{aligned}
    \tag{\ensuremath{QP}}
\]
где
\[
    Q=Q^{\mathsf T}\succ0.
\]

Функция Лагранжа равна
\[
    \mathcal L(x,\nu)
    =
    \frac12x^{\mathsf T}Qx
    +
    c^{\mathsf T}x
    +
    \nu^{\mathsf T}(Ax-b).
\]

Ограничений-неравенств нет, поэтому условия двойственной допустимости
и дополняющей нежёсткости отсутствуют.

Стационарность по \(x\) даёт
\[
    Qx^\star+c+A^{\mathsf T}\nu^\star=0.
\]
Прямая допустимость имеет вид
\[
    Ax^\star=b.
\]

Эти равенства объединяются в систему
\[
    \begin{pmatrix}
        Q & A^{\mathsf T}
        \\
        A & 0
    \end{pmatrix}
    \begin{pmatrix}
        x^\star
        \\
        \nu^\star
    \end{pmatrix}
    =
    \begin{pmatrix}
        -c
        \\
        b
    \end{pmatrix}.
\]

Матрица
\[
    \begin{pmatrix}
        Q & A^{\mathsf T}
        \\
        A & 0
    \end{pmatrix}
\]
называется KKT-матрицей. При \(Q\succ0\) и полном строчном ранге
\(A\) она невырождена, поэтому решение \((x^\star,\nu^\star)\)
единственно.

\begin{remark}
Для простой квадратичной задачи KKT-система непосредственно даёт
решение. В более сложных задачах условия KKT обычно используются как
критерий остановки итерационного алгоритма.
\end{remark}



Для строго линейно разделимой выборки рассмотрим прямую задачу
\[
    \begin{aligned}
        \minimize_{w\in\RR^n,\ b\in\RR}
        \quad&
        \frac12\norm{w}^2,
        \\
        \subjectto
        \quad&
        y_i(w^{\mathsf T}x_i+b)\ge1,
        \qquad i=1,\ldots,m.
    \end{aligned}
\]

Запишем ограничения в стандартной форме:
\[
    h_i(w,b)
    =
    1-y_i(w^{\mathsf T}x_i+b)
    \le0.
\]

Функция Лагранжа равна
\[
    \mathcal L(w,b,\lambda)
    =
    \frac12\norm{w}^2
    +
    \sum_{i=1}^{m}
    \lambda_i
    \left[
        1-y_i(w^{\mathsf T}x_i+b)
    \right].
\]

Условия стационарности:
\[
    \nabla_w\mathcal L=0
    \quad\Longleftrightarrow\quad
    w^\star
    =
    \sum_{i=1}^{m}
    \lambda_i^\star y_i x_i,
\]
\[
    \frac{\partial\mathcal L}{\partial b}=0
    \quad\Longleftrightarrow\quad
    \sum_{i=1}^{m}\lambda_i^\star y_i=0.
\]

Прямая допустимость:
\[
    y_i(w^{\star\mathsf T}x_i+b^\star)\ge1.
\]

Двойственная допустимость:
\[
    \lambda_i^\star\ge0.
\]

Дополняющая нежёсткость:
\[
    \lambda_i^\star
    \left[
        1-y_i(w^{\star\mathsf T}x_i+b^\star)
    \right]
    =
    0.
\]

Если
\[
    \lambda_i^\star>0,
\]
то
\[
    y_i(w^{\star\mathsf T}x_i+b^\star)=1.
\]
Такие точки лежат на границе максимального отступа.

\begin{definition}[Опорный вектор]
Точку \(x_i\) будем называть опорным вектором двойственного
представления, если
\[
    \lambda_i^\star>0.
\]
\end{definition}

В вырожденном случае точка может лежать на границе отступа и при этом
иметь нулевой множитель. Поэтому геометрическое множество граничных
точек и носитель двойственного решения не обязаны совпадать.

Для любого опорного вектора
\[
    b^\star
    =
    y_i-w^{\star\mathsf T}x_i.
\]
На практике значение \(b^\star\) можно усреднить по нескольким
опорным векторам.

Точки с \(\lambda_i^\star=0\) не входят в представление
\[
    w^\star
    =
    \sum_i\lambda_i^\star y_i x_i.
\]
Если большинство множителей нулевые, решение имеет разреженное
представление через опорные объекты. Однако малый размер этого
множества не следует из KKT-условий автоматически.

\section{Геометрия градиентного спуска и импульс}

Рассмотрим задачу без ограничений
\[
    \minimize_{x\in\RR^n}f(x).
\]
Градиентный спуск имеет вид
\[
    x^{k+1}
    =
    x^k-\eta\nabla f(x^k).
\]

Градиент ортогонален линии уровня функции и указывает направление
наибольшего локального возрастания. Поэтому вектор
\[
    -\nabla f(x^k)
\]
задаёт направление наискорейшего локального убывания в евклидовой
норме.

Для плохо обусловленных квадратичных функций линии уровня сильно
вытянуты. В этом случае градиентный спуск может многократно пересекать
узкую долину, образуя зигзагообразную траекторию.



Для уменьшения колебаний можно учитывать направление предыдущего
шага:
\[
    x^{k+1}
    =
    x^k
    -
    \eta\nabla f(x^k)
    +
    \beta(x^k-x^{k-1}),
\]
где
\[
    \beta\in[0,1)
\]
--- параметр импульса.

При согласованной инициализации
\[
    x^k-x^{k-1}=-\eta v^k
\]
эквивалентная запись использует дополнительную переменную:
\[
    v^{k+1}
    =
    \beta v^k+\nabla f(x^k),
\]
\[
    x^{k+1}
    =
    x^k-\eta v^{k+1}.
\]

Предыдущая скорость сглаживает резкие изменения направления
градиента и может ускорить движение вдоль вытянутой долины.

\begin{remark}[Классификация метода]
В транскрипте метод с импульсом сопоставляется с методами второго
порядка. В строгой терминологии это метод первого порядка с памятью:
он использует градиенты, но не использует Гессиан. Дополнительная
стоимость состоит в хранении ещё одного вектора размерности \(n\).
\end{remark}

Для модели с миллиардами параметров дополнительный вектор состояния
может требовать значительной памяти. Поэтому скорость по числу
итераций необходимо сравнивать с фактической стоимостью одной
итерации и объёмом памяти.

\section{Лемма об убывании и шаг градиентного спуска}

Далее предполагается, что \(f\) выпукла и \(L\)-гладка:
\[
    \norm{\nabla f(x)-\nabla f(y)}
    \le
    L\norm{x-y}.
\]

\begin{lemma}[Лемма об убывании]
Для любой \(L\)-гладкой функции
\[
    f(y)
    \le
    f(x)
    +
    \inner{\nabla f(x)}{y-x}
    +
    \frac{L}{2}\norm{y-x}^2.
\]
\end{lemma}

\begin{proof}
Положим
\[
    \gamma(t)=x+t(y-x),
    \qquad t\in[0,1].
\]
Тогда
\[
\begin{aligned}
    f(y)-f(x)
    &=
    \int_0^1
    \inner{
        \nabla f(\gamma(t))
    }{
        y-x
    }
    \dd t
    \\
    &=
    \inner{\nabla f(x)}{y-x}
    \\
    &\quad+
    \int_0^1
    \inner{
        \nabla f(\gamma(t))-\nabla f(x)
    }{
        y-x
    }
    \dd t.
\end{aligned}
\]
По неравенству Коши--Буняковского и \(L\)-гладкости
\[
\begin{aligned}
    \inner{
        \nabla f(\gamma(t))-\nabla f(x)
    }{
        y-x
    }
    &\le
    Lt\norm{y-x}^2.
\end{aligned}
\]
Интегрирование даёт
\[
    f(y)-f(x)
    \le
    \inner{\nabla f(x)}{y-x}
    +
    \frac{L}{2}\norm{y-x}^2.
\]
\end{proof}

Правая часть является квадратичной верхней моделью функции около
точки \(x\).



Зафиксируем текущую точку \(x^k\). Из леммы об убывании
\[
\begin{aligned}
    f(y)
    \le{}&
    f(x^k)
    +
    \inner{\nabla f(x^k)}{y-x^k}
    +
    \frac{L}{2}\norm{y-x^k}^2.
\end{aligned}
\]

Обозначим
\[
    g^k=\nabla f(x^k).
\]
Дополняя квадрат, получаем
\[
\begin{aligned}
    f(y)
    \le{}&
    f(x^k)
    -
    \frac{1}{2L}\norm{g^k}^2
    \\
    &+
    \frac{L}{2}
    \norm{
        y-
        \left(
            x^k-\frac1L g^k
        \right)
    }^2.
\end{aligned}
\]

Правая часть минимальна при
\[
    y
    =
    x^k-\frac1L\nabla f(x^k).
\]
Поэтому естественный шаг градиентного спуска равен
\[
    \eta=\frac1L,
\]
а обновление имеет вид
\[
    x^{k+1}
    =
    x^k-\frac1L\nabla f(x^k).
\]

\begin{theorem}[Гарантированное убывание]
Для шага \(\eta=1/L\)
\[
    f(x^{k+1})
    \le
    f(x^k)
    -
    \frac{1}{2L}
    \norm{\nabla f(x^k)}^2.
\]
\end{theorem}

Если
\[
    \nabla f(x^k)\ne0,
\]
то
\[
    f(x^{k+1})<f(x^k).
\]
Следовательно, последовательность значений целевой функции
невозрастает.

\begin{remark}[Исправление коэффициента]
При дополнении квадрата коэффициент перед нормой градиента равен
\[
    \frac{1}{2L},
\]
а не \(1/L\). Эта арифметическая неточность появляется в устном
выводе и затем исправляется.
\end{remark}



\begin{algorithm}[H]
\caption{Градиентный спуск для \(L\)-гладкой функции}
\label{alg:lecture11-gradient-descent}
\begin{algorithmic}[1]

\Require
функция \(f\), константа гладкости \(L>0\),
начальная точка \(x^0\)

\For{\(k=0,1,2,\ldots\)}

    \State
    \(g^k\gets\nabla f(x^k)\)

    \If{\(\norm{g^k}\) достаточно мала}

        \State
        \Return \(x^k\)

    \EndIf

    \State
    \(x^{k+1}\gets x^k-\dfrac1L g^k\)

\EndFor

\end{algorithmic}
\end{algorithm}

Если \(L\) неизвестна, фиксированный шаг \(1/L\) непосредственно
использовать нельзя. На практике шаг выбирают экспериментально или
определяют посредством поиска по направлению. Эти процедуры
рассматриваются отдельно.

\section{Скорость для выпуклой гладкой функции}

Пусть \(x^\star\) --- глобальный минимум выпуклой \(L\)-гладкой
функции.

\begin{theorem}[Сублинейная сходимость]
Для градиентного спуска
\[
    x^{k+1}
    =
    x^k-\frac1L\nabla f(x^k)
\]
выполняется
\[
    f(x^k)-f(x^\star)
    \le
    \frac{
        L\norm{x^0-x^\star}^2
    }{
        2k
    },
    \qquad k\ge1.
\]
\end{theorem}

\begin{proof}
Обозначим \(g^k=\nabla f(x^k)\). По выпуклости
\[
    f(x^k)-f(x^\star)
    \le
    \inner{g^k}{x^k-x^\star}.
\]
По гарантированному убыванию
\[
    f(x^{k+1})
    \le
    f(x^k)-\frac{1}{2L}\norm{g^k}^2.
\]
Следовательно,
\[
\begin{aligned}
    f(x^{k+1})-f(x^\star)
    \le{}&
    \inner{g^k}{x^k-x^\star}
    -
    \frac{1}{2L}\norm{g^k}^2
    \\
    ={}&
    \frac{L}{2}
    \left(
        \norm{x^k-x^\star}^2
        -
        \norm{x^{k+1}-x^\star}^2
    \right).
\end{aligned}
\]
Суммируя по \(k=0,\ldots,N-1\), получаем
\[
    \sum_{k=0}^{N-1}
    \bigl(
        f(x^{k+1})-f(x^\star)
    \bigr)
    \le
    \frac{L}{2}
    \norm{x^0-x^\star}^2.
\]
Поскольку \(f(x^k)\) невозрастает,
\[
    N\bigl(f(x^N)-f(x^\star)\bigr)
    \le
    \frac{L}{2}
    \norm{x^0-x^\star}^2.
\]
\end{proof}

Итак,
\[
    f(x^k)-f^\star
    =
    \BigO\left(\frac1k\right).
\]

Чтобы гарантировать
\[
    f(x^k)-f^\star\le\varepsilon,
\]
достаточно
\[
    k
    \ge
    \frac{
        L\norm{x^0-x^\star}^2
    }{
        2\varepsilon
    }.
\]
Итерационная сложность составляет
\[
    \BigO\left(\frac1\varepsilon\right).
\]

\section{Асимптотические обозначения и порядок сходимости}

\begin{definition}[Большое \(O\)]
Запись
\[
    a_k=\BigO(b_k)
\]
означает, что существуют \(C>0\) и \(k_0\), такие что
\[
    \abs{a_k}
    \le
    C\abs{b_k}
\]
для всех \(k\ge k_0\).
\end{definition}

Например,
\[
    k^2-2k+3k^3
    =
    \BigO(k^3).
\]

Оценка
\[
    f(x^k)-f^\star
    =
    \BigO(k^{-1})
\]
описывает скорость уменьшения ошибки по значению функции, а не
обязательно скорость уменьшения расстояния
\(\norm{x^k-x^\star}\).



Пусть
\[
    x^k\to x^\star,
    \qquad
    e_k=\norm{x^k-x^\star},
\]
и \(e_k>0\) для всех достаточно больших \(k\).

\begin{definition}[Порядок сходимости]
Последовательность имеет порядок сходимости \(q\ge1\), если
существует \(\rho\in(0,\infty)\), для которого
\[
    \lim_{k\to\infty}
    \frac{e_{k+1}}{e_k^q}
    =
    \rho.
\]
Число \(\rho\) называется асимптотической константой ошибки.
\end{definition}

При \(q=1\) говорят о линейной сходимости, если \(\rho<1\).
При \(q=2\) сходимость называется квадратичной.

Оценку
\[
    f(x^k)-f^\star=\BigO(k^{-1})
\]
называют сублинейной, поскольку она не задаёт геометрического
убывания ошибки. Из одной оценки большого \(O\) нельзя заключить, что
отношение последовательных ошибок обязательно стремится к единице.

Сильная выпуклость позволит получить более быструю линейную
сходимость градиентного метода. При дополнительных предположениях и подходящем выборе параметров
методы с импульсом позволяют улучшить зависимость скорости от числа
обусловленности.

\begin{lecturesummary}
Условия Каруша--Куна--Таккера состоят из стационарности, прямой и
двойственной допустимости и дополняющей нежёсткости. Для выпуклой
задачи они достаточны для глобальной оптимальности, а при условии
регулярности также необходимы. В квадратичной задаче с равенствами
KKT-условия образуют линейную систему с седловой KKT-матрицей. Для
SVM положительные двойственные множители выделяют опорные векторы и
позволяют восстановить смещение разделяющей гиперплоскости.

Для \(L\)-гладкой функции выполняется лемма об убывании
\[
    f(y)
    \le
    f(x)
    +
    \inner{\nabla f(x)}{y-x}
    +
    \frac{L}{2}\norm{y-x}^2.
\]
Минимизация этой верхней модели приводит к шагу
\[
    x^{k+1}
    =
    x^k-\frac1L\nabla f(x^k)
\]
и оценке
\[
    f(x^{k+1})
    \le
    f(x^k)
    -
    \frac{1}{2L}
    \norm{\nabla f(x^k)}^2.
\]
Для выпуклой \(L\)-гладкой функции градиентный спуск имеет скорость
\[
    f(x^k)-f^\star
    =
    \BigO\left(\frac1k\right)
\]
и требует порядка \(1/\varepsilon\) итераций для достижения ошибки
\(\varepsilon\) по значению функции.
\end{lecturesummary}
\lecture{12}{Ускорение при сильной выпуклости}

\sourcevideo
    {https://www.youtube.com/playlist?list=PLOzRYVm0a65fhKDS8cYIC7YCuVGJ4FOJx}
    {Week 4: Lecture 12: Acceleration under Strong Convexity}

Рассмотрим градиентный спуск с шагом \(1/L\):
\[
    x^{k+1}=x^k-\frac1L\nabla f(x^k).
\]
Для выпуклой \(L\)-гладкой функции он имеет скорость
\(\mathcal O(1/k)\). Если функция дополнительно удовлетворяет
неравенству Поляка--Лоясиевича, та же итерация сходится линейно.
Таким образом, в этой лекции ускоряется не алгоритм, а теоретическая
оценка за счёт более сильной геометрии целевой функции.

\section{Лемма об убывании}

Рассматривается задача
\[
    \minimize_{x\in\RR^n} f(x),
\]
где минимум достигается:
\[
    x^\star\in\operatorname*{argmin}_{x\in\RR^n}f(x),
    \qquad
    f^\star=f(x^\star).
\]
Гладкость с константой \(L\) означает
\[
    \norm{\nabla f(x)-\nabla f(y)}
    \le
    L\norm{x-y}.
\]

\begin{lemma}[Квадратичная верхняя оценка]
Если \(f\) является \(L\)-гладкой, то для любых \(x,y\in\RR^n\)
\[
    f(y)
    \le
    f(x)
    +
    \inner{\nabla f(x)}{y-x}
    +
    \frac L2\norm{y-x}^2.
\]
\end{lemma}

Подставляя
\[
    y=x-\frac1L\nabla f(x),
\]
получаем
\[
\begin{aligned}
    f\left(x-\frac1L\nabla f(x)\right)
    &\le
    f(x)
    -\frac1L\norm{\nabla f(x)}^2
    +\frac1{2L}\norm{\nabla f(x)}^2
    \\
    &=
    f(x)-\frac1{2L}\norm{\nabla f(x)}^2.
\end{aligned}
\]
Поэтому для градиентного спуска
\[
    f(x^{k+1})
    \le
    f(x^k)
    -
    \frac1{2L}\norm{\nabla f(x^k)}^2,
\]
и последовательность значений функции невозрастает.

\section{Выпуклый режим}

Для получения явной скорости используется более сильная форма
предыдущей оценки.

\begin{lemma}[Трёхточечное неравенство]
Пусть \(f\) выпукла и \(L\)-гладка, а
\[
    x^{k+1}=x^k-\frac1L\nabla f(x^k).
\]
Тогда для любого \(x\in\RR^n\)
\[
\begin{aligned}
    f(x^{k+1})
    \le
    f(x)
    +
    \frac L2
    \left(
        \norm{x^k-x}^2
        -
        \norm{x^{k+1}-x}^2
    \right).
\end{aligned}
\]
\end{lemma}

\begin{proof}
Из гладкости
\[
\begin{aligned}
    f(x^{k+1})
    \le{}&
    f(x^k)
    +
    \inner{\nabla f(x^k)}{x^{k+1}-x^k}
    \\
    &+
    \frac L2\norm{x^{k+1}-x^k}^2.
\end{aligned}
\]
По выпуклости
\[
    f(x^k)
    \le
    f(x)
    +
    \inner{\nabla f(x^k)}{x^k-x}.
\]
Складывая неравенства и используя
\[
    \nabla f(x^k)=L(x^k-x^{k+1}),
\]
получаем
\[
\begin{aligned}
    f(x^{k+1})
    \le{}&
    f(x)
    +
    L\inner{x^k-x^{k+1}}{x^{k+1}-x}
    \\
    &+
    \frac L2\norm{x^{k+1}-x^k}^2.
\end{aligned}
\]
Теперь применяется тождество
\[
    2\inner{a-b}{b-c}
    =
    \norm{a-c}^2
    -
    \norm{a-b}^2
    -
    \norm{b-c}^2.
\]
\end{proof}

Полагая \(x=x^\star\), имеем
\[
\begin{aligned}
    f(x^{k+1})-f^\star
    \le
    \frac L2
    \left(
        \norm{x^k-x^\star}^2
        -
        \norm{x^{k+1}-x^\star}^2
    \right).
\end{aligned}
\]
Суммирование по \(k=0,\ldots,K\) даёт
\[
    \sum_{k=0}^{K}
    \bigl(f(x^{k+1})-f^\star\bigr)
    \le
    \frac L2\norm{x^0-x^\star}^2.
\]
Поскольку \(f(x^k)\) невозрастает,
\[
\begin{aligned}
    (K+1)
    \bigl(f(x^{K+1})-f^\star\bigr)
    \le
    \frac L2\norm{x^0-x^\star}^2.
\end{aligned}
\]

\begin{theorem}[Скорость для выпуклой гладкой функции]
Пусть \(f\) выпукла и \(L\)-гладка. Тогда
\[
    f(x^{K+1})-f^\star
    \le
    \frac{L\norm{x^0-x^\star}^2}{2(K+1)}.
\]
\end{theorem}

Следовательно,
\[
    f(x^k)-f^\star
    =
    \mathcal O\!\left(\frac1k\right),
    \qquad
    k
    =
    \mathcal O\!\left(\frac1\varepsilon\right)
\]
для достижения ошибки не более \(\varepsilon\). Термин
«сублинейная скорость» относится к верхней оценке \(C/k\); он не
означает, что фактическое отношение соседних ошибок обязано равняться
\(k/(k+1)\).

\section{PL-условие и линейная сходимость}

\begin{definition}[Неравенство Поляка--Лоясиевича]
Дифференцируемая функция \(f\) удовлетворяет PL-неравенству с
константой \(\mu>0\), если
\[
    \frac12\norm{\nabla f(x)}^2
    \ge
    \mu\bigl(f(x)-f^\star\bigr)
\]
для всех \(x\).
\end{definition}

Всякая дифференцируемая \(\mu\)-сильно выпуклая функция
удовлетворяет PL-неравенству с той же константой, но дальнейшее
доказательство использует только PL-свойство.

Из леммы об убывании
\[
    f(x^{k+1})
    \le
    f(x^k)
    -
    \frac1{2L}\norm{\nabla f(x^k)}^2.
\]
По PL-неравенству
\[
    \frac12\norm{\nabla f(x^k)}^2
    \ge
    \mu\bigl(f(x^k)-f^\star\bigr).
\]
Следовательно,
\[
\begin{aligned}
    f(x^{k+1})-f^\star
    \le
    \left(1-\frac\mu L\right)
    \bigl(f(x^k)-f^\star\bigr).
\end{aligned}
\]

\begin{theorem}[Линейная сходимость при PL-условии]
Пусть \(f\) является \(L\)-гладкой и удовлетворяет
PL-неравенству с константой \(\mu>0\). Тогда
\[
    f(x^k)-f^\star
    \le
    \left(1-\frac\mu L\right)^k
    \bigl(f(x^0)-f^\star\bigr).
\]
\end{theorem}

Для \(L\)-гладкой PL-функции можно считать \(0<\mu\le L\), поэтому
коэффициент \(1-\mu/L\) принадлежит \([0,1)\). Ошибка уменьшается
геометрически: если она ещё велика, PL-неравенство не позволяет норме
градиента быть слишком малой.

Так как \(\nabla f(x^\star)=0\), из гладкости следует
\[
    f(x^0)-f^\star
    \le
    \frac L2\norm{x^0-x^\star}^2.
\]
Поэтому
\[
    f(x^k)-f^\star
    \le
    \frac L2
    \left(1-\frac\mu L\right)^k
    \norm{x^0-x^\star}^2.
\]
При нескольких минимизаторах \(x^\star\) можно выбирать произвольно;
однозначность гарантируется только сильной выпуклостью.

\section{Число обусловленности и сложность}

Для \(L\)-гладкой и \(\mu\)-сильно выпуклой функции вводится число
обусловленности
\[
    \kappa=\frac L\mu\ge1.
\]
Тогда коэффициент сходимости равен
\[
    1-\frac\mu L=1-\frac1\kappa.
\]
При \(\kappa\approx1\) кривизна функции почти одинакова во всех
направлениях. При \(\kappa\gg1\) линии уровня вытянуты, а коэффициент
сжатия близок к единице.

Для квадратичной функции
\[
    f(x)=\frac12x^{\mathsf T}Qx+c^{\mathsf T}x,
    \qquad
    Q\succ0,
\]
можно взять
\[
    L=\lambda_{\max}(Q),
    \qquad
    \mu=\lambda_{\min}(Q),
\]
поэтому
\[
    \kappa
    =
    \frac{\lambda_{\max}(Q)}{\lambda_{\min}(Q)}.
\]

Обозначим
\[
    \Delta_k=f(x^k)-f^\star.
\]
Из неравенства \(1-t\le e^{-t}\) получаем
\[
    \Delta_k
    \le
    \exp\left(-\frac\mu Lk\right)\Delta_0.
\]
Чтобы обеспечить \(\Delta_k\le\varepsilon\), достаточно
\[
    k
    \ge
    \frac L\mu
    \log\left(\frac{\Delta_0}{\varepsilon}\right).
\]
Следовательно,
\[
    k
    =
    \mathcal O\!\left(
        \kappa\log\frac1\varepsilon
    \right)
\]
при фиксированной начальной ошибке. В выпуклом случае зависимость была
\(\mathcal O(1/\varepsilon)\); PL-свойство заменяет её логарифмической
зависимостью от требуемой точности.

\section{PL-класс шире сильной выпуклости}

Сильная выпуклость влечёт PL-неравенство, но обратная импликация
неверна даже для выпуклых гладких функций. Рассмотрим
\[
    f(x_1,x_2)=\frac12x_1^2.
\]
Тогда
\[
    \nabla f(x_1,x_2)
    =
    \begin{pmatrix}
        x_1\\
        0
    \end{pmatrix},
    \qquad
    f^\star=0,
\]
и
\[
    \frac12\norm{\nabla f(x)}^2
    =
    \frac12x_1^2
    =
    f(x)-f^\star.
\]
Следовательно, функция удовлетворяет PL-неравенству с \(\mu=1\).
Однако она постоянна вдоль направления \((0,1)^{\mathsf T}\), не
является сильно выпуклой и имеет множество минимумов
\[
    \operatorname*{argmin}f
    =
    \setgiven{(0,x_2)}{x_2\in\RR}.
\]
Таким образом,
\[
    \text{сильная выпуклость}
    \Longrightarrow
    \text{PL-неравенство},
\]
но не наоборот.

\section{Интерпретация и ограничения}

В обоих режимах используется одна и та же итерация
\[
    x^{k+1}=x^k-\frac1L\nabla f(x^k).
\]
Для общего выпуклого класса градиент может быть малым без
количественной связи с ошибкой по значению, поэтому получается оценка
порядка \(1/k\). PL-неравенство устанавливает связь
\[
    \norm{\nabla f(x)}^2
    \ge
    2\mu\bigl(f(x)-f^\star\bigr),
\]
из-за которой каждая итерация уменьшает текущую ошибку как минимум в
фиксированное число раз.

Шаг \(1/L\) требует знания или оценки константы гладкости. Для
предсказания линейной скорости нужна также оценка \(\mu\). Если
\(\mu\ll L\), то большое число обусловленности делает даже линейную
сходимость медленной. Одна итерация использует один полный градиент и
хранит текущий вектор параметров. При стохастическом или
распределённом градиенте появляются дополнительные ошибки, связанные
с дисперсией и коммуникациями, поэтому оценки этой лекции напрямую к
ним не переносятся.

\lecture{13}{Инерционные методы}

\sourcevideo
    {https://www.youtube.com/playlist?list=PLOzRYVm0a65fhKDS8cYIC7YCuVGJ4FOJx}
    {Week 4: Lecture 13: Accelerate the Convergence Even Further}

Для выпуклой \(L\)-гладкой функции градиентный спуск имеет скорость
\(\mathcal O(1/k)\), а при PL-условии --- линейную скорость с
зависимостью от числа обусловленности \(\kappa=L/\mu\). В этой лекции
ускорение достигается не усилением предположений о функции, а
изменением самого алгоритма: к градиентному шагу добавляется память о
предыдущем движении.

\section{Мотивация импульса и исходные оценки}

Рассмотрим задачу
\[
    \min_{x\in\RR^n}f(x)
\]
и градиентный спуск
\[
    x^{k+1}=x^k-\frac1L\nabla f(x^k).
\]
Для выпуклой \(L\)-гладкой функции
\[
    f(x^k)-f^\star
    =
    \mathcal O\!\left(\frac1k\right),
\]
поэтому для достижения ошибки не более \(\varepsilon\) требуется
\[
    k
    =
    \mathcal O\!\left(\frac1\varepsilon\right).
\]
Если функция удовлетворяет PL-неравенству
\[
    \frac12\norm{\nabla f(x)}^2
    \ge
    \mu\bigl(f(x)-f^\star\bigr),
\]
то
\[
    f(x^k)-f^\star
    \le
    \left(1-\frac\mu L\right)^k
    \bigl(f(x^0)-f^\star\bigr).
\]
В сильно выпуклом случае
\[
    \kappa=\frac L\mu,
\]
и сложность равна
\[
    \mathcal O\!\left(
        \kappa\log\frac1\varepsilon
    \right).
\]

Причина замедления на плохо обусловленной функции видна по геометрии
линий уровня. Градиент ортогонален линии уровня, поэтому в вытянутой
долине последовательные шаги могут многократно пересекать её поперёк,
медленно продвигаясь к минимуму. Обычный градиентный спуск использует
только \(x^k\) и \(\nabla f(x^k)\); направление предыдущего
перемещения полностью теряется. Инерционный метод сохраняет часть
вектора
\[
    x^k-x^{k-1}
\]
и тем самым может ослаблять поперечные колебания, не уничтожая
накопленное движение вдоль долины.

\section{Метод тяжёлого шара}

Метод тяжёлого шара Поляка задаётся правилом
\[
    x^{k+1}
    =
    x^k
    -
    \alpha\nabla f(x^k)
    +
    \beta(x^k-x^{k-1}),
\]
где
\[
    \alpha>0,
    \qquad
    \beta\ge0.
\]
При \(\beta=0\) получается обычный градиентный спуск. Если ввести
скорость
\[
    v^k=x^k-x^{k-1},
\]
то алгоритм принимает вид
\[
    v^{k+1}
    =
    \beta v^k-\alpha\nabla f(x^k),
    \qquad
    x^{k+1}=x^k+v^{k+1}.
\]

\begin{algorithm}[H]
\caption{Метод тяжёлого шара}
\label{alg:heavy-ball}
\begin{algorithmic}[1]

\Require
\(x^0\in\RR^n\), \(v^0=0\),
параметры \(\alpha>0\), \(\beta\ge0\)

\For{\(k=0,1,2,\ldots\)}

    \State
    \(g^k\gets\nabla f(x^k)\)

    \State
    \(v^{k+1}\gets\beta v^k-\alpha g^k\)

    \State
    \(x^{k+1}\gets x^k+v^{k+1}\)

\EndFor

\end{algorithmic}
\end{algorithm}

Метод использует два последовательных состояния, но остаётся методом
первого порядка: Гессиан не вычисляется. Одна итерация требует одного
градиента, как и обычный градиентный спуск, однако дополнительно нужно
хранить один вектор размерности \(n\).

В лекции при надлежащем выборе \(\alpha\) и \(\beta\) приводятся
оценки
\[
    \mathcal O\!\left(\frac1\varepsilon\right)
\]
для общего выпуклого гладкого случая и
\[
    \mathcal O\!\left(
        \sqrt{\frac L\mu}
        \log\frac1\varepsilon
    \right)
    =
    \mathcal O\!\left(
        \sqrt\kappa
        \log\frac1\varepsilon
    \right)
\]
для сильно выпуклого случая. Подробное доказательство и точная область
допустимых параметров в лекции не приводятся; наиболее прозрачный
анализ метода тяжёлого шара получается для квадратичных функций.

\section{Метод Нестерова}

Метод Нестерова также использует импульс, но градиент вычисляется в
экстраполированной точке. Сначала строится
\[
    y^k
    =
    x^k+\beta_k(x^k-x^{k-1}),
\]
после чего выполняется шаг
\[
    x^{k+1}
    =
    y^k-\frac1L\nabla f(y^k).
\]
У тяжёлого шара градиент берётся в \(x^k\), а у метода Нестерова ---
в прогнозной точке \(y^k\).

Если \(f\) является \(L\)-гладкой и \(\mu\)-сильно выпуклой, можно
взять постоянный коэффициент
\[
    \beta
    =
    \frac{\sqrt L-\sqrt\mu}
         {\sqrt L+\sqrt\mu}
    =
    \frac{\sqrt\kappa-1}
         {\sqrt\kappa+1}.
\]
Тогда
\[
    y^k
    =
    x^k+\beta(x^k-x^{k-1}),
\]
\[
    x^{k+1}
    =
    y^k-\frac1L\nabla f(y^k),
\]
и для подходящей константы \(C>0\)
\[
    f(x^k)-f^\star
    \le
    C
    \left(
        1-\sqrt{\frac\mu L}
    \right)^k.
\]
Следовательно,
\[
    k
    =
    \mathcal O\!\left(
        \sqrt{\frac L\mu}
        \log\frac1\varepsilon
    \right).
\]

Удобно использовать стандартную эквивалентную параметризацию через \(t_k\).

\begin{algorithm}[H]
\caption{Ускоренный градиентный метод для выпуклой функции}
\label{alg:nesterov-convex}
\begin{algorithmic}[1]

\Require
\(x^{-1}=x^0\), константа \(L>0\), \(t_0=1\)

\For{\(k=0,1,2,\ldots\)}

    \If{\(k=0\)}

        \State
        \(y^0\gets x^0\)

    \Else

        \State
        \[
            y^k
            \gets
            x^k
            +
            \frac{t_{k-1}-1}{t_k}
            (x^k-x^{k-1})
        \]

    \EndIf

    \State
    \(x^{k+1}
        \gets
        y^k-\dfrac1L\nabla f(y^k)\)

    \State
    \[
        t_{k+1}
        \gets
        \frac{1+\sqrt{1+4t_k^2}}2
    \]

\EndFor

\end{algorithmic}
\end{algorithm}

\section{Скорости и оптимальность}

Для ускоренного метода Нестерова в выпуклом случае выполняется оценка
\[
    f(x^k)-f^\star
    \le
    \frac{
        2L\norm{x^0-x^\star}^2
    }{
        (k+1)^2
    }.
\]
Поэтому
\[
    f(x^k)-f^\star
    =
    \mathcal O\!\left(\frac1{k^2}\right),
\]
а для достижения точности \(\varepsilon\) достаточно
\[
    k
    =
    \mathcal O\!\left(
        \sqrt{
            \frac{
                L\norm{x^0-x^\star}^2
            }{
                \varepsilon
            }
        }
    \right).
\]
При фиксированных \(L\) и начальном расстоянии это даёт зависимость
\[
    \mathcal O\!\left(
        \frac1{\sqrt\varepsilon}
    \right)
\]
вместо
\[
    \mathcal O\!\left(
        \frac1\varepsilon
    \right)
\]
у градиентного спуска.

Основные оценки лекции имеют вид
\[
\begin{array}{c|c|c}
    \text{метод}
    &
    \text{выпуклая функция}
    &
    \text{сильно выпуклая функция}
    \\
    \hline
    \text{градиентный спуск}
    &
    \mathcal O(1/\varepsilon)
    &
    \mathcal O\!\left(
        \kappa\log(1/\varepsilon)
    \right)
    \\
    \text{тяжёлый шар}
    &
    \mathcal O(1/\varepsilon)
    &
    \mathcal O\!\left(
        \sqrt\kappa\log(1/\varepsilon)
    \right)
    \\
    \text{метод Нестерова}
    &
    \mathcal O(1/\sqrt\varepsilon)
    &
    \mathcal O\!\left(
        \sqrt\kappa\log(1/\varepsilon)
    \right)
\end{array}
\]

В модели оракула первого порядка существуют нижние оценки
\[
    \Omega\!\left(
        \sqrt{
            \frac{
                L\norm{x^0-x^\star}^2
            }{
                \varepsilon
            }
        }
    \right)
\]
для выпуклых \(L\)-гладких функций и
\[
    \Omega\!\left(
        \sqrt{\frac L\mu}
        \log\frac1\varepsilon
    \right)
\]
для \(L\)-гладких \(\mu\)-сильно выпуклых функций. Метод Нестерова
достигает этих порядков. Оптимальность относится именно к данной
модели доступа через значения функции и градиенты; дополнительная
структура задачи или использование информации второго порядка могут
изменить нижние оценки.

\section{Переход к непрерывному времени}

Градиентный спуск с шагом \(\eta>0\) можно записать как
\[
    \frac{x^{k+1}-x^k}{\eta}
    =
    -\nabla f(x^k).
\]
При формальном предельном переходе \(\eta\to0\) получается градиентный
поток
\[
    \dot x(t)
    =
    -\nabla f(x(t)).
\]

\begin{definition}[Градиентный поток]
Динамическая система
\[
    \dot x=-\nabla f(x)
\]
называется градиентным потоком функции \(f\).
\end{definition}

Для
\[
    f(x)=\frac12x^2
\]
имеем
\[
    \dot x=-x
\]
и
\[
    x(t)=e^{-t}x(0).
\]
Более общая система
\[
    \dot x=-ax,
    \qquad
    a>0,
\]
имеет решение
\[
    x(t)=e^{-at}x(0).
\]
В непрерывной модели увеличение \(a\) формально ускоряет затухание.
Однако дискретизация
\[
    x^{k+1}
    =
    (1-\eta a)x^k
\]
устойчива только при
\[
    0<\eta a<2.
\]
Следовательно, простое масштабирование непрерывного поля не даёт
бесплатного ускорения дискретного метода: рост коэффициента требует
уменьшения шага.

Глобальная \(L\)-гладкость особенно важна для анализа конечного шага,
поскольку она контролирует изменение градиента. В непрерывном анализе
некоторые оценки устойчивости можно получить без глобальной
константы \(L\), но условия регулярности поля всё равно нужны для
существования и единственности классического решения.

\section{Нормированный поток и практические ограничения}

Рассмотрим нормированный поток
\[
    \dot x
    =
    -
    \frac{
        \nabla f(x)
    }{
        \norm{\nabla f(x)}
    },
    \qquad
    \nabla f(x)\ne0.
\]
Его скорость удовлетворяет
\[
    \norm{\dot x}=1
\]
и не уменьшается около стационарной точки. Для
\[
    f(x)=\frac12x^2
\]
получаем
\[
    \dot x=-\operatorname{sign}(x),
\]
а решение до достижения нуля имеет вид
\[
    x(t)
    =
    \operatorname{sign}(x(0))
    \max\{
        \abs{x(0)}-t,
        0
    \}.
\]
Равновесие достигается за конечное время
\[
    T=\abs{x(0)}.
\]

В точке \(\nabla f(x)=0\) нормированное поле не определено, поэтому
нужно отдельно задать продолжение динамики или использовать
дифференциальное включение. После дискретизации возможны колебания
около минимума.

Все три дискретных метода используют один полный градиент за
итерацию. Градиентный спуск не требует дополнительного инерционного
состояния, а тяжёлый шар и метод Нестерова хранят ещё один вектор.
Ускорение по числу итераций не всегда совпадает с ускорением по
времени и памяти.

В распределённой задаче один градиентный шаг может включать
коммуникационный раунд, поэтому уменьшение числа итераций может
непосредственно сокращать число обменов, если стоимость одного шага
существенно не возрастает.

Непрерывная интерпретация связывает стационарные точки функции с
равновесиями динамической системы, сходимость алгоритма ---
с устойчивостью, а скорость оптимизации --- со скоростью затухания.
Следующая лекция вводит функции Ляпунова и формализует этот переход.

\lecture{14}{Основы теории устойчивости}

\sourcevideo
    {https://www.youtube.com/playlist?list=PLOzRYVm0a65fhKDS8cYIC7YCuVGJ4FOJx}
    {Week 4: Lecture 14: Stability Theory}

Оптимизационный алгоритм можно рассматривать как динамическую
систему, а его сходимость --- как устойчивость соответствующего
равновесия. Сначала это иллюстрируется сравнением градиентного спуска
и метода тяжёлого шара для плоской функции, затем вводятся основные
понятия устойчивости и прямой метод Ляпунова.

\section{Численный пример и метод тяжёлого шара}

Рассмотрим функцию
\[
    f(x)=x^4,
    \qquad
    f'(x)=4x^3,
    \qquad
    f''(x)=12x^2.
\]
Она строго выпукла и имеет единственный минимум
\[
    x^\star=0,
    \qquad
    f^\star=0,
\]
но не является глобально сильно выпуклой, поскольку
\[
    \inf_{x\in\RR}f''(x)=0,
\]
и не имеет глобально липшицева градиента, поскольку
\[
    \sup_{x\in\RR}f''(x)=+\infty.
\]
Вблизи минимума градиент \(4x^3\) быстро стремится к нулю, поэтому
градиентный спуск
\[
    x^{k+1}
    =
    x^k-4\alpha(x^k)^3
    =
    x^k\bigl(1-4\alpha(x^k)^2\bigr)
\]
замедляется. При большом шаге или большой начальной точке множитель
\(1-4\alpha(x^k)^2\) может иметь большую абсолютную величину, что
вызывает колебания или расходимость. Поэтому упоминаемую в лекции
граница шага около \(2\) относится только к конкретного эксперимента, а не как универсальную глобальную гарантию.

Метод тяжёлого шара имеет вид
\[
    x^{k+1}
    =
    x^k-4\alpha(x^k)^3
    +\beta(x^k-x^{k-1}),
\]
где \(\alpha>0\) --- шаг, а \(\beta\ge0\) --- коэффициент импульса.
Если
\[
    v^k=x^k-x^{k-1},
\]
то
\[
    v^{k+1}=\beta v^k-4\alpha(x^k)^3,
    \qquad
    x^{k+1}=x^k+v^{k+1}.
\]
В демонстрации используются, в частности,
\[
    x^0=-0.5,
    \qquad
    \alpha=0.5,
    \qquad
    \beta=0.4.
\]
Малое значение \(\beta\), например \(0.1\), почти не отличается от
обычного градиентного шага. Большое значение, например \(0.9\), может
ускорить движение в плоской области, но создаёт выраженные колебания
и немонотонность. Поэтому \(\alpha\) и \(\beta\) должны
согласовываться с гладкостью, кривизной и начальной точкой. Для
квадратичных сильно выпуклых функций параметры можно связать с
крайними собственными значениями гессиана; для общей функции обычно
нужны оценки \(L\) и \(\mu\) либо численная настройка.

Поскольку \(f^\star=0\), величина
\[
    f(x^k)=(x^k)^4
\]
непосредственно измеряет ошибку по значению функции. График
\(k\mapsto\log f(x^k)\) показывает различие в порядках величины и
выявляет временные увеличения ошибки, вызванные чрезмерным импульсом.

\section{Переход к непрерывной динамике}

Из дискретного обновления
\[
    x^{k+1}=x^k-\eta\nabla f(x^k)
\]
получаем
\[
    \frac{x^{k+1}-x^k}{\eta}
    =
    -\nabla f(x^k).
\]
В пределе малого шага возникает градиентный поток.

\begin{definition}[Градиентный поток]
Динамическая система
\[
    \dot x=-\nabla f(x)
\]
называется градиентным потоком функции \(f\).
\end{definition}

Его равновесия удовлетворяют
\[
    \nabla f(x^{\mathrm e})=0.
\]
Для выпуклой функции они являются глобальными минимумами, для строго
выпуклой --- единственным минимумом, а без выпуклости могут быть
локальными минимумами, максимумами или седловыми точками.

Непрерывная постановка позволяет проектировать и другие поля,
например
\[
    \dot x
    =
    -\frac{\nabla f(x)}
    {\lVert\nabla f(x)\rVert+\varepsilon},
    \qquad
    \varepsilon>0.
\]
Формальный выбор \(\varepsilon=0\) даёт единичную скорость вне
стационарных точек. Для
\[
    f(x)=\frac12x^2
\]
получается
\[
    \dot x=-\operatorname{sign}(x),
\]
и нуль достигается за конечное время. При \(\varepsilon>0\) поле
становится регулярнее, но точное конечно-временное достижение обычно
заменяется асимптотическим приближением.

Общая динамическая система записывается как
\[
    \dot x=g(t,x).
\]
Она называется неавтономной, если правая часть явно зависит от
времени, и автономной, если
\[
    \dot x=g(x).
\]
Точка \(x^{\mathrm e}\) является равновесием автономной системы, если
\[
    g(x^{\mathrm e})=0.
\]
Для неавтономной системы постоянное равновесие должно удовлетворять
\(g(t,x^{\mathrm e})=0\) для всех \(t\). После замены
\[
    z=x-x^{\mathrm e}
\]
равновесие переносится в начало координат. Например,
\[
    \dot x=-2(x-1)
\]
переходит в
\[
    \dot z=-2z.
\]

Непрерывные модели удобны для анализа и проектирования, но свойства
их дискретизации требуют отдельной проверки: устойчивое непрерывное
поле может потерять устойчивость при слишком крупном шаге.

\section{Устойчивость, притягательность и скорость}

Рассмотрим систему
\[
    \dot x=g(t,x),
    \qquad
    g(t,0)=0.
\]

\begin{definition}[Устойчивость по Ляпунову]
Равновесие \(x=0\) называется устойчивым, если для каждого
\(\varepsilon>0\) и каждого начального момента \(t_0\) существует
\(\delta(\varepsilon,t_0)>0\) такое, что
\[
    \lVert x(t_0)\rVert<\delta
    \quad\Longrightarrow\quad
    \lVert x(t)\rVert<\varepsilon
\]
для всех \(t\ge t_0\).
\end{definition}

Устойчивость означает сохранение близости, но не требует
\(x(t)\to0\). Она называется равномерной, если \(\delta\) можно
выбрать независимо от \(t_0\); для автономных систем эта независимость
обычно следует из инвариантности относительно сдвига времени.

\begin{definition}[Притягательность]
Равновесие локально притягивает траектории, если существует \(r>0\),
такое что
\[
    \lVert x(t_0)\rVert<r
    \quad\Longrightarrow\quad
    \lim_{t\to\infty}x(t)=0.
\]
Если предел выполняется для всех начальных состояний,
притягательность называется глобальной.
\end{definition}

Притягательность не контролирует промежуточные отклонения. Поэтому
равновесие может притягивать траектории, но не быть устойчивым.

\begin{definition}[Асимптотическая устойчивость]
Равновесие называется локально асимптотически устойчивым, если оно
устойчиво по Ляпунову и локально притягивает траектории. Если
притяжение глобально, говорят о глобальной асимптотической
устойчивости.
\end{definition}

\begin{definition}[Экспоненциальная устойчивость]
Равновесие локально экспоненциально устойчиво, если существуют
\(M\ge1\), \(\alpha>0\) и \(r>0\), для которых
\[
    \lVert x(t)\rVert
    \le
    M\exp\bigl(-\alpha(t-t_0)\bigr)
    \lVert x(t_0)\rVert
\]
при \(\lVert x(t_0)\rVert<r\) и всех \(t\ge t_0\). Если оценка
справедлива для любых начальных состояний, устойчивость глобальная.
\end{definition}

Экспоненциальная устойчивость влечёт асимптотическую и даёт время
достижения точности порядка
\[
    \frac1\alpha\log\frac1\varepsilon.
\]
Простейший пример ---
\[
    \dot x=-ax,
    \qquad
    a>0,
\]
для которого
\[
    x(t)=e^{-a(t-t_0)}x(t_0).
\]
Это одновременно градиентный поток функции
\[
    f(x)=\frac a2x^2.
\]

\section{Конечное и фиксированное время}

\begin{definition}[Конечно-временная сходимость]
Равновесие достигается за конечное время, если для каждого начального
состояния из некоторой области существует
\[
    T(x(t_0))<\infty,
\]
такое что
\[
    x(t)=0
\]
для всех \(t\ge t_0+T(x(t_0))\). Если равновесие при этом устойчиво
по Ляпунову, оно называется конечно-временно устойчивым.
\end{definition}

Для системы
\[
    \dot x=-\operatorname{sign}(x)
\]
имеем
\[
    T(x_0)=|x_0|,
\]
поэтому время зависит от начального состояния.

\begin{definition}[Фиксированно-временная сходимость]
Равновесие достигается за фиксированное время, если существует
\(T_{\max}<\infty\), такое что
\[
    T(x_0)\le T_{\max}
\]
для всех допустимых начальных состояний. При устойчивости по Ляпунову
говорят о фиксированно-временной устойчивости.
\end{definition}

Асимптотическая устойчивость гарантирует только существование
предела. Экспоненциальная устойчивость задаёт скорость затухания.
Конечно-временная сходимость обеспечивает точное достижение за время,
зависящее от \(x_0\). Фиксированно-временная сходимость даёт единую
верхнюю границу для всех запусков. В распределённых системах это
позволяет заранее ограничить момент завершения согласования.

Для градиентного потока соответствия имеют вид
\[
    \text{равновесие}
    \quad\Longleftrightarrow\quad
    \nabla f(x^\star)=0,
\]
\[
    \text{асимптотическая устойчивость}
    \quad\Longleftrightarrow\quad
    x(t)\to x^\star,
\]
\[
    \text{экспоненциальная устойчивость}
    \quad\Longleftrightarrow\quad
    \lVert x(t)-x^\star\rVert
    \le
    M e^{-\alpha(t-t_0)}
    \lVert x(t_0)-x^\star\rVert.
\]
Конечно- и фиксированно-временная устойчивость означают точное
достижение оптимума после конечного момента.

\section{Прямой метод Ляпунова}

Прямой метод исследует устойчивость без явного решения
дифференциального уравнения. Для этого строится скалярная функция
\(V\colon\RR^n\to\RR\).

\begin{definition}[Положительно определённая функция]
Функция \(V\) положительно определена относительно начала координат,
если
\[
    V(0)=0,
    \qquad
    V(x)>0
\]
для всех \(x\ne0\) в рассматриваемой области.
\end{definition}

Вдоль траектории автономной системы \(\dot x=g(x)\)
\[
    \dot V(x)
    =
    \nabla V(x)^{\mathsf T}g(x).
\]
Если \(\dot V\le0\), обобщённая энергия не возрастает; если
\(\dot V<0\) вне равновесия, она строго убывает.

\begin{theorem}[Прямой метод Ляпунова]
Пусть \(x=0\) --- равновесие системы \(\dot x=g(x)\) и в некоторой
его окрестности существует непрерывно дифференцируемая положительно
определённая функция \(V\). Если
\[
    \dot V(x)\le0,
\]
то равновесие устойчиво по Ляпунову. Если
\[
    \dot V(x)<0
\]
для всех \(x\ne0\), то равновесие локально асимптотически устойчиво.
\end{theorem}

Для глобального вывода обычно дополнительно требуется радиальная
неограниченность:
\[
    V(x)\to+\infty
    \qquad
    \text{при }
    \lVert x\rVert\to\infty.
\]

\begin{theorem}[Экспоненциальная оценка]
Пусть существуют \(c_1,c_2,c_3>0\), такие что
\[
    c_1\lVert x\rVert^2
    \le
    V(x)
    \le
    c_2\lVert x\rVert^2
\]
и
\[
    \dot V(x)
    \le
    -c_3\lVert x\rVert^2.
\]
Тогда равновесие экспоненциально устойчиво.
\end{theorem}

Действительно,
\[
    \dot V
    \le
    -\frac{c_3}{c_2}V,
\]
поэтому
\[
    V(x(t))
    \le
    \exp\left(
        -\frac{c_3}{c_2}(t-t_0)
    \right)
    V(x(t_0)).
\]
Двусторонняя оценка на \(V\) переносит это неравенство на норму
состояния. 

Для механической системы естественным кандидатом служит энергия. В
модели «масса--пружина--демпфер» можно взять
\[
    V(x,\dot x)
    =
    \frac12kx^2+
    \frac12m\dot x^2.
\]
Демпфер обеспечивает \(\dot V\le0\). Однако функция Ляпунова не
обязана иметь физический смысл и может быть специально
сконструированной математической величиной.

\section{Функции Ляпунова в оптимизации}

Для градиентного потока
\[
    \dot x=-\nabla f(x)
\]
естественно выбрать
\[
    V(x)=f(x)-f^\star.
\]
Тогда
\[
    \dot V(x)
    =
    \nabla f(x)^{\mathsf T}\dot x
    =
    -\lVert\nabla f(x)\rVert^2
    \le0.
\]
Если \(\nabla f(x)=0\) только в точке \(x^\star\), производная строго
отрицательна вне минимума; при дополнительных условиях регулярности и
компактности это ведёт к асимптотической устойчивости.

Если функция удовлетворяет PL-неравенству
\[
    \frac12\lVert\nabla f(x)\rVert^2
    \ge
    \mu\bigl(f(x)-f^\star\bigr),
\]
то
\[
    \dot V
    \le
    -2\mu V
\]
и
\[
    V(x(t))
    \le
    e^{-2\mu(t-t_0)}V(x(t_0)).
\]
Следовательно, ошибка по значению функции убывает экспоненциально.
Другой стандартный кандидат ---
\[
    V(x)=\frac12\lVert x-x^\star\rVert^2.
\]

Для алгоритмов с импульсом, ограничениями или сетевым взаимодействием
функция Ляпунова может одновременно включать ошибку по значению,
расстояние до решения, скорость, двойственные переменные и ошибку
согласования. Главная трудность прямого метода состоит именно в
выборе такой комбинации. Универсального простого правила, выдающего
удобную аналитическую функцию для любой нелинейной системы, нет.

Вместе с тем существуют обратные теоремы Ляпунова, гарантирующие
существование функции при различных видах устойчивости; сложной
остаётся её явная и удобная для расчётов конструкция.

\lecture{15}{Функции Ляпунова в оптимизации}

\sourcevideo
    {https://www.youtube.com/playlist?list=PLOzRYVm0a65fhKDS8cYIC7YCuVGJ4FOJx}
    {Week 4: Lecture 15: Connections to Optimization Problems}

Функция Ляпунова --- скалярная характеристика состояния, которая
принимает минимальное значение на исследуемом равновесии и не
возрастает вдоль траекторий. Скорость её убывания позволяет доказать
устойчивость, асимптотическую или экспоненциальную сходимость. В
оптимизации естественными кандидатами служат расстояние до множества
решений, ошибка по значению функции и норма градиента.

\section{Классы функций и определённость}

\begin{definition}[Функция класса \(\mathcal K\)]
Непрерывная функция
\[
    \alpha\colon[0,a)\to[0,\infty)
\]
принадлежит классу \(\mathcal K\), если
\[
    \alpha(0)=0
\]
и \(\alpha\) строго возрастает. Если область определения равна
\([0,\infty)\) и
\[
    \alpha(r)\to\infty
    \qquad
    \text{при }r\to\infty,
\]
то \(\alpha\in\mathcal K_\infty\).
\end{definition}

Функции \(r\), \(r^2\) и \(r^p\) при \(p>0\) принадлежат
\(\mathcal K_\infty\). Функция
\[
    \alpha(r)=1-e^{-r}
\]
принадлежит \(\mathcal K\), но не \(\mathcal K_\infty\), поскольку
ограничена сверху единицей.

Пусть \(V\colon\RR^n\to\RR\) непрерывна. Она называется положительно
полуопределённой относительно начала координат, если
\[
    V(0)=0,
    \qquad
    V(x)\ge0,
\]
и положительно определённой, если дополнительно
\[
    V(x)>0
    \qquad
    \text{при }x\ne0.
\]
Локальную положительную определённость можно выразить оценкой
\[
    V(x)\ge\alpha_1(\lVert x\rVert),
    \qquad
    \alpha_1\in\mathcal K.
\]
Если она справедлива глобально с
\(\alpha_1\in\mathcal K_\infty\), то \(V\) радиально неограничена:
\[
    V(x)\to\infty
    \qquad
    \text{при }\lVert x\rVert\to\infty.
\]

Например,
\[
    V(x)=\frac12\lVert x\rVert^2
\]
положительно определена и радиально неограничена, тогда как
\[
    V(x)=1-e^{-\lVert x\rVert^2}
\]
положительно определена, но ограничена. Поэтому положительная
определённость и радиальная неограниченность являются различными
свойствами.

Функция \(W\) называется отрицательно определённой, если
\[
    W(0)=0,
    \qquad
    W(x)<0
    \quad
    \text{при }x\ne0.
\]
Это эквивалентно положительной определённости \(-W\). Отрицательная
полуопределённость означает \(W(x)\le0\). В прямом методе положительно
определённой выбирают саму функцию Ляпунова, а отрицательной или
полуотрицательной --- её производную вдоль траекторий.

\section{Прямой метод и экспоненциальный критерий}

Рассмотрим автономную систему
\[
    \dot x=g(x),
    \qquad
    g(0)=0,
\]
и непрерывно дифференцируемую функцию \(V\). Вдоль решения
\[
\begin{aligned}
    \dot V(x)
    &=
    \frac{\mathrm d}{\mathrm dt}V(x(t))
    \\
    &=
    \nabla V(x)^{\mathsf T}\dot x
    \\
    &=
    \nabla V(x)^{\mathsf T}g(x).
\end{aligned}
\]

\begin{definition}[Функция Ляпунова]
Положительно определённая непрерывно дифференцируемая функция \(V\)
называется функцией Ляпунова для равновесия \(x=0\), если
\[
    \dot V(x)\le0
\]
в некоторой его окрестности.
\end{definition}

Функция Ляпунова не обязана совпадать с физической энергией. Это
может быть специально сконструированная величина, удобная для
доказательства требуемого типа устойчивости.

\begin{theorem}[Прямой метод Ляпунова]
Пусть в окрестности начала координат существует непрерывно
дифференцируемая положительно определённая функция \(V\). Если
\[
    \dot V(x)\le0,
\]
то равновесие устойчиво по Ляпунову. Если
\[
    \dot V(x)<0
\]
для всех \(x\ne0\), то равновесие локально асимптотически устойчиво.
Если условия выполняются глобально и \(V\) радиально неограничена, то
вывод является глобальным.
\end{theorem}

Условие \(\dot V\le0\) не гарантирует притяжения: обобщённая энергия
может сохраняться на ненулевых траекториях. Строгое убывание вне
равновесия устраняет эту возможность.

\begin{theorem}[Квадратичный критерий]
Пусть существуют \(c_1,c_2,c_3>0\), для которых
\[
    c_1\lVert x\rVert^2
    \le
    V(x)
    \le
    c_2\lVert x\rVert^2
\]
и
\[
    \dot V(x)
    \le
    -c_3\lVert x\rVert^2.
\]
Тогда
\[
    \lVert x(t)\rVert
    \le
    \sqrt{\frac{c_2}{c_1}}
    \exp\left(
        -\frac{c_3}{2c_2}(t-t_0)
    \right)
    \lVert x(t_0)\rVert,
\]
то есть равновесие экспоненциально устойчиво.
\end{theorem}

\begin{proof}
Из верхней оценки на \(V\) следует
\[
    \lVert x\rVert^2\ge\frac1{c_2}V(x),
\]
поэтому
\[
    \dot V\le-\frac{c_3}{c_2}V.
\]
Интегрирование даёт
\[
    V(x(t))
    \le
    e^{-\frac{c_3}{c_2}(t-t_0)}V(x(t_0)).
\]
Используя обе квадратичные оценки, получаем
\[
\begin{aligned}
    c_1\lVert x(t)\rVert^2
    &\le
    V(x(t))
    \\
    &\le
    e^{-\frac{c_3}{c_2}(t-t_0)}
    c_2\lVert x(t_0)\rVert^2.
\end{aligned}
\]
После извлечения квадратного корня получается требуемое неравенство.
\end{proof}

Таким образом, в стандартной оценке можно взять
\[
    M=\sqrt{\frac{c_2}{c_1}},
    \qquad
    \alpha=\frac{c_3}{2c_2}.
\]

\section{Механическая энергия и принцип ЛаСалля}

Рассмотрим систему «масса--пружина--демпфер»
\[
    m\ddot q+b\dot q+kq=0,
    \qquad
    m,b,k>0.
\]
Положим
\[
    x=
    \begin{pmatrix}
        q\\
        v
    \end{pmatrix},
    \qquad
    v=\dot q.
\]
Тогда
\[
    \dot x
    =
    \begin{pmatrix}
        v\\[1mm]
        -\dfrac{k}{m}q-\dfrac{b}{m}v
    \end{pmatrix},
\]
а единственным равновесием является \((q,v)=(0,0)\).

Естественный кандидат --- полная механическая энергия
\[
    V_0(q,v)
    =
    \frac12kq^2+\frac12mv^2.
\]
Она положительно определена и радиально неограничена. Её производная
равна
\[
\begin{aligned}
    \dot V_0
    &=
    kqv+mv\dot v
    \\
    &=
    kqv+mv
    \left(
        -\frac{k}{m}q-\frac{b}{m}v
    \right)
    \\
    &=
    -bv^2
    \le0.
\end{aligned}
\]
Отсюда следует устойчивость по Ляпунову, но производная лишь
отрицательно полуопределена: при \(v=0\) и \(q\ne0\) она равна нулю.

Рассмотрим множество
\[
    \mathcal Z
    =
    \{(q,v):\dot V_0(q,v)=0\}
    =
    \{(q,v):v=0\}.
\]
Чтобы траектория оставалась в \(\mathcal Z\), необходимо
\(v(t)\equiv0\) и, следовательно, \(\dot v(t)\equiv0\). Но при \(v=0\)
\[
    \dot v=-\frac{k}{m}q,
\]
поэтому инвариантность требует \(q=0\). Наибольшее инвариантное
подмножество \(\mathcal Z\) состоит только из равновесия.

\begin{theorem}[Следствие принципа ЛаСалля]
Равновесие системы «масса--пружина--демпфер» глобально
асимптотически устойчиво.
\end{theorem}

Производная функции Ляпунова может обращаться в нуль в отдельных
неравновесных точках. Это не мешает сходимости, если траектория не
может навсегда остаться в соответствующем множестве.

\section{Перекрёстный член и строгая оценка}

Физическая энергия доказывает асимптотическую устойчивость, но не
даёт непосредственно строгой квадратичной оценки производной. Для
этого добавим малый перекрёстный член:
\[
    V_\varepsilon(q,v)
    =
    \frac12kq^2
    +
    \frac12mv^2
    +
    \varepsilon m qv,
\]
где \(\varepsilon>0\). В матричной форме
\[
    V_\varepsilon(x)
    =
    \frac12x^{\mathsf T}P_\varepsilon x,
    \qquad
    P_\varepsilon
    =
    \begin{pmatrix}
        k&\varepsilon m\\
        \varepsilon m&m
    \end{pmatrix}.
\]
Матрица \(P_\varepsilon\) положительно определена, если
\[
    0<\varepsilon<\sqrt{\frac{k}{m}}.
\]

Вдоль траекторий
\[
\begin{aligned}
    \dot V_\varepsilon
    ={}&
    kqv+mv\dot v
    +
    \varepsilon m(v^2+q\dot v)
    \\
    ={}&
    -bv^2+
    \varepsilon mv^2-
    \varepsilon bqv-
    \varepsilon kq^2.
\end{aligned}
\]
Следовательно,
\[
    \dot V_\varepsilon
    =
    -x^{\mathsf T}R_\varepsilon x,
\]
где
\[
    R_\varepsilon
    =
    \begin{pmatrix}
        \varepsilon k&\dfrac{\varepsilon b}{2}\\[2mm]
        \dfrac{\varepsilon b}{2}&b-\varepsilon m
    \end{pmatrix}.
\]
Для достаточно малого \(\varepsilon\) матрица \(R_\varepsilon\)
положительно определена. Достаточно обеспечить
\[
    b-\varepsilon m>0
\]
и
\[
    \varepsilon k(b-\varepsilon m)
    -
    \frac{\varepsilon^2b^2}{4}
    >0.
\]
Тогда существуют \(c_1,c_2,c_3>0\), такие что
\[
    c_1\lVert x\rVert^2
    \le
    V_\varepsilon(x)
    \le
    c_2\lVert x\rVert^2,
    \qquad
    \dot V_\varepsilon(x)
    \le
    -c_3\lVert x\rVert^2.
\]
По квадратичному критерию равновесие глобально экспоненциально
устойчиво.

Одна и та же система допускает разные функции Ляпунова. Энергия
\(V_0\) вместе с принципом ЛаСалля даёт асимптотическую устойчивость,
а модифицированная энергия \(V_\varepsilon\) --- непосредственную
экспоненциальную оценку. Поэтому слабый результат иногда отражает не
свойство системы, а недостаточно точный выбор функции Ляпунова.

\section{Градиентный поток и три естественных кандидата}

Рассмотрим задачу
\[
    \min_{x\in\RR^n}f(x)
\]
и градиентный поток
\[
    \dot x=-\nabla f(x).
\]
Стационарная точка \(x^\star\), удовлетворяющая
\[
    \nabla f(x^\star)=0,
\]
является равновесием. Для выпуклой функции она является глобальным
минимумом, а для строго или сильно выпуклой функции минимум
единственен.

Пусть \(f\in C^2\) и
\[
    \nabla^2f(x)\succeq\mu I,
    \qquad
    \mu>0.
\]
Для функции
\[
    V_\nabla(x)
    =
    \frac12\lVert\nabla f(x)\rVert^2
\]
имеем
\[
\begin{aligned}
    \dot V_\nabla
    &=
    \nabla f(x)^{\mathsf T}\nabla^2f(x)\dot x
    \\
    &=
    -\nabla f(x)^{\mathsf T}
    \nabla^2f(x)\nabla f(x)
    \\
    &\le
    -\mu\lVert\nabla f(x)\rVert^2
    \\
    &=
    -2\mu V_\nabla.
\end{aligned}
\]
Следовательно,
\[
    \lVert\nabla f(x(t))\rVert
    \le
    e^{-\mu(t-t_0)}
    \lVert\nabla f(x(t_0))\rVert.
\]

Сильная выпуклость также даёт
\[
    \langle
        \nabla f(x)-\nabla f(x^\star),
        x-x^\star
    \rangle
    \ge
    \mu\lVert x-x^\star\rVert^2.
\]
Поскольку \(\nabla f(x^\star)=0\), получаем
\[
    \lVert\nabla f(x)\rVert
    \ge
    \mu\lVert x-x^\star\rVert.
\]
Если градиент \(L\)-липшицев, то
\[
    \lVert\nabla f(x(t_0))\rVert
    \le
    L\lVert x(t_0)-x^\star\rVert,
\]
и потому
\[
    \lVert x(t)-x^\star\rVert
    \le
    \frac{L}{\mu}
    e^{-\mu(t-t_0)}
    \lVert x(t_0)-x^\star\rVert.
\]

Более прямой выбор ---
\[
    V_x(x)
    =
    \frac12\lVert x-x^\star\rVert^2.
\]
Тогда
\[
\begin{aligned}
    \dot V_x
    &=
    -\langle x-x^\star,\nabla f(x)\rangle
    \\
    &\le
    -\mu\lVert x-x^\star\rVert^2
    \\
    &=
    -2\mu V_x.
\end{aligned}
\]
Следовательно,
\[
    \lVert x(t)-x^\star\rVert
    \le
    e^{-\mu(t-t_0)}
    \lVert x(t_0)-x^\star\rVert.
\]

Наконец, для ошибки по значению
\[
    V_f(x)=f(x)-f^\star
\]
всегда выполняется
\[
    \dot V_f
    =
    -\lVert\nabla f(x)\rVert^2.
\]
Если функция удовлетворяет неравенству Поляка--Лоясиевича
\[
    \frac12\lVert\nabla f(x)\rVert^2
    \ge
    \mu\bigl(f(x)-f^\star\bigr),
\]
то
\[
    \dot V_f\le-2\mu V_f
\]
и
\[
    f(x(t))-f^\star
    \le
    e^{-2\mu(t-t_0)}
    \bigl(f(x(t_0))-f^\star\bigr).
\]
PL-условие тем самым обеспечивает экспоненциальное убывание ошибки
по значению даже без выпуклости.

\section{Множество решений и выбор функции}

PL-неравенство не гарантирует единственности минимизатора. Например,
\[
    f(x_1,x_2)=\frac12x_1^2
\]
удовлетворяет
\[
    \frac12\lVert\nabla f(x)\rVert^2
    =
    f(x)-f^\star,
\]
но
\[
    \mathcal X^\star
    =
    \{(0,x_2):x_2\in\RR\}.
\]
В этом случае \(V_f=f-f^\star\) положительно определена относительно
множества решений:
\[
    V_f(x)=0
    \quad\Longleftrightarrow\quad
    x\in\mathcal X^\star.
\]
Экспоненциальная оценка гарантирует сходимость значения функции, но
сама по себе не выбирает единственный элемент множества
\(\mathcal X^\star\). Верное общее следствие PL-свойства состоит в том,
что каждая стационарная точка является глобальным минимумом.

Три основных кандидата отвечают разным целям. Функция
\[
    \frac12\lVert x-x^\star\rVert^2
\]
естественна при сильной выпуклости и непосредственно оценивает
расстояние до решения. Функция
\[
    f(x)-f^\star
\]
сочетается с PL-неравенством и контролирует ошибку по значению.
Функция
\[
    \frac12\lVert\nabla f(x)\rVert^2
\]
удобна при положительной определённости гессиана и даёт оценку нормы
градиента.

Для инерционных алгоритмов функция Ляпунова обычно включает ошибку по
функции, квадратичный член по скорости и перекрёстные слагаемые. Для
распределённых алгоритмов к ним добавляется ошибка согласования,
например
\[
    \sum_{i=1}^{N}\lVert x_i-\bar x\rVert^2.
\]
Прямой метод не требует явного решения системы: достаточно найти
равновесие, построить положительно определённую функцию, вычислить её
производную и оценить скорость убывания. Основная трудность состоит
в выборе \(V\). Обратные теоремы Ляпунова часто гарантируют
существование подходящей функции, но не дают автоматически её простой
аналитической формы.

\lecture{16}{\texorpdfstring{Экспоненциальная\\ устойчивость}{Экспоненциальная устойчивость}}

\sourcevideo
    {https://www.youtube.com/playlist?list=PLOzRYVm0a65fhKDS8cYIC7YCuVGJ4FOJx}
    {Week 5: Lecture 16: Exponential Stability}

Неравенство
\[
    \dot V\le-\gamma V
\]
гарантирует экспоненциальное убывание функции Ляпунова, но для
экспоненциальной устойчивости состояния нужно дополнительно сравнить
\(V(x)\) с расстоянием до равновесия. В этой лекции такой переход
выполняется для градиентного потока сильно выпуклой функции.

\section{Квадратичный критерий}

Рассмотрим автономную систему
\[
    \dot x=g(x)
\]
с равновесием \(x^\star\), то есть \(g(x^\star)=0\), и обозначим
\[
    e=x-x^\star.
\]

\begin{theorem}[Критерий экспоненциальной устойчивости]
Пусть существует непрерывно дифференцируемая функция
\(V\colon\RR^n\to\RR\) и положительные константы
\(c_1,c_2,c_3\), для которых
\[
    c_1\lVert e\rVert^2
    \le
    V(x)
    \le
    c_2\lVert e\rVert^2
\]
и
\[
    \dot V(x)
    \le
    -c_3\lVert e\rVert^2.
\]
Тогда
\[
    \lVert x(t)-x^\star\rVert
    \le
    \sqrt{\frac{c_2}{c_1}}
    \exp\left(
        -\frac{c_3}{2c_2}(t-t_0)
    \right)
    \lVert x(t_0)-x^\star\rVert.
\]
\end{theorem}

\begin{proof}
Из верхней оценки на \(V\) следует
\[
    \lVert e\rVert^2\ge\frac1{c_2}V(x),
\]
поэтому
\[
    \dot V\le-\frac{c_3}{c_2}V.
\]
Интегрирование даёт
\[
    V(x(t))
    \le
    \exp\left(
        -\frac{c_3}{c_2}(t-t_0)
    \right)
    V(x(t_0)).
\]
Используя обе квадратичные оценки, получаем
\[
\begin{aligned}
    c_1\lVert e(t)\rVert^2
    &\le
    V(x(t))
    \\
    &\le
    c_2
    \exp\left(
        -\frac{c_3}{c_2}(t-t_0)
    \right)
    \lVert e(t_0)\rVert^2.
\end{aligned}
\]
После извлечения квадратного корня получается требуемая оценка.
\end{proof}

Таким образом, в стандартной записи
\[
    \lVert e(t)\rVert
    \le
    M e^{-\alpha(t-t_0)}\lVert e(t_0)\rVert
\]
можно взять
\[
    M=\sqrt{\frac{c_2}{c_1}},
    \qquad
    \alpha=\frac{c_3}{2c_2}.
\]
Знак минус в условии \(\dot V\le-\gamma V\) принципиален: без него
функция Ляпунова может возрастать.

\section{Сильно выпуклый градиентный поток}

Рассмотрим задачу
\[
    \min_{x\in\RR^n}f(x)
\]
и градиентный поток
\[
    \dot x=-\nabla f(x).
\]
Пусть функция \(f\) является \(\mu\)-сильно выпуклой и
\(L\)-гладкой, где
\[
    0<\mu\le L.
\]
Тогда минимум \(x^\star\) единственен и
\[
    \nabla f(x^\star)=0.
\]

Из сильной выпуклости при выборе базовой точки \(x^\star\) следует
\[
    f(x)-f^\star
    \ge
    \frac\mu2\lVert x-x^\star\rVert^2.
\]
Из леммы об убывании и \(L\)-гладкости получаем
\[
    f(x)-f^\star
    \le
    \frac L2\lVert x-x^\star\rVert^2.
\]
Следовательно,
\[
    \frac\mu2\lVert x-x^\star\rVert^2
    \le
    f(x)-f^\star
    \le
    \frac L2\lVert x-x^\star\rVert^2.
\]
Именно эта двусторонняя оценка позволяет использовать ошибку по
значению функции как квадратично сравнимую функцию Ляпунова.

Сильная выпуклость также влечёт PL-неравенство
\[
    \frac12\lVert\nabla f(x)\rVert^2
    \ge
    \mu\bigl(f(x)-f^\star\bigr).
\]
В дальнейшем сильная выпуклость используется для нижней квадратичной
оценки и PL-неравенства, а гладкость — для верхней квадратичной
оценки.

\section{Ошибка по значению функции}

Положим
\[
    V_f(x)=f(x)-f^\star.
\]
При сильной выпуклости эта функция положительно определена
относительно \(x^\star\), а вдоль градиентного потока
\[
\begin{aligned}
    \dot V_f(x)
    &=
    \nabla f(x)^{\mathsf T}\dot x
    \\
    &=
    -\lVert\nabla f(x)\rVert^2.
\end{aligned}
\]
По PL-неравенству
\[
    \dot V_f\le-2\mu V_f,
\]
откуда
\[
    f(x(t))-f^\star
    \le
    e^{-2\mu(t-t_0)}
    \bigl(f(x(t_0))-f^\star\bigr).
\]
Значение целевой функции сходится к минимуму экспоненциально.

Чтобы получить оценку состояния по общему квадратичному критерию,
используем
\[
    \dot V_f
    \le
    -2\mu V_f
    \le
    -\mu^2\lVert x-x^\star\rVert^2.
\]
Можно выбрать
\[
    c_1=\frac\mu2,
    \qquad
    c_2=\frac L2,
    \qquad
    c_3=\mu^2.
\]
Следовательно,
\[
    \lVert x(t)-x^\star\rVert
    \le
    \sqrt{\frac L\mu}
    \exp\left(
        -\frac{\mu^2}{L}(t-t_0)
    \right)
    \lVert x(t_0)-x^\star\rVert.
\]
В терминах числа обусловленности
\[
    \kappa=\frac L\mu
\]
эта оценка принимает вид
\[
    \lVert x(t)-x^\star\rVert
    \le
    \sqrt\kappa
    \exp\left(
        -\frac\mu\kappa(t-t_0)
    \right)
    \lVert x(t_0)-x^\star\rVert.
\]
Чем больше \(\kappa\), тем слабее полученная граница. Это свойство
данного доказательства через \(V_f\), а не точная характеристика
самого градиентного потока.

\section{Расстояние и норма градиента}

Более прямой выбор функции Ляпунова равен
\[
    V_x(x)
    =
    \frac12\lVert x-x^\star\rVert^2.
\]
Тогда
\[
    \dot V_x
    =
    -\bigl\langle
        x-x^\star,
        \nabla f(x)
    \bigr\rangle.
\]
Сильная монотонность градиента даёт
\[
    \bigl\langle
        \nabla f(x)-\nabla f(x^\star),
        x-x^\star
    \bigr\rangle
    \ge
    \mu\lVert x-x^\star\rVert^2.
\]
Поскольку \(\nabla f(x^\star)=0\), получаем
\[
    \dot V_x
    \le
    -\mu\lVert x-x^\star\rVert^2
    =
    -2\mu V_x.
\]
Следовательно,
\[
    \lVert x(t)-x^\star\rVert
    \le
    e^{-\mu(t-t_0)}
    \lVert x(t_0)-x^\star\rVert.
\]
Эта оценка сильнее лекционного вывода через \(f-f^\star\) и не
требует \(L\)-гладкости для самой оценки, если траектория существует.
Гладкость нужна только в предыдущем доказательстве, где используется
верхняя квадратичная оценка на \(f-f^\star\).

Если \(f\in C^2\), можно также взять
\[
    V_{\nabla}(x)
    =
    \frac12\lVert\nabla f(x)\rVert^2.
\]
Тогда
\[
\begin{aligned}
    \dot V_{\nabla}
    &=
    \nabla f(x)^{\mathsf T}
    \nabla^2f(x)\dot x
    \\
    &=
    -\nabla f(x)^{\mathsf T}
    \nabla^2f(x)\nabla f(x)
    \\
    &\le
    -\mu\lVert\nabla f(x)\rVert^2
    =
    -2\mu V_{\nabla}.
\end{aligned}
\]
Поэтому
\[
    \lVert\nabla f(x(t))\rVert
    \le
    e^{-\mu(t-t_0)}
    \lVert\nabla f(x(t_0))\rVert.
\]
При \(L\)-гладкости и сильной выпуклости
\[
    \mu\lVert x-x^\star\rVert
    \le
    \lVert\nabla f(x)\rVert
    \le
    L\lVert x-x^\star\rVert,
\]
откуда следует ещё одна оценка состояния:
\[
    \lVert x(t)-x^\star\rVert
    \le
    \frac L\mu
    e^{-\mu(t-t_0)}
    \lVert x(t_0)-x^\star\rVert.
\]

Функции \(V_f\), \(V_x\) и \(V_{\nabla}\) контролируют разные
характеристики: ошибку по значению, расстояние до решения и норму
градиента. Выбор зависит от того, какую величину требуется оценить и
какие предположения доступны.

\section{PL-условие без сильной выпуклости}

Пусть функция удовлетворяет PL-неравенству
\[
    \frac12\lVert\nabla f(x)\rVert^2
    \ge
    \mu\bigl(f(x)-f^\star\bigr),
\]
но не обязательно является сильно выпуклой. Для
\[
    V_f(x)=f(x)-f^\star
\]
по-прежнему выполняется
\[
    \dot V_f
    =
    -\lVert\nabla f(x)\rVert^2
    \le
    -2\mu V_f,
\]
и потому
\[
    f(x(t))-f^\star
    \le
    e^{-2\mu(t-t_0)}
    \bigl(f(x(t_0))-f^\star\bigr).
\]

Однако PL-свойство само по себе не даёт двусторонней связи между
ошибкой по значению и расстоянием до единственной точки. Множество
глобальных минимумов может содержать несколько элементов, поэтому в
общем случае естественно исследовать устойчивость множества
\[
    \mathcal X^\star
    =
    \operatorname*{argmin}f.
\]
Для перехода от ошибки по функции к расстоянию
\[
    \operatorname{dist}(x,\mathcal X^\star)
\]
нужно дополнительное условие роста или регулярности.

Из PL-неравенства следует, что всякая стационарная точка является глобальным минимумом.

\section{Связь с дискретным методом}

Градиентный поток является непрерывным аналогом итерации
\[
    x^{k+1}
    =
    x^k-\eta\nabla f(x^k).
\]
Для \(\mu\)-сильно выпуклой \(L\)-гладкой функции и
\[
    \eta=\frac1L
\]
дискретный метод удовлетворяет оценке
\[
    f(x^k)-f^\star
    \le
    \left(1-\frac\mu L\right)^k
    \bigl(f(x^0)-f^\star\bigr).
\]
В непрерывном времени соответствующая функциональная оценка имеет
вид
\[
    f(x(t))-f^\star
    \le
    e^{-2\mu t}
    \bigl(f(x(0))-f^\star\bigr).
\]
Эти скорости нельзя сравнивать только по коэффициентам: один
дискретный шаг имеет вычислительную стоимость, а масштаб времени
непрерывной системы изменяется при умножении правой части на
положительную константу.

Практическое доказательство экспоненциальной устойчивости состоит из
четырёх действий: определить равновесие, выбрать функцию Ляпунова,
связать её с нужной мерой ошибки и оценить производную. Неравенство
\[
    \dot V\le-\gamma V
\]
непосредственно даёт экспоненциальное убывание выбранной величины.
Чтобы получить экспоненциальную устойчивость состояния, дополнительно
нужна количественная связь между \(V\) и расстоянием до равновесия.

\lecture{17}{Ньютоновский поток и дивергенция Брэгмана}

\sourcevideo
    {https://www.youtube.com/playlist?list=PLOzRYVm0a65fhKDS8cYIC7YCuVGJ4FOJx}
    {Week 5: Lecture 17: Bregman Divergence}

Обычный градиентный поток использует только направление
\(-\nabla f(x)\) и потому чувствителен к различию кривизны по
координатам. Ньютоновский поток компенсирует эту анизотропию обратным
Гессианом, а дивергенция Брэгмана заменяет евклидово расстояние мерой
расхождения, согласованной с геометрией задачи.

\section{Строгая выпуклость и скорость градиентного потока}

Пусть \(f\in C^2(\RR^n)\) и
\[
    \nabla^2f(x)\succ0
\]
во всей выпуклой области.

\begin{theorem}
Если Гессиан положительно определён в каждой точке выпуклой области,
то функция \(f\) строго выпукла на этой области.
\end{theorem}

Обратное утверждение неверно. Функция
\[
    f(x)=x^4
\]
строго выпукла, но
\[
    f''(0)=0.
\]
Строгая выпуклость гарантирует единственность минимума, если он
существует, но не гарантирует его существование. Например,
\(f(x)=e^x\) строго выпукла на \(\RR\), однако не достигает своего
инфимума.

Сильная выпуклость требует единой оценки
\[
    \nabla^2f(x)\succeq\mu I,
    \qquad
    \mu>0.
\]
При одной строгой выпуклости такая константа может отсутствовать.
Для \(f(x)=x^4\) градиентный поток
\[
    \dot x=-4x^3
\]
имеет при \(x(0)=x_0\ne0\) решение
\[
    x(t)
    =
    \frac{x_0}{\sqrt{1+8x_0^2t}},
\]
поэтому
\[
    |x(t)|=O(t^{-1/2}).
\]
Итак, строгая выпуклость обеспечивает единственность решения, но не
равномерную экспоненциальную скорость обычного градиентного потока.

\section{Ньютоновский поток}

Предположим, что
\[
    \nabla^2f(x)\succ0
\]
для всех рассматриваемых \(x\). Тогда можно определить динамику
\[
    \dot x
    =
    -\bigl(\nabla^2f(x)\bigr)^{-1}\nabla f(x).
\]
Она называется непрерывным методом Ньютона, или ньютоновским
потоком. В дискретном времени той же идее соответствует шаг
\[
    x^{k+1}
    =
    x^k-
    \bigl(\nabla^2f(x^k)\bigr)^{-1}\nabla f(x^k).
\]

Рассмотрим функцию Ляпунова
\[
    V(x)=\frac12\lVert\nabla f(x)\rVert^2.
\]
Если \(x^\star\) — единственная стационарная точка, то \(V\)
положительно определена относительно \(x^\star\). Вдоль
ньютоновского потока
\[
\begin{aligned}
    \dot V
    &=
    \nabla f(x)^{\mathsf T}\nabla^2f(x)\dot x
    \\
    &=
    -\nabla f(x)^{\mathsf T}\nabla^2f(x)
    \bigl(\nabla^2f(x)\bigr)^{-1}\nabla f(x)
    \\
    &=
    -\lVert\nabla f(x)\rVert^2
    =-2V.
\end{aligned}
\]
Следовательно,
\[
    V(x(t))
    =
    e^{-2(t-t_0)}V(x(t_0))
\]
и
\[
    \lVert\nabla f(x(t))\rVert
    =
    e^{-(t-t_0)}
    \lVert\nabla f(x(t_0))\rVert.
\]

Тот же результат получается непосредственно. Для
\[
    g(t)=\nabla f(x(t))
\]
имеем
\[
    \dot g
    =
    \nabla^2f(x)\dot x
    =-g,
\]
то есть нелинейная динамика состояния превращается в линейную
динамику градиента.

Экспоненциальное убывание градиента не автоматически даёт такую же
оценку расстояния до минимума. Для перехода нужна связь вида
\[
    \lVert\nabla f(x)\rVert
    \ge
    \mu\lVert x-x^\star\rVert.
\]
Она выполняется при \(\mu\)-сильной выпуклости, но может отсутствовать
при одной строгой выпуклости. Кроме того, для ньютоновского потока
нужно отдельно обеспечить существование траектории и невырожденность
Гессиана вдоль неё.

Если \(f\) только выпукла и
\[
    \nabla^2f(x)\succeq0,
\]
обратный Гессиан может не существовать. Для обычного градиентного
потока
\[
    \dot x=-\nabla f(x)
\]
та же функция даёт
\[
    \dot V
    =
    -\nabla f(x)^{\mathsf T}\nabla^2f(x)\nabla f(x)
    \le0.
\]
Отсюда следует невозрастание нормы градиента, но не скорость
сходимости и не обязательно сходимость к заранее выбранному
минимизатору. При нескольких решениях \(V\) положительно определена
относительно всего множества стационарных точек.

\section{Компенсация плохой обусловленности}

Рассмотрим квадратичную функцию
\[
    f(x,y)
    =
    \frac12x^2+
    \frac{10^{-4}}2y^2.
\]
Её градиент и Гессиан равны
\[
    \nabla f(x,y)
    =
    \begin{pmatrix}
        x\\
        10^{-4}y
    \end{pmatrix},
    \qquad
    H
    =
    \begin{pmatrix}
        1&0\\
        0&10^{-4}
    \end{pmatrix}.
\]
Число обусловленности составляет
\[
    \kappa(H)=10^4.
\]
Обычный градиентный поток имеет вид
\[
    \dot x=-x,
    \qquad
    \dot y=-10^{-4}y,
\]
поэтому
\[
    x(t)=e^{-t}x(0),
    \qquad
    y(t)=e^{-10^{-4}t}y(0).
\]
Медленная координата определяет общую скорость.

Поскольку
\[
    H^{-1}
    =
    \begin{pmatrix}
        1&0\\
        0&10^4
    \end{pmatrix},
\]
ньютоновский поток становится
\[
    \begin{pmatrix}
        \dot x\\
        \dot y
    \end{pmatrix}
    =
    -H^{-1}\nabla f(x,y)
    =
    -
    \begin{pmatrix}
        x\\
        y
    \end{pmatrix}.
\]
Обе координаты затухают как \(e^{-t}\). Обратный Гессиан выравнивает
различную кривизну направлений и преобразует вытянутую квадратичную
геометрию в изотропную.

Нормированный поток
\[
    \dot x
    =
    -\frac{\nabla f(x)}
    {\lVert\nabla f(x)\rVert+\varepsilon},
    \qquad
    \varepsilon>0,
\]
изменяет только длину градиентного вектора и сохраняет его
направление. Обратный Гессиан, напротив, масштабирует координаты
по-разному и в общем случае меняет также направление. Поэтому
скалярная нормировка не устраняет анизотропию так, как это делает
матричное преобразование второго порядка.

Основные евклидовы функции Ляпунова в оптимизации равны
\[
    f(x)-f^\star,
    \qquad
    \frac12\lVert x-x^\star\rVert^2,
    \qquad
    \frac12\lVert\nabla f(x)\rVert^2.
\]
Первая обращается в нуль на множестве глобальных минимумов, вторая
— только в выбранной точке \(x^\star\), третья — на множестве
стационарных точек. Для невыпуклой функции стационарность не
равносильна глобальной оптимальности.

\section{Дивергенция Брэгмана}

Евклидова норма не всегда соответствует структуре задачи. Пусть
\(h\colon\mathcal D\to\RR\) дифференцируема на выпуклом множестве
\(\mathcal D\).

\begin{definition}[Дивергенция Брэгмана]
Дивергенцией Брэгмана, порождённой функцией \(h\), называется
\[
    D_h(p,q)
    =
    h(p)-h(q)
    -\langle\nabla h(q),p-q\rangle.
\]
\end{definition}

Величина
\[
    h(q)+\langle\nabla h(q),p-q\rangle
\]
является аффинной касательной моделью \(h\) в точке \(q\),
вычисленной в \(p\). Поэтому \(D_h(p,q)\) измеряет ошибку этой
линейной аппроксимации.

Из выпуклости первого порядка следует
\[
    D_h(p,q)\ge0.
\]
Если \(h\) строго выпукла, то
\[
    D_h(p,q)=0
    \quad\Longleftrightarrow\quad
    p=q.
\]
При одной выпуклости нулевое значение возможно и для различных точек,
например на аффинном участке функции.

Если \(h\) \(\mu\)-сильно выпукла, то
\[
    D_h(p,q)
    \ge
    \frac\mu2\lVert p-q\rVert^2.
\]
Если она дополнительно \(L\)-гладка, то
\[
    D_h(p,q)
    \le
    \frac L2\lVert p-q\rVert^2.
\]
Таким образом,
\[
    \frac\mu2\lVert p-q\rVert^2
    \le
    D_h(p,q)
    \le
    \frac L2\lVert p-q\rVert^2.
\]
Сильная выпуклость даёт нижнюю оценку, а не точное равенство.
Равенство возникает лишь для специальных квадратичных функций.

При
\[
    h(x)=\frac12\lVert x\rVert^2
\]
получаем
\[
    D_h(p,q)=\frac12\lVert p-q\rVert^2.
\]
Для
\[
    h(x)=\frac12x^{\mathsf T}Px,
    \qquad
    P\succ0,
\]
имеем
\[
    D_h(p,q)
    =
    \frac12(p-q)^{\mathsf T}P(p-q).
\]
Матрица \(P\) задаёт взвешенную евклидову геометрию.

В общем случае дивергенция Брэгмана не является метрикой:
\[
    D_h(p,q)\ne D_h(q,p)
\]
и не обязана удовлетворять неравенству треугольника. Полезным
заменителем евклидовых тождеств служит трёхточечная формула
\[
\begin{aligned}
    D_h(p,q)+D_h(q,r)-D_h(p,r)
    =
    \langle\nabla h(r)-\nabla h(q),p-q\rangle.
\end{aligned}
\]
Она широко используется в анализе зеркальных методов.

\section{Брэгмановская геометрия в алгоритмах}

Если \(x^\star\) — оптимальное решение, то естественными
кандидатами на функцию Ляпунова являются
\[
    D_h(x^\star,x)
    \qquad\text{и}\qquad
    D_h(x,x^\star).
\]
Порядок аргументов существенен из-за несимметричности. При строгой
выпуклости \(h\) обе величины неотрицательны и обращаются в нуль
только при \(x=x^\star\); при сильной выпуклости они дополнительно
контролируют квадрат евклидова расстояния. Конкретная производная
вдоль траектории зависит от выбранного потока, поэтому функцию \(h\)
и динамику нужно проектировать совместно.

Евклидов градиентный шаг можно записать как
\[
    x^{k+1}
    =
    \operatorname*{argmin}_{x}
    \left\{
        \langle\nabla f(x^k),x-x^k\rangle
        +
        \frac1{2\eta}\lVert x-x^k\rVert^2
    \right\}.
\]
Замена квадрата нормы дивергенцией Брэгмана даёт зеркальный шаг
\[
    x^{k+1}
    =
    \operatorname*{argmin}_{x}
    \left\{
        \langle\nabla f(x^k),x-x^k\rangle
        +
        \frac1\eta D_h(x,x^k)
    \right\}.
\]
Функция \(h\) определяет геометрию регуляризующего члена и может быть
согласована с допустимым множеством и масштабом координат. Зеркальный
спуск подробно в лекции не выводится; эта запись поясняет практическую
роль дивергенции.

Квадратичная \(h\) приводит к взвешенной евклидовой геометрии, а
неквадратичная позволяет строить существенно неевклидовы шаги.
Вариационный взгляд на ускоренные методы использует подобные
дивергенции для построения непрерывных динамик и функций энергии. 

\section{Вычислительная стоимость и границы метода}

Обычный градиентный поток требует вычисления \(\nabla f(x)\).
Ньютоновский поток дополнительно требует Гессиан и решение системы
\[
    \nabla^2f(x)d=-\nabla f(x).
\]
Явно строить обратную матрицу не следует. Для плотного Гессиана
размерности \(n\) факторизация обычно требует порядка
\[
    O(n^3)
\]
операций и
\[
    O(n^2)
\]
памяти. Поэтому методы второго порядка могут использовать меньше
итераций, но каждая итерация существенно дороже.

Брэгмановские методы обычно не требуют Гессиана целевой функции, но
требуют решения вспомогательной задачи, определяемой \(h\). Их
преимущество проявляется, когда евклидова геометрия плохо отражает
структуру пространства или ограничений.

Ньютоновский поток особенно полезен при обратимом Гессиане, сильной
анизотропии и доступном решении линейных систем. Нормированный
градиент управляет общей длиной шага, но не выравнивает координатную
кривизну. Дивергенция Брэгмана позволяет согласовать функцию
Ляпунова и регуляризующий член с неевклидовой геометрией задачи.

\lecture{18}{Масштабированный градиентный поток}

\sourcevideo
    {https://www.youtube.com/playlist?list=PLOzRYVm0a65fhKDS8cYIC7YCuVGJ4FOJx}
    {Week 5: Lecture 18: Rescaled Gradient Flow}

Обычный градиентный поток замедляется около стационарной точки,
поскольку норма градиента стремится к нулю. Положительное скалярное
масштабирование сохраняет геометрические траектории непрерывной
системы, но изменяет скорость их прохождения. Подходящий степенной
множитель позволяет получить точное достижение оптимума сначала за
конечное, а затем за фиксированное время.

\section{Масштабирование градиентного потока}

Рассмотрим задачу
\[
    \min_{x\in\RR^n}f(x)
\]
и обычный градиентный поток
\[
    \dot x=-\nabla f(x).
\]
Его равновесия совпадают со стационарными точками функции. Если
\(f\) сильно выпукла, стационарная точка единственна и равна
глобальному минимуму \(x^\star\). Скорость обычного потока равна
\[
    \lVert\dot x\rVert=\lVert\nabla f(x)\rVert,
\]
поэтому движение замедляется при приближении к решению.

Пусть \(p>2\) и
\[
    \rho_p=\frac{p-2}{p-1}\in(0,1).
\]
Масштабированный поток определяется как
\[
    \dot x
    =
    -\frac{\nabla f(x)}
    {\lVert\nabla f(x)\rVert^{\rho_p}}
\]
при \(\nabla f(x)\ne0\), а в стационарных точках полагается
\(\dot x=0\). Его скорость равна
\[
    \lVert\dot x\rVert
    =
    \lVert\nabla f(x)\rVert^{1-\rho_p}
    =
    \lVert\nabla f(x)\rVert^{1/(p-1)}.
\]
Поскольку \(1/(p-1)<1\), скорость около минимума убывает медленнее,
чем у обычного градиентного потока.

Вне стационарных точек масштабированное поле имеет вид
\[
    \widetilde F(x)=a(x)F(x),
    \qquad
    a(x)=\lVert\nabla f(x)\rVert^{-\rho_p}>0,
\]
где \(F(x)=-\nabla f(x)\). Поэтому оба поля задают одинаковые
непараметризованные фазовые траектории и различаются только
перепараметризацией времени. После дискретизации это свойство в общем
случае теряется.

Величина правой части около равновесия имеет порядок
\[
    \lVert\nabla f(x)\rVert^{1/(p-1)}.
\]
Показатель меньше единицы, поэтому поле обычно гёльдерово, но не
локально липшицево в стационарной точке. Именно такая потеря
липшицевости допускает точное достижение равновесия за конечное время.

\section{Конечно-временной критерий Ляпунова}

\begin{definition}[Конечно-временная устойчивость]
Равновесие \(x^\star\) называется конечно-временно устойчивым, если
оно устойчиво по Ляпунову и для каждого допустимого начального
состояния \(x_0\) существует число \(T(x_0)<\infty\), такое что
\[
    x(t)=x^\star
\]
для всех \(t\ge T(x_0)\). Число \(T(x_0)\) называется временем
установления.
\end{definition}

\begin{theorem}[Конечно-временной критерий]
Пусть \(V\) положительно определена относительно \(x^\star\) и вдоль
траекторий выполняется
\[
    \dot V\le-cV^\alpha,
    \qquad
    c>0,
    \qquad
    0<\alpha<1.
\]
Тогда равновесие достигается за конечное время, причём
\[
    T(x_0)
    \le
    \frac{V(x_0)^{1-\alpha}}
    {c(1-\alpha)}.
\]
\end{theorem}

\begin{proof}
Пока \(V(t)>0\), имеем
\[
    V(t)^{-\alpha}\dot V(t)\le-c.
\]
Следовательно,
\[
    \frac{\mathrm d}{\mathrm dt}V(t)^{1-\alpha}
    \le
    -c(1-\alpha).
\]
После интегрирования получаем
\[
    V(t)^{1-\alpha}
    \le
    V(x_0)^{1-\alpha}-c(1-\alpha)t.
\]
Правая часть обращается в нуль не позднее указанного момента. Так как
\(V\ge0\), после этого должно выполняться \(V(t)=0\), а положительная
определённость даёт \(x(t)=x^\star\).
\end{proof}

Отличие от экспоненциального условия
\[
    \dot V\le-cV
\]
состоит в показателе \(\alpha<1\). При \(0<V<1\) выполняется
\(V^\alpha>V\), поэтому убывание около равновесия достаточно быстро
для точного достижения нуля.

\section{Конечное время масштабированного потока}

Предположим, что \(f\in C^2(\RR^n)\) и
\[
    \nabla^2f(x)\succeq\mu I,
    \qquad
    \mu>0.
\]
Тогда функция сильно выпукла, минимизатор \(x^\star\) единственен и
\[
    \nabla f(x)=0
    \quad\Longleftrightarrow\quad
    x=x^\star.
\]
Выберем функцию Ляпунова
\[
    V(x)=\frac12\lVert\nabla f(x)\rVert^2.
\]
Она положительно определена относительно \(x^\star\). Вдоль
масштабированного потока
\[
\begin{aligned}
    \dot V
    &=
    -\frac{
        \nabla f(x)^{\mathsf T}
        \nabla^2f(x)
        \nabla f(x)
    }{
        \lVert\nabla f(x)\rVert^{\rho_p}
    }
    \\
    &\le
    -\mu
    \lVert\nabla f(x)\rVert^{2-\rho_p}.
\end{aligned}
\]
Поскольку
\[
    2-\rho_p=\frac{p}{p-1},
\]
и \(\lVert\nabla f(x)\rVert^2=2V(x)\), получаем
\[
    \dot V
    \le
    -\mu2^{\alpha_p}V^{\alpha_p},
    \qquad
    \alpha_p=\frac{p}{2(p-1)}.
\]
При \(p>2\)
\[
    \frac12<\alpha_p<1,
\]
поэтому применим конечно-временной критерий.

\begin{theorem}[Конечно-временная сходимость]
Пусть \(f\in C^2(\RR^n)\) является \(\mu\)-сильно выпуклой и
\(p>2\). Тогда равновесие \(x^\star\) масштабированного потока
конечно-временно устойчиво, причём
\[
    T(x_0)
    \le
    \frac{p-1}{\mu(p-2)}
    \lVert\nabla f(x_0)\rVert^{\frac{p-2}{p-1}}.
\]
\end{theorem}

Время конечно, но зависит от начального состояния. В конечно-временном
случае после момента \(T\) выполняется \(V(t)=0\). При сильной
выпуклости это непосредственно означает
\[
    \nabla f(x(t))=0
    \quad\Longleftrightarrow\quad
    x(t)=x^\star,
\]
поэтому дополнительная двусторонняя оценка между \(V\) и расстоянием
до решения не требуется.

В лекции также упоминаются полиномиальные оценки вида
\[
    f(x(t))-f^\star=O(t^{-p}).
\]
Их нельзя напрямую сравнивать с полученным конечно-временным
результатом: такие оценки могут относиться к более широкому классу
выпуклых функций, тогда как приведённый вывод использует равномерную
нижнюю границу Гессиана.

\section{Фиксированное время и два режима}

\begin{definition}[Фиксированно-временная устойчивость]
Равновесие называется фиксированно-временно устойчивым, если оно
конечно-временно устойчиво и существует число \(T_{\max}<\infty\), не
зависящее от начального состояния, для которого
\[
    T(x_0)\le T_{\max}
\]
для всех допустимых \(x_0\).
\end{definition}

\begin{theorem}[Фиксированно-временной критерий]
Пусть \(V\) положительно определена относительно \(x^\star\) и
\[
    \dot V
    \le
    -c_1V^{\alpha_1}
    -c_2V^{\alpha_2},
\]
где
\[
    c_1,c_2>0,
    \qquad
    0<\alpha_1<1<\alpha_2.
\]
Тогда
\[
    T(x_0)
    \le
    T_{\max}
    =
    \frac1{c_1(1-\alpha_1)}
    +
    \frac1{c_2(\alpha_2-1)}.
\]
\end{theorem}

Доказательство разбивает движение на два этапа. При \(V\ge1\)
оставляем только второе слагаемое:
\[
    \dot V\le-c_2V^{\alpha_2}.
\]
Интегрирование от \(V(x_0)\) до единицы даёт время не больше
\[
    \frac1{c_2(\alpha_2-1)}.
\]
После входа в область \(V\le1\) используется
\[
    \dot V\le-c_1V^{\alpha_1},
\]
и дополнительное время до нуля не превосходит
\[
    \frac1{c_1(1-\alpha_1)}.
\]
Сумма не содержит начального значения. Степень выше единицы
контролирует движение вдали от решения, а степень ниже единицы
обеспечивает точное достижение около него.

\section{Двухрежимный масштабированный поток}

Выберем
\[
    p>2,
    \qquad
    1<q<2,
\]
и положим
\[
    \rho_p=\frac{p-2}{p-1},
    \qquad
    \rho_q=\frac{q-2}{q-1}.
\]
Тогда \(0<\rho_p<1\), а \(\rho_q<0\). Рассмотрим систему
\[
\begin{aligned}
    \dot x
    ={}&
    -\frac{\nabla f(x)}
    {\lVert\nabla f(x)\rVert^{\rho_p}}
    \\
    &-
    \frac{\nabla f(x)}
    {\lVert\nabla f(x)\rVert^{\rho_q}}
\end{aligned}
\]
при \(\nabla f(x)\ne0\), а в стационарной точке положим
\(\dot x=0\).

Первое слагаемое имеет величину
\[
    \lVert\nabla f(x)\rVert^{1/(p-1)}
\]
с показателем меньше единицы и действует около решения. Второе имеет
величину
\[
    \lVert\nabla f(x)\rVert^{1/(q-1)}
\]
с показателем больше единицы и доминирует вдали от решения.

Для
\[
    V(x)=\frac12\lVert\nabla f(x)\rVert^2
\]
и \(\nabla^2f(x)\succeq\mu I\) получаем
\[
    \dot V
    \le
    -\mu2^{\alpha_1}V^{\alpha_1}
    -\mu2^{\alpha_2}V^{\alpha_2},
\]
где
\[
    \alpha_1=\frac{p}{2(p-1)}<1,
    \qquad
    \alpha_2=\frac{q}{2(q-1)}>1.
\]

\begin{theorem}[Фиксированное время]
Пусть \(f\in C^2(\RR^n)\) является \(\mu\)-сильно выпуклой,
\(p>2\) и \(1<q<2\). Тогда оптимальное равновесие двухрежимного
потока фиксированно-временно устойчиво и
\[
\begin{aligned}
    T(x_0)
    \le T_{\max}
    ={}&
    \frac1{
        \mu2^{\alpha_1}(1-\alpha_1)
    }
    \\
    &+
    \frac1{
        \mu2^{\alpha_2}(\alpha_2-1)
    }.
\end{aligned}
\]
Правая часть не зависит от \(x_0\).
\end{theorem}

Эта конструкция иллюстрирует проектирование динамики по требуемому
неравенству Ляпунова. Сначала задаются желаемые степени
\(V^{\alpha_1}\) и \(V^{\alpha_2}\), затем выбирается функция
Ляпунова, после чего правая часть системы строится так, чтобы её
производная содержала эти степени.

\section{Регуляризация, дискретизация и стоимость}

В численной реализации знаменатель часто заменяют на
\[
    \lVert\nabla f(x)\rVert+\varepsilon,
    \qquad
    \varepsilon>0.
\]
Например,
\[
    \dot x
    =
    -\frac{\nabla f(x)}
    {\bigl(\lVert\nabla f(x)\rVert+\varepsilon\bigr)^{\rho_p}}.
\]
Это устраняет деление на нуль, но около оптимума поток ведёт себя как
обычный градиентный поток с постоянным масштабом
\(\varepsilon^{-\rho_p}\). Поэтому при фиксированном
\(\varepsilon>0\) точная конечно-временная гарантия обычно теряется;
остаётся асимптотическая сходимость или быстрое попадание в малую
окрестность.

Явная дискретизация одночленного потока имеет вид
\[
    x^{k+1}
    =
    x^k
    -\eta
    \frac{\nabla f(x^k)}
    {\lVert\nabla f(x^k)\rVert^{\rho_p}}.
\]
Для двухрежимной системы к ней добавляется аналогичный член с
\(\rho_q\). Непрерывная конечно- или фиксированно-временная
устойчивость не переносится автоматически на эту рекурсию. Около
решения возможны перескакивание через минимум, колебания и высокая
чувствительность к округлению. Поэтому шаг, регуляризация и критерий
остановки требуют отдельного анализа.

По сравнению с обычным градиентным методом дополнительно вычисляются
норма градиента и одна или две скалярные степени. Для вектора
размерности \(n\) вычисление нормы и масштабирование требуют
\(O(n)\) операций, тогда как основная стоимость обычно связана с
самим градиентом. В распределённой системе глобальная норма может
потребовать отдельного агрегирования.

Теоремы предполагают точный непрерывный поток, точный градиент,
\(f\in C^2\) и равномерную сильную выпуклость. Стохастический шум,
задержки, коммуникационные ошибки и конечный шаг могут уничтожить
точное достижение решения. Эмпирические сравнения с Adam и RMSProp зависят от задачи и параметров и сами по себе не устанавливают универсального превосходства одного метода.

\lecture{19}{\texorpdfstring{PL-неравенство\\ и конечное время}{PL-неравенство и конечное время}}

\sourcevideo
    {https://www.youtube.com/playlist?list=PLOzRYVm0a65fhKDS8cYIC7YCuVGJ4FOJx}
    {Week 6: Lecture 19: Advanced Results on PL Inequality: Part 1}

Масштабирование положительным скаляром сохраняет геометрические
траектории непрерывной системы и изменяет только скорость движения по
ним. Для ньютоновского поля это позволяет получить фиксированное время
сходимости без равномерной константы сильной выпуклости. Во второй
части лекции доказывается геометрическое следствие неравенства
Поляка--Лоясиевича: квадратичный рост относительно всего множества
глобальных минимумов.

\section{Критерии времени и масштабирование поля}

Пусть \(V\) положительно определена относительно равновесия
\(x^\star\). Если
\[
    \dot V\le-cV^\alpha,
    \qquad
    c>0,
    \qquad
    0<\alpha<1,
\]
то равновесие достигается за конечное время и
\[
    T(x_0)
    \le
    \frac{V(x_0)^{1-\alpha}}{c(1-\alpha)}.
\]
При экспоненциальной оценке \(\dot V\le-cV\) функция обычно лишь
стремится к нулю. Степень \(\alpha<1\) не позволяет убыванию стать
слишком медленным около равновесия.

Если
\[
    \dot V
    \le
    -c_1V^{\alpha_1}
    -c_2V^{\alpha_2},
\]
где
\[
    c_1,c_2>0,
    \qquad
    0<\alpha_1<1<\alpha_2,
\]
то время установления равномерно ограничено:
\[
    T(x_0)
    \le
    \frac1{c_1(1-\alpha_1)}
    +
    \frac1{c_2(\alpha_2-1)}.
\]
Степень выше единицы отвечает за быстрое движение вдали от решения,
а степень ниже единицы --- за точное достижение после входа в его
окрестность.

Пусть система
\[
    \dot x=F(x)
\]
имеет устойчивое равновесие и \(a(x)>0\) вне него. Система
\[
    \dot x=a(x)F(x)
\]
задаёт те же непараметризованные фазовые траектории и отличается
перепараметризацией времени. Чтобы получить конечно- или
фиксированно-временной результат, необходимо отдельно проверить
определение поля в равновесии, существование решений, устойчивость и
нужное неравенство для функции Ляпунова.

После явной дискретизации это геометрическое совпадение теряется.
Например, фиксированно-временной градиентный поток формально приводит
к рекурсии
\[
\begin{aligned}
    x^{k+1}
    ={}&
    x^k
    -\eta
    \frac{\nabla f(x^k)}{
        \lVert\nabla f(x^k)\rVert^{\frac{p-2}{p-1}}
    }
    \\
    &-\eta
    \frac{\nabla f(x^k)}{
        \lVert\nabla f(x^k)\rVert^{\frac{q-2}{q-1}}
    },
\end{aligned}
\]
где \(p>2\) и \(1<q<2\). В дискретном времени возможны
перескакивание через решение и колебания. Поэтому непрерывная
фиксированно-временная устойчивость не означает точного достижения за
фиксированное число шагов. Для дискретной схемы требуется отдельный выбор достаточно малого шага; универсальной границы на \(\eta\) для общего нелинейного случая нет.

\section{Фиксированно-временной ньютоновский поток}

Пусть
\[
    f\in C^2(\RR^n),
    \qquad
    \nabla^2f(x)\succ0
\]
во всех рассматриваемых точках, а минимум \(x^\star\) существует.
Тогда \(f\) строго выпукла, минимум единственен, а Гессиан обратим.
Обозначим
\[
    g(x)=\nabla f(x),
    \qquad
    H(x)=\nabla^2f(x).
\]
Обычный ньютоновский поток имеет вид
\[
    \dot x=-H(x)^{-1}g(x).
\]
Для
\[
    V(x)=\frac12\lVert g(x)\rVert^2
\]
получаем
\[
    \dot V
    =
    g(x)^{\mathsf T}H(x)\dot x
    =
    -\lVert g(x)\rVert^2
    =
    -2V.
\]
Следовательно,
\[
    \lVert\nabla f(x(t))\rVert
    =
    e^{-(t-t_0)}
    \lVert\nabla f(x(t_0))\rVert.
\]

Выберем параметры
\[
    p>2,
    \qquad
    1<q<2,
\]
и положим
\[
    \rho_p=\frac{p-2}{p-1},
    \qquad
    \rho_q=\frac{q-2}{q-1}.
\]
Масштабированный поток задаётся уравнением
\[
\begin{aligned}
    \dot x
    ={}&
    -H(x)^{-1}
    \frac{g(x)}{\lVert g(x)\rVert^{\rho_p}}
    \\
    &-
    H(x)^{-1}
    \frac{g(x)}{\lVert g(x)\rVert^{\rho_q}},
\end{aligned}
\]
а при \(g(x)=0\) полагается \(\dot x=0\). Обычный и
масштабированный ньютоновские потоки имеют одинаковые геометрические
траектории. Они в общем случае не совпадают с траекториями обычного
градиентного потока.

\section{Функция Ляпунова и оценка времени}

Поскольку строгая выпуклость и существование минимума дают
\[
    g(x)=0
    \quad\Longleftrightarrow\quad
    x=x^\star,
\]
функция
\[
    V(x)=\frac12\lVert g(x)\rVert^2
\]
положительно определена относительно \(x^\star\). Вдоль
масштабированного ньютоновского потока
\[
\begin{aligned}
    \dot V
    ={}&
    -\lVert g(x)\rVert^{2-\rho_p}
    \\
    &-
    \lVert g(x)\rVert^{2-\rho_q}.
\end{aligned}
\]
Положим
\[
    \alpha_1=\frac{p}{2(p-1)},
    \qquad
    \alpha_2=\frac{q}{2(q-1)}.
\]
Так как \(\lVert g(x)\rVert^2=2V(x)\), имеем
\[
    \dot V
    =
    -2^{\alpha_1}V^{\alpha_1}
    -
    2^{\alpha_2}V^{\alpha_2}.
\]
При \(p>2\) выполняется \(0<\alpha_1<1\), а при
\(1<q<2\) --- \(\alpha_2>1\).

\begin{theorem}[Фиксированное время ньютоновского потока]
Пусть \(f\in C^2(\RR^n)\), минимум \(x^\star\) существует,
\(\nabla^2f(x)\succ0\) вдоль рассматриваемых траекторий,
\(p>2\) и \(1<q<2\). Если соответствующее решение существует на
всём промежутке до достижения равновесия, то
\(x^\star\) фиксированно-временно устойчиво и
\[
    T(x_0)
    \le
    \frac1{2^{\alpha_1}(1-\alpha_1)}
    +
    \frac1{2^{\alpha_2}(\alpha_2-1)}.
\]
\end{theorem}

В оценке не появляется константа сильной выпуклости: в производной
\(V\) Гессиан сокращается с обратным Гессианом. Однако условие
\(H(x)\succ0\) не даёт равномерной оценки
\[
    H(x)\succeq\mu I.
\]
Поэтому \(\lVert H(x)^{-1}\rVert\) может быть большой в областях
малой кривизны. Полный результат требует контролировать существование
траектории, отсутствие ухода на бесконечность, корректность поля при
малых собственных значениях Гессиана и продолжение решения после
достижения равновесия.

Можно ввести коэффициенты \(a,b>0\):
\[
\begin{aligned}
    \dot x
    ={}&
    -aH(x)^{-1}
    \frac{g(x)}{\lVert g(x)\rVert^{\rho_p}}
    \\
    &-
    bH(x)^{-1}
    \frac{g(x)}{\lVert g(x)\rVert^{\rho_q}}.
\end{aligned}
\]
Тогда
\[
    T_{\max}
    \le
    \frac1{a2^{\alpha_1}(1-\alpha_1)}
    +
    \frac1{b2^{\alpha_2}(\alpha_2-1)}.
\]
Коэффициенты можно выбрать так, чтобы эта верхняя граница не
превышала заранее заданного времени. Параметры \(p,q\) задают степени,
а \(a,b\) непосредственно масштабируют длительность движения.

\section{PL-неравенство и квадратичный рост}

Пусть дифференцируемая функция достигает значения \(f^\star\) и
удовлетворяет неравенству Поляка--Лоясиевича
\[
    \frac12\lVert\nabla f(x)\rVert^2
    \ge
    \mu\bigl(f(x)-f^\star\bigr),
    \qquad
    \mu>0.
\]
Выпуклость здесь не предполагается. Из условия следует
\[
    \nabla f(x)=0
    \quad\Longrightarrow\quad
    f(x)=f^\star,
\]
то есть неоптимальных стационарных точек нет. Множество минимумов
\[
    \mathcal X^\star
    =
    \operatorname*{argmin}_{x\in\RR^n}f(x)
\]
при этом может содержать несколько точек.

\begin{definition}[Квадратичный рост]
Функция удовлетворяет условию квадратичного роста с константой
\(\mu>0\), если
\[
    f(x)-f^\star
    \ge
    \frac\mu2
    \operatorname{dist}(x,\mathcal X^\star)^2,
\]
где
\[
    \operatorname{dist}(x,\mathcal X^\star)
    =
    \inf_{z\in\mathcal X^\star}\lVert x-z\rVert.
\]
\end{definition}

\begin{theorem}[PL-неравенство влечёт квадратичный рост]
Пусть \(f\) дифференцируема, достигает минимума и удовлетворяет
PL-неравенству с константой \(\mu>0\). При условиях регулярности,
обеспечивающих существование используемого ниже градиентного потока до
достижения \(\mathcal X^\star\), выполняется
\[
    f(x)-f^\star
    \ge
    \frac\mu2
    \operatorname{dist}(x,\mathcal X^\star)^2
\]
для всех \(x\).
\end{theorem}

Таким образом, малая ошибка по значению возможна только вблизи
множества глобальных решений. При единственном минимизаторе условие
принимает привычный вид
\[
    f(x)-f^\star
    \ge
    \frac\mu2\lVert x-x^\star\rVert^2,
\]
но ни PL-условие, ни квадратичный рост сами по себе не требуют
выпуклости.

\section{Доказательство квадратичного роста}

Возьмём \(x_0\notin\mathcal X^\star\) и определим
\[
    G(x)=\sqrt{f(x)-f^\star}.
\]
Вне множества решений
\[
    \nabla G(x)
    =
    \frac{\nabla f(x)}
    {2\sqrt{f(x)-f^\star}}.
\]
Поэтому
\[
    \lVert\nabla G(x)\rVert^2
    =
    \frac{\lVert\nabla f(x)\rVert^2}
    {4(f(x)-f^\star)}
    \ge
    \frac\mu2.
\]
Рассмотрим поток
\[
    \dot x=-\nabla G(x)
\]
до первого момента достижения \(\mathcal X^\star\). Вдоль него
\[
\begin{aligned}
    \frac{\mathrm d}{\mathrm dt}G(x(t))
    &=
    -\lVert\nabla G(x(t))\rVert^2
    \\
    &\le
    -\frac\mu2.
\end{aligned}
\]
Следовательно, множество решений достигается не позднее времени
\[
    T\le\frac{2G(x_0)}\mu.
\]
Обозначим конечную точку через \(x_T\in\mathcal X^\star\).

С другой стороны,
\[
\begin{aligned}
    G(x_0)
    &=
    \int_0^T
    \lVert\nabla G(x(t))\rVert^2\,\mathrm dt
    \\
    &\ge
    \sqrt{\frac\mu2}
    \int_0^T
    \lVert\nabla G(x(t))\rVert\,\mathrm dt.
\end{aligned}
\]
Поскольку
\[
    \lVert\dot x(t)\rVert
    =
    \lVert\nabla G(x(t))\rVert,
\]
последний интеграл является длиной траектории и не меньше расстояния
между её концами:
\[
    \int_0^T
    \lVert\dot x(t)\rVert\,\mathrm dt
    \ge
    \lVert x_0-x_T\rVert
    \ge
    \operatorname{dist}(x_0,\mathcal X^\star).
\]
Значит,
\[
    G(x_0)
    \ge
    \sqrt{\frac\mu2}
    \operatorname{dist}(x_0,\mathcal X^\star).
\]
После возведения в квадрат получаем
\[
    f(x_0)-f^\star
    \ge
    \frac\mu2
    \operatorname{dist}(x_0,\mathcal X^\star)^2.
\]
Поскольку начальная точка произвольна, теорема доказана.

\section{Следствия, примеры и границы}

Из квадратичного роста следует
\[
    f(x)-f^\star\le\varepsilon
    \quad\Longrightarrow\quad
    \operatorname{dist}(x,\mathcal X^\star)
    \le
    \sqrt{\frac{2\varepsilon}{\mu}}.
\]
Из PL-неравенства получаем
\[
    f(x)-f^\star
    \le
    \frac1{2\mu}
    \lVert\nabla f(x)\rVert^2.
\]
Объединение двух оценок даёт
\[
    \operatorname{dist}(x,\mathcal X^\star)
    \le
    \frac1\mu
    \lVert\nabla f(x)\rVert.
\]
Поэтому малая норма градиента у PL-функции означает близость к
множеству глобальных решений, а не только почти стационарность.

PL-условие не гарантирует единственности минимума. Для
\[
    f(x_1,x_2)=\frac12x_1^2
\]
имеем
\[
    \frac12\lVert\nabla f(x)\rVert^2
    =
    f(x)-f^\star,
\]
но
\[
    \mathcal X^\star
    =
    \{(0,s):s\in\RR\}.
\]
При этом
\[
    f(x)-f^\star
    =
    \frac12
    \operatorname{dist}(x,\mathcal X^\star)^2.
\]
Следовательно, утверждение, будто PL-функция либо постоянна, либо
имеет единственный минимум, неверно.

PL-свойство исключает неоптимальные стационарные точки; инвексность является отдельным структурным условием. Функция
\[
    f(x)=x^2+3\sin^2x
\]
иллюстрирует невыпуклую геометрию; без отдельной проверки PL-неравенства она не используется как доказательный пример.

Для дифференцируемых функций сохраняется цепочка
\[
    \text{сильная выпуклость}
    \Longrightarrow
    \text{PL}
    \Longrightarrow
    \text{квадратичный рост},
\]
при указанных предпосылках существования минимума и регулярности.
Обратные импликации в общем случае неверны.

Ньютоновская конструкция требует решения системы
\[
    H(x)d=-g(x),
\]
а не явного вычисления обратной матрицы, но такая операция всё равно
может быть дорогой. PL-результат требует
\(\mathcal X^\star\ne\varnothing\), а доказательство через \(G\) ---
существования траектории до достижения множества решений.
Непрерывные фиксированно-временные оценки не переносятся автоматически
на методы с конечным шагом.

\lecture{20}{PL-свойство композиций}

\sourcevideo
    {https://www.youtube.com/playlist?list=PLOzRYVm0a65fhKDS8cYIC7YCuVGJ4FOJx}
    {Week 6: Lecture 20: Advanced Results on PL Inequality: Part 2}

Лекция состоит из двух частей. Сначала доказывается, что композиция
сильно выпуклой функции с линейным отображением может потерять
сильную выпуклость, но сохраняет неравенство Поляка--Лоясиевича.
Затем исследуется робастность фиксированно-временного градиентного
потока к возмущению, исчезающему около множества решений.

\section{PL-условие и композиция}

Пусть дифференцируемая функция
\[
    f\colon\mathbb R^n\to\mathbb R
\]
достигает минимального значения
\[
    f^\star=\min_x f(x),
    \qquad
    \mathcal X^\star=\operatorname*{argmin}_x f(x).
\]

\begin{definition}[Неравенство Поляка--Лоясиевича]
Функция \(f\) удовлетворяет PL-неравенству с константой \(\mu>0\),
если
\[
    \frac12\lVert\nabla f(x)\rVert^2
    \ge
    \mu\bigl(f(x)-f^\star\bigr)
\]
для всех \(x\in\mathbb R^n\).
\end{definition}

PL-условие не требует выпуклости. Оно исключает неоптимальные
стационарные точки, поскольку
\[
    \nabla f(x)=0
    \quad\Longrightarrow\quad
    f(x)=f^\star,
\]
но не гарантирует единственности минимизатора.

Пусть
\[
    g\colon\mathbb R^m\to\mathbb R
\]
является \(\sigma\)-сильно выпуклой и
\[
    f(x)=g(Ax-b),
    \qquad
    A\in\mathbb R^{m\times n},
    \quad
    b\in\mathbb R^m.
\]
Тогда
\[
    \nabla f(x)=A^{\mathsf T}\nabla g(Ax-b).
\]
Из сильной выпуклости \(g\) следует
\[
\begin{aligned}
    f(y)
    \ge{}&
    f(x)
    +\langle\nabla f(x),y-x\rangle
    \\
    &+
    \frac\sigma2
    \lVert A(y-x)\rVert^2.
\end{aligned}
\]
Последнее слагаемое содержит \(\lVert A(y-x)\rVert\), поэтому
сильная выпуклость внешней функции не всегда переносится на
композицию.

Если \(A\) имеет полный столбцовый ранг, то
\[
    \lVert Az\rVert
    \ge
    \sigma_{\min}(A)\lVert z\rVert,
\]
и \(f\) является сильно выпуклой с константой
\[
    \sigma\sigma_{\min}(A)^2.
\]
Если же \(\ker A\ne\{0\}\), то
\[
    f(x+z)=f(x),
    \qquad
    z\in\ker A,
\]
поэтому глобальная сильная выпуклость невозможна.

\section{Сингулярное число, оценка Хоффмана и PL-свойство}

Пусть \(A\ne0\), но матрица может быть вырожденной. Обозначим её
наименьшее положительное сингулярное число через
\[
    \theta_A=\sigma_{\min}^{+}(A).
\]
Для каждого \(z\in(\ker A)^\perp\)
\[
    \lVert Az\rVert\ge\theta_A\lVert z\rVert.
\]
Сильная выпуклость \(g\) обеспечивает единственность решения задачи
\[
    \min_{u\in\operatorname{range}(A)-b}g(u).
\]
Обозначим его через \(u^\star\). Тогда
\[
    \mathcal X^\star
    =
    \{x\in\mathbb R^n:Ax-b=u^\star\}.
\]

Для заданного \(x\) выберем ближайший минимизатор
\[
    x^\star
    \in
    \operatorname*{argmin}_{z\in\mathcal X^\star}
    \lVert x-z\rVert.
\]
Тогда \(x-x^\star\in(\ker A)^\perp\), и потому
\[
    \lVert A(x-x^\star)\rVert
    \ge
    \theta_A\operatorname{dist}(x,\mathcal X^\star).
\]
Это частный случай оценки Хоффмана для линейной системы
\(Ax=Ax^\star\).

Подставляя \(y=x^\star\) в неравенство для композиции, получаем
\[
\begin{aligned}
    f^\star
    \ge{}&
    f(x)
    +\langle\nabla f(x),x^\star-x\rangle
    \\
    &+
    \frac{\sigma\theta_A^2}{2}
    \lVert x^\star-x\rVert^2.
\end{aligned}
\]
Это ограниченная сильная выпуклость относительно ближайшей точки
множества решений, а не обычная сильная выпуклость для произвольной
пары точек.

Положим
\[
    \mu_A=\sigma\theta_A^2.
\]
Для любых \(a,d\) выполняется
\[
    \langle a,d\rangle
    +\frac{\mu_A}{2}\lVert d\rVert^2
    \ge
    -\frac1{2\mu_A}\lVert a\rVert^2.
\]
Применяя это неравенство при
\[
    a=\nabla f(x),
    \qquad
    d=x^\star-x,
\]
получаем
\[
    f^\star
    \ge
    f(x)-\frac1{2\mu_A}\lVert\nabla f(x)\rVert^2.
\]
Следовательно,
\[
    \frac12\lVert\nabla f(x)\rVert^2
    \ge
    \mu_A\bigl(f(x)-f^\star\bigr).
\]

\begin{theorem}[PL-свойство линейной композиции]
Пусть \(g\) является \(\sigma\)-сильно выпуклой, \(A\ne0\), а
функция
\[
    f(x)=g(Ax-b)
\]
достигает минимума. Тогда \(f\) удовлетворяет PL-неравенству с
константой
\[
    \mu_A
    =
    \sigma
    \bigl(\sigma_{\min}^{+}(A)\bigr)^2.
\]
\end{theorem}

Матрица \(A\) может быть вырожденной. В этом случае функция не
является сильно выпуклой, но остаётся PL-функцией. При \(A=0\)
композиция постоянна и PL-неравенство выполняется тривиально.

\section{Наименьшие квадраты и границы переноса}

Для задачи
\[
    f(x)=\frac12\lVert Ax-b\rVert^2
\]
внешняя функция
\[
    g(u)=\frac12\lVert u\rVert^2
\]
является \(1\)-сильно выпуклой. Поэтому
\[
    \frac12\lVert\nabla f(x)\rVert^2
    \ge
    \bigl(\sigma_{\min}^{+}(A)\bigr)^2
    \bigl(f(x)-f^\star\bigr),
\]
где
\[
    \nabla f(x)=A^{\mathsf T}(Ax-b).
\]
Если \(A\) имеет полный столбцовый ранг, решение единственно и
функция сильно выпукла. При ненулевом ядре минимизаторы имеют вид
\[
    x^\star+z,
    \qquad
    z\in\ker A,
\]
но PL-свойство сохраняется. Поэтому градиентные методы могут линейно
сходиться по значению функции даже для вырожденной задачи наименьших
квадратов.

Тот же аргумент применим к любой композиции \(g(Ax-b)\), если
внешняя функция действительно сильно выпукла на рассматриваемой
области. Для стандартной нерегуляризованной логистической потери это
условие глобально обычно не выполнено; требуются дополнительные
ограничения или регуляризация.

\section{Возмущённый поток}

Пусть \(f\) удовлетворяет PL-неравенству с константой \(\mu>0\).
Сначала предположим, что минимизатор \(x^\star\) единственен. Выберем
\[
    p>2,
    \qquad
    1<q<2,
\]
и обозначим
\[
    \rho_p=\frac{p-2}{p-1},
    \qquad
    \rho_q=\frac{q-2}{q-1}.
\]
Рассмотрим возмущённую систему
\[
\begin{aligned}
    \dot x
    ={}&
    -c_1
    \frac{\nabla f(x)}{
        \lVert\nabla f(x)\rVert^{\rho_p}
    }
    \\
    &-c_2
    \frac{\nabla f(x)}{
        \lVert\nabla f(x)\rVert^{\rho_q}
    }
    +\varepsilon(x),
\end{aligned}
\]
где \(c_1,c_2>0\), а в стационарной точке правая часть определяется
равной нулю. Предполагается
\[
    \lVert\varepsilon(x)\rVert
    \le
    \ell\lVert x-x^\star\rVert^2.
\]
Отсюда \(\varepsilon(x^\star)=0\), поэтому оптимум сохраняется как
равновесие. При ненулевом остаточном возмущении точная сходимость к
\(x^\star\) в общем случае невозможна; можно ожидать лишь попадание
в окрестность.

Из предыдущей лекции PL-условие влечёт квадратичный рост:
\[
    f(x)-f^\star
    \ge
    \frac\mu2\lVert x-x^\star\rVert^2.
\]
Совместно с PL-неравенством это даёт
\[
    \lVert x-x^\star\rVert
    \le
    \frac1\mu\lVert\nabla f(x)\rVert.
\]
Следовательно,
\[
    \lVert\varepsilon(x)\rVert
    \le
    \frac\ell{\mu^2}
    \lVert\nabla f(x)\rVert^2.
\]
Обозначим
\[
    \overline\ell=\frac\ell{\mu^2},
    \qquad
    V(x)=f(x)-f^\star,
    \qquad
    s=\lVert\nabla f(x)\rVert.
\]
Вдоль возмущённого потока
\[
    \dot V
    \le
    -c_1s^{2-\rho_p}
    -c_2s^{2-\rho_q}
    +\overline\ell s^3.
\]
Положим
\[
    r_1=2-\rho_p=\frac{p}{p-1},
    \qquad
    r_2=2-\rho_q=\frac{q}{q-1}.
\]
Тогда
\[
    1<r_1<2,
    \qquad
    r_2>2.
\]

\section{Поглощение возмущения и робастная теорема}

При \(0\le s\le1\) из \(r_1<3\) следует
\[
    s^{r_1}\ge s^3.
\]
При \(s\ge1\) требуется
\[
    s^{r_2}\ge s^3,
\]
то есть \(r_2\ge3\). Это условие эквивалентно
\[
    1<q\le\frac32.
\]
Таким образом, для всех \(s\ge0\)
\[
    s^3\le s^{r_1}+s^{r_2},
\]
и потому
\[
    \dot V
    \le
    -\bigl(c_1-\overline\ell\bigr)s^{r_1}
    -\bigl(c_2-\overline\ell\bigr)s^{r_2}.
\]
Если
\[
    c_1>\overline\ell,
    \qquad
    c_2>\overline\ell,
\]
оба коэффициента положительны.

Введём
\[
    \alpha_1=\frac{r_1}{2}=\frac{p}{2(p-1)},
    \qquad
    \alpha_2=\frac{r_2}{2}=\frac{q}{2(q-1)}.
\]
Тогда
\[
    0<\alpha_1<1<\alpha_2.
\]
Из PL-неравенства \(s^2\ge2\mu V\) следует
\[
    \dot V
    \le
    -a_1V^{\alpha_1}
    -a_2V^{\alpha_2},
\]
где
\[
    a_1
    =
    \left(c_1-\frac\ell{\mu^2}\right)
    (2\mu)^{\alpha_1},
\]
\[
    a_2
    =
    \left(c_2-\frac\ell{\mu^2}\right)
    (2\mu)^{\alpha_2}.
\]

\begin{theorem}[Робастность к исчезающему возмущению]
Пусть \(f\) удовлетворяет PL-неравенству с константой \(\mu>0\) и
имеет единственный минимизатор \(x^\star\). Предположим, что
\[
    \lVert\varepsilon(x)\rVert
    \le
    \ell\lVert x-x^\star\rVert^2,
\]
а параметры удовлетворяют
\[
    p>2,
    \qquad
    1<q\le\frac32,
    \qquad
    c_1,c_2>\frac\ell{\mu^2}.
\]
Тогда \(x^\star\) является фиксированно-временно устойчивым
равновесием, причём
\[
    T(x_0)
    \le
    \frac1{a_1(1-\alpha_1)}
    +
    \frac1{a_2(\alpha_2-1)}.
\]
Правая часть не зависит от начальной точки.
\end{theorem}

Первое слагаемое потока поглощает возмущение при малом градиенте, а
второе --- при большом. Увеличение \(c_1,c_2\) уменьшает
теоретическую границу времени, но после дискретизации чрезмерное
усиление может ухудшить численную устойчивость.

\section{Несколько решений, дискретизация и дальнейший переход}

Если \(\mathcal X^\star\) содержит несколько точек, квадратичный рост
записывается как
\[
    f(x)-f^\star
    \ge
    \frac\mu2
    \operatorname{dist}(x,\mathcal X^\star)^2.
\]
Естественное условие на возмущение принимает вид
\[
    \lVert\varepsilon(x)\rVert
    \le
    \ell
    \operatorname{dist}(x,\mathcal X^\star)^2.
\]
Тогда исследуется фиксированно-временная устойчивость множества
решений, а не сходимость к заранее выбранному его элементу.

Явная дискретизация имеет вид
\[
\begin{aligned}
    x^{k+1}
    ={}&
    x^k
    -\eta c_1
    \frac{\nabla f(x^k)}{
        \lVert\nabla f(x^k)\rVert^{\rho_p}
    }
    \\
    &-\eta c_2
    \frac{\nabla f(x^k)}{
        \lVert\nabla f(x^k)\rVert^{\rho_q}
    }
    +\eta\varepsilon(x^k).
\end{aligned}
\]
Непрерывная фиксированно-временная гарантия не переносится на эту
схему автоматически. При слишком большом шаге или коэффициентах
возможны перескакивание через решение, колебания и расходимость. В
лекции утверждается существование достаточно малого шага, при котором
последовательность попадает в заданную окрестность, но универсальная
явная формула для общего нелинейного случая не приводится.

Далее курс переходит от задач без ограничений к постановкам
\[
    \min_x f(x)
    \qquad
    \text{при условии}
    \qquad
    Ax=b.
\]
Основными инструментами станут функция Лагранжа, дополненная функция
Лагранжа, метод множителей, двойственный подъём и ADMM.

\lecture{21}{Двойственные и седловые методы}

\sourcevideo
    {https://www.youtube.com/playlist?list=PLOzRYVm0a65fhKDS8cYIC7YCuVGJ4FOJx}
    {Week 6: Lecture 21: Constrained Optimization Problem}

Задача с ограничениями может быть преобразована в двойственную задачу
по множителям Лагранжа или в задачу поиска седловой точки. Для
линейных равенств двойственная функция выражается через сопряжённую
функцию Фенхеля, а её градиент совпадает с невязкой ограничения.

\section{Функция Лагранжа и двойственная задача}

Рассмотрим задачу
\[
    \min_{x\in\mathbb R^n} f_0(x)
\]
при условиях
\[
    f_i(x)\le0,
    \qquad
    i=1,\ldots,m,
\]
\[
    h_j(x)=0,
    \qquad
    j=1,\ldots,r.
\]
Переменная \(x\) называется прямой. Для неравенств вводятся множители
\(\lambda_i\ge0\), а для равенств --- свободные множители
\(\nu_j\in\mathbb R\).

\begin{definition}[Функция Лагранжа]
Функцией Лагранжа называется
\[
    \mathcal L(x,\lambda,\nu)
    =
    f_0(x)
    +
    \sum_{i=1}^{m}\lambda_i f_i(x)
    +
    \sum_{j=1}^{r}\nu_j h_j(x).
\]
Переменные \(\lambda\) и \(\nu\) называются двойственными.
\end{definition}

\begin{definition}[Двойственная функция]
Двойственная функция определяется формулой
\[
    g(\lambda,\nu)
    =
    \inf_{x\in\mathbb R^n}
    \mathcal L(x,\lambda,\nu).
\]
\end{definition}

При каждом фиксированном \(x\) функция
\[
    (\lambda,\nu)
    \longmapsto
    \mathcal L(x,\lambda,\nu)
\]
аффинна. Поэтому её поточечный инфимум \(g\) всегда вогнут, даже если
исходная задача невыпукла. Двойственная задача имеет вид
\[
    \max_{\lambda\ge0,\,\nu}g(\lambda,\nu).
\]

Для любой допустимой пары \((\lambda,\nu)\) выполняется слабая
двойственность
\[
    g(\lambda,\nu)\le p^\star,
\]
где \(p^\star\) --- оптимальное значение прямой задачи. Если
\[
    d^\star=p^\star,
\]
двойственный разрыв отсутствует. Для выпуклых задач это обычно
обеспечивается условием регулярности, например условием Слейтера.

\section{Сопряжённая функция Фенхеля}

\begin{definition}[Сопряжённая функция]
Для функции
\[
    f\colon\mathbb R^n
    \to
    \mathbb R\cup\{+\infty\}
\]
сопряжённая функция Фенхеля определяется как
\[
    f^\ast(y)
    =
    \sup_{x\in\mathbb R^n}
    \left\{
        \langle y,x\rangle-f(x)
    \right\}.
\]
\end{definition}

Для каждого \(x\) выражение
\[
    \langle y,x\rangle-f(x)
\]
аффинно по \(y\), поэтому \(f^\ast\) всегда выпукла. Геометрически
\(f^\ast(y)\) является максимальным разрывом между линейной функцией
\[
    x\longmapsto\langle y,x\rangle
\]
и функцией \(f\).

Для замкнутых собственных выпуклых функций сильная выпуклость и
гладкость переходят к сопряжённой функции с обратными константами.

\begin{theorem}
Если \(f\) является \(\mu\)-сильно выпуклой, то \(f^\ast\)
дифференцируема и
\[
    \left\lVert
        \nabla f^\ast(y_1)
        -
        \nabla f^\ast(y_2)
    \right\rVert
    \le
    \frac1\mu
    \lVert y_1-y_2\rVert.
\]
Кроме того,
\[
    \nabla f^\ast(y)
    =
    \operatorname*{argmax}_{x}
    \left\{
        \langle y,x\rangle-f(x)
    \right\}.
\]
\end{theorem}

\begin{theorem}
Если \(f\) выпукла и \(L\)-гладка, то \(f^\ast\) является
\(1/L\)-сильно выпуклой на своей области определения.
\end{theorem}

Таким образом, сильная выпуклость исходной функции обеспечивает
гладкость сопряжённой, а гладкость исходной функции --- сильную
выпуклость сопряжённой.

\section{Линейные равенства и двойственный подъём}

Рассмотрим задачу
\[
    \min_{x\in\mathbb R^n}f(x)
    \qquad
    \text{при условии}
    \qquad
    Ax=b,
\]
где
\[
    A\in\mathbb R^{m\times n},
    \qquad
    b\in\mathbb R^m.
\]
Её функция Лагранжа равна
\[
    \mathcal L(x,\nu)
    =
    f(x)+\nu^{\mathsf T}(Ax-b).
\]

Двойственная функция вычисляется следующим образом:
\[
\begin{aligned}
    g(\nu)
    &=
    \inf_x
    \left\{
        f(x)+\nu^{\mathsf T}(Ax-b)
    \right\}
    \\
    &=
    -b^{\mathsf T}\nu
    +
    \inf_x
    \left\{
        f(x)+\langle A^{\mathsf T}\nu,x\rangle
    \right\}
    \\
    &=
    -f^\ast(-A^{\mathsf T}\nu)
    -
    b^{\mathsf T}\nu.
\end{aligned}
\]
Следовательно, двойственная задача не имеет ограничений по \(\nu\):
\[
    \max_{\nu\in\mathbb R^m}
    \left\{
        -f^\ast(-A^{\mathsf T}\nu)
        -
        b^{\mathsf T}\nu
    \right\}.
\]

Пусть \(f\) сильно выпукла. Тогда
\[
    x(\nu)
    =
    \operatorname*{argmin}_x
    \mathcal L(x,\nu)
\]
единственен и
\[
    x(\nu)
    =
    \nabla f^\ast(-A^{\mathsf T}\nu).
\]
По теореме Данскина
\[
    \nabla g(\nu)=Ax(\nu)-b.
\]
Таким образом, градиент двойственной функции совпадает с невязкой
прямого ограничения.

Обычный двойственный подъём имеет вид
\[
    \nu^{k+1}
    =
    \nu^k
    +
    \eta
    \bigl(
        Ax(\nu^k)-b
    \bigr).
\]
На каждой итерации сначала решается внутренняя задача по \(x\), затем
множитель изменяется в направлении нарушения ограничения.

Пусть вогнутая функция \(g\) удовлетворяет PL-условию для
максимизации:
\[
    \frac12
    \lVert\nabla g(\nu)\rVert^2
    \ge
    \mu_d
    \bigl(
        g^\star-g(\nu)
    \bigr).
\]
Для
\[
    p>2,
    \qquad
    1<q<2
\]
положим
\[
    \rho_p=\frac{p-2}{p-1},
    \qquad
    \rho_q=\frac{q-2}{q-1}.
\]
Фиксированно-временной двойственный поток записывается как
\[
\begin{aligned}
    \dot\nu
    ={}&
    c_1
    \frac{\nabla g(\nu)}
    {\lVert\nabla g(\nu)\rVert^{\rho_p}}
    \\
    &+
    c_2
    \frac{\nabla g(\nu)}
    {\lVert\nabla g(\nu)\rVert^{\rho_q}},
\end{aligned}
\]
где \(c_1,c_2>0\). Для
\[
    V(\nu)=g^\star-g(\nu)
\]
имеем
\[
    \dot V
    \le
    -c_1(2\mu_d)^{\alpha_1}V^{\alpha_1}
    -
    c_2(2\mu_d)^{\alpha_2}V^{\alpha_2},
\]
где
\[
    \alpha_1
    =
    \frac{p}{2(p-1)}
    <1,
    \qquad
    \alpha_2
    =
    \frac{q}{2(q-1)}
    >1.
\]
По критерию фиксированного времени двойственный оптимум достигается за
время, равномерно ограниченное по начальному множителю.

\section{Квадратичная задача и система KKT}

Рассмотрим
\[
    \min_x
    \left\{
        \frac12x^{\mathsf T}Qx+c^{\mathsf T}x
    \right\}
    \qquad
    \text{при условии}
    \qquad
    Ax=b,
\]
где
\[
    Q\succ0.
\]
Для
\[
    f(x)
    =
    \frac12x^{\mathsf T}Qx+c^{\mathsf T}x
\]
сопряжённая функция равна
\[
    f^\ast(y)
    =
    \frac12
    (y-c)^{\mathsf T}
    Q^{-1}
    (y-c).
\]
Действительно, условие оптимальности во внутренней задаче даёт
\[
    y-Qx-c=0,
    \qquad
    x=Q^{-1}(y-c).
\]

Поэтому двойственная функция имеет вид
\[
\begin{aligned}
    g(\nu)
    ={}&
    -\frac12
    \bigl(
        c+A^{\mathsf T}\nu
    \bigr)^{\mathsf T}
    Q^{-1}
    \bigl(
        c+A^{\mathsf T}\nu
    \bigr)
    \\
    &-
    b^{\mathsf T}\nu,
\end{aligned}
\]
а минимизатор функции Лагранжа равен
\[
    x(\nu)
    =
    -Q^{-1}
    \bigl(
        c+A^{\mathsf T}\nu
    \bigr).
\]
Следовательно,
\[
    \nabla g(\nu)
    =
    -AQ^{-1}
    \bigl(
        c+A^{\mathsf T}\nu
    \bigr)
    -
    b
\]
и
\[
    \nabla^2g(\nu)
    =
    -AQ^{-1}A^{\mathsf T}
    \preceq0.
\]
Если \(A\) имеет полный строчный ранг, то
\[
    AQ^{-1}A^{\mathsf T}\succ0,
\]
поэтому \(g\) строго вогнута.

Условия оптимальности равны
\[
    Qx^\star+c+A^{\mathsf T}\nu^\star=0,
    \qquad
    Ax^\star=b.
\]
Они образуют систему KKT
\[
\begin{pmatrix}
    Q&A^{\mathsf T}\\
    A&0
\end{pmatrix}
\begin{pmatrix}
    x^\star\\
    \nu^\star
\end{pmatrix}
=
\begin{pmatrix}
    -c\\
    b
\end{pmatrix}.
\]
При \(Q\succ0\) и полном строчном ранге \(A\) пара
\[
    (x^\star,\nu^\star)
\]
единственна.

\section{Седловая динамика и ньютоновский поток}

Для задачи
\[
    \min_x\max_z\Phi(x,z)
\]
пара \((x^\star,z^\star)\) называется седловой точкой, если
\[
    \Phi(x^\star,z)
    \le
    \Phi(x^\star,z^\star)
    \le
    \Phi(x,z^\star)
\]
для всех допустимых \(x,z\). Для линейного ограничения функция
Лагранжа задаёт именно такую постановку:
\[
    \min_x\max_\nu\mathcal L(x,\nu).
\]

Естественная динамика спуска--подъёма имеет вид
\[
    \dot x=-\nabla_x\Phi(x,z),
    \qquad
    \dot z=\nabla_z\Phi(x,z).
\]
Введём
\[
    w=
    \begin{pmatrix}
        x\\
        z
    \end{pmatrix},
    \qquad
    F(w)=
    \begin{pmatrix}
        \nabla_x\Phi(x,z)\\
        -\nabla_z\Phi(x,z)
    \end{pmatrix}.
\]
Тогда седловая точка удовлетворяет
\[
    F(w^\star)=0,
\]
а спуск--подъём записывается как
\[
    \dot w=-F(w).
\]

Якобиан седлового оператора равен
\[
    J_F(w)
    =
    \begin{pmatrix}
        \nabla^2_{xx}\Phi
        &
        \nabla^2_{xz}\Phi
        \\
        -\nabla^2_{zx}\Phi
        &
        -\nabla^2_{zz}\Phi
    \end{pmatrix}.
\]
Если \(\Phi\) строго выпукла по \(x\), строго вогнута по \(z\), а
смешанные производные согласованы, то
\[
    \frac{
        J_F(w)+J_F(w)^{\mathsf T}
    }2
    =
    \begin{pmatrix}
        \nabla^2_{xx}\Phi&0\\
        0&-\nabla^2_{zz}\Phi
    \end{pmatrix}
    \succ0.
\]
Следовательно, \(J_F(w)\) невырожден.

Ньютоновский седловой поток имеет вид
\[
    \dot w=-J_F(w)^{-1}F(w).
\]
Для
\[
    V(w)=\frac12\lVert F(w)\rVert^2
\]
получаем
\[
\begin{aligned}
    \dot V
    &=
    F(w)^{\mathsf T}J_F(w)\dot w
    \\
    &=
    -\lVert F(w)\rVert^2
    =
    -2V.
\end{aligned}
\]
Поэтому невязка седловых условий затухает экспоненциально.

Фиксированно-временная модификация задаётся уравнением
\[
\begin{aligned}
    \dot w
    ={}&
    -c_1J_F(w)^{-1}
    \frac{F(w)}
    {\lVert F(w)\rVert^{\rho_p}}
    \\
    &-
    c_2J_F(w)^{-1}
    \frac{F(w)}
    {\lVert F(w)\rVert^{\rho_q}},
\end{aligned}
\]
а при \(F(w)=0\) полагается \(\dot w=0\). Тогда
\[
    \dot V
    =
    -c_1 2^{\alpha_1}V^{\alpha_1}
    -
    c_2 2^{\alpha_2}V^{\alpha_2}
\]
и
\[
    T_{\max}
    =
    \frac1{
        c_1 2^{\alpha_1}(1-\alpha_1)
    }
    +
    \frac1{
        c_2 2^{\alpha_2}(\alpha_2-1)
    }.
\]
Если корень \(F\) единственен, то
\[
    w(t)=w^\star
\]
после времени \(T_{\max}\).

\section{Предпосылки и переход к дополненной функции}

Строгая выпуклость по \(x\) и строгая вогнутость по \(z\)
обеспечивают не более одной седловой точки, но не всегда гарантируют
её существование на неограниченной области. Дополнительно могут
потребоваться компактность, коэрцитивность, антикоэрцитивность или
условия сильной двойственности. Для ньютоновской динамики также
предполагаются существование траектории и обратимость \(J_F(w)\)
вдоль неё.

Двойственный метод требует на каждой итерации вычислять
\[
    x(\nu)
    =
    \operatorname*{argmin}_x
    \mathcal L(x,\nu)
\]
или, что эквивалентно при сильной выпуклости, значение
\[
    \nabla f^\ast(-A^{\mathsf T}\nu).
\]
В общем случае это самостоятельная задача оптимизации.

Ньютоновский седловой поток требует решать линейную систему
\[
    J_F(w)d=-F(w).
\]
Матрицу явно не обращают. После дискретизации фиксированно-временные
гарантии требуют отдельного анализа шага и устойчивости.

Обычная функция Лагранжа
\[
    \mathcal L(x,\nu)
    =
    f(x)+\nu^{\mathsf T}(Ax-b)
\]
может иметь недостаточную кривизну по прямой переменной. В следующей
лекции к ней добавляется квадратичный штраф:
\[
    \mathcal L_\rho(x,\nu)
    =
    f(x)
    +
    \nu^{\mathsf T}(Ax-b)
    +
    \frac\rho2
    \lVert Ax-b\rVert^2.
\]
Эта дополненная функция Лагранжа лежит в основе метода множителей и
ADMM.

\lecture{22}{Дополненная функция Лагранжа}

\sourcevideo
    {https://www.youtube.com/playlist?list=PLOzRYVm0a65fhKDS8cYIC7YCuVGJ4FOJx}
    {Week 6: Lecture 22: Augmented Lagrangian}

Рассмотрим задачу с ограничениями-равенствами
\[
    \min_{x\in\mathbb R^n} f(x)
    \qquad
    \text{при условии}
    \qquad
    h(x)=0,
\]
где
\[
    h\colon\mathbb R^n\to\mathbb R^m.
\]
Обычная функция Лагранжа равна
\[
    \mathcal L(x,\nu)
    =
    f(x)+\nu^{\mathsf T}h(x),
\]
а дополненная функция Лагранжа определяется формулой
\[
    \mathcal L_c(x,\nu)
    =
    f(x)
    +
    \nu^{\mathsf T}h(x)
    +
    \frac c2\lVert h(x)\rVert^2,
    \qquad
    c>0.
\]
Квадратичный член не изменяет значение функции на допустимом
множестве, но добавляет кривизну в направлениях, нарушающих
ограничение. Эта конструкция лежит в основе метода множителей и ADMM.

\section{Штрафная и дополненная постановки}

Следует различать исходную задачу
\[
    \min_x f(x)
    \qquad
    \text{при условии}
    \qquad
    h(x)=0,
\]
эквивалентную задачу
\[
    \min_x
    \left\{
        f(x)+\frac c2\lVert h(x)\rVert^2
    \right\}
    \qquad
    \text{при условии}
    \qquad
    h(x)=0
\]
и безусловную штрафную задачу
\[
    \min_x
    \left\{
        f(x)+\frac c2\lVert h(x)\rVert^2
    \right\}.
\]
Первые две задачи имеют одинаковые допустимые точки и одинаковые
значения целевой функции на них. В третьей задаче ограничение удалено,
поэтому при конечном \(c\) её решение в общем случае не обязано быть
допустимым. При подходящих предпосылках выполняется лишь
\[
    h(x_c)\to0
    \qquad
    \text{при }c\to\infty.
\]

Дополненная функция Лагранжа сочетает квадратичный штраф с линейным
членом \(\nu^{\mathsf T}h(x)\). Именно обновление множителя позволяет
методу множителей работать при конечном и умеренном штрафном
параметре.

Для линейного ограничения
\[
    h(x)=Ax-b,
    \qquad
    A\in\mathbb R^{m\times n},
\]
штраф имеет градиент и Гессиан
\[
    \nabla_x\frac c2\lVert Ax-b\rVert^2
    =
    cA^{\mathsf T}(Ax-b),
\]
\[
    \nabla_x^2\frac c2\lVert Ax-b\rVert^2
    =
    cA^{\mathsf T}A.
\]
Матрица \(A^{\mathsf T}A\) положительно полуопределена, поскольку
\[
    y^{\mathsf T}A^{\mathsf T}Ay
    =
    \lVert Ay\rVert^2
    \ge0.
\]
Она положительно определена тогда и только тогда, когда
\[
    \ker A=\{0\},
\]
то есть \(A\) имеет полный столбцовый ранг. В этом случае штраф
сильно выпуклый с константой
\[
    c\sigma_{\min}(A)^2.
\]
Если \(\ker A\ne\{0\}\), штраф не создаёт кривизну вдоль ядра.
Для положительной определённости \(A^{\mathsf T}A\) необходим полный столбцовый ранг \(A\).

\section{Роль множителя и квадратичный пример}

На допустимом множестве
\[
    \mathcal L_c(x,\nu)
    =
    \mathcal L(x,\nu)
    =
    f(x).
\]
При фиксированном произвольном \(\nu\) минимизаторы
\(\mathcal L(\cdot,\nu)\) и \(\mathcal L_c(\cdot,\nu)\) могут
различаться и не обязаны совпадать с решением исходной задачи.
Совпадение обеспечивается при правильном множителе и соответствующих
условиях оптимальности.

Рассмотрим задачу
\[
    \min_{x_1,x_2}
    \frac12(x_1^2+x_2^2)
    \qquad
    \text{при условии}
    \qquad
    x_1=1.
\]
Её решение равно
\[
    x^\star=
    \begin{pmatrix}
        1\\
        0
    \end{pmatrix}.
\]
Функция Лагранжа имеет вид
\[
    \mathcal L(x_1,x_2,\nu)
    =
    \frac12(x_1^2+x_2^2)+\nu(x_1-1).
\]
Из условий стационарности и допустимости
\[
    x_1+\nu=0,
    \qquad
    x_2=0,
    \qquad
    x_1=1
\]
получаем
\[
    \nu^\star=-1.
\]
Дополненная функция равна
\[
\begin{aligned}
    \mathcal L_c(x_1,x_2,\nu)
    ={}&
    \frac12(x_1^2+x_2^2)
    +\nu(x_1-1)
    \\
    &+
    \frac c2(x_1-1)^2.
\end{aligned}
\]
При фиксированном \(\nu\) её минимизатор удовлетворяет
\[
    x_2=0,
    \qquad
    x_1+\nu+c(x_1-1)=0,
\]
то есть
\[
    x_1(c,\nu)
    =
    \frac{c-\nu}{c+1}.
\]
При \(\nu=\nu^\star=-1\) получаем \(x_1=1\) для любого \(c\ge0\),
тогда как при произвольном фиксированном множителе
\[
    x_1(c,\nu)\to1
    \qquad
    \text{при }c\to\infty.
\]
Тем самым к решению можно приближаться либо увеличением штрафа, либо
уточнением двойственной переменной. Метод множителей использует второй
механизм.

\section{KKT-условия и локальная кривизна}

Пусть \(x^\star\) является локальным решением, а якобиан ограничения
равен
\[
    J_h(x)
    =
    \begin{pmatrix}
        \nabla h_1(x)^{\mathsf T}\\
        \vdots\\
        \nabla h_m(x)^{\mathsf T}
    \end{pmatrix}.
\]
При условии регулярности существует множитель \(\nu^\star\), для
которого выполняются KKT-условия
\[
    h(x^\star)=0,
\]
\[
    \nabla f(x^\star)
    +
    J_h(x^\star)^{\mathsf T}\nu^\star
    =0.
\]
Эквивалентно,
\[
    \nabla_x\mathcal L(x^\star,\nu^\star)=0.
\]
Поскольку \(h(x^\star)=0\), добавление штрафа сохраняет
стационарность:
\[
    \nabla_x\mathcal L_c(x^\star,\nu^\star)=0.
\]

Для произвольной точки
\[
\begin{aligned}
    \nabla_{xx}^2\mathcal L_c(x,\nu)
    ={}&
    \nabla_{xx}^2\mathcal L(x,\nu)
    +
    cJ_h(x)^{\mathsf T}J_h(x)
    \\
    &+
    c\sum_{i=1}^{m}h_i(x)\nabla^2h_i(x).
\end{aligned}
\]
В допустимой точке последний член исчезает:
\[
\begin{aligned}
    \nabla_{xx}^2
    \mathcal L_c(x^\star,\nu^\star)
    ={}&
    \nabla_{xx}^2
    \mathcal L(x^\star,\nu^\star)
    \\
    &+
    cJ_h(x^\star)^{\mathsf T}J_h(x^\star).
\end{aligned}
\]

Касательное пространство линеаризованного ограничения равно
\[
    \mathcal T(x^\star)
    =
    \ker J_h(x^\star).
\]
Если \(d\in\mathcal T(x^\star)\), то
\[
    d^{\mathsf T}
    J_h(x^\star)^{\mathsf T}J_h(x^\star)d
    =
    \lVert J_h(x^\star)d\rVert^2
    =
    0.
\]
Следовательно, квадратичный штраф не добавляет кривизны в касательных
направлениях. Там положительность должна обеспечиваться исходной
функцией Лагранжа. В нормальных направлениях штраф добавляет
\[
    c\lVert J_h(x^\star)d\rVert^2.
\]

\begin{definition}[Достаточное условие второго порядка]
Пара \((x^\star,\nu^\star)\) удовлетворяет достаточному условию
второго порядка, если выполняются KKT-условия и
\[
    d^{\mathsf T}
    \nabla_{xx}^2\mathcal L(x^\star,\nu^\star)d
    >0
\]
для каждого ненулевого
\[
    d\in\ker J_h(x^\star).
\]
\end{definition}

\section{Локальная точность дополненной функции}

\begin{lemma}
Пусть \(P\) и \(Q\) симметричны, \(Q\succeq0\), а
\[
    y^{\mathsf T}Py>0
\]
для каждого ненулевого \(y\in\ker Q\). Тогда существует
\(\overline c\ge0\), такое что
\[
    P+cQ\succ0
\]
для всех \(c\ge\overline c\).
\end{lemma}

Матрица \(P\) уже положительна в направлениях, где \(Q\) обращается
в нуль. В остальных направлениях достаточно большой положительный
член \(cQ\) компенсирует возможную отрицательную кривизну \(P\). В
рассматриваемой задаче
\[
    P=
    \nabla_{xx}^2\mathcal L(x^\star,\nu^\star),
    \qquad
    Q=
    J_h(x^\star)^{\mathsf T}J_h(x^\star),
\]
причём
\[
    \ker Q=\ker J_h(x^\star).
\]

\begin{theorem}[Локальная точность]
Пусть \(f\) и \(h\) дважды непрерывно дифференцируемы, а пара
\((x^\star,\nu^\star)\) удовлетворяет KKT-условиям и достаточному
условию второго порядка. Тогда существует \(\overline c\ge0\), такое
что для каждого \(c\ge\overline c\) точка \(x^\star\) является
строгим локальным минимумом функции
\[
    x\longmapsto\mathcal L_c(x,\nu^\star).
\]
\end{theorem}

\begin{proof}
KKT-условия и допустимость дают
\[
    \nabla_x\mathcal L_c(x^\star,\nu^\star)=0.
\]
В этой точке
\[
\begin{aligned}
    \nabla_{xx}^2
    \mathcal L_c(x^\star,\nu^\star)
    ={}&
    \nabla_{xx}^2
    \mathcal L(x^\star,\nu^\star)
    \\
    &+
    cJ_h(x^\star)^{\mathsf T}J_h(x^\star).
\end{aligned}
\]
По достаточному условию второго порядка первая матрица положительна на
ядре второй. Лемма обеспечивает положительную определённость полного
Гессиана при всех достаточно больших \(c\). Поэтому \(x^\star\) —
строгий локальный минимум дополненной функции при фиксированном
множителе \(\nu^\star\).
\end{proof}

Теорема локальна. Она не утверждает, что для произвольного \(\nu\)
глобальный минимум \(\mathcal L_c(\cdot,\nu)\) совпадает с
\(x^\star\).

В выпуклой задаче с аффинным ограничением результат усиливается. Если
\((x^\star,\nu^\star)\) является седловой точкой обычной функции
Лагранжа, то
\[
    \mathcal L(x^\star,\nu)
    \le
    \mathcal L(x^\star,\nu^\star)
    \le
    \mathcal L(x,\nu^\star).
\]
Поскольку штраф неотрицателен и равен нулю в \(x^\star\),
\[
    \mathcal L_c(x^\star,\nu)
    \le
    \mathcal L_c(x^\star,\nu^\star)
    \le
    \mathcal L_c(x,\nu^\star).
\]
Следовательно, та же пара остаётся седловой точкой дополненной
функции Лагранжа.

\section{Метод штрафов и метод множителей}

Метод квадратичных штрафов решает последовательность задач
\[
    x^{k+1}
    \in
    \operatorname*{argmin}_x
    \left\{
        f(x)
        +
        \frac{c_k}{2}\lVert h(x)\rVert^2
    \right\},
    \qquad
    c_k\to\infty.
\]
При подходящих предпосылках \(h(x^k)\to0\), но рост \(c_k\) ухудшает
обусловленность. Для линейного ограничения Гессиан содержит
\(c_kA^{\mathsf T}A\): кривизна в нормальных направлениях быстро
растёт, а в касательных может почти не изменяться.

Метод множителей использует итерации
\[
    x^{k+1}
    \in
    \operatorname*{argmin}_x
    \mathcal L_c(x,\nu^k),
\]
\[
    \nu^{k+1}
    =
    \nu^k+c\,h(x^{k+1}).
\]
Второе равенство корректирует множитель пропорционально невязке.
Для \(h(x)=Ax-b\)
\[
    \nu^{k+1}
    =
    \nu^k+c(Ax^{k+1}-b).
\]
Двойственная переменная аккумулирует нарушение допустимости и меняет
следующую прямую подзадачу. В отличие от чистого метода штрафов,
параметр \(c\) часто можно оставить фиксированным. Постоянное увеличение \(c\) возможно, но не является обязательным свойством метода множителей.

Слишком большой штраф нежелателен. Гессиан прямой подзадачи содержит
\[
    cJ_h(x)^{\mathsf T}J_h(x),
\]
поэтому растут число обусловленности и липшицева константа градиента,
уменьшается допустимый шаг явных методов и усиливается чувствительность
к ошибкам вычисления невязки. На практике выбирают умеренный \(c\) и
корректируют его по прогрессу допустимости.

Логарифмический барьер для неравенства \(g_i(x)\le0\) имеет вид
\[
    -\tau\sum_i\log\bigl(-g_i(x)\bigr).
\]
Он определён только при строгой допустимости и не позволяет пересекать
границу. Квадратичный штраф определён по обе стороны равенства и
допускает недопустимые промежуточные точки. Для равенства
\(h(x)=0\) логарифмический барьер непосредственно неприменим, поскольку
допустимое множество не имеет внутренности.

\section{Границы модели и переход к ADMM}

Содержательная идея состоит в следующем: достаточно большой
штраф компенсирует отрицательную кривизну в направлениях, меняющих
ограничение, но не влияет на касательное пространство.

При анализе дополненной функции сначала проверяются KKT-условия, затем
определяется
\[
    \mathcal T(x^\star)=\ker J_h(x^\star),
\]
и на нём проверяется условие второго порядка. После этого формула
\[
\begin{aligned}
    \nabla_{xx}^2
    \mathcal L_c(x^\star,\nu^\star)
    ={}&
    \nabla_{xx}^2
    \mathcal L(x^\star,\nu^\star)
    \\
    &+
    cJ_h(x^\star)^{\mathsf T}J_h(x^\star)
\end{aligned}
\]
позволяет выбрать достаточно большое, но численно приемлемое значение
\(c\).

Если прямая переменная разбита на несколько блоков и минимизация
дополненной функции выполняется по ним последовательно, возникает
метод чередующихся направлений множителей, или ADMM. Блочная структура
разделяет исходную задачу на подзадачи и далее используется для
построения распределённых алгоритмов.

\lecture{23}{Метод множителей}

\sourcevideo
    {https://www.youtube.com/playlist?list=PLOzRYVm0a65fhKDS8cYIC7YCuVGJ4FOJx}
    {Week 7: Lecture 23: Method of Multipliers}

Рассмотрим задачу
\[
    \min_{x\in\mathbb R^n}f(x)
    \qquad
    \text{при условии}
    \qquad
    h(x)=0,
\]
где
\[
    h\colon\mathbb R^n\to\mathbb R^m.
\]
Дополненная функция Лагранжа имеет вид
\[
    \mathcal L_c(x,\nu)
    =
    f(x)
    +
    \nu^{\mathsf T}h(x)
    +
    \frac c2\lVert h(x)\rVert^2,
    \qquad
    c>0.
\]
Метод множителей чередует минимизацию по прямой переменной и
коррекцию двойственного множителя по невязке ограничения.

\section{Схема метода и стационарность}

При заданных \(\nu^k\) и \(c_k>0\) выполняется прямой шаг
\[
    x^{k+1}
    \in
    \operatorname*{argmin}_{x\in\mathbb R^n}
    \mathcal L_{c_k}(x,\nu^k),
\]
после чего множитель обновляется по правилу
\[
    \nu^{k+1}
    =
    \nu^k+c_kh(x^{k+1}).
\]
В лекции дополнительно рассматривается увеличение штрафа
\[
    c_{k+1}=\beta c_k,
    \qquad
    \beta>1,
\]
и упоминаются значения \(\beta\in[5,10]\). Такое быстрое увеличение
может ухудшать обусловленность, поэтому на практике часто используют
умеренный рост или фиксируют \(c_k=c\).

Прямая подзадача является безусловной:
\[
    \min_x
    \left\{
        f(x)
        +
        (\nu^k)^{\mathsf T}h(x)
        +
        \frac{c_k}{2}\lVert h(x)\rVert^2
    \right\}.
\]
Её можно решать градиентным, инерционным, ньютоновским или
квазиньютоновским методом. В теоретической записи минимизатор
находится точно, а в реализации допускается приближённое решение.

Предположим сначала, что прямой шаг точен. Тогда
\[
    \nabla_x\mathcal L_{c_k}(x^{k+1},\nu^k)=0.
\]
Поскольку
\[
\begin{aligned}
    \nabla_x\mathcal L_{c_k}(x,\nu^k)
    ={}&
    \nabla f(x)
    +
    J_h(x)^{\mathsf T}\nu^k
    \\
    &+
    c_kJ_h(x)^{\mathsf T}h(x),
\end{aligned}
\]
получаем
\[
    \nabla f(x^{k+1})
    +
    J_h(x^{k+1})^{\mathsf T}\nu^{k+1}
    =0.
\]
Таким образом, прямой шаг и обновление множителя автоматически дают
стационарность обычной функции Лагранжа. Для выполнения KKT-условий
остаётся обеспечить допустимость
\[
    h(x^{k+1})\to0.
\]
Из двойственного обновления непосредственно следует
\[
    h(x^{k+1})
    =
    \frac{\nu^{k+1}-\nu^k}{c_k}.
\]
При фиксированном \(c_k=c>0\) достаточно затухания приращений
множителя. При ограниченной последовательности \(\{\nu^k\}\) и
\(c_k\to\infty\) невязка также стремится к нулю при дополнительных
условиях существования предельных точек.

\section{Предельные точки и неточный прямой шаг}

При стандартных предпосылках сходимости, если \(f\) и \(h\)
непрерывны, допустимое множество непусто, прямые подзадачи решаются
глобально, последовательность множителей ограничена, а
\(c_k\to\infty\), то при стандартных предпосылках компактности всякая
предельная точка \(\{x^k\}\) является глобальным решением исходной
задачи. Возрастающий штраф не позволяет предельной точке сохранять
ненулевую невязку.

Для невыпуклой задачи глобальное решение прямой подзадачи обычно
недоступно. Если внутренний алгоритм находит лишь локальный минимум,
то в общем случае сохраняются только локальные гарантии.

При приближённой минимизации вместо точного равенства используется
условие
\[
    \left\lVert
        \nabla_x
        \mathcal L_{c_k}(x^{k+1},\nu^k)
    \right\rVert
    \le
    \varepsilon_k,
    \qquad
    \varepsilon_k\to0.
\]
На ранних внешних итерациях допустим грубый внутренний допуск, а по
мере сходимости точность увеличивается. Для сходимости требуется согласовать скорость убывания \(\varepsilon_k\) с выбором \(c_k\). При полном ранге якобиана
ограничения и подходящем убывании ошибки можно получить также
сходимость приближённых множителей
\[
    \nu^k+c_kh(x^{k+1})\to\nu^\star.
\]

Большой штраф улучшает допустимость, но может сделать внутреннюю
задачу плохо обусловленной. Для линейного ограничения
\[
    h(x)=Ax-b
\]
Гессиан прямой подзадачи содержит
\[
    \nabla^2f(x)+c_kA^{\mathsf T}A.
\]
Кривизна растёт только в направлениях, на которые действует \(A\),
и не меняется вдоль \(\ker A\). Поэтому увеличение \(c_k\) может
резко увеличить число обусловленности. Практически используют тёплый
запуск из предыдущей прямой точки, предобусловливание, ньютоновские
или квазиньютоновские методы и не увеличивают штраф без необходимости.

\section{Выпуклый квадратичный пример}

Рассмотрим задачу
\[
    \min_{x_1,x_2}
    \frac12(x_1^2+x_2^2)
    \qquad
    \text{при условии}
    \qquad
    x_1=1.
\]
Её решение и оптимальный множитель равны
\[
    x^\star=
    \begin{pmatrix}
        1\\
        0
    \end{pmatrix},
    \qquad
    \nu^\star=-1.
\]
Дополненная функция Лагранжа имеет вид
\[
\begin{aligned}
    \mathcal L_c(x,\nu)
    ={}&
    \frac12(x_1^2+x_2^2)
    +
    \nu(x_1-1)
    \\
    &+
    \frac c2(x_1-1)^2.
\end{aligned}
\]
Условия минимума по прямым переменным дают
\[
    x_1^{k+1}
    =
    \frac{c_k-\nu^k}{1+c_k},
    \qquad
    x_2^{k+1}=0.
\]
Следовательно,
\[
    x_1^{k+1}-1
    =
    -\frac{\nu^k+1}{1+c_k}.
\]
После двойственного обновления
\[
    \nu^{k+1}
    =
    \nu^k+c_k(x_1^{k+1}-1)
\]
получаем точную рекурсию
\[
    \nu^{k+1}-\nu^\star
    =
    \frac{\nu^k-\nu^\star}{1+c_k}.
\]
Для любого \(c_k>0\) коэффициент лежит в \((0,1)\), поэтому ошибка
множителя монотонно сокращается. При постоянном \(c_k=c\)
\[
    |\nu^k-\nu^\star|
    =
    \frac{|\nu^0-\nu^\star|}{(1+c)^k}.
\]
Рост штрафа до бесконечности в этом примере не нужен.

Гессиан прямой подзадачи равен
\[
    \nabla^2_{xx}\mathcal L_c
    =
    \begin{pmatrix}
        1+c&0\\
        0&1
    \end{pmatrix},
\]
поэтому
\[
    \kappa=1+c.
\]
Большое \(c\) уменьшает внешний коэффициент сжатия \(1/(1+c)\), но
одновременно ухудшает внутреннюю обусловленность. Это основной
компромисс метода множителей.

\section{Невыпуклый квадратичный пример}

Теперь рассмотрим задачу
\[
    \min_{x_1,x_2}
    \left\{
        -\frac12x_1^2+\frac12x_2^2
    \right\}
    \qquad
    \text{при условии}
    \qquad
    x_1=1.
\]
Хотя целевая функция невыпукла по \(x_1\), ограничение фиксирует эту
координату. Поэтому
\[
    x^\star=
    \begin{pmatrix}
        1\\
        0
    \end{pmatrix},
    \qquad
    \nu^\star=1.
\]
Дополненная функция Лагранжа равна
\[
\begin{aligned}
    \mathcal L_c(x,\nu)
    ={}&
    -\frac12x_1^2
    +
    \frac12x_2^2
    +
    \nu(x_1-1)
    \\
    &+
    \frac c2(x_1-1)^2.
\end{aligned}
\]
Её Гессиан имеет вид
\[
    \nabla^2_{xx}\mathcal L_c
    =
    \begin{pmatrix}
        c-1&0\\
        0&1
    \end{pmatrix}.
\]
Следовательно, прямая подзадача сильно выпукла только при
\[
    c>1.
\]
Для \(c_k>1\) прямой шаг равен
\[
    x_1^{k+1}
    =
    \frac{c_k-\nu^k}{c_k-1},
    \qquad
    x_2^{k+1}=0,
\]
а невязка имеет вид
\[
    x_1^{k+1}-1
    =
    -\frac{\nu^k-1}{c_k-1}.
\]
После двойственного обновления получаем
\[
    \nu^{k+1}-\nu^\star
    =
    -\frac{\nu^k-\nu^\star}{c_k-1}.
\]
Знак минус означает
чередование направления ошибки. Для сжатия по модулю необходимо
\[
    \frac1{c_k-1}<1,
    \qquad
    \text{то есть}
    \qquad
    c_k>2.
\]
При \(1<c_k<2\) прямая подзадача уже сильно выпукла, но внешняя
итерация расходится; при \(c_k=2\) ошибка только меняет знак; при
\(c_k>2\) возникает сжатие. Улучшение кривизны внутренней задачи само
по себе поэтому не гарантирует сходимость всего метода.

В общих задачах порог заранее может быть неизвестен. Штраф начинают с
умеренного значения и увеличивают только при недостаточном уменьшении
невязки. Слишком малый \(c_k\) слабо обеспечивает допустимость и может
не устранить отрицательную кривизну, а слишком большой ухудшает
обусловленность, уменьшает допустимый внутренний шаг и усиливает
численные ошибки.

\section{Неравенства и проекционное обновление}

Для задачи
\[
    \min_x f(x)
\]
при условиях
\[
    h(x)=0,
    \qquad
    g_j(x)\le0,
    \quad
    j=1,\ldots,r,
\]
каждое неравенство можно заменить равенством
\[
    g_j(x)+z_j^2=0,
    \qquad
    z_j\in\mathbb R.
\]
Такое равенство разрешимо по \(z_j\) тогда и только тогда, когда
\(g_j(x)\le0\). В векторной форме используется
\[
    g(x)+z\circ z=0.
\]
Дополненная функция Лагранжа принимает вид
\[
\begin{aligned}
    \mathcal L_c(x,z,\nu,\lambda)
    ={}&
    f(x)
    +
    \nu^{\mathsf T}h(x)
    +
    \frac c2\lVert h(x)\rVert^2
    \\
    &+
    \lambda^{\mathsf T}(g(x)+z\circ z)
    \\
    &+
    \frac c2
    \lVert g(x)+z\circ z\rVert^2.
\end{aligned}
\]
Прямой шаг выполняется по \((x,z)\), а множители обновляются как
\[
    \nu^{k+1}
    =
    \nu^k+c_kh(x^{k+1}),
\]
\[
\begin{aligned}
    \lambda^{k+1}
    ={}&
    \lambda^k
    +
    c_k
    \bigl(
        g(x^{k+1})
        +
        z^{k+1}\circ z^{k+1}
    \bigr).
\end{aligned}
\]

Переменные \(z_j\) можно исключить. Для одной компоненты минимизируется
по \(t=z_j^2\ge0\) выражение
\[
    \lambda_j(g_j(x)+t)
    +
    \frac c2(g_j(x)+t)^2.
\]
После исключения получается проекционное обновление
\[
    \lambda_j^{k+1}
    =
    \max\left\{
        0,
        \lambda_j^k+c_kg_j(x^{k+1})
    \right\},
\]
или в векторной форме
\[
    \lambda^{k+1}
    =
    \bigl[
        \lambda^k+c_kg(x^{k+1})
    \bigr]_+.
\]
Оно сохраняет \(\lambda^k\ge0\). В оптимальной точке выполняются
\[
    g_j(x^\star)\le0,
    \qquad
    \lambda_j^\star\ge0,
    \qquad
    \lambda_j^\star g_j(x^\star)=0.
\]
Если ограничение неактивно, то \(g_j(x^\star)<0\) и обязательно
\(\lambda_j^\star=0\). При активном ограничении множитель может быть
положительным или нулевым; в устном пояснении лекции эти случаи были
перепутаны.

\section{Реализация и двойственная интерпретация}

Практический алгоритм имеет внешний и внутренний циклы. При
фиксированных \(\nu^k\), \(\lambda^k\) и \(c_k\) приближённо
минимизируется дополненная функция, причём предыдущая прямая точка
используется как тёплый старт. Для равенств контролируются прямая и
стационарная невязки
\[
    r_{\mathrm p}^k=h(x^k),
\]
\[
    r_{\mathrm s}^k
    =
    \nabla f(x^k)
    +
    J_h(x^k)^{\mathsf T}\nu^k.
\]
При неравенствах дополнительно проверяются допустимость,
неотрицательность множителей и дополнительная нежёсткость
\[
    \lVert\lambda^k\circ g(x^k)\rVert
    \le
    \varepsilon_{\mathrm c}.
\]

Если
\[
    g_c(\nu)
    =
    \inf_x\mathcal L_c(x,\nu),
\]
то при единственности внутреннего минимизатора
\[
    \nabla g_c(\nu)=h(x(\nu)).
\]
Поэтому обновление
\[
    \nu^{k+1}
    =
    \nu^k+ch(x^{k+1})
\]
можно рассматривать как градиентный подъём по дополненной двойственной
функции с шагом \(c\). Квадратичное дополнение сглаживает двойственную
задачу, но может связать блоки прямой переменной. Последовательное
разрешение этих блоков приводит к ADMM.

\lecture{24}{Двойственный подъём и декомпозиция}

\sourcevideo
    {https://www.youtube.com/playlist?list=PLOzRYVm0a65fhKDS8cYIC7YCuVGJ4FOJx}
    {Week 7: Lecture 24: Dual Ascent and Dual Decomposition}

Рассматривается выпуклая задача с линейным равенством
\[
    \min_{x\in\mathbb R^n} f(x)
    \qquad
    \text{при условии}
    \qquad
    Ax=b.
\]
Двойственный подъём вычисляет градиент двойственной функции через
вспомогательную минимизацию по прямой переменной. Если функция и
переменные разделяются на блоки, эта минимизация распадается на
независимые локальные задачи.

\section{Сопряжённая функция и субградиенты}

Для собственной функции
\[
    f\colon\mathbb R^n\to\mathbb R\cup\{+\infty\}
\]
сопряжённая функция Фенхеля определяется формулой
\[
    f^*(y)
    =
    \sup_{x\in\mathbb R^n}
    \bigl\{\langle y,x\rangle-f(x)\bigr\}.
\]
Она всегда выпукла как поточечный супремум аффинных функций.
Эпиграф функции равен
\[
    \operatorname{epi}f
    =
    \bigl\{(x,t)\in\mathbb R^{n+1}:f(x)\le t\bigr\}.
\]
Функция называется замкнутой, если её эпиграф замкнут; для собственной
выпуклой функции это эквивалентно нижней полунепрерывности.

Повторное сопряжение задаётся формулой
\[
    f^{**}(x)
    =
    \sup_y\bigl\{\langle x,y\rangle-f^*(y)\bigr\}.
\]
По теореме Фенхеля--Моро для собственной замкнутой выпуклой функции
\[
    f^{**}=f.
\]
В общем случае \(f^{**}\) является замкнутой выпуклой оболочкой
исходной функции.

Из определения сопряжения следует неравенство Фенхеля--Юнга
\[
    f(x)+f^*(y)\ge \langle x,y\rangle.
\]
Для собственной замкнутой выпуклой функции эквивалентны условия
\[
    y\in\partial f(x),
    \qquad
    x\in\partial f^*(y),
    \qquad
    f(x)+f^*(y)=\langle x,y\rangle.
\]
При дифференцируемости обеих функций это даёт
\[
    y=\nabla f(x)
    \quad\Longleftrightarrow\quad
    x=\nabla f^*(y).
\]

\section{Двойственная функция}

Пусть
\[
    A\in\mathbb R^{m\times n},
    \qquad
    b\in\mathbb R^m.
\]
Функция Лагранжа и двойственная функция имеют вид
\[
    \mathcal L(x,\nu)
    =
    f(x)+\nu^{\mathsf T}(Ax-b),
    \qquad
    g(\nu)=\inf_x\mathcal L(x,\nu).
\]
Используя определение сопряжённой функции, получаем
\[
\begin{aligned}
    g(\nu)
    &=
    -b^{\mathsf T}\nu
    +
    \inf_x
    \bigl\{f(x)+\langle A^{\mathsf T}\nu,x\rangle\bigr\}
    \\
    &=
    -f^*(-A^{\mathsf T}\nu)-b^{\mathsf T}\nu.
\end{aligned}
\]
Двойственная задача состоит в максимизации \(g\) по
\(\nu\in\mathbb R^m\).

Предположим, что для каждого \(\nu\) минимум достигается в единственной
точке
\[
    x(\nu)
    =
    \operatorname*{argmin}_x
    \bigl\{f(x)+\nu^{\mathsf T}(Ax-b)\bigr\}.
\]
По теореме Данскина
\[
    \nabla g(\nu)=Ax(\nu)-b.
\]
Таким образом, двойственный градиент совпадает с невязкой прямого
ограничения. Явную формулу для \(f^*\) вычислять не требуется:
достаточно решить задачу минимизации по \(x\).

При нескольких минимизаторах функция \(g\) может быть
недифференцируемой. Для любого
\[
    x\in\operatorname*{argmin}_z\mathcal L(z,\nu)
\]
вектор \(Ax-b\) является суперградиентом вогнутой функции \(g\).
Сильная выпуклость \(f\) гарантирует единственность \(x(\nu)\) и
дифференцируемость двойственной функции.

\section{Двойственный подъём и скорость сходимости}

Одна итерация двойственного подъёма состоит из прямой минимизации
\[
    x^{k+1}
    =
    \operatorname*{argmin}_x
    \bigl\{f(x)+(\nu^k)^{\mathsf T}(Ax-b)\bigr\}
\]
и обновления
\[
    \nu^{k+1}
    =
    \nu^k+\alpha_k\bigl(Ax^{k+1}-b\bigr),
    \qquad
    \alpha_k>0.
\]
Поскольку максимизируется вогнутая функция, это именно подъём; запись
как спуска относится к минимизации \(-g\).

Если \(f\) является \(\mu\)-сильно выпуклой, то \(f^*\) имеет
\(1/\mu\)-липшицев градиент. Поэтому
\[
    \lVert\nabla g(\nu_1)-\nabla g(\nu_2)\rVert
    \le
    \frac{\lVert A\rVert^2}{\mu}
    \lVert\nu_1-\nu_2\rVert.
\]
Обозначим
\[
    L_d=\frac{\lVert A\rVert^2}{\mu}.
\]
Для шага \(\alpha=1/L_d\) и существующего двойственного решения
\(\nu^*\) стандартная оценка гладкого градиентного подъёма даёт
\[
    g^*-g(\nu^k)
    \le
    \frac{L_d\lVert\nu^0-\nu^*\rVert^2}{2k}.
\]
Следовательно, достижение двойственной точности \(\varepsilon\)
требует порядка \(O(\varepsilon^{-1})\) итераций.

Если \(f\) дополнительно \(L\)-гладка, а \(A\) имеет полный строчный
ранг, то \(-g\) сильно выпукла с константой
\[
    \mu_d
    =
    \frac{\sigma_{\min}(A)^2}{L}.
\]
При шаге
\[
    \alpha=\frac{2}{L_d+\mu_d}
\]
получаем линейное сжатие
\[
    \lVert\nu^{k+1}-\nu^*\rVert
    \le
    \frac{L_d-\mu_d}{L_d+\mu_d}
    \lVert\nu^k-\nu^*\rVert.
\]
Двойственное число обусловленности равно
\[
    \kappa_d
    =
    \frac{L_d}{\mu_d}
    =
    \frac{L\lVert A\rVert^2}
         {\mu\sigma_{\min}(A)^2},
\]
поэтому скорость зависит и от геометрии функции, и от обусловленности
линейного ограничения.

На практике прямая подзадача часто решается приближённо. Если вместо
\(x(\nu^k)\) используется \(\widetilde x^{k+1}\), то обновление
строится по вектору
\[
    A\widetilde x^{k+1}-b.
\]
Для сохранения сходимости погрешности внутренних решений должны
убывать или быть суммируемыми; точное условие зависит от выбранной
внешней схемы.

\section{Блочно-разделимая задача}

Рассмотрим задачу
\[
\begin{aligned}
    \min_{x_1,\ldots,x_B}\quad
    &\sum_{i=1}^{B}f_i(x_i),
    \\
    \text{при условии}\quad
    &\sum_{i=1}^{B}A_ix_i=b,
\end{aligned}
\]
где
\[
    x_i\in\mathbb R^{n_i},
    \qquad
    A_i\in\mathbb R^{m\times n_i}.
\]
Её функция Лагранжа разделяется:
\[
    \mathcal L(x_1,\ldots,x_B,\nu)
    =
    \sum_{i=1}^{B}
    \bigl(f_i(x_i)+\nu^{\mathsf T}A_ix_i\bigr)
    -\nu^{\mathsf T}b.
\]
Поэтому при фиксированном \(\nu\) минимизация по блокам независима,
а двойственная функция равна
\[
    g(\nu)
    =
    -\nu^{\mathsf T}b
    -
    \sum_{i=1}^{B}f_i^*(-A_i^{\mathsf T}\nu).
\]

На итерации \(k\) каждый агент параллельно вычисляет
\[
    x_i^{k+1}
    =
    \operatorname*{argmin}_{x_i}
    \bigl\{f_i(x_i)+(\nu^k)^{\mathsf T}A_ix_i\bigr\}.
\]
После этого координатор формирует невязку
\[
    r^{k+1}
    =
    \sum_{i=1}^{B}A_ix_i^{k+1}-b
\]
и обновляет общий множитель
\[
    \nu^{k+1}
    =
    \nu^k+\alpha_k r^{k+1}.
\]
Именно распад прямой минимизации после перехода к двойственной задаче
называется двойственной декомпозицией.

Разделимость нарушается, если целевая функция содержит член
\(\psi(x_i,x_j)\), связывающий разные блоки. Тогда связанные
переменные объединяют в один блок, вводят локальные копии с условиями
согласования или применяют другой метод декомпозиции. Например,
\[
    \frac12x_1^2+\frac12x_2^2+x_2x_3
    =
    f_1(x_1)+f_2(x_2,x_3),
\]
поэтому \(x_2\) и \(x_3\) образуют один блок. Сам этот пример
иллюстрирует только структуру разделения: без дополнительных условий
локальная функция не обязана быть выпуклой.

\section{Архитектура и коммуникации}

Алгоритм имеет схему «рассылка--сбор». Координатор передаёт всем
агентам \(\nu^k\), агенты параллельно решают локальные задачи и
возвращают векторы \(A_ix_i^{k+1}\), после чего координатор вычисляет
общую невязку и новый множитель. При ограничении размерности \(m\)
объём передачи в каждом направлении имеет порядок \(O(Bm)\) скалярных
чисел за синхронную итерацию. Размер локальной переменной \(n_i\) не
обязан совпадать с размером сообщения.

Передача только \(A_ix_i\), а не полной функции \(f_i\) или вектора
\(x_i\), уменьшает объём раскрываемой информации, но не обеспечивает
конфиденциальность. При известной матрице \(A_i\) и малом ядре
локальная переменная может быть частично восстановлена; повторные
сообщения также несут информацию о функции. Формальные гарантии
требуют дополнительных механизмов защиты.

Несмотря на параллельные вычисления, схема остаётся централизованной:
координатор агрегирует сообщения, обновляет множитель и является
единой точкой отказа. Синхронная версия ждёт самого медленного агента;
асинхронные варианты допускают устаревшие сообщения, но требуют
отдельного анализа. Двойственную переменную можно интерпретировать как
цену общего ресурса: положительный дисбаланс увеличивает цену, а
недоиспользование уменьшает её.

\section{Связь с дополненной функцией Лагранжа}

Обычный двойственный подъём использует
\[
    \mathcal L(x,\nu)
    =
    f(x)+\nu^{\mathsf T}(Ax-b).
\]
Метод множителей заменяет её дополненной функцией
\[
    \mathcal L_c(x,\nu)
    =
    f(x)+\nu^{\mathsf T}(Ax-b)
    +
    \frac c2\lVert Ax-b\rVert^2.
\]
Квадратичный штраф улучшает свойства двойственного шага, но для
блочного ограничения создаёт перекрёстные члены
\[
    x_i^{\mathsf T}A_i^{\mathsf T}A_jx_j,
    \qquad
    i\ne j,
\]
и тем самым обычно разрушает независимость локальных подзадач. ADMM
использует дополненную функцию Лагранжа, но минимизирует её по блокам
последовательно, сохраняя часть преимуществ декомпозиции.

Двойственная декомпозиция требует существования решений локальных
задач и подходящего выбора шага. При неединственности локальных
минимизаторов сходимость двойственных переменных не обязательно сразу
влечёт сходимость каждого прямого вектора. Условие консенсуса также
можно записать как линейное равенство, однако полностью
децентрализованные методы используют структуру графа и обходятся без
единого координатора.

\lecture{25}{Метод чередующихся направлений множителей}

\sourcevideo
    {https://www.youtube.com/playlist?list=PLOzRYVm0a65fhKDS8cYIC7YCuVGJ4FOJx}
    {Week 7: Lecture 25: ADMM Algorithm}

Метод чередующихся направлений множителей (ADMM) применяется к
задачам, в которых целевая функция разделяется по блокам, а переменные
связаны линейным равенством. Метод сохраняет квадратичный штраф
дополненной функции Лагранжа, но вместо совместной минимизации по всем
переменным выполняет последовательные блочные шаги.

\section{Двухблочная задача и дополненная функция Лагранжа}

Рассмотрим выпуклую задачу
\[
\begin{aligned}
    \min_{x\in\mathbb R^n,\,z\in\mathbb R^p}\quad
    &f(x)+g(z),
    \\
    \text{при условии}\quad
    &Ax+Bz=c,
\end{aligned}
\]
где
\[
    A\in\mathbb R^{m\times n},
    \qquad
    B\in\mathbb R^{m\times p},
    \qquad
    c\in\mathbb R^m.
\]
Прямая невязка равна
\[
    r(x,z)=Ax+Bz-c.
\]
Для двойственной переменной \(y\in\mathbb R^m\) обычная функция
Лагранжа имеет вид
\[
    \mathcal L(x,z,y)
    =
    f(x)+g(z)+y^{\mathsf T}(Ax+Bz-c).
\]
При параметре \(\rho>0\) дополненная функция Лагранжа определяется
формулой
\[
\begin{aligned}
    \mathcal L_\rho(x,z,y)
    ={}&
    f(x)+g(z)+y^{\mathsf T}(Ax+Bz-c)
    \\
    &+
    \frac{\rho}{2}\lVert Ax+Bz-c\rVert^2.
\end{aligned}
\]
На допустимом множестве штраф равен нулю, поэтому он не меняет
исходную целевую функцию. При этом квадрат невязки улучшает кривизну
в направлениях, видимых отображением
\((x,z)\mapsto Ax+Bz\), но не создаёт кривизну в его ядре и потому не
делает произвольную невыпуклую задачу глобально сильно выпуклой.

Метод множителей совместно вычисляет
\[
    (x^{k+1},z^{k+1})
    \in
    \operatorname*{argmin}_{x,z}
    \mathcal L_\rho(x,z,y^k)
\]
и затем обновляет
\[
    y^{k+1}
    =
    y^k+\rho(Ax^{k+1}+Bz^{k+1}-c).
\]
Совместная подзадача обычно не разделяется, поскольку раскрытие
квадрата порождает перекрёстный член
\(2x^{\mathsf T}A^{\mathsf T}Bz\). В обычной двойственной
декомпозиции такого штрафа нет, поэтому блоки разделяются, но
теряется регуляризация нарушения ограничения.

\section{Итерация ADMM}

ADMM заменяет совместную минимизацию двумя последовательными шагами:
\[
\begin{aligned}
    x^{k+1}
    &\in
    \operatorname*{argmin}_{x}
    \mathcal L_\rho(x,z^k,y^k),
    \\
    z^{k+1}
    &\in
    \operatorname*{argmin}_{z}
    \mathcal L_\rho(x^{k+1},z,y^k),
    \\
    y^{k+1}
    &=
    y^k+\rho(Ax^{k+1}+Bz^{k+1}-c).
\end{aligned}
\]
После удаления слагаемых, не зависящих от оптимизируемого блока,
прямые подзадачи принимают вид
\[
\begin{aligned}
    x^{k+1}
    \in
    \operatorname*{argmin}_{x}
    \Bigl\{
        f(x)
        +(y^k)^{\mathsf T}Ax
        +\frac{\rho}{2}
        \lVert Ax+Bz^k-c\rVert^2
    \Bigr\},
\end{aligned}
\]
\[
\begin{aligned}
    z^{k+1}
    \in
    \operatorname*{argmin}_{z}
    \Bigl\{
        g(z)
        +(y^k)^{\mathsf T}Bz
        +\frac{\rho}{2}
        \lVert Ax^{k+1}+Bz-c\rVert^2
    \Bigr\}.
\end{aligned}
\]
Каждая задача содержит только один переменный блок, однако второй шаг
использует уже вычисленное значение \(x^{k+1}\). Поэтому классический
двухблочный ADMM является схемой Гаусса--Зейделя, а не параллельной
схемой Якоби. Распределённая реализация возможна за счёт размещения
блоков на разных узлах и внутреннего распараллеливания каждой
подзадачи, но сами два блочных обновления остаются последовательными.

\section{Масштабированная форма}

Введём масштабированную двойственную переменную
\[
    u^k=\frac{y^k}{\rho}.
\]
Для \(r=Ax+Bz-c\) справедливо тождество
\[
    y^{\mathsf T}r+\frac{\rho}{2}\lVert r\rVert^2
    =
    \frac{\rho}{2}\lVert r+u\rVert^2
    -
    \frac{\rho}{2}\lVert u\rVert^2.
\]
Последнее слагаемое не зависит от прямых переменных. Поэтому ADMM
можно записать как
\[
\begin{aligned}
    x^{k+1}
    &\in
    \operatorname*{argmin}_{x}
    \Bigl\{
        f(x)
        +\frac{\rho}{2}
        \lVert Ax+Bz^k-c+u^k\rVert^2
    \Bigr\},
    \\
    z^{k+1}
    &\in
    \operatorname*{argmin}_{z}
    \Bigl\{
        g(z)
        +\frac{\rho}{2}
        \lVert Ax^{k+1}+Bz-c+u^k\rVert^2
    \Bigr\},
    \\
    u^{k+1}
    &=
    u^k+Ax^{k+1}+Bz^{k+1}-c.
\end{aligned}
\]
Если
\[
    r^{k+1}=Ax^{k+1}+Bz^{k+1}-c,
\]
то
\[
    u^{k+1}=u^k+r^{k+1},
    \qquad
    r^{k+1}=u^{k+1}-u^k.
\]
Таким образом, \(u^k\) аккумулирует предыдущие нарушения ограничения,
но обычно сходится не к нулю, а к \(y^\star/\rho\). Поэтому качество
итерации оценивают по невязкам, а не по норме самой масштабированной
переменной.

\section{Условия оптимальности и невязки}

Для выпуклой задачи седловая точка
\((x^\star,z^\star,y^\star)\) удовлетворяет условиям
\[
    0\in\partial f(x^\star)+A^{\mathsf T}y^\star,
\]
\[
    0\in\partial g(z^\star)+B^{\mathsf T}y^\star,
\]
\[
    Ax^\star+Bz^\star-c=0.
\]
Условие оптимальности \(z\)-подзадачи вместе с двойственным
обновлением даёт точное соотношение
\[
    0
    \in
    \partial g(z^{k+1})+B^{\mathsf T}y^{k+1}.
\]
Для \(x\)-подзадачи получаем
\[
\begin{aligned}
    0
    \in{}&
    \partial f(x^{k+1})+A^{\mathsf T}y^{k+1}
    \\
    &-
    \rho A^{\mathsf T}B(z^{k+1}-z^k).
\end{aligned}
\]
Поэтому двойственную невязку определяют как
\[
    s^{k+1}
    =
    \rho A^{\mathsf T}B(z^{k+1}-z^k).
\]
Прямая невязка измеряет нарушение линейного ограничения, а
двойственная --- отклонение от стационарности по блоку \(x\). При
\[
    r^k\to0,
    \qquad
    s^k\to0
\]
условия KKT выполняются асимптотически.

Практический критерий остановки имеет вид
\[
    \lVert r^k\rVert\le\varepsilon_{\mathrm p},
    \qquad
    \lVert s^k\rVert\le\varepsilon_{\mathrm d}.
\]
Если штраф \(\rho\) изменяется, масштабированную переменную следует
пересчитать так, чтобы немасштабированный множитель
\(y=\rho u\) сохранился.

\section{Сходимость и выбор штрафа}

\begin{theorem}[Базовая сходимость ADMM]
Пусть \(f\) и \(g\) --- собственные замкнутые выпуклые функции,
седловая точка обычной функции Лагранжа существует, обе прямые
подзадачи решаются точно и \(\rho>0\) фиксировано. Тогда
\[
    r^k\to0,
\]
\[
    f(x^k)+g(z^k)\to p^\star,
\]
а двойственная последовательность сходится к двойственному решению:
\[
    y^k\to y^\star.
\]
\end{theorem}

Для этих выводов полный ранг \(A\) и \(B\) не требуется. Однако
сходимость самих последовательностей \(x^k\) и \(z^k\),
единственность решений подзадач и единственность прямого предела могут
потребовать сильной выпуклости, коэрцитивности, ранговых условий или
единственности оптимального решения. Классическая теорема не
распространяется автоматически на произвольные невыпуклые задачи и
многоблочные схемы с тремя и более последовательными обновлениями.

Большое \(\rho\) сильнее подавляет прямую невязку, но может ухудшить
обусловленность подзадач. Малое \(\rho\) делает прямые шаги мягче, но
может замедлить достижение допустимости. В отличие от чистого метода
квадратичных штрафов, для выпуклого двухблочного ADMM не требуется
предел \(\rho_k\to\infty\): базовая теория работает при любом
фиксированном \(\rho>0\). Адаптивное изменение штрафа допустимо лишь
в рамках правила, для которого проверены условия сходимости.

Если подзадачи решаются приближённо, ошибки их оптимальности должны
контролироваться. Одно из достаточных условий имеет вид
\[
    \sum_{k=0}^{\infty}\lVert e_x^{k+1}\rVert<\infty,
    \qquad
    \sum_{k=0}^{\infty}\lVert e_z^{k+1}\rVert<\infty,
\]
хотя точные требования зависят от варианта метода.

\section{Разбиение задачи и распределённая реализация}

Преимущество ADMM определяется выбранным разбиением. Оно полезно,
если обе блочные подзадачи существенно проще исходной совместной
задачи. При размещении \(x\) и \(z\) на разных узлах первый узел
вычисляет \(x^{k+1}\), после чего передаёт второму необходимый вклад,
например \(Ax^{k+1}\); второй узел вычисляет \(z^{k+1}\), а затем
обновляется двойственная переменная. Объём сообщений зависит от
структуры \(A\) и \(B\), а последовательность шагов требует
синхронизации.

В частном случае
\[
    A=I,
    \qquad
    B=-I,
    \qquad
    c=0
\]
ограничение имеет вид \(x=z\), а масштабированные обновления
становятся проксимальными:
\[
    x^{k+1}
    =
    \operatorname{prox}_{f/\rho}(z^k-u^k),
\]
\[
    z^{k+1}
    =
    \operatorname{prox}_{g/\rho}(x^{k+1}+u^k).
\]
Поэтому ADMM особенно эффективен, когда проксимальные операторы
отдельных частей вычисляются просто.

По сравнению с двойственным подъёмом ADMM сохраняет квадратичный
штраф, но теряет полную параллельность прямых обновлений. По сравнению
с методом множителей он заменяет одну связанную минимизацию двумя
блочными. Метод не гарантирует хорошую скорость при неудачном
разбиении, не определяет автоматически оптимальное значение \(\rho\)
и требует отдельной теории для невыпуклых, многоблочных и полностью
параллельных модификаций.

\lecture{26}{Основы теории графов}

\sourcevideo
    {https://www.youtube.com/playlist?list=PLOzRYVm0a65fhKDS8cYIC7YCuVGJ4FOJx}
    {Week 8: Lecture 26: Basics of Graph Theory}

В распределённой оптимизации агент использует собственную функцию и
сообщения, полученные от разрешённых соседей. Структуру обмена задаёт
граф: его вершины соответствуют агентам, а рёбра --- доступным каналам
связи. Ниже вводятся понятия, необходимые для записи консенсусных
ограничений и построения матрицы Лапласа.

\section{Распределённая постановка}

Пусть агент \(i\) знает функцию
\[
    f_i\colon\mathbb R^d\to\mathbb R,
    \qquad
    i=1,\ldots,N.
\]
Централизованная задача имеет вид
\[
    \min_{x\in\mathbb R^d}
    F(x),
    \qquad
    F(x)=\sum_{i=1}^{N}f_i(x).
\]
В машинном обучении локальная функция может быть эмпирическим риском
\[
    f_i(x)
    =
    \frac1{|\mathcal D_i|}
    \sum_{(\xi,y)\in\mathcal D_i}
    \ell(x;\xi,y).
\]
Локальные минимизаторы в общем случае различаются и не совпадают с
минимизатором суммы.

В распределённой формулировке агент хранит копию
\(x_i\in\mathbb R^d\):
\[
\begin{aligned}
    \min_{x_1,\ldots,x_N}\quad
    &\sum_{i=1}^{N}f_i(x_i),
    \\
    \text{при условии}\quad
    &x_1=\cdots=x_N.
\end{aligned}
\]
При выполнении консенсусного ограничения существует общий вектор
\(x\), для которого \(x_i=x\), поэтому задача эквивалентна исходной.
Одна итерация распределённого метода сочетает локальный
оптимизационный шаг, например
\[
    x_i^k-\alpha_k\nabla f_i(x_i^k),
\]
и обмен векторами с соседями. Стоимость алгоритма определяется не
только вычислениями, но и числом коммуникационных раундов, размером
сообщений и задержками.

Передача локальных данных может не требоваться, однако обмен
параметрами и градиентами сам по себе не обеспечивает формальной
конфиденциальности. Для неё нужны дополнительные механизмы, например
защищённая агрегация или дифференциальная приватность.

\section{Графы, рёбра и соседи}

\begin{definition}[Граф]
Графом называется пара
\[
    \mathcal G=(\mathcal V,\mathcal E),
\]
где \(\mathcal V=\{1,\ldots,N\}\) --- множество вершин, а
\(\mathcal E\) --- множество рёбер.
\end{definition}

В неориентированном графе ребро \(\{i,j\}\) разрешает двусторонний
обмен между агентами \(i\) и \(j\). В ориентированном графе ребро
\((i,j)\) разрешает передачу от \(i\) к \(j\), тогда как обратное
ребро может отсутствовать. Во взвешенном графе ребру сопоставляется
вес \(a_{ij}>0\), характеризующий интенсивность взаимодействия,
доверие или коэффициент усреднения. Для неориентированного графа
обычно предполагается
\[
    a_{ij}=a_{ji}.
\]

\begin{definition}[Соседи]
Для неориентированного графа множество соседей вершины \(i\) равно
\[
    \mathcal N_i
    =
    \{j\in\mathcal V:\{i,j\}\in\mathcal E\}.
\]
Для ориентированного графа определяются
\[
    \mathcal N_i^{\mathrm{out}}
    =
    \{j:(i,j)\in\mathcal E\},
    \qquad
    \mathcal N_i^{\mathrm{in}}
    =
    \{j:(j,i)\in\mathcal E\}.
\]
\end{definition}

Агент может непосредственно использовать состояния только своих
соседей. Информация от удалённой вершины поступает через цепочку
промежуточных узлов. Для ориентированных графов необходимо заранее
зафиксировать соглашение о матрице смежности. Далее используется
\[
    a_{ij}>0
    \quad\Longleftrightarrow\quad
    (i,j)\in\mathcal E,
\]
то есть строка \(i\) описывает исходящие рёбра вершины \(i\).

\section{Пути и связность}

\begin{definition}[Путь]
Путём из \(i\) в \(j\) называется последовательность
\[
    i=v_0,v_1,\ldots,v_\ell=j,
\]
в которой каждая соседняя пара соединена ребром. Число \(\ell\)
называется длиной пути. В ориентированном графе направления рёбер
должны совпадать с направлением движения.
\end{definition}

\begin{definition}[Связность]
Неориентированный граф называется связным, если путь существует между
любыми двумя вершинами. Максимальный связный подграф называется
компонентой связности.
\end{definition}

Если граф несвязен, информация не передаётся между его компонентами.
Поэтому агенты могут согласовать значения внутри каждой компоненты,
но в общем случае не могут решить задачу для полной суммы
\(\sum_i f_i\). Связность является минимальным условием глобального
консенсуса в статической неориентированной сети.

Ориентированный граф называется сильно связным, если ориентированный
путь существует для каждой упорядоченной пары вершин. Он называется
слабо связным, если после удаления направлений получается связный
неориентированный граф. Для двустороннего распространения информации
слабой связности обычно недостаточно.

В изменяющейся сети используется последовательность
\[
    \mathcal G^k=(\mathcal V,\mathcal E^k).
\]
Отдельный граф может быть несвязным. Одно из стандартных условий
требует существования \(T\ge1\), для которого объединённый граф
\[
    \left(
        \mathcal V,
        \bigcup_{t=k}^{k+T-1}\mathcal E^t
    \right)
\]
связен для каждого \(k\). Тогда информация распространяется по всей
сети в течение каждого окна длины \(T\).

\section{Консенсус по рёбрам}

Для графа локальное согласование записывается условиями
\[
    x_i=x_j
    \qquad
    \text{для всех }\{i,j\}\in\mathcal E.
\]

\begin{theorem}[Локальные равенства и глобальный консенсус]
Пусть неориентированный граф \(\mathcal G\) связен. Тогда равенства
на всех рёбрах эквивалентны условию
\[
    x_1=\cdots=x_N.
\]
\end{theorem}

\begin{proof}
Глобальный консенсус очевидно влечёт равенство на каждом ребре.
Обратно, возьмём любые вершины \(i\) и \(j\). По связности существует
путь
\[
    i=v_0,v_1,\ldots,v_\ell=j.
\]
На каждом его ребре выполнено
\(x_{v_{r-1}}=x_{v_r}\), поэтому
\[
    x_i=x_{v_1}=\cdots=x_j.
\]
Поскольку пара \(i,j\) произвольна, все локальные переменные равны.
\end{proof}

Следовательно, на связном графе распределённую задачу можно записать
как
\[
\begin{aligned}
    \min_{x_1,\ldots,x_N}\quad
    &\sum_{i=1}^{N}f_i(x_i),
    \\
    \text{при условиях}\quad
    &x_i=x_j,
    \qquad
    \{i,j\}\in\mathcal E.
\end{aligned}
\]
Каждое ограничение связывает только соседних агентов, что делает
возможной полностью децентрализованную реализацию. Если граф несвязен,
такие равенства обеспечивают консенсус лишь внутри компонент.

\section{Матрицы смежности и степеней}

\begin{definition}[Матрица смежности]
Матрицей смежности называется матрица
\[
    A=(a_{ij})_{i,j=1}^{N},
\]
элементы которой кодируют наличие и вес рёбер. Для простого
неориентированного невзвешенного графа
\[
    a_{ij}
    =
    \begin{cases}
        1,& \{i,j\}\in\mathcal E,\\
        0,& \text{иначе}.
    \end{cases}
\]
При отсутствии петель \(a_{ii}=0\), а неориентированность даёт
\[
    A=A^{\mathsf T}.
\]
\end{definition}

\begin{definition}[Степень]
В неориентированном графе степень вершины равна
\[
    d_i=|\mathcal N_i|.
\]
Во взвешенном случае используется
\[
    d_i^{\mathrm w}=\sum_{j=1}^{N}a_{ij}.
\]
Матрица степеней определяется формулой
\[
    D=\operatorname{diag}(d_1,\ldots,d_N).
\]
\end{definition}

Для неориентированного графа степень равна сумме строки и, благодаря
симметрии, сумме столбца:
\[
    d_i=\sum_j a_{ij}=\sum_j a_{ji}.
\]
Для ориентированного графа при принятом соглашении
\[
    d_i^{\mathrm{out}}=\sum_j a_{ij},
    \qquad
    d_i^{\mathrm{in}}=\sum_j a_{ji}.
\]
Соответственно вводятся диагональные матрицы
\(D^{\mathrm{out}}\) и \(D^{\mathrm{in}}\).

Рассмотрим неориентированный граф с рёбрами
\[
\begin{aligned}
    \mathcal E
    =
    \bigl\{
        &\{1,2\},\{2,3\},\{2,4\},
        \{3,4\},\\
        &\{4,5\},\{4,6\}
    \bigr\}.
\end{aligned}
\]
Его степени равны
\[
    (d_1,\ldots,d_6)=(1,3,2,4,1,1),
\]
а матрицы имеют вид
\[
    D
    =
    \begin{pmatrix}
        1&0&0&0&0&0\\
        0&3&0&0&0&0\\
        0&0&2&0&0&0\\
        0&0&0&4&0&0\\
        0&0&0&0&1&0\\
        0&0&0&0&0&1
    \end{pmatrix},
\]
\[
    A
    =
    \begin{pmatrix}
        0&1&0&0&0&0\\
        1&0&1&1&0&0\\
        0&1&0&1&0&0\\
        0&1&1&0&1&1\\
        0&0&0&1&0&0\\
        0&0&0&1&0&0
    \end{pmatrix}.
\]
Суммы строк \(A\) совпадают с диагональными элементами \(D\).

\section{Топология и коммуникационная стоимость}

В полном неориентированном графе
\[
    d_i=N-1,
    \qquad
    |\mathcal E|=\frac{N(N-1)}2.
\]
Он обеспечивает прямой обмен между любыми агентами, но требует
квадратичного числа каналов. Разреженный связный граф уменьшает
стоимость одной итерации, сохраняя возможность распространения
информации по всей сети.

Если агент \(i\) передаёт соседям вектор из \(\mathbb R^d\), объём
его обмена за раунд имеет порядок
\[
    O(d_i d),
\]
а суммарный объём по сети --- порядок
\[
    O(|\mathcal E|d),
\]
с точностью до соглашения о подсчёте двух направлений передачи по
неориентированному ребру. Более плотная сеть может ускорять
распространение информации, но повышает стоимость раунда.

Матрица смежности и матрица степеней образуют матрицу Лапласа
\[
    L=D-A.
\]
Она кодирует структуру рёбер, степени, связность и консенсусное
подпространство. Её спектральные свойства и применение в
распределённых алгоритмах рассматриваются в следующей лекции.

\lecture{27}{Матрица Лапласа и спектральная связность}

\sourcevideo
    {https://www.youtube.com/playlist?list=PLOzRYVm0a65fhKDS8cYIC7YCuVGJ4FOJx}
    {Week 8: Lecture 27: Basics of Graph Theory-2}

Пусть неориентированная взвешенная коммуникационная сеть задана
графом
\[
    \mathcal G=(\mathcal V,\mathcal E),
    \qquad
    \mathcal V=\{1,\ldots,N\}.
\]
Её матрица смежности равна
\[
    A=(a_{ij})_{i,j=1}^{N},
    \qquad
    a_{ij}=a_{ji}\ge0,
\]
а матрица степеней имеет вид
\[
    D=\operatorname{diag}(d_1,\ldots,d_N),
    \qquad
    d_i=\sum_{j=1}^{N}a_{ij}.
\]
Матрица Лапласа объединяет локальную структуру рёбер и глобальные
свойства связности. Её спектр определяет скорость консенсусных и
распределённых алгоритмов.

\section{Лапласиан и его квадратичная форма}

\begin{definition}[Матрица Лапласа]
Матрицей Лапласа неориентированного взвешенного графа называется
\[
    L=D-A.
\]
\end{definition}

Её элементы равны
\[
    L_{ij}
    =
    \begin{cases}
        d_i, & i=j,\\
        -a_{ij}, & i\ne j.
    \end{cases}
\]
Для простого невзвешенного графа ненулевые внедиагональные элементы
равны \(-1\). Поскольку \(A=A^{\mathsf T}\), матрица \(L\)
симметрична. Сумма элементов каждой строки равна нулю:
\[
    \sum_{j=1}^{N}L_{ij}
    =
    d_i-\sum_{j=1}^{N}a_{ij}
    =
    0.
\]
Следовательно,
\[
    L\mathbf1=0,
    \qquad
    \mathbf1^{\mathsf T}L=0,
\]
где \(\mathbf1\in\mathbb R^N\) --- вектор из единиц. Поэтому нуль
всегда является собственным значением Лапласиана.

Для произвольного \(x\in\mathbb R^N\)
\[
\begin{aligned}
    (Lx)_i
    &=
    d_ix_i-\sum_{j=1}^{N}a_{ij}x_j
    \\
    &=
    \sum_{j=1}^{N}a_{ij}(x_i-x_j).
\end{aligned}
\]
Таким образом, компоненту \((Lx)_i\) можно вычислить по состоянию
вершины \(i\) и состояниям её соседей.

Квадратичная форма имеет вид
\[
\begin{aligned}
    x^{\mathsf T}Lx
    &=
    \frac12
    \sum_{i,j=1}^{N}
    a_{ij}(x_i-x_j)^2
    \\
    &=
    \sum_{\{i,j\}\in\mathcal E}
    a_{ij}(x_i-x_j)^2,
\end{aligned}
\]
где во второй сумме каждое неориентированное ребро учитывается один
раз. Все слагаемые неотрицательны.

\begin{theorem}
Матрица Лапласа неориентированного графа положительно
полуопределена:
\[
    L\succeq0.
\]
\end{theorem}

Следовательно, все собственные значения \(L\) вещественны и
неотрицательны.

\section{Ядро и компоненты связности}

Из формулы для квадратичной формы следует
\[
    x^{\mathsf T}Lx=0
\]
тогда и только тогда, когда
\[
    x_i=x_j
\]
на каждом ребре положительного веса. Поскольку \(L\succeq0\), это
эквивалентно условию \(Lx=0\). Значит, любой вектор из ядра
Лапласиана постоянен на каждой компоненте связности.

Если граф связен, то
\[
    \ker L=\operatorname{span}\{\mathbf1\}.
\]
Если граф состоит из \(r\) компонент
\[
    \mathcal V_1,\ldots,\mathcal V_r,
\]
то
\[
    \ker L
    =
    \operatorname{span}
    \left\{
        \mathbf1_{\mathcal V_1},
        \ldots,
        \mathbf1_{\mathcal V_r}
    \right\},
\]
где \(\mathbf1_{\mathcal V_s}\) --- индикатор компоненты
\(\mathcal V_s\). Поэтому
\[
    \dim\ker L=r.
\]

\begin{theorem}
Кратность нулевого собственного значения Лапласиана равна числу
компонент связности графа.
\end{theorem}

После подходящей перенумерации вершин матрицы смежности, степеней и
Лапласа принимают блочно-диагональный вид:
\[
    L=
    \begin{pmatrix}
        L_1&0&\cdots&0\\
        0&L_2&\cdots&0\\
        \vdots&\vdots&\ddots&\vdots\\
        0&0&\cdots&L_r
    \end{pmatrix}.
\]
Каждый блок соответствует одной связной компоненте и имеет ровно одно
нулевое собственное значение.

\section{Алгебраическая связность и консенсус}

Упорядочим собственные значения:
\[
    0=\lambda_1(L)
    \le
    \lambda_2(L)
    \le
    \cdots
    \le
    \lambda_N(L).
\]
Граф связен тогда и только тогда, когда нулевое собственное значение
простое:
\[
    \mathcal G\text{ связен}
    \quad\Longleftrightarrow\quad
    \lambda_2(L)>0.
\]

\begin{definition}[Алгебраическая связность]
Второе наименьшее собственное значение
\[
    \lambda_2(L)
\]
называется алгебраической связностью, или собственным значением
Фидлера.
\end{definition}

Для симметричного Лапласиана справедлива вариационная формула Рэлея:
\[
    \lambda_2(L)
    =
    \min_{\substack{x\ne0\\x^{\mathsf T}\mathbf1=0}}
    \frac{x^{\mathsf T}Lx}{x^{\mathsf T}x}.
\]
Следовательно, для каждого \(x\perp\mathbf1\)
\[
    x^{\mathsf T}Lx
    \ge
    \lambda_2(L)\lVert x\rVert^2.
\]

Для произвольного \(x\in\mathbb R^N\) определим среднее
\[
    \overline x
    =
    \frac1N\mathbf1^{\mathsf T}x
\]
и рассогласование
\[
    x_\perp=x-\overline x\,\mathbf1.
\]
Тогда \(x_\perp\perp\mathbf1\), а из \(L\mathbf1=0\) следует
\[
    Lx=Lx_\perp,
    \qquad
    x^{\mathsf T}Lx
    =
    x_\perp^{\mathsf T}Lx_\perp.
\]
Для связного графа получаем основную оценку
\[
    x^{\mathsf T}Lx
    \ge
    \lambda_2(L)
    \lVert
        x-\overline x\,\mathbf1
    \rVert^2.
\]
Она показывает, насколько сумма локальных различий на рёбрах
контролирует глобальное расстояние до консенсуса. Малое
\(\lambda_2(L)\) указывает на слабую связь между частями сети, а
большое значение означает более сильное спектральное перемешивание.

Если каждый агент хранит вектор \(x_i\in\mathbb R^d\), то для
составного состояния
\[
    \mathbf x=
    \begin{pmatrix}
        x_1\\
        \vdots\\
        x_N
    \end{pmatrix}
    \in\mathbb R^{Nd}
\]
используется оператор \(L\otimes I_d\). При связном графе
\[
    \ker(L\otimes I_d)
    =
    \{
        \mathbf1\otimes v:
        v\in\mathbb R^d
    \},
\]
поэтому
\[
    (L\otimes I_d)\mathbf x=0
    \quad\Longleftrightarrow\quad
    x_1=\cdots=x_N.
\]

\section{Нормализованные Лапласианы}

Матрицу
\[
    L=D-A
\]
называют ненормализованным Лапласианом. Если все степени положительны,
определяются симметричный нормализованный Лапласиан
\[
    L_{\mathrm{sym}}
    =
    D^{-1/2}LD^{-1/2}
    =
    I-D^{-1/2}AD^{-1/2}
\]
и Лапласиан случайного блуждания
\[
    L_{\mathrm{rw}}
    =
    D^{-1}L
    =
    I-D^{-1}A.
\]
Они связаны преобразованием подобия
\[
    L_{\mathrm{rw}}
    =
    D^{-1/2}
    L_{\mathrm{sym}}
    D^{1/2}
\]
и потому имеют одинаковые собственные значения.

Ненормализованный Лапласиан учитывает абсолютные веса рёбер, поэтому
вершины большой степени сильнее влияют на квадратичную форму.
Нормализация сравнивает связи относительно степеней вершин. По этой
причине \(L_{\mathrm{sym}}\) и \(L_{\mathrm{rw}}\) используются в
спектральной кластеризации и моделях случайного блуждания.

При наличии изолированных вершин обычные матрицы \(D^{-1}\) и
\(D^{-1/2}\) не определены. Тогда нормализация задаётся только на
вершинах положительной степени или используется соглашение с
псевдообратной матрицей.

\section{Расстояние и диаметр графа}

\begin{definition}[Графовое расстояние]
Расстоянием между вершинами \(i\) и \(j\) связного графа называется
длина кратчайшего пути между ними:
\[
    \operatorname{dist}_{\mathcal G}(i,j).
\]
\end{definition}

\begin{definition}[Диаметр]
Диаметром связного графа называется величина
\[
    \operatorname{diam}(\mathcal G)
    =
    \max_{i,j\in\mathcal V}
    \operatorname{dist}_{\mathcal G}(i,j).
\]
\end{definition}

Для пути \(P_N\)
\[
    \operatorname{diam}(P_N)=N-1,
\]
а для звезды с \(N\ge3\)
\[
    \operatorname{diam}(\mathcal G)=2.
\]
Диаметр определяет минимальное число последовательных передач,
необходимых для доставки одного сообщения между наиболее удалёнными
вершинами.

Малый диаметр не означает, что стандартный итерационный алгоритм
усреднения достигает точного консенсуса за такое же число итераций.
Скорость численного алгоритма зависит также от весов и спектра
матрицы обновления. Диаметр является комбинаторной характеристикой, а
\(\lambda_2(L)\) --- спектральной. У хорошо связанных графов диаметр
обычно мал, а алгебраическая связность относительно велика, однако
универсального равенства между этими величинами нет.

\section{Динамика консенсуса и локальная реализация}

Рассмотрим непрерывную систему
\[
    \dot x=-Lx.
\]
Среднее состояние сохраняется, поскольку
\[
\begin{aligned}
    \frac{\mathrm d}{\mathrm dt}
    \left(
        \frac1N\mathbf1^{\mathsf T}x
    \right)
    &=
    -\frac1N\mathbf1^{\mathsf T}Lx
    =
    0.
\end{aligned}
\]
Следовательно, рассогласование
\[
    x_\perp=x-\overline x\,\mathbf1
\]
удовлетворяет уравнению
\[
    \dot x_\perp=-Lx_\perp.
\]
Для функции Ляпунова
\[
    V(x_\perp)
    =
    \frac12\lVert x_\perp\rVert^2
\]
получаем
\[
\begin{aligned}
    \dot V
    &=
    -x_\perp^{\mathsf T}Lx_\perp
    \\
    &\le
    -\lambda_2(L)
    \lVert x_\perp\rVert^2
    \\
    &=
    -2\lambda_2(L)V.
\end{aligned}
\]
Поэтому для связного графа
\[
    \lVert x_\perp(t)\rVert
    \le
    e^{-\lambda_2(L)t}
    \lVert x_\perp(0)\rVert.
\]
Собственное значение Фидлера тем самым непосредственно задаёт
гарантированную экспоненциальную скорость непрерывного консенсуса.

Явно хранить плотную матрицу \(L\) не требуется. Каждый агент
вычисляет
\[
    (Lx)_i
    =
    \sum_{j\in\mathcal N_i}
    a_{ij}(x_i-x_j)
\]
по сообщениям соседей. Для разреженного графа стоимость вычисления
всех компонент имеет порядок
\[
    O(\lvert\mathcal E\rvert),
\]
а для состояний размерности \(d\) ---
\[
    O(\lvert\mathcal E\rvert d).
\]
Именно сочетание локальной реализуемости и спектрального контроля
глобального рассогласования делает Лапласиан основным оператором
децентрализованной оптимизации.

\lecture{28}{Консенсус и средний консенсус}

\sourcevideo
    {https://www.youtube.com/playlist?list=PLOzRYVm0a65fhKDS8cYIC7YCuVGJ4FOJx}
    {Week 8: Lecture 28: Consensus and Average Consensus}

Пусть агенты обмениваются состояниями только с соседями по
коммуникационному графу. Консенсус означает сходимость всех состояний
к одному значению, а средний консенсус дополнительно требует, чтобы
это значение совпадало со средним начальных состояний. Ниже
рассматриваются распространение информации по графу, максимальный
консенсус и линейная модель Де Гроота.

\section{Распространение информации по графу}

Рассмотрим ориентированный граф
\[
    \mathcal G=(\mathcal V,\mathcal E),
    \qquad
    \mathcal V=\{1,\ldots,N\}.
\]
Зафиксируем соглашение: \(a_{ij}>0\), если агент \(i\) получает
информацию от агента \(j\). Тогда строка \(i\) матрицы смежности
\(A=(a_{ij})\) описывает входящие связи вершины \(i\), а
\[
    (Ax)_i=\sum_{j=1}^{N}a_{ij}x_j
\]
является взвешенной суммой состояний, доступных за один шаг. При
противоположном соглашении о направлении рёбер входящее
агрегирование записывается через \(A^{\mathsf T}x\).

Элемент степени матрицы имеет вид
\[
    (A^k)_{ij}
    =
    \sum_{i_1,\ldots,i_{k-1}}
    a_{ii_{k-1}}a_{i_{k-1}i_{k-2}}\cdots a_{i_1j}.
\]
Он равен суммарному весу ориентированных блужданий длины \(k\) от
\(j\) к \(i\); в невзвешенном графе это число таких блужданий.
Поэтому при линейном обновлении
\[
    x^{k+1}=Ax^k,
    \qquad
    x^k=A^kx^0,
\]
через \(k\) шагов на состояние вершины могут влиять начальные
значения из её \(k\)-шаговой входящей окрестности. Тот же принцип
лежит в основе линейных графовых слоёв
\[
    H^{\ell+1}
    =
    \sigma\!\left(\widetilde A H^\ell W^\ell\right),
\]
где между операциями агрегирования добавляются обучаемые матрицы,
нормировка и нелинейность.

\section{Консенсус и максимальный консенсус}

\begin{definition}[Консенсус]
Последовательность \(x^k\in\mathbb R^N\) достигает консенсуса, если
существует \(c\in\mathbb R\), для которого
\[
    x^k\longrightarrow c\mathbf1.
\]
\end{definition}

Обычный консенсус не фиксирует значение \(c\). Например, итерация
\(x^{k+1}=\tfrac12x^k\) приводит все состояния к нулю, но не сохраняет
информацию о начальных данных.

\begin{definition}[Средний консенсус]
Алгоритм достигает среднего консенсуса, если
\[
    x^k
    \longrightarrow
    \frac1N\mathbf1\mathbf1^{\mathsf T}x^0.
\]
\end{definition}

Для связного неориентированного графа максимальный консенсус задаётся
обновлением
\[
    x_i^{k+1}
    =
    \max\bigl\{x_i^k,\,x_j^k:j\in\mathcal N_i\bigr\}.
\]
Каждая последовательность \(x_i^k\) не убывает и не превосходит
начальный глобальный максимум. Если максимум первоначально находится
в вершине \(j^\star\), то после \(k\) шагов его знают все вершины на
расстоянии не более \(k\) от \(j^\star\).

\begin{theorem}[Конечное время максимального консенсуса]
На связном неориентированном графе
\[
    x_i^k=\max_j x_j^0
\]
для всех \(i\) и всех
\[
    k\ge\operatorname{diam}(\mathcal G).
\]
\end{theorem}

Аналогичное обновление с минимумом достигает минимального консенсуса.
Среднее устроено сложнее: необходимо одновременно распространить
информацию, сохранить её суммарную массу и устранить рассогласование.
Передача всех начальных значений позволяет получить точный ответ за
конечное время, но требует памяти и сообщений, растущих с \(N\).
Линейные алгоритмы используют сообщения фиксированной размерности и
обычно сходятся лишь асимптотически.

\section{Модель Де Гроота и стохастические матрицы}

В модели Де Гроота агент обновляет состояние по правилу
\[
    x_i^{k+1}=\sum_{j=1}^{N}w_{ij}x_j^k,
    \qquad
    x^{k+1}=Wx^k.
\]
Матрица \(W\) согласована с графом: \(w_{ij}=0\), если агент \(i\) не
получает сообщение от \(j\).

\begin{definition}[Строчная стохастичность]
Неотрицательная матрица \(W\) называется строчно-стохастической, если
\[
    W\mathbf1=\mathbf1.
\]
\end{definition}

Тогда каждое новое состояние является выпуклой комбинацией
предыдущих, поэтому
\[
    \min_jx_j^k
    \le x_i^{k+1}\le
    \max_jx_j^k.
\]
Однако строчной стохастичности недостаточно для сходимости. Матрица
\[
    W=
    \begin{pmatrix}
        0&1\\
        1&0
    \end{pmatrix}
\]
дважды стохастична, но \(W^2=I\), поэтому состояния периодически
меняются местами.

\begin{definition}[Примитивность]
Неотрицательная матрица \(W\) называется примитивной, если существует
\(m\ge1\), для которого \(W^m>0\) покомпонентно.
\end{definition}

Для стохастической матрицы примитивность исключает разложение сети и
периодические колебания. В графовых терминах она соответствует
сильной связности и апериодичности; на связном неориентированном графе
апериодичность обычно обеспечивается положительными диагональными
весами.

\begin{theorem}[Предел степеней]
Пусть \(W\) --- примитивная строчно-стохастическая матрица. Тогда
существует единственный вектор \(\pi>0\), удовлетворяющий
\[
    \pi^{\mathsf T}W=\pi^{\mathsf T},
    \qquad
    \mathbf1^{\mathsf T}\pi=1,
\]
и
\[
    W^k\longrightarrow\mathbf1\pi^{\mathsf T}.
\]
Следовательно,
\[
    x^k\longrightarrow
    \mathbf1\pi^{\mathsf T}x^0.
\]
\end{theorem}

Предел является взвешенным средним начальных состояний. Вес агента
\(i\) равен \(\pi_i\).

\section{Средний консенсус}

\begin{definition}[Двойная стохастичность]
Неотрицательная матрица \(W\) называется дважды стохастической, если
\[
    W\mathbf1=\mathbf1,
    \qquad
    \mathbf1^{\mathsf T}W=\mathbf1^{\mathsf T}.
\]
\end{definition}

Столбцовая стохастичность сохраняет сумму состояний:
\[
    \mathbf1^{\mathsf T}x^{k+1}
    =
    \mathbf1^{\mathsf T}Wx^k
    =
    \mathbf1^{\mathsf T}x^k.
\]
Если одновременно достигается консенсус \(x^k\to c\mathbf1\), то
\[
    Nc=\mathbf1^{\mathsf T}x^0,
\]
то есть \(c\) равно начальному среднему.

\begin{theorem}[Средний консенсус]
Если \(W\) примитивна и дважды стохастична, то
\[
    W^k
    \longrightarrow
    J
    :=
    \frac1N\mathbf1\mathbf1^{\mathsf T},
\]
и для любого \(x^0\)
\[
    x^k\longrightarrow Jx^0.
\]
\end{theorem}

\begin{proof}
Примитивность и строчная стохастичность дают
\(W^k\to\mathbf1\pi^{\mathsf T}\). По столбцовой стохастичности
\(N^{-1}\mathbf1\) является нормированным левым собственным вектором
для собственного значения \(1\). Единственность стационарного вектора
даёт \(\pi=N^{-1}\mathbf1\).
\end{proof}

Двойная стохастичность без примитивности недостаточна, как показывает
матрица перестановки. Если \(W=W^{\mathsf T}\) и
\(W\mathbf1=\mathbf1\), то после транспонирования получаем
\(\mathbf1^{\mathsf T}W=\mathbf1^{\mathsf T}\); поэтому симметричная
строчно-стохастическая матрица автоматически дважды стохастична.

Проекторы
\[
    J=\frac1N\mathbf1\mathbf1^{\mathsf T},
    \qquad
    J_\perp=I-J
\]
выделяют консенсусную часть \(Jx=\bar x\mathbf1\) и рассогласование
\(J_\perp x=x-\bar x\mathbf1\). Для дважды стохастической матрицы
\(WJ=JW=J\), поэтому
\[
    J_\perp x^{k+1}
    =
    (W-J)J_\perp x^k.
\]
Средний консенсус достигается тогда, когда оператор \(W-J\) сжимает
подпространство рассогласования.

\section{Спектральная скорость и лапласовский шаг}

Если \(W\) диагонализируема,
\[
    W=P\Lambda P^{-1},
    \qquad
    W^k=P\Lambda^kP^{-1}.
\]
Поэтому при простом собственном значении \(1\) и условии
\[
    |\lambda_i(W)|<1,
    \qquad i\ge2,
\]
степени матрицы сходятся к проектору на консенсусное подпространство.
Для вещественной симметричной дважды стохастической матрицы введём
\[
    \rho_\perp(W)
    =
    \max_{i\ge2}|\lambda_i(W)|.
\]
Тогда
\[
    \lVert W^k-J\rVert_2
    =
    \rho_\perp(W)^k
\]
и
\[
    \lVert x^k-Jx^0\rVert
    \le
    \rho_\perp(W)^k
    \lVert x^0-Jx^0\rVert.
\]
Чем меньше \(\rho_\perp(W)\), тем быстрее затухает рассогласование.

Для связного неориентированного графа можно взять
\[
    W=I-\alpha L,
\]
где \(L\) --- матрица Лапласа. Локальное обновление имеет вид
\[
    x_i^{k+1}
    =
    x_i^k
    -
    \alpha
    \sum_{j\in\mathcal N_i}
    a_{ij}(x_i^k-x_j^k).
\]
Поскольку \(L\mathbf1=0\) и \(L=L^{\mathsf T}\), матрица \(W\)
сохраняет среднее. Для неотрицательности весов достаточно, например,
\[
    0<\alpha\le\frac1{d_{\max}}
\]
в невзвешенном графе. Если
\[
    0=\lambda_1(L)<\lambda_2(L)\le\cdots\le\lambda_N(L),
\]
то собственные значения \(W\) равны
\[
    1-\alpha\lambda_i(L).
\]
Спектральная сходимость обеспечивается условием
\[
    0<\alpha<\frac{2}{\lambda_N(L)},
\]
а её скорость определяется величиной
\[
    \max_{i\ge2}|1-\alpha\lambda_i(L)|.
\]
В частности, малое значение алгебраической связности
\(\lambda_2(L)\) замедляет консенсус.

Если матрица не диагонализируема, используется жорданова форма.
Степень блока \(J_\lambda=\lambda I+N\) размера \(r\) равна
\[
    J_\lambda^k
    =
    \sum_{s=0}^{r-1}
    \binom{k}{s}\lambda^{k-s}N^s.
\]
При \(|\lambda|<1\) экспоненциальное убывание подавляет
полиномиальные множители. Для сходимости собственное значение \(1\)
должно быть полупростым; у примитивной стохастической матрицы оно
простое, а остальные собственные значения лежат внутри единичного
круга.

\section{Ориентированные и векторные модели}

На ориентированном графе локально естественно строить
строчно-стохастическую матрицу. При сильной связности и
апериодичности она даёт
\[
    x^k\longrightarrow
    \mathbf1\pi^{\mathsf T}x^0,
\]
но в общем случае \(\pi\ne N^{-1}\mathbf1\), поэтому предел является
взвешенным средним. Для устранения этого смещения применяют, например,
push-sum и алгоритмы отношения масс.

Если состояние агента векторное,
\[
    x_i^k\in\mathbb R^d,
\]
то объединённое обновление записывается как
\[
    \mathbf x^{k+1}
    =
    (W\otimes I_d)\mathbf x^k.
\]
При \(W^k\to J\) каждая координата усредняется независимо:
\[
    \mathbf x^k
    \longrightarrow
    \mathbf1\otimes
    \left(
        \frac1N\sum_{i=1}^{N}x_i^0
    \right).
\]

В одном синхронном раунде агент отправляет соседям текущее состояние,
получает их состояния и вычисляет взвешенную сумму. Для векторов
размерности \(d\) суммарный объём передачи по активным рёбрам имеет
порядок
\[
    O(|\mathcal E|d),
\]
а число раундов до заданной точности определяется спектральным
коэффициентом смешивания. Базовый анализ предполагает фиксированный
граф, синхронные итерации, точную передачу сообщений и отсутствие
злонамеренных агентов; задержки, потери пакетов, асинхронность и
изменяющаяся топология требуют отдельных условий.

\lecture{29}{Спектральный анализ консенсуса}

\sourcevideo
    {https://www.youtube.com/playlist?list=PLOzRYVm0a65fhKDS8cYIC7YCuVGJ4FOJx}
    {Week 8: Lecture 29: Consensus and Average Consensus-2}

Рассматривается линейная динамика
\[
    x^{k+1}=Wx^k,
\]
где \(W\in\mathbb R^{N\times N}\) --- неотрицательная матрица
весов. Спектр \(W\) определяет предельное значение и скорость
консенсуса. Во второй части лекции выводятся два тождества для
нечётных функций относительных состояний, необходимые для анализа
непрерывных нелинейных алгоритмов.

\section{Стохастические матрицы и теорема Перрона--Фробениуса}

\begin{definition}[Строчная стохастичность]
Неотрицательная матрица \(W=(w_{ij})\) называется
строчно-стохастической, если
\[
    \sum_{j=1}^{N}w_{ij}=1,
    \qquad i=1,\ldots,N,
\]
или, эквивалентно,
\[
    W\mathbf1=\mathbf1.
\]
\end{definition}

Следовательно, число \(1\) является собственным значением \(W\), а
\(\mathbf1\) --- соответствующим правым собственным вектором.

\begin{theorem}[Спектральный радиус]
Если \(W\) строчно-стохастична, то каждое её собственное значение
\(\lambda\) удовлетворяет
\[
    |\lambda|\le1,
\]
и потому
\[
    \rho(W)=1.
\]
\end{theorem}

\begin{proof}
Для индуцированной нормы
\[
    \lVert W\rVert_\infty
    =
    \max_i\sum_{j=1}^{N}|w_{ij}|
    =1.
\]
Если \(Wv=\lambda v\) и \(v\ne0\), то
\[
    |\lambda|\lVert v\rVert_\infty
    =
    \lVert Wv\rVert_\infty
    \le
    \lVert W\rVert_\infty\lVert v\rVert_\infty,
\]
откуда \(|\lambda|\le1\). Поскольку \(1\) уже принадлежит спектру,
спектральный радиус равен единице.
\end{proof}

Граф матрицы задаётся соглашением
\[
    w_{ij}>0
    \quad\Longleftrightarrow\quad
    \text{агент }i\text{ получает сообщение от агента }j.
\]
Сильная связность этого графа эквивалентна неприводимости
неотрицательной матрицы. По теореме Перрона--Фробениуса для
неприводимой строчно-стохастической матрицы собственное значение
\(1\) алгебраически простое и существует единственный вектор
\(\pi>0\), удовлетворяющий
\[
    \pi^{\mathsf T}W=\pi^{\mathsf T},
    \qquad
    \mathbf1^{\mathsf T}\pi=1.
\]
Он называется стационарным распределением.

Одной сильной связности недостаточно для сходимости. Например,
\[
    W=
    \begin{pmatrix}
        0&1\\
        1&0
    \end{pmatrix}
\]
имеет собственные значения \(1\) и \(-1\), поэтому состояния
периодически меняются местами.

\begin{definition}[Примитивность]
Неотрицательная матрица \(W\) называется примитивной, если существует
\(m\ge1\), для которого \(W^m>0\) покомпонентно.
\end{definition}

Для неприводимой стохастической матрицы примитивность эквивалентна
апериодичности. В этом случае
\[
    \lambda_1(W)=1,
    \qquad
    |\lambda_i(W)|<1,
    \quad i=2,\ldots,N.
\]
На связном неориентированном графе апериодичность, в частности,
обеспечивается положительными диагональными весами \(w_{ii}>0\).

\section{Взвешенный и средний консенсус}

Пусть матрица \(W\) примитивна и строчно-стохастична. Общий
спектральный результат для простого доминирующего собственного
значения даёт
\[
    W^k\longrightarrow \mathbf1\pi^{\mathsf T}.
\]
Следовательно,
\[
    x^k=W^kx^0
    \longrightarrow
    \mathbf1\pi^{\mathsf T}x^0.
\]
Все агенты сходятся к одному значению
\[
    x_\infty=\pi^{\mathsf T}x^0,
\]
то есть достигается взвешенный консенсус.

\begin{definition}[Двойная стохастичность]
Матрица \(W\) называется дважды стохастической, если
\[
    W\mathbf1=\mathbf1,
    \qquad
    \mathbf1^{\mathsf T}W=\mathbf1^{\mathsf T}.
\]
\end{definition}

Столбцовая стохастичность означает сохранение суммы:
\[
    \mathbf1^{\mathsf T}x^{k+1}
    =
    \mathbf1^{\mathsf T}x^k.
\]
В этом случае нормированный левый собственный вектор равен
\[
    \pi=\frac1N\mathbf1.
\]

\begin{theorem}[Средний консенсус]
Если \(W\) примитивна и дважды стохастична, то
\[
    W^k\longrightarrow
    J:=\frac1N\mathbf1\mathbf1^{\mathsf T}
\]
и
\[
    x^k\longrightarrow Jx^0.
\]
Иными словами,
\[
    x_i^k\longrightarrow
    \frac1N\sum_{j=1}^{N}x_j^0
\]
для каждого агента \(i\).
\end{theorem}

Если \(W=W^{\mathsf T}\) и \(W\mathbf1=\mathbf1\), то матрица
автоматически дважды стохастична. Однако симметричность и связность
без апериодичности недостаточны: приведённая выше матрица перестановки
симметрична и дважды стохастична, но не сходится. Поэтому стандартное
достаточное условие включает неотрицательность, симметричность,
связность графа и положительные диагональные веса.

В симметричном случае существует ортогональное разложение
\[
    W=Q\Lambda Q^{\mathsf T},
\]
причём собственный вектор для \(\lambda_1=1\) равен
\[
    q_1=\frac1{\sqrt N}\mathbf1.
\]
Отсюда
\[
    W^k
    =
    q_1q_1^{\mathsf T}
    +
    \sum_{i=2}^{N}\lambda_i^kq_iq_i^{\mathsf T}
    \longrightarrow
    \frac1N\mathbf1\mathbf1^{\mathsf T}.
\]

\section{Скорость сходимости и топология}

Для симметричной матрицы смешивания положим
\[
    \rho_\perp(W)
    =
    \max_{i=2,\ldots,N}|\lambda_i(W)|.
\]
При среднем консенсусе \(\rho_\perp(W)<1\), и
\[
    \lVert W^k-J\rVert_2
    =
    \rho_\perp(W)^k.
\]
Поэтому
\[
    \lVert x^k-Jx^0\rVert
    \le
    \rho_\perp(W)^k
    \lVert x^0-Jx^0\rVert.
\]
Для относительной точности \(\varepsilon\) достаточно
\[
    k
    \ge
    \frac{\log(1/\varepsilon)}{-\log\rho_\perp(W)}.
\]

Топология определяет одновременно скорость смешивания и цену одного
раунда. Для пути \(P_N\)
\[
    |\mathcal E|=N-1,
    \qquad
    \operatorname{diam}(P_N)=N-1,
\]
поэтому обмен дешёв, но информация распространяется медленно. Для
полного графа \(K_N\)
\[
    |\mathcal E|=\frac{N(N-1)}2,
    \qquad
    \operatorname{diam}(K_N)=1,
\]
поэтому смешивание быстрое, но каждый раунд дорог. У кольца
\[
    |\mathcal E|=N,
    \qquad
    \operatorname{diam}(C_N)
    =
    \left\lfloor\frac N2\right\rfloor,
\]
а степень каждой вершины равна двум.

В экспоненциальных топологиях вершина соединяется с узлами на
циклических расстояниях
\[
    1,2,4,8,\ldots.
\]
И степень, и диаметр таких графов могут иметь порядок \(O(\log N)\).
Они дают промежуточный компромисс между кольцом и полным графом.

Для лапласовского консенсуса скорость связана с алгебраической
связностью \(\lambda_2(L)\). Узкие места уменьшают
\(\lambda_2(L)\), а рёбра между удалёнными частями сети обычно её
увеличивают. Поэтому сравнивать алгоритмы следует по полной
коммуникационной стоимости: числу раундов, умноженному на объём
сообщений за один раунд.

\section{Нечётные взаимодействия и сохранение среднего}

Пусть \(A=(a_{ij})\) --- симметричная матрица весов
неориентированного графа и
\[
    \phi\colon\mathbb R^d\to\mathbb R^d,
    \qquad
    \phi(-z)=-\phi(z).
\]
Примерами служат \(\phi(z)=z\), покоординатный знак и в скалярном
случае
\[
    \phi_\alpha(z)
    =
    |z|^\alpha\operatorname{sign}(z),
    \qquad \alpha>0.
\]

\begin{lemma}[Первая формула симметризации]
Для любых состояний \(x_1,\ldots,x_N\)
\[
    \sum_{i=1}^{N}\sum_{j=1}^{N}
    a_{ij}\phi(x_i-x_j)=0.
\]
\end{lemma}

\begin{proof}
Обозначим левую часть через \(S\). После перестановки индексов
\(i\) и \(j\) получаем
\[
    S
    =
    \sum_{i,j}a_{ji}\phi(x_j-x_i)
    =
    -\sum_{i,j}a_{ij}\phi(x_i-x_j)
    =-S.
\]
Здесь использованы симметрия \(a_{ij}=a_{ji}\) и нечётность
\(\phi\). Следовательно, \(S=0\).
\end{proof}

Рассмотрим непрерывную динамику
\[
    \dot x_i
    =
    -\sum_{j\in\mathcal N_i}
    a_{ij}\phi(x_i-x_j).
\]
По первой формуле
\[
    \frac{\mathrm d}{\mathrm dt}
    \sum_{i=1}^{N}x_i
    =0,
\]
а значит,
\[
    \bar x(t)
    =
    \frac1N\sum_{i=1}^{N}x_i(t)
    =
    \bar x(0).
\]
Если динамика достигает консенсуса, общее значение неизбежно равно
начальному среднему.

\section{Симметризация и функция Ляпунова}

\begin{lemma}[Вторая формула симметризации]
Для произвольных векторов \(e_i\in\mathbb R^d\)
\[
\begin{aligned}
    &\sum_{i,j}
    a_{ij}e_i^{\mathsf T}\phi(x_i-x_j)
    \\
    &\qquad=
    \frac12
    \sum_{i,j}
    a_{ij}(e_i-e_j)^{\mathsf T}\phi(x_i-x_j).
\end{aligned}
\]
\end{lemma}

\begin{proof}
Пусть
\[
    T=\sum_{i,j}a_{ij}e_i^{\mathsf T}\phi(x_i-x_j).
\]
После перестановки индексов и использования симметрии и нечётности
\[
    T
    =
    -\sum_{i,j}a_{ij}e_j^{\mathsf T}\phi(x_i-x_j).
\]
Сложение двух представлений даёт требуемое тождество.
\end{proof}

При \(e_i=x_i\)
\[
\begin{aligned}
    &\sum_{i,j}
    a_{ij}x_i^{\mathsf T}\phi(x_i-x_j)
    \\
    &\qquad=
    \frac12
    \sum_{i,j}
    a_{ij}(x_i-x_j)^{\mathsf T}\phi(x_i-x_j).
\end{aligned}
\]
Для \(\phi(z)=z\) правая часть равна
\[
    \frac12\sum_{i,j}a_{ij}\lVert x_i-x_j\rVert^2,
\]
а для покоординатного знака ---
\[
    \frac12\sum_{i,j}a_{ij}\lVert x_i-x_j\rVert_1.
\]

Положим
\[
    e_i=x_i-\bar x,
    \qquad
    V=\frac12\sum_{i=1}^{N}\lVert e_i\rVert^2.
\]
Поскольку среднее сохраняется, \(\dot e_i=\dot x_i\). Тогда
\[
    \dot V
    =
    -\frac12
    \sum_{i,j}
    a_{ij}(x_i-x_j)^{\mathsf T}\phi(x_i-x_j).
\]
Если
\[
    z^{\mathsf T}\phi(z)>0
    \qquad\text{при }z\ne0,
\]
то \(\dot V\le0\), и функция рассогласования не возрастает.

\section{Негладкие динамики и ограничения анализа}

Для знакового взаимодействия
\[
    \dot x_i
    =
    -\sum_{j\in\mathcal N_i}
    a_{ij}\operatorname{sign}(x_i-x_j)
\]
получаем
\[
    \dot V
    =
    -\frac12
    \sum_{i,j}a_{ij}|x_i-x_j|.
\]
Зависимость от первой степени рассогласований создаёт возможность
конечно-временной сходимости. Однако правая часть разрывна на
множестве \(x_i=x_j\), поэтому строгий анализ проводится в смысле
дифференциальных включений Филиппова.

Для степенного взаимодействия
\[
    \phi_\alpha(z)
    =
    |z|^\alpha\operatorname{sign}(z)
\]
имеем
\[
    \dot V
    =
    -\frac12
    \sum_{i,j}a_{ij}|x_i-x_j|^{1+\alpha}.
\]
Показатель \(0<\alpha<1\) используется в конечно-временных схемах,
а взаимодействие с \(\alpha>1\) усиливает движение вдали от
консенсусного множества. Сочетание
\[
    0<\alpha<1<\beta
\]
лежит в основе фиксированно-временных алгоритмов.

Обе формулы симметризации опираются на условие
\(a_{ij}=a_{ji}\). В ориентированных сетях парное сокращение в общем
случае исчезает; тогда применяются весовая сбалансированность,
специальные левые собственные векторы, push-sum или дополнительные
нормирующие переменные. Первая формула отвечает за сохранение
среднего, а вторая переводит производную функции Ляпунова в сумму
относительных рассогласований на рёбрах.

\lecture{30}{Алгоритмы консенсуса}

\sourcevideo
    {https://www.youtube.com/playlist?list=PLOzRYVm0a65fhKDS8cYIC7YCuVGJ4FOJx}
    {Week 9: Lecture 30: Consensus Algorithms}

В дискретном среднем консенсусе требуется построить матрицу смешивания,
которая сохраняет сумму состояний и подавляет рассогласование. Для
статических неориентированных графов это можно сделать локально с
помощью весов Метрополиса. В непрерывном времени ту же структуру
задаёт лапласовская динамика
\[
    \dot x=-Lx.
\]

\section{Дважды стохастические веса}

Пусть
\[
    x^{k+1}=Wx^k,
    \qquad
    W=(w_{ij})_{i,j=1}^{N}\ge0.
\]
Строчная и столбцовая стохастичность записываются как
\[
    W\mathbf1=\mathbf1,
    \qquad
    \mathbf1^{\mathsf T}W=\mathbf1^{\mathsf T}.
\]
При выполнении второго равенства сумма состояний сохраняется:
\[
    \mathbf1^{\mathsf T}x^{k+1}
    =
    \mathbf1^{\mathsf T}x^k.
\]
Если одновременно достигается консенсус
\(x^k\to c\mathbf1\), то
\[
    c=\frac1N\sum_{i=1}^{N}x_i^0.
\]
Следовательно, двойная стохастичность превращает обычный консенсус в
средний. Для сходимости дополнительно требуется апериодичность:
матрица
\[
    \begin{pmatrix}
        0&1\\
        1&0
    \end{pmatrix}
\]
дважды стохастична и имеет связный граф, но порождает периодические
колебания. На связном неориентированном графе апериодичность обычно
обеспечивается положительными диагональными весами.

Один общий способ построения дважды стохастической матрицы ---
масштабирование Синкхорна--Кноппа. Для неотрицательной матрицы
\(M^0\) последовательно нормируют строки и столбцы:
\[
    r_i^k=\sum_jM_{ij}^k,
    \qquad
    \widetilde M_{ij}^k=\frac{M_{ij}^k}{r_i^k},
\]
\[
    c_j^k=\sum_i\widetilde M_{ij}^k,
    \qquad
    M_{ij}^{k+1}=\frac{\widetilde M_{ij}^k}{c_j^k}.
\]
В матричной форме это диагональное масштабирование
\[
    M^k\longmapsto R_kM^k\longmapsto R_kM^kC_k.
\]
Сходимость справедлива не для любой неотрицательной матрицы: носитель
должен допускать дважды стохастическое масштабирование; стандартное
достаточное условие формулируется через полный носитель, а полная
неразложимость обеспечивает единственность масштабирования с точностью
до общего множителя.

Алгоритм требует повторного вычисления нормировок по всей матрице и
поэтому не является непосредственно локальным. Он также не решает
дискретную задачу назначения: матрицы перестановок являются частным
случаем дважды стохастических матриц, но масштабирование лишь переводит
положительные оценки в соответствующее непрерывное множество.

\section{Веса Метрополиса}

Пусть \(\mathcal G=(\mathcal V,\mathcal E)\) --- связный
неориентированный граф, а
\[
    d_i=|\mathcal N_i|
\]
--- степень вершины \(i\). Для ребра \(\{i,j\}\in\mathcal E\)
положим
\[
    w_{ij}
    =
    \frac1{1+\max\{d_i,d_j\}},
\]
для несмежных различных вершин --- \(w_{ij}=0\), а диагональный вес
определим формулой
\[
    w_{ii}
    =
    1-\sum_{j\in\mathcal N_i}w_{ij}.
\]
Для вычисления этих коэффициентов агенту нужны только собственная
степень и степени соседей, которые можно передать за один
предварительный раунд.

По симметрии выражения
\[
    w_{ij}=w_{ji},
\]
поэтому \(W=W^{\mathsf T}\). По определению диагонального веса
\[
    W\mathbf1=\mathbf1,
\]
а симметричность даёт
\[
    \mathbf1^{\mathsf T}W=\mathbf1^{\mathsf T}.
\]
Кроме того,
\[
    w_{ij}
    \le
    \frac1{1+d_i},
    \qquad
    \sum_{j\in\mathcal N_i}w_{ij}
    \le
    \frac{d_i}{1+d_i}<1,
\]
следовательно, \(w_{ii}>0\). Поэтому на связном графе матрица
примитивна и
\[
    W^k
    \longrightarrow
    \frac1N\mathbf1\mathbf1^{\mathsf T}.
\]

Например, для графа с рёбрами
\[
    \mathcal E
    =
    \{\{1,2\},\{2,3\},\{2,4\},\{3,4\}\}
\]
имеем
\[
    (d_1,d_2,d_3,d_4)=(1,3,2,2)
\]
и
\[
    W=
    \begin{pmatrix}
        \frac34&\frac14&0&0\\[1mm]
        \frac14&\frac14&\frac14&\frac14\\[1mm]
        0&\frac14&\frac5{12}&\frac13\\[1mm]
        0&\frac14&\frac13&\frac5{12}
    \end{pmatrix}.
\]
Локальная итерация принимает вид
\[
    x_i^{k+1}
    =
    w_{ii}x_i^k
    +
    \sum_{j\in\mathcal N_i}w_{ij}x_j^k,
\]
то есть использует только собственное состояние и состояния соседей.

\section{Непрерывный лапласовский консенсус}

Пусть \(A=(a_{ij})\) --- симметричная матрица неотрицательных весов
неориентированного графа,
\[
    D=\operatorname{diag}(d_1,\ldots,d_N),
    \qquad
    d_i=\sum_ja_{ij},
\]
а
\[
    L=D-A
\]
--- матрица Лапласа. Непрерывный алгоритм задаётся системой
\[
    \dot x=-Lx.
\]
Её локальная форма равна
\[
    \dot x_i
    =
    \sum_{j\in\mathcal N_i}a_{ij}(x_j-x_i).
\]
Таким образом, матричная запись нужна для анализа, а реализация
остаётся полностью локальной.

Для пути из трёх вершин
\[
    A=
    \begin{pmatrix}
        0&1&0\\
        1&0&1\\
        0&1&0
    \end{pmatrix},
    \qquad
    D=
    \begin{pmatrix}
        1&0&0\\
        0&2&0\\
        0&0&1
    \end{pmatrix},
\]
поэтому
\[
    L=
    \begin{pmatrix}
        1&-1&0\\
        -1&2&-1\\
        0&-1&1
    \end{pmatrix}
\]
и
\[
    \dot x_1=x_2-x_1,
    \qquad
    \dot x_2=(x_1-x_2)+(x_3-x_2),
    \qquad
    \dot x_3=x_2-x_3.
\]

Для связного неориентированного графа
\[
    \ker L=\operatorname{span}\{\mathbf1\},
\]
поэтому равновесия имеют вид \(x=c\mathbf1\). Поскольку
\[
    \mathbf1^{\mathsf T}L=0,
\]
сумма и среднее состояний сохраняются:
\[
    \sum_{i=1}^{N}x_i(t)
    =
    \sum_{i=1}^{N}x_i(0),
    \qquad
    \bar x(t)=\bar x(0).
\]
Следовательно, любой достигнутый консенсус является средним.

\section{Сходимость и её скорость}

Решение линейной системы имеет вид
\[
    x(t)=\exp(-Lt)x(0).
\]
Для связного графа спектр симметричного Лапласиана удовлетворяет
\[
    0=\lambda_1(L)
    <
    \lambda_2(L)
    \le\cdots\le
    \lambda_N(L),
\]
а нормированный собственный вектор для нулевого собственного значения
равен
\[
    q_1=\frac1{\sqrt N}\mathbf1.
\]
Из спектрального разложения
\[
    \exp(-Lt)
    =
    \sum_{i=1}^{N}
    e^{-\lambda_i(L)t}q_iq_i^{\mathsf T}
\]
получаем
\[
    \lim_{t\to\infty}\exp(-Lt)
    =
    \frac1N\mathbf1\mathbf1^{\mathsf T}
\]
и, следовательно,
\[
    x(t)
    \longrightarrow
    \frac1N\mathbf1\mathbf1^{\mathsf T}x(0).
\]

Положим
\[
    e(t)=x(t)-\bar x\mathbf1,
    \qquad
    V(e)=\frac12\lVert e\rVert^2.
\]
Так как \(e\perp\mathbf1\), то
\[
    \dot V
    =
    -e^{\mathsf T}Le
    \le
    -\lambda_2(L)\lVert e\rVert^2
    =
    -2\lambda_2(L)V.
\]
Отсюда
\[
    V(t)
    \le
    e^{-2\lambda_2(L)t}V(0)
\]
и
\[
    \lVert x(t)-\bar x\mathbf1\rVert
    \le
    e^{-\lambda_2(L)t}
    \lVert x(0)-\bar x\mathbf1\rVert.
\]
Скорость непрерывного консенсуса определяется алгебраической
связностью \(\lambda_2(L)\).

Знак минус в динамике существенен. При \(\dot x=Lx\) компоненты в
собственных направлениях с \(\lambda_i(L)>0\) растут как
\(e^{\lambda_i(L)t}\), поэтому рассогласование увеличивается.

\section{Дискретизация и векторные состояния}

Явная схема Эйлера с шагом \(h>0\) даёт
\[
    x^{k+1}
    =
    x^k-hLx^k
    =
    (I-hL)x^k.
\]
Следовательно,
\[
    W=I-hL.
\]
Матрица \(W\) симметрична и сохраняет среднее. Локально
\[
    x_i^{k+1}
    =
    (1-hd_i)x_i^k
    +
    h\sum_{j\in\mathcal N_i}a_{ij}x_j^k.
\]
Для неотрицательности весов достаточно
\[
    0<h\le\frac1{d_{\max}},
    \qquad
    d_{\max}=\max_i d_i,
\]
а спектральная устойчивость требует
\[
    0<h<\frac2{\lambda_N(L)}.
\]
Первое условие обеспечивает стохастическую интерпретацию, второе ---
сжатие всех неконсенсусных мод; при стандартных невзвешенных графах
первое условие автоматически является безопасным и для устойчивости.

Если агент хранит вектор \(x_i(t)\in\mathbb R^d\), то после
объединения состояний
\[
    \mathbf x(t)
    =
    (x_1(t)^{\mathsf T},\ldots,x_N(t)^{\mathsf T})^{\mathsf T}
\]
динамика записывается как
\[
    \dot{\mathbf x}
    =
    -(L\otimes I_d)\mathbf x.
\]
Для связного графа
\[
    \mathbf x(t)
    \longrightarrow
    \mathbf1\otimes
    \left(
        \frac1N\sum_{i=1}^{N}x_i(0)
    \right),
\]
то есть каждая координата усредняется независимо.

\section{Реализация и границы модели}

Для локального вычисления
\[
    \dot x_i
    =
    \sum_{j\in\mathcal N_i}a_{ij}(x_j-x_i)
\]
агенту нужны только собственное состояние, состояния соседей и веса
инцидентных рёбер. Полная матрица Лапласа и глобальная топология ему
не требуются.

Полученные гарантии относятся к статическому связному
неориентированному графу с симметричными неотрицательными весами. В
ориентированной сети лапласиан обычно несимметричен и среднее может не
сохраняться. Для изменяющихся графов требуется совместная связность на
временных окнах, а задержки, потери сообщений и асинхронность требуют
отдельного анализа. Линейная динамика достигает консенсуса лишь
асимптотически; конечно-временные и фиксированно-временные варианты
строятся с помощью нелинейных функций относительных состояний.

\lecture{31}{Фиксированно-временной консенсус}

\sourcevideo
    {https://www.youtube.com/playlist?list=PLOzRYVm0a65fhKDS8cYIC7YCuVGJ4FOJx}
    {Week 9: Lecture 31: Consensus Algorithms-Fixed time}

Линейная динамика
\[
    \dot x=-Lx
\]
достигает среднего консенсуса только асимптотически. В этой лекции
строится нелинейный алгоритм, для которого существует единая верхняя
граница времени установления, не зависящая от начальных состояний.

\section{Критерий фиксированного времени}

Пусть множество равновесий системы задаётся условием \(V=0\), где
\(V\ge0\) --- функция Ляпунова. Сходимость называется
конечно-временной, если для каждого начального состояния существует
конечное время \(T(x(0))\), после которого \(V=0\). Она называется
фиксированно-временной, если
\[
    T(x(0))\le T_{\max}<\infty
\]
для всех начальных состояний.

\begin{theorem}[Критерий фиксированного времени]
Если вдоль траекторий
\[
    \dot V
    \le
    -c_1V^{\alpha_1}
    -c_2V^{\alpha_2},
\]
где
\[
    c_1,c_2>0,
    \qquad
    0<\alpha_1<1<\alpha_2,
\]
то множество \(\{V=0\}\) достигается не позднее времени
\[
    T_{\max}
    \le
    \frac1{c_1(1-\alpha_1)}
    +
    \frac1{c_2(\alpha_2-1)}.
\]
\end{theorem}

Степень меньше единицы обеспечивает конечное время вблизи
равновесия, а степень больше единицы ограничивает время движения из
сколь угодно удалённой начальной точки.

Для \(z\in\mathbb R^d\) и \(\mu>0\) определим знаковую степень
\[
    \operatorname{sig}^{\mu}(z)
    =
    \begin{cases}
        z\lVert z\rVert^{\mu-1},&z\ne0,\\[1mm]
        0,&z=0.
    \end{cases}
\]
В скалярном случае
\[
    \operatorname{sig}^{\mu}(z)
    =
    |z|^\mu\operatorname{sign}(z).
\]
Для любого \(\mu>0\)
\[
    \operatorname{sig}^{\mu}(-z)
    =
    -\operatorname{sig}^{\mu}(z),
    \qquad
    z^{\mathsf T}\operatorname{sig}^{\mu}(z)
    =
    \lVert z\rVert^{1+\mu}.
\]
Определение через норму корректно для произвольного вещественного
\(\mu>0\); ограничивать показатель отношением нечётных целых не
требуется.

\section{Алгоритм и сохранение среднего}

Пусть \(\mathcal G=(\mathcal V,\mathcal E)\) --- связный
неориентированный невзвешенный граф, а агент \(i\) хранит состояние
\(x_i(t)\in\mathbb R^d\). Выберем
\[
    k_1,k_2>0,
    \qquad
    0<\mu_1<1<\mu_2.
\]
Локальная динамика имеет вид
\[
\begin{aligned}
    \dot x_i
    ={}&
    -k_1
    \sum_{j\in\mathcal N_i}
    \operatorname{sig}^{\mu_1}(x_i-x_j)
    \\
    &-
    k_2
    \sum_{j\in\mathcal N_i}
    \operatorname{sig}^{\mu_2}(x_i-x_j).
\end{aligned}
\]
Первое слагаемое отвечает за конечное время около консенсусного
множества, второе --- за равномерную по начальным условиям оценку
вдали от него. При \(\mu_1=\mu_2=1\) получается линейный
лапласовский консенсус с изменённым коэффициентом.

Просуммируем уравнения по всем агентам. Каждому неориентированному
ребру \(\{i,j\}\) соответствуют два противоположных слагаемых,
поскольку
\[
    \operatorname{sig}^{\mu}(x_j-x_i)
    =
    -\operatorname{sig}^{\mu}(x_i-x_j).
\]
Следовательно,
\[
    \frac{\mathrm d}{\mathrm dt}
    \sum_{i=1}^{N}x_i=0.
\]
Среднее
\[
    \bar x
    =
    \frac1N\sum_{i=1}^{N}x_i(t)
\]
сохраняется. Поэтому любой достигнутый консенсус совпадает с
начальным средним.

Введём рассогласования
\[
    \widetilde x_i=x_i-\bar x,
    \qquad
    \widetilde{\mathbf x}
    =
    (\widetilde x_1^{\mathsf T},\ldots,
      \widetilde x_N^{\mathsf T})^{\mathsf T}.
\]
Тогда
\[
    (\mathbf1^{\mathsf T}\otimes I_d)
    \widetilde{\mathbf x}=0,
    \qquad
    \dot{\widetilde x}_i=\dot x_i.
\]

\section{Функция Ляпунова}

Рассмотрим
\[
    V
    =
    \frac12
    \sum_{i=1}^{N}\lVert\widetilde x_i\rVert^2
    =
    \frac12\lVert\widetilde{\mathbf x}\rVert^2.
\]
Равенство \(V=0\) эквивалентно среднему консенсусу. Производная
равна
\[
    \dot V
    =
    \sum_{i=1}^{N}
    \widetilde x_i^{\mathsf T}\dot x_i.
\]
Для симметричного графа и нечётной функции \(\phi\) выполняется
тождество
\[
\begin{aligned}
    &\sum_i\sum_{j\in\mathcal N_i}
    e_i^{\mathsf T}\phi(x_i-x_j)
    \\
    &\qquad=
    \frac12
    \sum_i\sum_{j\in\mathcal N_i}
    (e_i-e_j)^{\mathsf T}\phi(x_i-x_j).
\end{aligned}
\]
Применяя его с \(e_i=\widetilde x_i\), получаем
\[
\begin{aligned}
    \dot V
    ={}&
    -k_1
    \sum_{\{i,j\}\in\mathcal E}
    \lVert\widetilde x_i-
            \widetilde x_j\rVert^{1+\mu_1}
    \\
    &-
    k_2
    \sum_{\{i,j\}\in\mathcal E}
    \lVert\widetilde x_i-
            \widetilde x_j\rVert^{1+\mu_2}.
\end{aligned}
\]
Каждое ребро здесь учитывается один раз.

\section{Степенные и спектральные оценки}

Положим
\[
    m=|\mathcal E|,
    \qquad
    \eta_{ij}
    =
    \lVert\widetilde x_i-
            \widetilde x_j\rVert^2,
\]
\[
    p_1=\frac{1+\mu_1}{2},
    \qquad
    p_2=\frac{1+\mu_2}{2}.
\]
Тогда
\[
    0<p_1<1<p_2
\]
и
\[
    \dot V
    =
    -k_1\sum_{\{i,j\}\in\mathcal E}\eta_{ij}^{p_1}
    -k_2\sum_{\{i,j\}\in\mathcal E}\eta_{ij}^{p_2}.
\]
Для неотрицательных чисел
\[
    \sum_e\eta_e^{p_1}
    \ge
    \left(\sum_e\eta_e\right)^{p_1},
\]
а по неравенству между степенными средними
\[
    \sum_e\eta_e^{p_2}
    \ge
    m^{1-p_2}
    \left(\sum_e\eta_e\right)^{p_2}.
\]

Для неориентированного графа
\[
    \sum_{\{i,j\}\in\mathcal E}
    \lVert\widetilde x_i-
            \widetilde x_j\rVert^2
    =
    \widetilde{\mathbf x}^{\mathsf T}
    (L\otimes I_d)
    \widetilde{\mathbf x}.
\]
Поскольку граф связен и
\(\widetilde{\mathbf x}\) ортогонален консенсусному подпространству,
\[
    \widetilde{\mathbf x}^{\mathsf T}
    (L\otimes I_d)
    \widetilde{\mathbf x}
    \ge
    \lambda_2(L)
    \lVert\widetilde{\mathbf x}\rVert^2
    =
    2\lambda_2(L)V.
\]
Следовательно,
\[
    \dot V
    \le
    -c_1V^{p_1}
    -c_2V^{p_2},
\]
где
\[
    c_1
    =
    k_1\bigl(2\lambda_2(L)\bigr)^{p_1},
\]
\[
    c_2
    =
    k_2m^{1-p_2}
    \bigl(2\lambda_2(L)\bigr)^{p_2}.
\]

\section{Фиксированно-временная сходимость}

\begin{theorem}[Средний консенсус за фиксированное время]
Пусть \(\mathcal G\) --- связный неориентированный невзвешенный
граф и
\[
    k_1,k_2>0,
    \qquad
    0<\mu_1<1<\mu_2.
\]
Тогда приведённая динамика достигает среднего консенсуса:
\[
    x_i(t)
    =
    \frac1N\sum_{\ell=1}^{N}x_\ell(0)
\]
для всех \(i\) и всех \(t\ge T_{\max}\), где
\[
    T_{\max}
    \le
    \frac1{c_1(1-p_1)}
    +
    \frac1{c_2(p_2-1)}.
\]
Эта граница не зависит от начального состояния.
\end{theorem}

С учётом определений показателей
\[
\begin{aligned}
    T_{\max}
    \le{}&
    \frac{2}{c_1(1-\mu_1)}
    \\
    &+
    \frac{2}{c_2(\mu_2-1)}.
\end{aligned}
\]
Увеличение \(\lambda_2(L)\) повышает обе константы \(c_1,c_2\) и
уменьшает гарантированное время. Узкие места в графе уменьшают
алгебраическую связность и ухудшают оценку. Число рёбер входит через
множитель \(m^{1-p_2}\); при отсутствии точного значения можно
использовать грубую границу
\[
    m\le\frac{N(N-1)}2.
\]
Запись по упорядоченным парам вместо неориентированных рёбер меняет
только константы, поскольку каждое взаимодействие учитывается дважды.

\section{Расширения и ограничения}

Для симметричного взвешенного графа алгоритм записывается как
\[
\begin{aligned}
    \dot x_i
    ={}&
    -k_1
    \sum_{j\in\mathcal N_i}
    a_{ij}\operatorname{sig}^{\mu_1}(x_i-x_j)
    \\
    &-
    k_2
    \sum_{j\in\mathcal N_i}
    a_{ij}\operatorname{sig}^{\mu_2}(x_i-x_j),
\end{aligned}
\]
где \(a_{ij}=a_{ji}>0\). Симметрия сохраняет среднее, а явная оценка
времени дополнительно зависит от весов и спектра взвешенного
лапласиана.

Для изменяющегося графа то же точечное доказательство сохраняется,
если каждый мгновенный граф неориентирован и связен, а
\[
    \lambda_2(L(t))
    \ge
    \underline\lambda>0
\]
равномерно по времени. Одной совместной связности на временных окнах
недостаточно для приведённой оценки, поскольку в отдельные моменты
\(\lambda_2(L(t))\) может обращаться в нуль.

Алгоритм полностью децентрализован: агент использует только своё
состояние, состояния соседей и параметры
\(k_1,k_2,\mu_1,\mu_2\). Однако вычисление общей априорной границы
\(T_{\max}\) требует знания или нижней оценки
\(\lambda_2(L)\).

Полученная гарантия относится к непрерывному времени. Явная
дискретизация Эйлера не обязана сохранять точное достижение
консенсуса за число шагов, не зависящее от начальных условий. Большой
шаг может вызвать колебания и потерю убывания функции Ляпунова, а для
строгой дискретной фиксированно-временной гарантии нужен отдельный
анализ.

Наконец, консенсус обеспечивает только равенство локальных
переменных. Для распределённой оптимизации его необходимо совместить
с локальными градиентами и отдельно доказать сходимость общего
состояния к минимуму суммы локальных функций.

\lecture{32}{Распределённая задача экономического диспетчирования}

\sourcevideo
    {https://www.youtube.com/playlist?list=PLOzRYVm0a65fhKDS8cYIC7YCuVGJ4FOJx}
    {Week 9: Lecture 32: Distributed Economic Dispatch Problem}

Экономическое диспетчирование состоит в распределении заданной
нагрузки между генераторами с минимальной суммарной стоимостью. Каждый
генератор знает собственную функцию затрат и обменивается с соседями
только локальными переменными. В этой лекции выводятся условия
оптимальности и формулируются требования к распределённому алгоритму;
конкретная динамика строится в следующей лекции.

\section{Постановка задачи}

Пусть в сети имеется \(N\) генераторов. Через \(P_i\) обозначим
мощность генератора \(i\), а через \(P_{\mathrm D}\) --- суммарную
нагрузку. В модели без потерь должен выполняться баланс
\[
    \sum_{i=1}^{N}P_i=P_{\mathrm D}.
\]
Генератору \(i\) соответствует функция стоимости \(C_i(P_i)\). В
лекции используется квадратичная модель
\[
    C_i(P_i)
    =
    \alpha_iP_i^2+\beta_iP_i+\gamma_i,
    \qquad
    \alpha_i>0,
\]
для которой
\[
    C_i'(P_i)=2\alpha_iP_i+\beta_i.
\]
Производная \(C_i'(P_i)\) называется предельной стоимостью генерации.
Постоянный коэффициент \(\gamma_i\) меняет значение стоимости, но не
оптимальное распределение мощности.

Полная централизованная задача имеет вид
\[
\begin{aligned}
    \min_{P_1,\ldots,P_N}\quad
    &\sum_{i=1}^{N}C_i(P_i),
    \\
    \text{при условиях}\quad
    &\sum_{i=1}^{N}P_i=P_{\mathrm D},
    \\
    &P_i^{\min}\le P_i\le P_i^{\max},
    \qquad i=1,\ldots,N.
\end{aligned}
\]
Необходимое условие допустимости равно
\[
    \sum_{i=1}^{N}P_i^{\min}
    \le
    P_{\mathrm D}
    \le
    \sum_{i=1}^{N}P_i^{\max}.
\]

Сначала ограничения мощности отбрасываются. Полученная задача
называется некапацитированной:
\[
\begin{aligned}
    \min_{P_1,\ldots,P_N}\quad
    &\sum_{i=1}^{N}C_i(P_i),
    \\
    \text{при условии}\quad
    &\sum_{i=1}^{N}P_i=P_{\mathrm D}.
\end{aligned}
\]
При строгой выпуклости функций \(C_i\) её оптимальный вектор
единственен.

В отличие от консенсусной оптимизации, здесь обычно
\[
    P_i^\star\ne P_j^\star.
\]
Согласовываются не мощности, а двойственная переменная, связанная с
общим балансом.

\section{Условия оптимальности}

Для ограничения
\[
    P_{\mathrm D}-\sum_{i=1}^{N}P_i=0
\]
введём множитель \(\lambda\in\mathbb R\). Функция Лагранжа равна
\[
    \mathcal L(P,\lambda)
    =
    \sum_{i=1}^{N}C_i(P_i)
    +
    \lambda
    \left(
        P_{\mathrm D}-\sum_{i=1}^{N}P_i
    \right).
\]
Для некапацитированной выпуклой задачи условия стационарности дают
\[
    C_i'(P_i^\star)=\lambda^\star,
    \qquad i=1,\ldots,N,
\]
а производная по \(\lambda\) восстанавливает баланс
\[
    \sum_{i=1}^{N}P_i^\star=P_{\mathrm D}.
\]
Следовательно,
\[
    C_1'(P_1^\star)
    =
    \cdots
    =
    C_N'(P_N^\star)
    =
    \lambda^\star.
\]

Если у двух свободных генераторов предельные стоимости различны, то
перенос малой мощности от более дорогого генератора к более дешёвому
сохраняет баланс и уменьшает общую стоимость. Поэтому равенство
предельных стоимостей является необходимым условием оптимальности.

\section{Квадратичная модель}

Для квадратичных затрат условие стационарности принимает вид
\[
    2\alpha_iP_i^\star+\beta_i=\lambda^\star,
\]
откуда
\[
    P_i^\star
    =
    \frac{\lambda^\star-\beta_i}{2\alpha_i}.
\]
Подстановка в условие баланса даёт
\[
    \sum_{i=1}^{N}
    \frac{\lambda^\star-\beta_i}{2\alpha_i}
    =
    P_{\mathrm D},
\]
и поэтому
\[
    \lambda^\star
    =
    \frac{
        P_{\mathrm D}
        +
        \displaystyle\sum_{i=1}^{N}
        \frac{\beta_i}{2\alpha_i}
    }{
        \displaystyle\sum_{i=1}^{N}
        \frac1{2\alpha_i}
    }.
\]
Зная \(\lambda^\star\), каждый генератор вычисляет собственную
оптимальную мощность только по своим коэффициентам. Однако прямое
вычисление общего множителя требует централизованного сбора
\(P_{\mathrm D}\), всех \(\alpha_i\) и всех \(\beta_i\). Именно от
этого сбора отказывается распределённая постановка.

\section{Распределённая формулировка}

Обмен сообщениями описывается связным графом
\[
    \mathcal G=(\mathcal V,\mathcal E),
    \qquad
    \mathcal V=\{1,\ldots,N\}.
\]
Ребро \(\{i,j\}\in\mathcal E\) означает наличие канала связи между
генераторами \(i\) и \(j\). Коммуникационный граф может совпадать с
физической сетью, но это не требуется.

Вместо одного глобального множителя генератор \(i\) хранит локальную
оценку \(\lambda_i(t)\). Целевое состояние удовлетворяет
\[
    \lambda_1(t)
    =
    \cdots
    =
    \lambda_N(t)
    \longrightarrow
    \lambda^\star,
\]
а текущая мощность связывается с ней формулой
\[
    P_i(t)
    =
    \frac{\lambda_i(t)-\beta_i}{2\alpha_i}.
\]
При \(\lambda_i(t)\to\lambda^\star\) получаем
\[
    P_i(t)\to P_i^\star.
\]

Обычного среднего консенсуса по \(\lambda_i\) недостаточно: его
предел определяется начальными значениями и не обязан совпадать с
\(\lambda^\star\). Распределённая динамика должна одновременно
согласовывать предельные стоимости и обеспечивать баланс мощности.

Обмен локальными двойственными оценками не требует явной передачи
коэффициентов \(\alpha_i,\beta_i,\gamma_i\), но сам по себе не даёт
формальной гарантии конфиденциальности.

\section{Сохранение баланса и ограничения мощности}

Баланс удобно обеспечить уже при инициализации. Пусть \(D_i\) ---
нагрузка, локально закреплённая за генератором \(i\), причём
\[
    \sum_{i=1}^{N}D_i=P_{\mathrm D}.
\]
Выбор
\[
    P_i(0)=D_i
\]
даёт допустимое начальное состояние. Если алгоритм удовлетворяет
\[
    \frac{\mathrm d}{\mathrm dt}
    \sum_{i=1}^{N}P_i(t)=0,
\]
то
\[
    \sum_{i=1}^{N}P_i(t)=P_{\mathrm D}
\]
для всех \(t\). Полная нагрузка при этом не обязана быть известна
каждому генератору отдельно.

В капацитированной задаче дополнительно требуется
\[
    P_i^{\min}\le P_i\le P_i^{\max}.
\]
Формула
\[
    P_i^\star
    =
    \frac{\lambda^\star-\beta_i}{2\alpha_i}
\]
применима только к генераторам с неактивными ограничениями. Если это
значение выходит за допустимый интервал, генератор работает на
соответствующей границе. Поэтому активные ограничения изменяют
условия оптимальности и требуют отдельного алгоритма.

\section{Требования к алгоритму и границы модели}

Для некапацитированной задачи целевое состояние определяется двумя
условиями:
\[
    C_i'(P_i^\star)=\lambda^\star,
    \qquad i=1,\ldots,N,
\]
и
\[
    \sum_{i=1}^{N}P_i^\star=P_{\mathrm D}.
\]
Следующая лекция строит распределённую динамику, которая сохраняет
баланс и за фиксированное время согласовывает предельные стоимости.
Генератор использует собственную функцию стоимости, текущую мощность,
локальную двойственную оценку и сообщения соседей; полные функции
стоимости остальных генераторов ему не нужны.

Рассматриваемая модель не учитывает потери в линиях, ограничения их
пропускной способности, динамику частоты и напряжения, стоимость
запуска и остановки, невыпуклые режимы, неопределённость
возобновляемой генерации и изменение нагрузки во времени. Поэтому она
служит базовой постановкой для введения распределённых методов, а не
полной моделью энергосистемы.

\lecture{33}{Некапацитированное экономическое диспетчирование}

\sourcevideo
    {https://www.youtube.com/playlist?list=PLOzRYVm0a65fhKDS8cYIC7YCuVGJ4FOJx}
    {Week 9: Lecture 33: Algorithm for Uncapacitated EDP}

Рассматривается задача
\[
\begin{aligned}
    \min_{P_1,\ldots,P_N}\quad
    &\sum_{i=1}^{N}C_i(P_i),
    \\
    \text{при условии}\quad
    &\sum_{i=1}^{N}P_i=P_{\mathrm D},
\end{aligned}
\]
где
\[
    C_i(P_i)=\alpha_iP_i^2+\beta_iP_i+\gamma_i,
    \qquad \alpha_i>0.
\]
В оптимуме
\[
    2\alpha_iP_i^\star+\beta_i=\lambda^\star,
    \qquad
    \sum_{i=1}^{N}P_i^\star=P_{\mathrm D}.
\]
Следовательно, распределённый алгоритм должен сохранять баланс
мощности, установить локальную связь между мощностью и предельной
стоимостью и затем согласовать предельные стоимости.

\section{Локальные переменные и алгоритм}

Пусть генераторы соединены связным неориентированным графом
\[
    \mathcal G=(\mathcal V,\mathcal E),
    \qquad
    \mathcal V=\{1,\ldots,N\}.
\]
Через \(\mathcal N_i\) обозначим множество соседей генератора \(i\).
Для \(z\in\mathbb R\) и \(\mu>0\) положим
\[
    \operatorname{sig}^{\mu}(z)
    =
    |z|^\mu\operatorname{sign}(z).
\]
Эта функция нечётна и удовлетворяет равенству
\[
    z\operatorname{sig}^{\mu}(z)=|z|^{1+\mu}.
\]
Выберем
\[
    k_1,k_2,r_1,r_2>0,
    \qquad
    0<\mu_1<1<\mu_2.
\]

Каждый генератор хранит мощность \(P_i\) и локальную предельную
стоимость \(\lambda_i\). Ошибка локального условия оптимальности
равна
\[
    e_i
    =
    P_i-
    \frac{\lambda_i-\beta_i}{2\alpha_i}.
\]
Равенство \(e_i=0\) эквивалентно
\[
    \lambda_i=2\alpha_iP_i+\beta_i=C_i'(P_i).
\]
Определим консенсусный сигнал
\[
\begin{aligned}
    \Phi_i(\lambda)
    ={}&
    k_1
    \sum_{j\in\mathcal N_i}
    \operatorname{sig}^{\mu_1}(\lambda_j-\lambda_i)
    \\
    &+
    k_2
    \sum_{j\in\mathcal N_i}
    \operatorname{sig}^{\mu_2}(\lambda_j-\lambda_i).
\end{aligned}
\]
Распределённая динамика задаётся уравнениями
\[
    \dot P_i=\Phi_i(\lambda),
\]
\[
\begin{aligned}
    \dot\lambda_i
    =
    2\alpha_i
    \Bigl(
        \Phi_i(\lambda)
        &+r_1\operatorname{sig}^{\mu_1}(e_i)
        \\
        &+r_2\operatorname{sig}^{\mu_2}(e_i)
    \Bigr).
\end{aligned}
\]
Для реализации агенту нужны только собственные параметры
\(\alpha_i,\beta_i\), собственные переменные \(P_i,\lambda_i\) и
значения \(\lambda_j\) соседей.

\section{Сохранение баланса и динамика ошибки}

По нечётности знаковой степени и симметрии графа каждому слагаемому,
соответствующему ребру \(\{i,j\}\), отвечает противоположное
слагаемое. Поэтому
\[
    \sum_{i=1}^{N}\Phi_i(\lambda)=0.
\]
Следовательно,
\[
    \frac{\mathrm d}{\mathrm dt}
    \sum_{i=1}^{N}P_i=0
\]
и
\[
    \sum_{i=1}^{N}P_i(t)
    =
    \sum_{i=1}^{N}P_i(0).
\]
Если
\[
    \sum_{i=1}^{N}P_i(0)=P_{\mathrm D},
\]
то баланс мощности выполняется при всех \(t\ge0\).

Поскольку \(\beta_i\) постоянно,
\[
    \dot e_i
    =
    \dot P_i-
    \frac{\dot\lambda_i}{2\alpha_i}.
\]
После подстановки алгоритма консенсусные сигналы сокращаются:
\[
    \dot e_i
    =
    -r_1\operatorname{sig}^{\mu_1}(e_i)
    -r_2\operatorname{sig}^{\mu_2}(e_i).
\]
Таким образом, ошибки локальной оптимальности подчиняются автономной
динамике, не зависящей от графа и состояний соседей.

\section{Локальная оптимальность}

Рассмотрим функцию Ляпунова
\[
    V_e=\frac12\sum_{i=1}^{N}e_i^2.
\]
Её производная равна
\[
    \dot V_e
    =
    -r_1\sum_{i=1}^{N}|e_i|^{1+\mu_1}
    -r_2\sum_{i=1}^{N}|e_i|^{1+\mu_2}.
\]
Положим
\[
    p_1=\frac{1+\mu_1}{2},
    \qquad
    p_2=\frac{1+\mu_2}{2}.
\]
Тогда \(0<p_1<1<p_2\). Для неотрицательных чисел выполняются
оценки
\[
    \sum_{i=1}^{N}|e_i|^{1+\mu_1}
    \ge
    \left(
        \sum_{i=1}^{N}e_i^2
    \right)^{p_1},
\]
\[
    \sum_{i=1}^{N}|e_i|^{1+\mu_2}
    \ge
    N^{1-p_2}
    \left(
        \sum_{i=1}^{N}e_i^2
    \right)^{p_2}.
\]
Так как \(\sum_i e_i^2=2V_e\), получаем
\[
    \dot V_e
    \le
    -c_{e,1}V_e^{p_1}
    -c_{e,2}V_e^{p_2},
\]
где
\[
    c_{e,1}=r_1 2^{p_1},
    \qquad
    c_{e,2}=r_2N^{1-p_2}2^{p_2}.
\]
По критерию фиксированного времени все ошибки обращаются в нуль не
позднее момента
\[
    T_1
    \le
    \frac1{c_{e,1}(1-p_1)}
    +
    \frac1{c_{e,2}(p_2-1)}.
\]
Следовательно, при \(t\ge T_1\)
\[
    P_i
    =
    \frac{\lambda_i-\beta_i}{2\alpha_i},
\]
и это равенство сохраняется при дальнейшем движении.

\section{Консенсус предельных стоимостей}

После первого этапа
\[
    \frac{\dot\lambda_i}{2\alpha_i}
    =
    \dot P_i
    =
    \Phi_i(\lambda).
\]
Положим
\[
    w_i=\frac1{2\alpha_i}.
\]
Тогда
\[
    \frac{\mathrm d}{\mathrm dt}
    \sum_{i=1}^{N}w_i\lambda_i
    =
    \sum_{i=1}^{N}\Phi_i(\lambda)
    =0.
\]
В момент \(T_1\)
\[
\begin{aligned}
    \sum_{i=1}^{N}\frac{\lambda_i(T_1)}{2\alpha_i}
    &=
    \sum_{i=1}^{N}P_i(T_1)
    +
    \sum_{i=1}^{N}\frac{\beta_i}{2\alpha_i}
    \\
    &=
    P_{\mathrm D}
    +
    \sum_{i=1}^{N}\frac{\beta_i}{2\alpha_i}.
\end{aligned}
\]
Если \(\lambda_i\) сходятся к общему значению
\(\lambda_{\mathrm c}\), то
\[
    \lambda_{\mathrm c}
    \sum_{i=1}^{N}\frac1{2\alpha_i}
    =
    P_{\mathrm D}
    +
    \sum_{i=1}^{N}\frac{\beta_i}{2\alpha_i}.
\]
Отсюда
\[
    \lambda_{\mathrm c}
    =
    \frac{
        P_{\mathrm D}
        +
        \displaystyle\sum_{i=1}^{N}
        \frac{\beta_i}{2\alpha_i}
    }{
        \displaystyle\sum_{i=1}^{N}
        \frac1{2\alpha_i}
    }
    =
    \lambda^\star.
\]
Поэтому предельные мощности равны
\[
    P_i^\star
    =
    \frac{\lambda^\star-\beta_i}{2\alpha_i}.
\]
Консенсус достигается по предельным стоимостям, а мощности в общем
случае остаются различными.

\section{Фиксированное время второго этапа}

Положим
\[
    \widetilde\lambda_i=\lambda_i-\lambda^\star
\]
и введём диагональную матрицу
\[
    W
    =
    \operatorname{diag}
    \left(
        \frac1{2\alpha_1},\ldots,
        \frac1{2\alpha_N}
    \right).
\]
После момента \(T_1\)
\[
    \sum_{i=1}^{N}
    \frac{\widetilde\lambda_i}{2\alpha_i}=0.
\]
Рассмотрим
\[
    V_\lambda
    =
    \frac12
    \widetilde\lambda^{\mathsf T}
    W\widetilde\lambda.
\]
Для связного графа существует \(\kappa_W>0\), такое что на
указанном подпространстве
\[
    \widetilde\lambda^{\mathsf T}
    L\widetilde\lambda
    \ge
    \kappa_W
    \widetilde\lambda^{\mathsf T}
    W\widetilde\lambda.
\]
Используя ту же симметризацию по рёбрам и степенные неравенства, что
и в доказательстве фиксированно-временного консенсуса, получаем
\[
    \dot V_\lambda
    \le
    -c_{\lambda,1}V_\lambda^{p_1}
    -c_{\lambda,2}V_\lambda^{p_2}
\]
с некоторыми положительными константами
\(c_{\lambda,1}\) и \(c_{\lambda,2}\). Поэтому консенсус по
\(\lambda_i\) достигается за время
\[
    T_2
    \le
    \frac1{c_{\lambda,1}(1-p_1)}
    +
    \frac1{c_{\lambda,2}(p_2-1)},
\]
не зависящее от значений \(\lambda_i(T_1)\).

\begin{theorem}[Фиксированно-временное диспетчирование]
Пусть граф связен и неориентирован, функции стоимости квадратичны и
строго выпуклы, а начальные мощности удовлетворяют балансу
\[
    \sum_{i=1}^{N}P_i(0)=P_{\mathrm D}.
\]
Тогда приведённый распределённый алгоритм достигает единственного
решения некапацитированной задачи экономического диспетчирования не
позднее времени
\[
    T_{\max}\le T_1+T_2.
\]
При всех \(t\ge T_{\max}\)
\[
    P_i(t)=P_i^\star,
    \qquad
    \lambda_i(t)=\lambda^\star.
\]
\end{theorem}

\section{Интерпретация и ограничения}

Переменная \(\lambda_i\) является локальной предельной, или
инкрементальной, стоимостью. Классические дискретные методы,
согласовывающие такие переменные, называются алгоритмами согласования
инкрементальных стоимостей; рассмотренная схема является их
непрерывным фиксированно-временным вариантом.

Результат относится к некапацитированной модели. Ограничения
\[
    P_i^{\min}\le P_i\le P_i^{\max}
\]
алгоритмом автоматически не поддерживаются. Анализ также предполагает
квадратичные строго выпуклые стоимости, статический связный
неориентированный граф, симметричный и точный обмен, непрерывное
выполнение динамики, отсутствие потерь мощности и сетевых
ограничений. Прямая дискретизация Эйлера не обязана сохранять точную
фиксированно-временную гарантию. Капацитированная задача требует
изменения условий оптимальности и рассматривается далее.

\lecture{34}{Капацитированное экономическое диспетчирование}

\sourcevideo
    {https://www.youtube.com/playlist?list=PLOzRYVm0a65fhKDS8cYIC7YCuVGJ4FOJx}
    {Week 9: Lecture 34: Capacitated EDP}

Учитываются физические ограничения генераторов
\[
    P_i^{\min}\le P_i\le P_i^{\max}.
\]
Они изменяют условия оптимальности: предельные стоимости свободных
генераторов совпадают, а генераторы на границах удовлетворяют
односторонним условиям KKT.

\section{Задача и условия KKT}

Пусть
\[
    C_i(P_i)=\alpha_iP_i^2+\beta_iP_i+\gamma_i,
    \qquad \alpha_i>0.
\]
Рассматривается выпуклая задача
\[
\begin{aligned}
    \min_{P_1,\ldots,P_N}\quad
    &\sum_{i=1}^{N}C_i(P_i),
    \\
    \text{при условиях}\quad
    &\sum_{i=1}^{N}P_i=P_{\mathrm D},
    \\
    &P_i^{\min}\le P_i\le P_i^{\max}.
\end{aligned}
\]
Допустимость требует
\[
    \sum_{i=1}^{N}P_i^{\min}
    \le P_{\mathrm D}\le
    \sum_{i=1}^{N}P_i^{\max}.
\]
Запишем ограничения мощности как
\[
    P_i-P_i^{\max}\le0,
    \qquad
    P_i^{\min}-P_i\le0.
\]
Для баланса введём множитель \(\lambda\in\mathbb R\), а для верхнего
и нижнего ограничений --- множители \(r_i,s_i\ge0\). Функция Лагранжа
равна
\[
\begin{aligned}
    \mathcal L(P,\lambda,r,s)
    ={}&
    \sum_{i=1}^{N}C_i(P_i)
    +\lambda\left(P_{\mathrm D}-\sum_{i=1}^{N}P_i\right)
    \\
    &+
    \sum_{i=1}^{N}r_i(P_i-P_i^{\max})
    +
    \sum_{i=1}^{N}s_i(P_i^{\min}-P_i).
\end{aligned}
\]
Условия KKT имеют вид
\[
    2\alpha_iP_i^\star+\beta_i
    -\lambda^\star+r_i^\star-s_i^\star=0,
\]
\[
    \sum_{i=1}^{N}P_i^\star=P_{\mathrm D},
    \qquad
    P_i^{\min}\le P_i^\star\le P_i^{\max},
\]
\[
    r_i^\star,s_i^\star\ge0,
    \qquad
    r_i^\star(P_i^\star-P_i^{\max})=0,
    \qquad
    s_i^\star(P_i^{\min}-P_i^\star)=0.
\]
Следовательно, для свободного генератора
\[
    P_i^{\min}<P_i^\star<P_i^{\max}
    \quad\Longrightarrow\quad
    \lambda^\star=2\alpha_iP_i^\star+\beta_i,
\]
на верхней границе
\[
    P_i^\star=P_i^{\max}
    \quad\Longrightarrow\quad
    \lambda^\star\ge2\alpha_iP_i^{\max}+\beta_i,
\]
а на нижней границе
\[
    P_i^\star=P_i^{\min}
    \quad\Longrightarrow\quad
    \lambda^\star\le2\alpha_iP_i^{\min}+\beta_i.
\]

\section{Проекция мощности и скалярное уравнение}

Для фиксированного \(\lambda\) локальная оптимальная реакция
генератора равна
\[
    P_i(\lambda)
    =
    \min\left\{
        P_i^{\max},
        \max\left\{
            P_i^{\min},
            \frac{\lambda-\beta_i}{2\alpha_i}
        \right\}
    \right\}.
\]
Тем самым некапацитированное значение проецируется на интервал
\([P_i^{\min},P_i^{\max}]\). Оптимальный множитель определяется
уравнением
\[
    \sum_{i=1}^{N}P_i(\lambda^\star)=P_{\mathrm D}.
\]
Левая часть непрерывна и не убывает по \(\lambda\). При наличии хотя
бы одного свободного генератора она строго возрастает в окрестности
решения, поэтому \(\lambda^\star\) определяется однозначно. Если все
генераторы находятся на границах, мощности уже фиксированы, а сам
множитель может быть неоднозначен на интервале, совместимом с KKT.

Уравнение можно решать централизованно бисекцией. В распределённой
схеме вместо прямого поиска корня используется активное множество и
несколько раундов среднего консенсуса.

\section{Активное множество и формула множителя}

Пусть \(\Theta\) --- множество генераторов, работающих на границах, и
\[
    \widehat P_i\in\{P_i^{\min},P_i^{\max}\},
    \qquad i\in\Theta.
\]
Для свободных генераторов
\[
    P_i^\star
    =
    \frac{\lambda^\star-\beta_i}{2\alpha_i},
    \qquad i\notin\Theta.
\]
Из баланса мощности получаем
\[
    \lambda^\star
    =
    \frac{
        P_{\mathrm D}
        -\displaystyle\sum_{i\in\Theta}\widehat P_i
        +\displaystyle\sum_{i\notin\Theta}
        \frac{\beta_i}{2\alpha_i}
    }{
        \displaystyle\sum_{i\notin\Theta}
        \frac1{2\alpha_i}
    }.
\]

Обозначим через \(\lambda_{\mathrm u}^\star\) множитель
некапацитированной задачи:
\[
    \lambda_{\mathrm u}^\star
    =
    \frac{
        P_{\mathrm D}
        +\displaystyle\sum_{i=1}^{N}
        \frac{\beta_i}{2\alpha_i}
    }{
        \displaystyle\sum_{i=1}^{N}
        \frac1{2\alpha_i}
    },
\]
\[
    P_{i,\mathrm u}^\star
    =
    \frac{\lambda_{\mathrm u}^\star-\beta_i}{2\alpha_i}.
\]
Тогда для фиксированного активного множества
\[
    \lambda^\star
    =
    \lambda_{\mathrm u}^\star
    +
    \frac{
        \displaystyle\sum_{i\in\Theta}
        \left(P_{i,\mathrm u}^\star-\widehat P_i\right)
    }{
        \displaystyle\sum_{i\notin\Theta}
        \frac1{2\alpha_i}
    }.
\]
Числитель показывает суммарное изменение мощности из-за насыщения, а
знаменатель --- чувствительность свободных генераторов к изменению
общего множителя.

\section{Распределённое вычисление поправки}

Для заданного множества \(\Theta\) агент \(i\) инициализирует
\[
    y_i(0)
    =
    \begin{cases}
        P_{i,\mathrm u}^\star-\widehat P_i,
        &i\in\Theta,\\[1mm]
        0,
        &i\notin\Theta,
    \end{cases}
\]
и
\[
    z_i(0)
    =
    \begin{cases}
        0,
        &i\in\Theta,\\[1mm]
        \dfrac1{2\alpha_i},
        &i\notin\Theta.
    \end{cases}
\]
Два параллельных алгоритма среднего консенсуса дают всем агентам
\[
    \overline y
    =
    \frac1N
    \sum_{i\in\Theta}
    \left(P_{i,\mathrm u}^\star-\widehat P_i\right),
    \qquad
    \overline z
    =
    \frac1N
    \sum_{i\notin\Theta}
    \frac1{2\alpha_i}.
\]
Множители \(1/N\) сокращаются, поэтому каждый агент вычисляет
\[
    \lambda^\star
    =
    \lambda_{\mathrm u}^\star
    +
    \frac{\overline y}{\overline z}.
\]
После этого свободные мощности пересчитываются по формуле
\[
    P_i
    =
    \frac{\lambda^\star-\beta_i}{2\alpha_i},
    \qquad i\notin\Theta,
\]
а активные остаются на назначенных границах. Коэффициенты
\(\alpha_i,\beta_i,\gamma_i\) напрямую не передаются, однако обмен
состояниями \(y_i,z_i\) не является криптографической гарантией
приватности.

\section{Активно-множественная процедура}

Сначала распределённо решается некапацитированная задача и строятся
\(\lambda_{\mathrm u}^\star\) и \(P_{i,\mathrm u}^\star\). Затем при
текущем \(\lambda\) каждый свободный генератор вычисляет
\[
    P_i=\frac{\lambda-\beta_i}{2\alpha_i}.
\]
Нарушители образуют множество
\[
\begin{aligned}
    \Omega
    =
    \bigl\{
        i\notin\Theta:
        P_i<P_i^{\min}
        \text{ или }
        P_i>P_i^{\max}
    \bigr\}.
\end{aligned}
\]
Для \(i\in\Omega\) мощность фиксируется на ближайшей нарушенной
границе:
\[
    \widehat P_i
    =
    \begin{cases}
        P_i^{\min},&P_i<P_i^{\min},\\[1mm]
        P_i^{\max},&P_i>P_i^{\max}.
    \end{cases}
\]
После обновления \(\Theta\) через два консенсуса пересчитывается
\(\lambda\), затем пересчитываются мощности свободных генераторов.
Процедура завершается, когда нарушителей нет и одновременно выполнены
односторонние условия KKT для активных индексов:
\[
    C_i'(P_i^{\max})\le\lambda^\star
    \quad\text{на верхней границе},
\]
\[
    C_i'(P_i^{\min})\ge\lambda^\star
    \quad\text{на нижней границе}.
\]
Одной проверки допустимости мощностей недостаточно.

В лекции активное множество только увеличивается. Это корректно, если
однажды насыщенный генератор после каждого пересчёта остаётся
активным. В общей задаче при одновременных нарушениях верхних и нижних
границ такое свойство требует отдельного доказательства. Надёжная
реализация должна разрешать удаление индексов из \(\Theta\) либо
проверять конечную точку через проекционное уравнение
\[
    \sum_{i=1}^{N}P_i(\lambda)=P_{\mathrm D}.
\]

\section{Завершение и границы модели}

Если на каждом раунде к корректному активному множеству добавляется
хотя бы один новый индекс и удаление не требуется, число обновлений не
превосходит \(N\). Если некапацитированная задача решается за время
\(T_{\mathrm u}\), а один раунд двух параллельных консенсусов занимает
не более \(T_{\mathrm c}\), то
\[
    T_{\mathrm{total}}
    \le
    T_{\mathrm u}+NT_{\mathrm c}.
\]
Эта оценка не зависит от начальных значений непрерывных консенсусных
переменных, но зависит от числа генераторов и коммуникационного графа.

Рассмотренная модель не учитывает потери мощности, ограничения потоков
по линиям, динамику генераторов, невыпуклые затраты, стоимость запуска
и изменение нагрузки во времени. Точная фиксированно-временная
гарантия относится к непрерывным консенсусным динамикам; после
дискретизации требуется отдельный анализ устойчивости и времени
сходимости.

\lecture{35}{Алгоритмы распределённой оптимизации}

\sourcevideo
    {https://www.youtube.com/playlist?list=PLOzRYVm0a65fhKDS8cYIC7YCuVGJ4FOJx}
    {Week 10: Lecture 35}

Пусть \(N\) агентов связаны коммуникационным графом, а агент \(i\)
знает только собственную функцию
\(f_i\colon\mathbb R^d\to\mathbb R\). Требуется минимизировать сумму
локальных функций, не собирая их в одном узле. Основной метод лекции
--- распределённый градиентный спуск (distributed gradient descent,
DGD), сочетающий локальное усреднение и градиентный шаг.

\section{Постановка задачи}

Централизованная задача имеет вид
\[
    \min_{x\in\mathbb R^d} F(x),
    \qquad
    F(x)=\sum_{i=1}^{N}f_i(x).
\]
При доступе ко всем функциям обычный градиентный спуск записывается как
\[
    x^{k+1}
    =
    x^k-\eta_k\nabla F(x^k),
    \qquad
    \nabla F(x)=\sum_{i=1}^{N}\nabla f_i(x).
\]
В распределённой сети агент не может вычислить эту сумму. Он хранит
локальную копию \(x_i\in\mathbb R^d\), вычисляет только
\(\nabla f_i(x_i)\) и обменивается сообщениями с соседями.
Эквивалентная формулировка имеет вид
\[
\begin{aligned}
    \min_{x_1,\ldots,x_N}\quad
    &\sum_{i=1}^{N}f_i(x_i),
    \\
    \text{при условии}\quad
    &x_1=\cdots=x_N.
\end{aligned}
\]
Ограничение консенсуса возвращает одну общую переменную и тем самым
восстанавливает исходную задачу.

Для примера рассмотрим
\[
    f_i(x)=\frac12(x-i)^2,
    \qquad i=1,\ldots,5.
\]
Тогда
\[
    F'(x)=\sum_{i=1}^{5}(x-i)=5x-15,
\]
поэтому глобальный минимум равен \(x^\star=3\). Изолированный агент
получил бы локальный минимум \(x_i^\star=i\), так что без коммуникации
глобальное решение недостижимо.

\section{Матрица смешивания}

Обмен сообщениями задаётся связным неориентированным графом
\[
    \mathcal G=(\mathcal V,\mathcal E),
    \qquad
    \mathcal V=\{1,\ldots,N\},
\]
и матрицей смешивания
\(W=(w_{ij})\in\mathbb R^{N\times N}\). Для базового анализа
предполагается
\[
    w_{ij}\ge0,
    \qquad
    w_{ij}=0
    \quad\text{при }j\notin\mathcal N_i\cup\{i\},
\]
а также двойная стохастичность
\[
    W\mathbf1=\mathbf1,
    \qquad
    \mathbf1^{\mathsf T}W=\mathbf1^{\mathsf T}.
\]
На неориентированном графе этим требованиям удовлетворяют, например,
симметричные веса Метрополиса:
\[
    w_{ij}
    =
    \frac{1}{1+\max\{d_i,d_j\}}
\]
для соседних \(i\ne j\) и
\[
    w_{ii}
    =
    1-\sum_{j\in\mathcal N_i}w_{ij}.
\]
Строчная стохастичность делает локальное обновление выпуклой
комбинацией соседних состояний, а столбцовая сохраняет их среднее.

\section{Распределённый градиентный спуск}

На итерации \(k\) агент сначала смешивает локальные переменные,
\[
    x_i^{k+\frac12}
    =
    \sum_{j=1}^{N}w_{ij}x_j^k,
\]
а затем выполняет градиентный шаг по своей функции:
\[
    x_i^{k+1}
    =
    x_i^{k+\frac12}
    -
    \eta_k\nabla f_i\!\left(x_i^{k+\frac12}\right).
\]
Смешивание уменьшает рассогласование, тогда как локальные градиенты
двигают агентов к минимумам разных функций. Возможен и обратный
порядок:
\[
    x_i^{k+\frac12}
    =
    x_i^k-\eta_k\nabla f_i(x_i^k),
    \qquad
    x_i^{k+1}
    =
    \sum_{j=1}^{N}w_{ij}x_j^{k+\frac12}.
\]
Обе схемы сочетают оптимизационный и консенсусный шаги, но имеют
разные рекуррентные соотношения и требуют отдельных доказательств.

Обозначим среднее состояние через
\[
    \bar x^k
    =
    \frac1N\sum_{i=1}^{N}x_i^k.
\]
Для первого порядка шагов столбцовая стохастичность даёт
\[
\begin{aligned}
    \frac1N\sum_{i=1}^{N}x_i^{k+\frac12}
    &=
    \frac1N\sum_{j=1}^{N}
    \left(\sum_{i=1}^{N}w_{ij}\right)x_j^k
    \\
    &=
    \bar x^k.
\end{aligned}
\]
После локального градиентного шага
\[
    \bar x^{k+1}
    =
    \bar x^k
    -
    \frac{\eta_k}{N}
    \sum_{i=1}^{N}
    \nabla f_i\!\left(x_i^{k+\frac12}\right).
\]
Если
\[
    \bar F(x)=\frac1N\sum_{i=1}^{N}f_i(x),
\]
то централизованный шаг использовал бы
\(\nabla\bar F(\bar x^k)\). Различие возникает потому, что локальные
градиенты вычисляются в разных точках.

\section{Средняя динамика и рассогласование}

Добавляя и вычитая градиенты в средней точке, получаем
\[
    \bar x^{k+1}
    =
    \bar x^k
    -
    \eta_k\nabla\bar F(\bar x^k)
    +
    e^k,
\]
где
\[
    e^k
    =
    \frac{\eta_k}{N}
    \sum_{i=1}^{N}
    \left[
        \nabla f_i(\bar x^k)
        -
        \nabla f_i\!\left(x_i^{k+\frac12}\right)
    \right].
\]
Следовательно, среднее состояние подчиняется возмущённому
централизованному градиентному спуску. Если каждый градиент
\(L_i\)-липшицев и \(L=\max_iL_i\), то
\[
    \lVert e^k\rVert
    \le
    \frac{\eta_kL}{N}
    \sum_{i=1}^{N}
    \left\lVert x_i^{k+\frac12}-\bar x^k\right\rVert.
\]
Ошибка исчезает, когда шаг и рассогласование становятся достаточно
малыми.

Для матричного описания положим
\[
    \mathbf x^k
    =
    \begin{pmatrix}
        x_1^k\\
        \vdots\\
        x_N^k
    \end{pmatrix},
    \qquad
    J=\frac1N\mathbf1\mathbf1^{\mathsf T}.
\]
Консенсусная часть равна
\(\mathbf1\otimes\bar x^k\), а рассогласование ---
\[
    \mathbf x^k-\mathbf1\otimes\bar x^k.
\]
Для симметричной примитивной дважды стохастической матрицы
\[
    \rho
    =
    \lVert W-J\rVert_2
    <1.
\]
Чистое смешивание сжимает рассогласование с коэффициентом не больше
\(\rho\), однако следующий локальный градиентный шаг снова его создаёт.
Поэтому анализ DGD должен одновременно контролировать оптимизационную
ошибку среднего состояния и расстояние локальных переменных до
консенсусного подпространства.

\section{Сходимость и выбор шага}

Классический анализ точной сходимости использует убывающий шаг
\[
    \eta_k>0,
    \qquad
    \sum_{k=0}^{\infty}\eta_k=\infty,
    \qquad
    \sum_{k=0}^{\infty}\eta_k^2<\infty.
\]
Первое условие не позволяет суммарному движению оказаться конечным до
достижения решения, а второе делает квадратичные возмущения
суммируемыми. Типичный выбор --- \(\eta_k=\eta_0/(k+1)\).

При связном неориентированном графе, примитивной дважды
стохастической матрице \(W\), выпуклых гладких функциях, непустом
множестве решений и стандартных дополнительных условиях
ограниченности можно доказать
\[
    \lVert x_i^k-x_j^k\rVert\to0
\]
и
\[
    \operatorname{dist}
    \left(
        \bar x^k,
        \operatorname*{arg\,min}F
    \right)
    \to0.
\]
Если минимизатор единственен, то
\[
    x_i^k\to x^\star
\]
для всех агентов.

При постоянном шаге \(\eta_k=\eta>0\) обычный DGD в общем случае
сходится лишь к окрестности решения. Её размер уменьшается вместе с
\(\eta\) и зависит от спектральных свойств графа и различия локальных
градиентов. Для точной сходимости с постоянным шагом вводят
корректирующие механизмы. Упомянутый в лекции алгоритм EXTRA устраняет
стационарное смещение DGD и при подходящих предположениях сходится к
точному решению; его формула и доказательство здесь не выводятся.

\section{Коммуникационная стоимость и ограничения}

На каждой итерации агент передаёт соседям вектор
\(x_i^k\in\mathbb R^d\). Поэтому суммарный объём обмена по активным
рёбрам имеет порядок
\[
    O\bigl(|\mathcal E|d\bigr),
\]
а локальная вычислительная стоимость определяется вычислением
\(\nabla f_i(x_i)\). Передача дополнительных градиентных или
инерционных переменных сохраняет порядок \(d\) на сообщение, но
увеличивает постоянный множитель.

Базовая схема предполагает синхронные итерации, фиксированный связный
граф, точную передачу сообщений и отсутствие злонамеренных агентов.
Для ориентированных графов двойная стохастичность может быть
недоступна, для изменяющейся топологии требуется совместная связность,
а задержки, асинхронность и стохастические градиенты приводят к другим
оценкам. Следующая лекция вводит градиентный трекинг, позволяющий
устранить смещение постоянного шага за счёт дополнительной локальной
переменной.

\lecture{36}{Градиентный трекинг}

\sourcevideo
    {https://www.youtube.com/playlist?list=PLOzRYVm0a65fhKDS8cYIC7YCuVGJ4FOJx}
    {Week 10: Lecture 36}

Распределённый градиентный спуск сочетает смешивание локальных
переменных с шагом по локальному градиенту, но при постоянном шаге в
общем случае сохраняет смещение. Метод градиентного трекинга вводит
дополнительную переменную, которая отслеживает средний градиент сети и
позволяет получить точную сходимость при постоянном шаге.

\section{Постановка и смещение DGD}

Пусть \(N\) агентов решают задачу
\[
    \min_{x\in\mathbb R^d} F(x),
    \qquad
    F(x)=\sum_{i=1}^{N}f_i(x),
\]
где агент \(i\) знает только функцию
\(f_i\colon\mathbb R^d\to\mathbb R\) и обменивается сообщениями с
соседями. Каждый агент хранит локальную оценку
\(x_i^k\in\mathbb R^d\), и требуется обеспечить
\[
    x_i^k\longrightarrow x^\star,
    \qquad
    x^\star\in\operatorname*{arg\,min}_{x}F(x).
\]

Распределённый градиентный спуск имеет вид
\[
    x_i^{k+1}
    =
    \sum_{j=1}^{N}w_{ij}x_j^k
    -\eta\nabla f_i(x_i^k).
\]
Первое слагаемое уменьшает рассогласование, второе --- локальную
функцию. При постоянном шаге эти действия конфликтуют: в глобальном
оптимуме обычно
\[
    \nabla f_i(x^\star)\ne0,
\]
хотя
\[
    \sum_{i=1}^{N}\nabla f_i(x^\star)=0.
\]
Поэтому локальные градиенты продолжают раздвигать агентов даже около
решения, и DGD может сходиться лишь к его окрестности.

\section{Алгоритм градиентного трекинга}

Каждый агент дополнительно хранит \(y_i^k\in\mathbb R^d\), которое
должно приближать средний текущий градиент сети. Пусть
\(W=(w_{ij})\in\mathbb R^{N\times N}\) --- примитивная дважды
стохастическая матрица смешивания:
\[
    W\mathbf1=\mathbf1,
    \qquad
    \mathbf1^{\mathsf T}W=\mathbf1^{\mathsf T}.
\]
При постоянном шаге \(\eta>0\) метод задаётся рекурсиями
\[
    x_i^{k+1}
    =
    \sum_{j=1}^{N}w_{ij}x_j^k
    -\eta y_i^k,
\]
\[
\begin{aligned}
    y_i^{k+1}
    ={}&
    \sum_{j=1}^{N}w_{ij}y_j^k
    \\
    &+
    \nabla f_i(x_i^{k+1})
    -
    \nabla f_i(x_i^k),
\end{aligned}
\]
с инициализацией
\[
    y_i^0=\nabla f_i(x_i^0).
\]

Обновление \(x_i^k\) использует не локальный градиент, а его
отслеживаемую сетевую оценку. В обновлении \(y_i^k\) смешивание
сочетается с приращением локального градиента. Без разности
\[
    \nabla f_i(x_i^{k+1})-\nabla f_i(x_i^k)
\]
обычный консенсус усреднял бы только начальные градиенты и не учитывал
бы движение точек \(x_i^k\).

На каждой итерации агент передаёт соседям два вектора: \(x_i^k\) и
\(y_i^k\). Наличие двух связанных рекурсий не делает метод методом
второго порядка: в нём нет ни гессианов, ни обратных гессианов, а
переменная \(y_i^k\) не является классическим импульсом.

\section{Инвариант и средняя динамика}

Главное тождество метода имеет вид
\[
    \sum_{i=1}^{N}y_i^k
    =
    \sum_{i=1}^{N}\nabla f_i(x_i^k)
\]
для всех \(k\). При \(k=0\) оно следует из инициализации. Если
равенство выполнено на итерации \(k\), то
\[
\begin{aligned}
    \sum_{i=1}^{N}y_i^{k+1}
    ={}&
    \sum_{i=1}^{N}\sum_{j=1}^{N}w_{ij}y_j^k
    \\
    &+
    \sum_{i=1}^{N}
    \bigl(
        \nabla f_i(x_i^{k+1})
        -
        \nabla f_i(x_i^k)
    \bigr).
\end{aligned}
\]
Столбцовая стохастичность даёт
\[
    \sum_{i=1}^{N}\sum_{j=1}^{N}w_{ij}y_j^k
    =
    \sum_{j=1}^{N}y_j^k,
\]
после чего индукционное предположение приводит к тому же тождеству на
итерации \(k+1\).

Определим средние
\[
    \overline x^k
    =
    \frac1N\sum_{i=1}^{N}x_i^k,
    \qquad
    \overline y^k
    =
    \frac1N\sum_{i=1}^{N}y_i^k.
\]
Тогда
\[
    \overline y^k
    =
    \frac1N
    \sum_{i=1}^{N}\nabla f_i(x_i^k)
\]
и, после усреднения первого обновления,
\[
    \overline x^{k+1}
    =
    \overline x^k-\eta\overline y^k
    =
    \overline x^k
    -
    \frac{\eta}{N}
    \sum_{i=1}^{N}\nabla f_i(x_i^k).
\]
Если локальные точки близки к среднему, то
\[
    \overline y^k
    \approx
    \frac1N\nabla F(\overline x^k),
\]
поэтому средняя переменная движется приблизительно как
централизованный градиентный спуск для функции \(F/N\).

\section{Матричная форма и рассогласование}

Объединим локальные переменные в столбцы
\[
    \mathbf x^k
    =
    \begin{pmatrix}
        x_1^k\\
        \vdots\\
        x_N^k
    \end{pmatrix},
    \qquad
    \mathbf y^k
    =
    \begin{pmatrix}
        y_1^k\\
        \vdots\\
        y_N^k
    \end{pmatrix},
\]
и положим
\[
    \mathbf g(\mathbf x)
    =
    \begin{pmatrix}
        \nabla f_1(x_1)\\
        \vdots\\
        \nabla f_N(x_N)
    \end{pmatrix}.
\]
Тогда
\[
    \mathbf x^{k+1}
    =
    (W\otimes I_d)\mathbf x^k
    -\eta\mathbf y^k,
\]
\[
\begin{aligned}
    \mathbf y^{k+1}
    ={}&
    (W\otimes I_d)\mathbf y^k
    \\
    &+
    \mathbf g(\mathbf x^{k+1})
    -
    \mathbf g(\mathbf x^k).
\end{aligned}
\]

Пусть
\[
    J=\frac1N\mathbf1\mathbf1^{\mathsf T}.
\]
Ошибки консенсуса равны
\[
    \mathbf e_x^k
    =
    \mathbf x^k-
    \mathbf1\otimes\overline x^k,
    \qquad
    \mathbf e_y^k
    =
    \mathbf y^k-
    \mathbf1\otimes\overline y^k.
\]
Матрица \(W-J\) действует только на неконсенсусные компоненты. Для
симметричной примитивной дважды стохастической матрицы
\[
    \rho=\lVert W-J\rVert_2<1,
\]
поэтому смешивание сжимает обе ошибки. В динамике
\(\mathbf e_y^k\) возникает возмущение, связанное с изменением
градиентов, но при их липшицевости оно контролируется изменением
\(\mathbf x^k\).

\section{Предельная точка и сходимость}

Предположим, что последовательности сходятся и локальные переменные
достигают консенсуса:
\[
    x_i^k\to x^\infty,
    \qquad
    y_i^k\to y_i^\infty.
\]
В пределе первое обновление даёт
\[
    \mathbf x^\infty
    =
    (W\otimes I_d)\mathbf x^\infty
    -\eta\mathbf y^\infty.
\]
На консенсусном подпространстве
\[
    (W\otimes I_d)\mathbf x^\infty=\mathbf x^\infty,
\]
поэтому
\[
    \mathbf y^\infty=0.
\]
Инвариант суммы отсюда влечёт
\[
    0
    =
    \sum_{i=1}^{N}\nabla f_i(x^\infty)
    =
    \nabla F(x^\infty).
\]
Для выпуклой функции \(F\) это условие означает глобальную
оптимальность. Именно здесь исчезает смещение DGD: в глобальном
оптимуме отслеживаемый средний градиент равен нулю, хотя отдельные
локальные градиенты могут оставаться ненулевыми.

\begin{theorem}[Линейная сходимость]
Пусть граф связен и неориентирован, \(W\) примитивна и дважды
стохастична, локальные градиенты липшицевы, глобальная функция \(F\)
сильно выпукла, а постоянный шаг \(\eta\) достаточно мал. Тогда
существуют \(C>0\) и \(q\in(0,1)\), для которых
\[
    \lVert x_i^k-x^\star\rVert
    \le
    Cq^k
\]
для всех агентов \(i\).
\end{theorem}

Допустимый шаг зависит от гладкости локальных функций и качества
смешивания, характеризуемого величиной \(\lVert W-J\rVert_2\).
Слишком большой шаг может вызвать колебания или расходимость даже при
сильной выпуклости.

\section{Пример, реализация и стоимость}

Для функций
\[
    f_i(x)=\frac12(x-i)^2,
    \qquad
    i=1,\ldots,5,
\]
имеем \(\nabla f_i(x)=x-i\) и
\[
    x^\star
    =
    \frac{1+2+3+4+5}{5}
    =3.
\]
При произвольных \(x_i^0\) и инициализации
\[
    y_i^0=x_i^0-i
\]
рекурсии принимают вид
\[
    x_i^{k+1}
    =
    \sum_jw_{ij}x_j^k-
    \eta y_i^k,
\]
\[
    y_i^{k+1}
    =
    \sum_jw_{ij}y_j^k
    +x_i^{k+1}-x_i^k.
\]
При подходящем шаге все локальные точки сходятся к \(3\).

Для неориентированного графа можно использовать веса Метрополиса:
\[
    w_{ij}
    =
    \frac1{1+\max\{d_i,d_j\}}
\]
для соседних \(i\ne j\) и
\[
    w_{ii}
    =
    1-
    \sum_{j\in\mathcal N_i}w_{ij}.
\]
Полученная матрица симметрична и дважды стохастична. В NumPy
смешивание записывается как \texttt{W @ x}; выражение
\texttt{W * x} выполняет покомпонентное умножение и не реализует
матричный шаг.

По каждому ребру на итерации передаются два вектора размерности \(d\),
поэтому суммарный объём обмена имеет порядок
\[
    O(|\mathcal E|d),
\]
но с примерно вдвое большей постоянной, чем у DGD. Локально агент
вычисляет один новый градиент и два смешивания; предыдущий градиент
можно хранить.

Базовая модель предполагает синхронные итерации, фиксированный граф и
точную передачу сообщений. Для ориентированных графов требуются
push-sum или другие корректирующие механизмы. При стохастических
градиентах переменная трекинга отслеживает шумную величину, поэтому в
анализе появляется дополнительная дисперсия.

\lecture{37}{Непрерывная распределённая оптимизация}

\sourcevideo
    {https://www.youtube.com/playlist?list=PLOzRYVm0a65fhKDS8cYIC7YCuVGJ4FOJx}
    {Week 10: Lecture 37}

Рассматривается задача
\[
    \min_{x\in\mathbb R^d}F(x),
    \qquad
    F(x)=\sum_{i=1}^{N}f_i(x),
\]
в которой агент \(i\) знает только функцию \(f_i\) и обменивается
сообщениями с соседями. Цель --- построить непрерывную систему,
достигающую единственного минимума за время, равномерно ограниченное
по всем начальным состояниям.

\section{Централизованный поток}

Предположим, что функции \(f_i\) дифференцируемы и выпуклы, а сумма
\(F\) является \(\mu\)-сильно выпуклой. Тогда минимум \(x^\star\)
единственен и
\[
    \lVert\nabla F(x)\rVert^2
    \ge
    2\mu\bigl(F(x)-F(x^\star)\bigr).
\]
Для \(z\in\mathbb R^d\) и \(\nu>0\) положим
\[
    \operatorname{sig}^{\nu}(z)
    =
    \begin{cases}
        z\lVert z\rVert^{\nu-1},&z\ne0,\\[1mm]
        0,&z=0.
    \end{cases}
\]
Тогда
\[
    z^{\mathsf T}\operatorname{sig}^{\nu}(z)
    =
    \lVert z\rVert^{1+\nu}.
\]
Выберем
\[
    a_1,a_2>0,
    \qquad
    0<\nu_1<1<\nu_2,
\]
и определим
\[
    \Psi(z)
    =
    a_1\operatorname{sig}^{\nu_1}(z)
    +
    a_2\operatorname{sig}^{\nu_2}(z).
\]
Централизованный поток имеет вид
\[
    \dot x=-\Psi\bigl(\nabla F(x)\bigr).
\]
Для
\[
    V(x)=F(x)-F(x^\star)
\]
получаем
\[
\begin{aligned}
    \dot V
    ={}&
    -a_1\lVert\nabla F(x)\rVert^{1+\nu_1}
    \\
    &-
    a_2\lVert\nabla F(x)\rVert^{1+\nu_2}.
\end{aligned}
\]
Положим
\[
    p_1=\frac{1+\nu_1}{2},
    \qquad
    p_2=\frac{1+\nu_2}{2}.
\]
Тогда \(0<p_1<1<p_2\) и
\[
    \dot V
    \le
    -c_1V^{p_1}-c_2V^{p_2},
\]
где
\[
    c_1=a_1(2\mu)^{p_1},
    \qquad
    c_2=a_2(2\mu)^{p_2}.
\]
Следовательно,
\[
    T_F
    \le
    \frac1{c_1(1-p_1)}
    +
    \frac1{c_2(p_2-1)}.
\]
Эта граница не зависит от начальной точки.

\section{Недоступность глобального градиента}

В распределённой сети агенты начинают из произвольных точек
\[
    x_i(0)\ne x_j(0)
\]
и агент \(i\) может вычислить только \(\nabla f_i(x_i)\), но не
\[
    \nabla F(x)=\sum_{j=1}^{N}\nabla f_j(x).
\]
Поэтому требуется одновременно согласовать состояния \(x_i\) и
восстановить текущий средний локальный градиент
\[
    \theta_{\mathrm c}(t)
    =
    \frac1N
    \sum_{i=1}^{N}\nabla f_i(x_i(t)).
\]
Каждый агент хранит оценку \(\theta_i(t)\in\mathbb R^d\). Если
локальная функция дважды дифференцируема, производная её градиента
равна
\[
    h_i(t)
    =
    \frac{\mathrm d}{\mathrm dt}\nabla f_i(x_i(t))
    =
    \nabla^2f_i(x_i(t))\dot x_i(t).
\]
Таким образом, точный непрерывный трекинг требует вычисления
произведения гессиана на вектор либо иного способа получить
производную локального градиента.

\section{Фиксированно-временной трекинг градиента}

Инициализация задаётся условием
\[
    \theta_i(0)=\nabla f_i(x_i(0)).
\]
Далее
\[
    \dot\theta_i=h_i+\omega_i,
\]
где для симметричных весов \(a_{ij}=a_{ji}>0\)
\[
\begin{aligned}
    \omega_i
    ={}&
    -k_0\sum_{j\in\mathcal N_i}
    a_{ij}\operatorname{sign}(\theta_i-\theta_j)
    \\
    &-
    k_1\sum_{j\in\mathcal N_i}
    a_{ij}\operatorname{sig}^{\gamma_1}(\theta_i-\theta_j)
    \\
    &-
    k_2\sum_{j\in\mathcal N_i}
    a_{ij}\operatorname{sig}^{\gamma_2}(\theta_i-\theta_j),
\end{aligned}
\]
с параметрами
\[
    k_0,k_1,k_2>0,
    \qquad
    0<\gamma_1<1<\gamma_2.
\]
Нечётность парных взаимодействий даёт
\[
    \sum_{i=1}^{N}\omega_i=0.
\]
Следовательно,
\[
\begin{aligned}
    \frac{\mathrm d}{\mathrm dt}
    \sum_{i=1}^{N}\theta_i
    &=
    \sum_{i=1}^{N}h_i
    \\
    &=
    \frac{\mathrm d}{\mathrm dt}
    \sum_{i=1}^{N}\nabla f_i(x_i).
\end{aligned}
\]
С учётом инициализации сохраняется инвариант
\[
    \sum_{i=1}^{N}\theta_i(t)
    =
    \sum_{i=1}^{N}\nabla f_i(x_i(t)).
\]
Поэтому после достижения консенсуса оценки обязаны совпасть с
\(\theta_{\mathrm c}(t)\).

Пусть для всех рёбер и всех \(t\)
\[
    \lVert h_i(t)-h_j(t)\rVert
    \le
    \rho_\theta.
\]
При достаточно большом \(k_0\), зависящем от этой границы и графа,
знаковый член подавляет возмущение, а два степенных члена обеспечивают
существование времени \(T_\theta\), не зависящего от начальных
оценок, такого что
\[
    \theta_i(t)=\theta_{\mathrm c}(t),
    \qquad
    t\ge T_\theta.
\]
Порог \(k_0\) выбирается так, чтобы используемые далее оценки выполнялись для всех \(k\ge k_0\).

\section{Согласование локальных переменных}

Используя оценку среднего градиента, агент строит оптимизационное
направление
\[
    g_i=-\Psi(N\theta_i).
\]
Для согласования переменных добавляется сигнал
\[
\begin{aligned}
    u_i
    ={}&
    -\ell_0\sum_{j\in\mathcal N_i}
    a_{ij}\operatorname{sign}(x_i-x_j)
    \\
    &-
    \ell_1\sum_{j\in\mathcal N_i}
    a_{ij}\operatorname{sig}^{\delta_1}(x_i-x_j)
    \\
    &-
    \ell_2\sum_{j\in\mathcal N_i}
    a_{ij}\operatorname{sig}^{\delta_2}(x_i-x_j),
\end{aligned}
\]
где
\[
    \ell_0,\ell_1,\ell_2>0,
    \qquad
    0<\delta_1<1<\delta_2.
\]
Полная динамика равна
\[
    \dot x_i=g_i+u_i.
\]
До завершения трекинга направления \(g_i\) могут различаться и
действуют как возмущение консенсусной подсистемы. Предположим, что
\[
    \lVert g_i(t)-g_j(t)\rVert\le\rho_x
\]
на каждом ребре. Если \(\ell_0\) достаточно велик относительно
\(\rho_x\) и параметров графа, то существует равномерное время
\(T_x\), после которого
\[
    x_1(t)=\cdots=x_N(t).
\]
Знаковый член также обеспечивает инвариантность консенсусного
множества при продолжающемся воздействии различающихся направлений.

\section{Сведение к централизованному потоку}

Положим
\[
    T_0=\max\{T_\theta,T_x\}.
\]
При \(t\ge T_0\) одновременно выполнены
\[
    \theta_i
    =
    \frac1N\sum_{j=1}^{N}\nabla f_j(x_j)
\]
и
\[
    x_1=\cdots=x_N=x.
\]
Поэтому
\[
    N\theta_i
    =
    \sum_{j=1}^{N}\nabla f_j(x)
    =
    \nabla F(x),
\]
а консенсусный сигнал удовлетворяет \(u_i=0\). Все агенты следуют
одному уравнению
\[
    \dot x=-\Psi\bigl(\nabla F(x)\bigr).
\]
После момента \(T_0\) до достижения \(x^\star\) требуется не более
\(T_F\). Следовательно,
\[
    T_{\mathrm{total}}
    \le
    \max\{T_\theta,T_x\}+T_F.
\]

\begin{theorem}[Фиксированно-временная распределённая оптимизация]
Пусть граф связен, неориентирован и имеет симметричные положительные
веса. Пусть \(F=\sum_i f_i\) сильно выпукла, локальные функции
достаточно гладки, различия \(h_i-h_j\) и \(g_i-g_j\) равномерно
ограничены, а робастные коэффициенты выбраны выше соответствующих
порогов. Тогда существует число \(T_{\max}<\infty\), не зависящее от
начальных состояний, такое что
\[
    x_i(t)=x^\star
\]
для всех агентов и всех \(t\ge T_{\max}\).
\end{theorem}

\section{Реализация и границы модели}

Каждый агент передаёт соседям два вектора размерности \(d\):
\(x_i(t)\) и \(\theta_i(t)\). Поэтому объём одного полного обмена
имеет порядок
\[
    O(|\mathcal E|d),
\]
но с примерно вдвое большей постоянной, чем при передаче только
состояний. Локально требуются \(\nabla f_i(x_i)\) и
\[
    h_i=\nabla^2f_i(x_i)\dot x_i.
\]

Функция \(\operatorname{sign}\) разрывна в нуле, поэтому строгий
анализ ведётся в смысле решений Филиппова. Гладкое насыщение уменьшает
дребезг, но обычно заменяет точный консенсус попаданием в малую
окрестность.

Результат предполагает непрерывное время, статический связный
неориентированный граф, точные сообщения без задержек, доступность
производной локального градиента и известные границы возмущений. Для
ориентированных, асинхронных и изменяющихся сетей требуется другая
схема трекинга. Прямая дискретизация также не обязана сохранять
точную фиксированно-временную гарантию.

\lecture{38}{Нейронные сети и их обучение}

\sourcevideo
    {https://www.youtube.com/playlist?list=PLOzRYVm0a65fhKDS8cYIC7YCuVGJ4FOJx}
    {Week 10: Lecture 38}

Нейронная сеть задаёт параметрическое нелинейное отображение, а её
обучение сводится к минимизации эмпирического риска. При распределении
данных между агентами эта задача снова принимает форму суммы локальных
функций, но размерность параметра и стоимость передачи градиентов
становятся существенными.

\section{Обучение с учителем и один слой}

Пусть задан набор наблюдений
\[
    \mathcal D
    =
    \{(\xi_s,y_s)\}_{s=1}^{M},
    \qquad
    \xi_s\in\mathbb R^{n_0}.
\]
В обучении с учителем требуется построить отображение
\[
    \xi\longmapsto \widehat y=f_\theta(\xi),
\]
где \(\theta\) объединяет все параметры модели. В регрессии выход
обычно принадлежит \(\mathbb R^q\), а в классификации задаёт оценки
вероятностей классов.

Один полносвязный слой преобразует вход \(x\in\mathbb R^n\) по
формуле
\[
    h=\sigma(Wx+b),
\]
где
\[
    W\in\mathbb R^{m\times n},
    \qquad
    b\in\mathbb R^m,
\]
а функция \(\sigma\) применяется покоординатно. Число параметров
слоя равно \(mn+m\).

Смещение можно включить в матрицу весов. Для
\[
    \widetilde x
    =
    \begin{pmatrix}
        x\\
        1
    \end{pmatrix},
    \qquad
    \widetilde W
    =
    \begin{pmatrix}
        W & b
    \end{pmatrix}
\]
имеем
\[
    Wx+b=\widetilde W\widetilde x.
\]
Поэтому термин «веса» иногда используют для всех параметров слоя.

\section{Многослойная сеть и нелинейность}

Полносвязная сеть с \(L\) слоями задаётся рекурсией
\[
    h^0=\xi,
\]
\[
    h^\ell
    =
    \sigma_\ell
    \left(
        W^\ell h^{\ell-1}+b^\ell
    \right),
    \qquad
    \ell=1,\ldots,L-1,
\]
и выходным слоем
\[
    z=W^Lh^{L-1}+b^L.
\]
Векторы \(h^1,\ldots,h^{L-1}\) называются скрытыми
представлениями, а полный набор параметров равен
\[
    \theta
    =
    (W^1,b^1,\ldots,W^L,b^L).
\]

Без функций активации композиция нескольких слоёв остаётся аффинной.
Например,
\[
    W^2(W^1x+b^1)+b^2
    =
    W^2W^1x+W^2b^1+b^2.
\]
Следовательно, увеличение числа линейных слоёв не расширяет класс
представимых функций. Нелинейность вводится именно для устранения
этого вырождения.

Основной пример — функция
\[
    \operatorname{ReLU}(t)=\max\{0,t\}.
\]
Она непрерывна и кусочно-линейна, но не дифференцируема в нуле; при
обучении там выбирают допустимое значение субградиента.

Ширина определяется числом нейронов в слое, а глубина — числом
последовательных преобразований. Глубокая архитектура в некоторых
классах задач может представлять функцию существенно меньшим числом
параметров, чем неглубокая, однако это зависит от класса функций и
требуемой точности. Теоремы универсальной аппроксимации утверждают,
что при подходящей нелинейной активации и достаточном числе нейронов
можно для непрерывной функции \(g\) на компактном множестве \(K\)
добиться
\[
    \sup_{x\in K}
    \lVert f_\theta(x)-g(x)\rVert
    <\varepsilon.
\]
Точная версия результата зависит от архитектуры и выбранной функции
активации.

\section{Выходы регрессии и классификации}

В скалярной регрессии последний слой обычно имеет вид
\[
    \widehat y=w^{\mathsf T}h+b.
\]
Для многомерной регрессии используется несколько выходных координат.

В классификации по \(K\) классам последний линейный слой выдаёт
логиты
\[
    z=(z_1,\ldots,z_K)^{\mathsf T}.
\]
Они не обязаны быть неотрицательными и не суммируются в единицу.
Вероятности получают с помощью softmax:
\[
    p_r
    =
    \frac{\exp(z_r)}
    {\displaystyle\sum_{q=1}^{K}\exp(z_q)},
    \qquad
    r=1,\ldots,K.
\]
Тогда
\[
    p_r>0,
    \qquad
    \sum_{r=1}^{K}p_r=1,
\]
а предсказываемый класс равен
\[
    \widehat r
    =
    \operatorname*{arg\,max}_{r}p_r.
\]

Для численно устойчивого вычисления положим
\[
    m=\max_r z_r.
\]
Поскольку общий множитель \(\exp(-m)\) сокращается,
\[
    p_r
    =
    \frac{\exp(z_r-m)}
    {\displaystyle\sum_{q=1}^{K}\exp(z_q-m)}.
\]
Теперь \(z_r-m\le0\), поэтому наибольшая экспонента равна единице и
переполнение предотвращается.

\section{Эмпирический риск и функции потерь}

Потеря на одном объекте записывается как
\[
    \ell_s(\theta)
    =
    \ell\bigl(f_\theta(\xi_s),y_s\bigr).
\]
Эмпирический риск равен
\[
    \mathcal L(\theta)
    =
    \frac1M\sum_{s=1}^{M}\ell_s(\theta),
\]
а обучение формулируется задачей
\[
    \min_\theta \mathcal L(\theta).
\]

Для регрессии часто используется среднеквадратичная потеря
\[
    \ell_s(\theta)
    =
    \frac12
    \lVert f_\theta(\xi_s)-y_s\rVert^2,
\]
откуда
\[
    \mathcal L_{\mathrm{MSE}}(\theta)
    =
    \frac1{2M}
    \sum_{s=1}^{M}
    \lVert f_\theta(\xi_s)-y_s\rVert^2.
\]

В классификации правильный ответ задаётся индикаторным вектором
\[
    t=(t_1,\ldots,t_K)^{\mathsf T},
\]
а потеря определяется кросс-энтропией
\[
    \ell_{\mathrm{CE}}
    =
    -\sum_{r=1}^{K}t_r\log p_r.
\]
Если правильный класс имеет номер \(r^\star\), то
\[
    \ell_{\mathrm{CE}}=-\log p_{r^\star}.
\]
Минимизация этой величины увеличивает вероятность правильного класса.

Квадратичная потеря выпукла по непосредственному предсказанию, а
кросс-энтропия с softmax имеет хорошие свойства по логитам. Однако
отображение \(\theta\mapsto f_\theta(\xi)\) в многослойной сети
нелинейно и содержит композиции параметров. Поэтому
\(\mathcal L(\theta)\) в общем случае невыпукла, и градиентные методы
не гарантируют нахождение глобального минимума произвольной сети.

\section{Режимы вычисления градиента}

Полный градиент эмпирического риска равен
\[
    \nabla\mathcal L(\theta)
    =
    \frac1M
    \sum_{s=1}^{M}\nabla_\theta\ell_s(\theta),
\]
а шаг полного градиентного спуска имеет вид
\[
    \theta^{k+1}
    =
    \theta^k
    -
    \eta_k
    \frac1M
    \sum_{s=1}^{M}
    \nabla_\theta\ell_s(\theta^k).
\]
Одна итерация использует весь набор данных и потому даёт стабильное,
но дорогое направление.

В SGD выбирается случайный индекс \(s_k\):
\[
    \theta^{k+1}
    =
    \theta^k
    -
    \eta_k\nabla_\theta\ell_{s_k}(\theta^k).
\]
При равномерной выборке оценка несмещённая:
\[
    \mathbb E
    \left[
        \nabla_\theta\ell_{s_k}(\theta)
    \right]
    =
    \nabla\mathcal L(\theta).
\]
Итерация дешева, но дисперсия направления может быть большой.

Для мини-пакета \(B_k\subset\{1,\ldots,M\}\)
\[
    g_k
    =
    \frac1{|B_k|}
    \sum_{s\in B_k}
    \nabla_\theta\ell_s(\theta^k),
    \qquad
    \theta^{k+1}=\theta^k-\eta_kg_k.
\]
Увеличение \(|B_k|\) обычно уменьшает дисперсию, но повышает стоимость
итерации. Полный градиент соответствует \(|B_k|=M\), а чистый SGD —
\(|B_k|=1\). Мини-пакеты дают компромисс и позволяют эффективно
использовать параллельные вычислители.

Шум SGD иногда помогает не задерживаться в неглубоких локальных
областях, но это не является общей гарантией нахождения лучшего
минимума. Поведение зависит от геометрии функции, шага, размера
мини-пакета и структуры шума; слишком большая дисперсия вызывает
колебания и замедляет сходимость.

\section{Распределённое обучение}

Пусть данные разделены между агентами:
\[
    \mathcal D
    =
    \bigcup_{i=1}^{N}\mathcal D_i.
\]
Локальный риск имеет вид
\[
    f_i(\theta)
    =
    \frac1{|\mathcal D_i|}
    \sum_{(\xi,y)\in\mathcal D_i}
    \ell\bigl(f_\theta(\xi),y\bigr),
\]
а глобальная задача записывается как
\[
    \min_\theta
    \sum_{i=1}^{N}\omega_i f_i(\theta),
\]
где веса \(\omega_i\) могут учитывать размеры локальных наборов. Это
та же структура суммы локальных функций, что и в распределённой
оптимизации, но градиенты обычно вычисляются стохастически, а
размерность параметра может быть огромной.

Пусть число параметров равно \(d_\theta\). Передача полного вектора
параметров или градиента требует сообщения размерности \(d_\theta\).
При обмене центрального сервера с \(N\) агентами объём одного раунда
имеет порядок
\[
    O(Nd_\theta).
\]
Поэтому при миллионах или миллиардах параметров коммуникация может
оказаться дороже локального вычисления. Архитектура распределённого
обучения должна учитывать расположение данных, способ агрегации,
синхронность, топологию сети и объём сообщений. Эти вопросы являются
предметом следующих лекций.

\lecture{39}{Крупномасштабное машинное обучение}

\sourcevideo
    {https://www.youtube.com/playlist?list=PLOzRYVm0a65fhKDS8cYIC7YCuVGJ4FOJx}
    {Week 10: Lecture 39: Large Scale Machine Learning}

Пусть параметр модели равен \(w\in\mathbb R^d\), а данные распределены
между \(N\) вычислительными работниками. При большой размерности
модели передача параметров и градиентов может занимать больше времени,
чем локальное вычисление. Ниже рассматриваются две архитектуры
синхронного обучения --- сервер параметров и кольцевой all-reduce, ---
а затем приводится базовая оценка для стохастического градиентного
спуска.

\section{Распределённый градиентный шаг}

Работник \(i\) хранит локальный набор \(\mathcal D_i\) мощности
\(n_i=\lvert\mathcal D_i\rvert\). При
\[
    M=\sum_{i=1}^{N}n_i
\]
глобальная и локальные эмпирические функции имеют вид
\[
    F(w)
    =
    \frac1M
    \sum_{i=1}^{N}
    \sum_{\zeta\in\mathcal D_i}
    \ell(w;\zeta),
    \qquad
    F_i(w)
    =
    \frac1{n_i}
    \sum_{\zeta\in\mathcal D_i}
    \ell(w;\zeta).
\]
Следовательно,
\[
    F(w)
    =
    \sum_{i=1}^{N}\frac{n_i}{M}F_i(w),
    \qquad
    \nabla F(w)
    =
    \sum_{i=1}^{N}\frac{n_i}{M}\nabla F_i(w).
\]
Данные остаются у работников; для общего шага требуется агрегировать
только локальные градиенты.

На итерации \(k\) все работники используют один параметр \(w^k\) и
вычисляют
\[
    g_i^k=\nabla F_i(w^k).
\]
После агрегации
\[
    g^k
    =
    \sum_{i=1}^{N}\frac{n_i}{M}g_i^k,
    \qquad
    w^{k+1}=w^k-\eta_k g^k.
\]
При точных локальных градиентах имеем \(g^k=\nabla F(w^k)\), поэтому
это обычный централизованный градиентный шаг, распределённый между
несколькими устройствами.

\section{Сервер параметров}

В архитектуре сервера параметров центральный узел хранит \(w^k\),
рассылает его работникам, принимает локальные градиенты, вычисляет их
взвешенную сумму и формирует \(w^{k+1}\). Получаемые векторы можно
сразу добавлять к накопленной сумме, поэтому без учёта хранения самой
модели дополнительная память сервера имеет порядок \(O(d)\).

За один раунд сервер принимает \(N\) векторов градиента и отправляет
\(N\) копий параметра. Его коммуникационная нагрузка имеет порядок
\[
    O(Nd).
\]
При росте \(N\) центральный узел становится узким местом и единой
точкой отказа. Преимущество этой архитектуры состоит в простоте и
малом числе логических стадий одного глобального шага.

\section{Кольцевой all-reduce}

Для синхронного шага каждому работнику достаточно получить сумму
локальных вкладов. При равных размерах выборок это
\[
    G=\sum_{i=1}^{N}g_i,
    \qquad
    w^{k+1}=w^k-\eta_k\frac1N G.
\]
При различных \(n_i\) работник заранее умножает свой градиент на
\(n_i/M\), после чего all-reduce суммирует уже взвешенные векторы.

В кольцевой реализации каждый участник передаёт сообщения следующему
работнику и получает их от предыдущего. Предположим сначала, что
\(d\) делится на \(N\), и разобьём
\[
    g_i
    =
    \bigl(g_i^{(0)},\ldots,g_i^{(N-1)}\bigr),
    \qquad
    g_i^{(r)}\in\mathbb R^{d/N}.
\]
При неделимости используются блоки близких размеров или дополнение
нулями. На блоки делится передаваемый вектор, а не локальный набор
данных.

Первая фаза, scatter-reduce, длится \(N-1\) раундов. В каждом раунде
участник отправляет один блок вперёд по кольцу, принимает другой блок
и прибавляет его к соответствующей локальной сумме. После завершения
фазы каждый работник хранит один полностью просуммированный блок
\[
    G^{(r)}=\sum_{i=1}^{N}g_i^{(r)},
\]
а разные блоки распределены между разными работниками.

Например, для
\[
    g_0=(A_0,B_0,C_0),\qquad
    g_1=(A_1,B_1,C_1),\qquad
    g_2=(A_2,B_2,C_2)
\]
после двух раундов получены суммы
\[
    A=A_0+A_1+A_2,\qquad
    B=B_0+B_1+B_2,\qquad
    C=C_0+C_1+C_2,
\]
причём каждая из них пока находится у одного участника.

Во второй фазе, all-gather, готовые блоки циркулируют по кольцу без
дальнейшего суммирования. Ещё через \(N-1\) раундов каждый работник
получает полный вектор
\[
    G
    =
    \bigl(G^{(0)},\ldots,G^{(N-1)}\bigr)
    =
    \sum_{i=1}^{N}g_i
\]
и независимо выполняет то же обновление параметра.

\section{Коммуникационная стоимость}

В одном раунде передаётся примерно \(d/N\) компонент. Поэтому за
scatter-reduce один работник передаёт \((N-1)d/N\) компонент и столько
же --- за all-gather. Полный объём передачи на одного работника равен
\[
    2\frac{N-1}{N}d
    \longrightarrow
    2d
    \qquad
    \text{при }N\to\infty.
\]
Таким образом, объём данных на одного участника имеет порядок
\(O(d)\), но один глобальный шаг требует \(2(N-1)\) последовательных
коммуникационных раундов. Ring all-reduce устраняет центральный
агрегатор и равномерно распределяет трафик, однако его задержка может
расти с числом работников.

Сервер параметров выгоден малым числом стадий, но концентрирует
нагрузку порядка \(Nd\) в одном узле. Кольцо распределяет нагрузку и
оставляет на каждом работнике объём порядка \(d\), но требует
линейного числа раундов. Выбор зависит от размерности модели, числа
работников, пропускной способности, задержки и надёжности сети.
Путь также имеет малую степень вершин, но большой диаметр; звезда
имеет диаметр два, но перегружает центр. Статические экспоненциальные
графы позволяют получить промежуточный компромисс между степенью и
диаметром.

Обе рассмотренные архитектуры предполагают, что до коллективной
операции все работники имеют одинаковый \(w^k\). После получения
одного агрегированного градиента они снова приходят к общему
\(w^{k+1}\). Поэтому эти схемы распределяют вычисление
централизованного шага, но не согласуют изначально различные локальные
модели.

\section{Стохастический градиентный спуск}

Пусть на итерации \(k\) используется случайный градиент
\[
    g_k=g(w^k;\zeta_k),
    \qquad
    \mathbb E[g_k\mid w^k]=\nabla F(w^k),
\]
и его условная дисперсия ограничена:
\[
    \mathbb E\!\left[
        \lVert g_k-\nabla F(w^k)\rVert^2
        \,\middle|\,
        w^k
    \right]
    \le
    \sigma^2.
\]
Для обновления \(w^{k+1}=w^k-\eta g_k\), функции с
\(L\)-липшицевым градиентом и нижней гранью \(F_{\inf}\), при
\(0<\eta\le 1/L\) выполняется стандартная оценка
\[
    \frac1T
    \sum_{k=0}^{T-1}
    \mathbb E\lVert\nabla F(w^k)\rVert^2
    \le
    \frac{2\bigl(F(w^0)-F_{\inf}\bigr)}{\eta T}
    +
    L\eta\sigma^2.
\]
Первое слагаемое уменьшается при росте шага, а второе отражает влияние
шума. При \(\eta\asymp T^{-1/2}\) получаем
\[
    \frac1T
    \sum_{k=0}^{T-1}
    \mathbb E\lVert\nabla F(w^k)\rVert^2
    =
    O\!\left(T^{-1/2}\right).
\]
Константа зависит от начального функционального разрыва, гладкости и
дисперсии. Для невыпуклой функции эта оценка характеризует
приближение к стационарности, а не к глобальному минимуму. Чтобы
средний квадрат нормы градиента не превосходил \(\varepsilon\), базовая
оценка требует \(T=O(\varepsilon^{-2})\) итераций.

\section{Переход к децентрализованному SGD}

Сервер параметров и ring all-reduce воспроизводят один
централизованный SGD-шаг. В полностью децентрализованной схеме агент
\(i\) хранит собственную переменную \(w_i^k\) и обменивается только с
соседями. Тогда к стохастической ошибке градиента добавляются
рассогласование локальных моделей и зависимость скорости смешивания
от спектра коммуникационного графа. Эти эффекты рассматриваются в
следующей лекции.

\lecture{40}{Децентрализованный SGD}

\sourcevideo
    {https://www.youtube.com/playlist?list=PLOzRYVm0a65fhKDS8cYIC7YCuVGJ4FOJx}
    {Week 11: Lecture 40: Decentralized Stochastic Gradient Descent}

В параллельном обучении все работники используют одну модель и после
каждого шага глобально усредняют градиенты. Децентрализованный SGD
отказывается от центрального сервера: агент хранит собственную модель,
обрабатывает локальные данные и обменивается параметрами только с
соседями. Поэтому алгоритм должен одновременно решать задачу
оптимизации и подавлять рассогласование локальных моделей.

\section{Задача и параллельный SGD}

Пусть агент \(i\) хранит распределение или набор данных
\(\mathcal D_i\) и минимизирует локальную функцию
\[
    f_i(x)
    =
    \mathbb E_{\xi\sim\mathcal D_i}
    \bigl[\ell(x;\xi)\bigr].
\]
Глобальная задача имеет вид
\[
    \min_{x\in\mathbb R^d}F(x),
    \qquad
    F(x)=\frac1N\sum_{i=1}^{N}f_i(x).
\]
Её децентрализованная форма записывается как
\[
\begin{aligned}
    \min_{x_1,\ldots,x_N}\quad
    &\frac1N\sum_{i=1}^{N}f_i(x_i),
    \\
    \text{при условии}\quad
    &x_1=\cdots=x_N.
\end{aligned}
\]
Каждый агент знает только свою функцию и не обязан передавать
исходные данные другим участникам.

На одном узле SGD записывается как
\[
    x^{k+1}=x^k-\eta_k g(x^k;\xi_k),
    \qquad
    \mathbb E[g(x;\xi)]=\nabla F(x).
\]
При несмещённости и ограниченной дисперсии градиента типичная оценка
стационарности для гладкой функции содержит стохастический член
порядка
\[
    \frac1T\sum_{k=0}^{T-1}
    \mathbb E\lVert\nabla F(x^k)\rVert^2
    =
    O\!\left(
        \frac{\sigma}{\sqrt T}
    \right),
\]
наряду с начальным функциональным разрывом и константой гладкости.

В параллельном SGD все \(N\) работников используют одну модель
\(x^k\), независимо вычисляют
\[
    g_i^k=g_i(x^k;\xi_i^k)
\]
и выполняют глобальное обновление
\[
    x^{k+1}
    =
    x^k-
    \eta_k\frac1N\sum_{i=1}^{N}g_i^k.
\]
Если шумы независимы и имеют дисперсию не более \(\sigma^2\), то
дисперсия среднего градиента уменьшается примерно до
\(\sigma^2/N\), а стохастический член принимает порядок
\[
    O\!\left(
        \frac{\sigma}{\sqrt{NT}}
    \right).
\]
Это линейное ускорение по числу работников требует достаточной
независимости шумов и пренебрежимо малой стоимости синхронизации.

\section{Коммуникационный граф и алгоритм}

Глобальное усреднение поддерживает равенство всех моделей после
каждой итерации, но создаёт узкое место. У сервера параметров нагрузка
центрального узла имеет порядок \(O(Nd)\), а ring all-reduce требует
\(2(N-1)\) последовательных коммуникационных стадий. В
децентрализованной схеме обмен задаётся связным неориентированным
графом
\[
    \mathcal G=(\mathcal V,\mathcal E).
\]
Агент \(i\) взаимодействует только с вершинами
\(j\in\mathcal N_i\). Матрица смешивания
\(W=(w_{ij})\in\mathbb R^{N\times N}\) удовлетворяет
\[
    w_{ij}\ge0,
    \qquad
    w_{ij}=0
    \quad
    \text{при }
    j\notin\mathcal N_i\cup\{i\},
\]
\[
    W\mathbf1=\mathbf1,
    \qquad
    \mathbf1^{\mathsf T}W=\mathbf1^{\mathsf T}.
\]
Таким образом, \(W\) является дважды стохастической.

На итерации \(k\) агент вычисляет локальный стохастический градиент
\[
    g_i^k=g_i(x_i^k;\xi_i^k),
\]
делает локальный шаг
\[
    x_i^{k+\frac12}
    =
    x_i^k-\eta_k g_i^k
\]
и смешивает промежуточные модели соседей:
\[
    x_i^{k+1}
    =
    \sum_{j=1}^{N}
    w_{ij}x_j^{k+\frac12}.
\]
Итоговая рекурсия равна
\[
    x_i^{k+1}
    =
    \sum_{j=1}^{N}
    w_{ij}
    \bigl(
        x_j^k-\eta_k g_j^k
    \bigr).
\]
После объединения локальных векторов
\[
    \mathbf x^k
    =
    \begin{pmatrix}
        x_1^k\\
        \vdots\\
        x_N^k
    \end{pmatrix},
    \qquad
    \mathbf g^k
    =
    \begin{pmatrix}
        g_1^k\\
        \vdots\\
        g_N^k
    \end{pmatrix},
\]
получаем матричную форму
\[
    \mathbf x^{k+1}
    =
    (W\otimes I_d)
    \bigl(
        \mathbf x^k-\eta_k\mathbf g^k
    \bigr).
\]
Можно сначала смешать модели, а затем вычислить локальный градиентный
шаг. Обе схемы используют те же два механизма, но порождают разные
траектории; далее рассматривается порядок «локальный шаг, затем
смешивание».

\section{Средняя модель и рассогласование}

Определим среднюю модель
\[
    \overline x^k
    =
    \frac1N
    \sum_{i=1}^{N}x_i^k.
\]
Благодаря столбцовой стохастичности
\[
\begin{aligned}
    \overline x^{k+1}
    &=
    \frac1N
    \sum_{i,j}
    w_{ij}
    \bigl(
        x_j^k-\eta_k g_j^k
    \bigr)
    \\
    &=
    \overline x^k
    -
    \frac{\eta_k}{N}
    \sum_{j=1}^{N}g_j^k.
\end{aligned}
\]
Следовательно,
\[
    \overline x^{k+1}
    =
    \overline x^k-\eta_k\overline g^k,
    \qquad
    \overline g^k
    =
    \frac1N
    \sum_{i=1}^{N}g_i^k.
\]
Это похоже на параллельный SGD, но градиенты вычисляются в различных
точках \(x_i^k\). Поэтому \(\overline g^k\) не является точным
стохастическим градиентом в точке \(\overline x^k\).

Пусть
\[
    J
    =
    \frac1N
    \mathbf1\mathbf1^{\mathsf T}
\]
и
\[
    \mathbf e^k
    =
    \mathbf x^k
    -
    \mathbf1\otimes\overline x^k.
\]
Чистый смешивающий шаг действует на неконсенсусную компоненту
матрицей \(W-J\). Введём коэффициент смешивания
\[
    \rho
    =
    \left\lVert
        W-
        \frac1N
        \mathbf1\mathbf1^{\mathsf T}
    \right\rVert_2.
\]
Для связной апериодической сети
\[
    0\le\rho<1,
\]
и без градиентного возмущения
\[
    \lVert\mathbf e^{k+1}\rVert
    \le
    \rho\lVert\mathbf e^k\rVert.
\]
Обычно спектральным разрывом называют \(1-\rho\), а не саму
величину \(\rho\). Для симметричной дважды стохастической матрицы
\[
    \rho
    =
    \max_{r\ge2}
    |\lambda_r(W)|.
\]

\section{Топология и скорость смешивания}

Для полного графа можно выбрать
\[
    W
    =
    \frac1N
    \mathbf1\mathbf1^{\mathsf T},
\]
поэтому \(\rho=0\) и точное среднее достигается за один раунд. Цена
состоит в обмене каждого агента с \(N-1\) соседями.

В пути и кольце степень вершины ограничена константой, поэтому один
раунд дешёв, но информация распространяется медленно. Для стандартных
локальных весов на кольце
\[
    \rho
    =
    1-
    \Theta\!\left(
        \frac1{N^2}
    \right),
\]
следовательно,
\[
    1-\rho
    =
    \Theta\!\left(
        \frac1{N^2}
    \right).
\]
Число раундов, необходимое для существенного уменьшения
рассогласования, поэтому растёт квадратично по порядку \(N\).

Диаметр графа равен \(N-1\) для пути,
\(\lfloor N/2\rfloor\) для кольца и единице для полного графа.
Малый диаметр обычно помогает распространению информации, но не
определяет скорость смешивания полностью: точная характеристика
задаётся спектром \(W\).

При выборе сети приходится одновременно контролировать максимальную
степень, диаметр и коэффициент \(\rho\). Полный граф минимизирует
число раундов, а путь и кольцо минимизируют локальную степень.
Экспоненциальные графы дают промежуточный вариант: небольшую степень и
существенно более быстрое распространение информации, чем в кольце.

\section{Ошибка и коммуникационная стоимость}

При липшицевости локальных градиентов
\[
    \left\lVert
        \nabla f_i(x_i^k)
        -
        \nabla f_i(\overline x^k)
    \right\rVert
    \le
    L_i
    \lVert
        x_i^k-\overline x^k
    \rVert.
\]
Поэтому рассогласование моделей непосредственно искажает направление
среднего градиентного шага. Анализ децентрализованного SGD должен
одновременно контролировать оптимизационную ошибку средней модели,
стохастический шум и ошибку консенсуса. Первая уменьшается за счёт
градиентных шагов, вторая зависит от дисперсии и размеров
мини-пакетов, а третья определяется графом и матрицей смешивания.

На одной итерации агент \(i\) передаёт вектор размерности \(d\)
каждому соседу. Его локальная коммуникационная нагрузка имеет порядок
\[
    O(d_i d),
\]
а максимальная нагрузка по сети --- порядок
\[
    O(d_{\max}d),
    \qquad
    d_{\max}
    =
    \max_i d_i.
\]
Разреженный граф уменьшает стоимость одного раунда, но обычно
увеличивает число итераций из-за медленного смешивания. Поэтому
сравнивать архитектуры следует по полной стоимости достижения заданной
точности, а не только по числу передаваемых компонент за один шаг.

\section{Предпосылки и границы модели}

Базовый анализ предполагает синхронные итерации, фиксированный
связный граф, дважды стохастическую матрицу смешивания, точную передачу
параметров и несмещённые локальные градиенты. Линейное ускорение
параллельного SGD дополнительно требует независимости шумов. При
неодинаковых распределениях данных, задержках, потерях сообщений,
отказах и асинхронности возникают дополнительные ошибки и нужны
модифицированные алгоритмы.

Децентрализация устраняет центральный узел и глобальную
синхронизацию, но не делает коммуникацию бесплатной. Итоговый темп
определяется совместным действием стохастического шума, локальной
геометрии задачи и спектральных свойств сети.

\lecture{41}{Децентрализованный SGD: скорость и устойчивость}

\sourcevideo
    {https://www.youtube.com/playlist?list=PLOzRYVm0a65fhKDS8cYIC7YCuVGJ4FOJx}
    {Week 11: Lecture 41: Decentralized Stochastic Gradient Descent -2}

В этой лекции изучается влияние коммуникационного графа на скорость
децентрализованного стохастического градиентного спуска. Отдельно
рассматривается переходный период до появления линейного ускорения и
способ ограничить влияние агента, отклоняющегося от протокола.

\section{Модель и предположения}

Пусть \(N\) агентов решают задачу
\[
    \min_{x\in\mathbb R^d}F(x),
    \qquad
    F(x)=\frac1N\sum_{i=1}^{N}f_i(x),
\]
где
\[
    f_i(x)
    =
    \mathbb E_{\xi\sim\mathcal D_i}
    \bigl[\ell(x;\xi)\bigr].
\]
Агент \(i\) хранит локальную модель \(x_i^k\) и вычисляет
стохастический градиент
\[
    g_i^k=g_i(x_i^k;\xi_i^k).
\]
Один шаг децентрализованного SGD имеет вид
\[
    x_i^{k+\frac12}
    =
    x_i^k-\gamma g_i^k,
    \qquad
    x_i^{k+1}
    =
    \sum_{j=1}^{N}w_{ij}x_j^{k+\frac12}.
\]
Матрица смешивания \(W=(w_{ij})\) согласована с графом и предполагается
дважды стохастической.

Каждая \(f_i\) имеет \(L\)-липшицев градиент:
\[
    \lVert\nabla f_i(x)-\nabla f_i(y)\rVert
    \le
    L\lVert x-y\rVert.
\]
Стохастические градиенты несмещены,
\[
    \mathbb E[g_i^k\mid x_i^k]
    =
    \nabla f_i(x_i^k),
\]
и имеют ограниченную условную дисперсию:
\[
    \mathbb E
    \left[
        \lVert
            g_i^k-\nabla f_i(x_i^k)
        \rVert^2
        \,\middle|\,
        x_i^k
    \right]
    \le
    \sigma^2.
\]
Шумы различных агентов считаются независимыми.

Неоднородность локальных данных измеряется условием
\[
    \frac1N
    \sum_{i=1}^{N}
    \lVert
        \nabla f_i(x)-\nabla F(x)
    \rVert^2
    \le
    b^2.
\]
При \(b=0\) локальные градиенты совпадают с глобальным. Для конечных
эмпирических выборок это является идеализированным случаем даже тогда,
когда данные получены из одного распределения.

\section{Смешивание и оценка скорости}

Положим
\[
    J=\frac1N\mathbf1\mathbf1^{\mathsf T},
    \qquad
    \rho=\lVert W-J\rVert_2.
\]
Для связной апериодической сети
\[
    0\le\rho<1.
\]
Малое \(\rho\) соответствует быстрому смешиванию, а величина
\(1-\rho\) играет роль спектрального разрыва.

Обозначим среднюю модель через
\[
    \bar x^k
    =
    \frac1N\sum_{i=1}^{N}x_i^k.
\]
Для шага порядка \(\gamma\asymp T^{-1/2}\) получаем оценку
следующей структуры:
\[
\begin{aligned}
    \frac1T
    \sum_{k=0}^{T-1}
    \mathbb E
    \lVert\nabla F(\bar x^k)\rVert^2
    \lesssim{}&
    \frac{\sigma}{\sqrt{NT}}
    \\
    &+
    \frac{
        \rho^{2/3}\sigma^{2/3}
    }{
        (1-\rho)^{1/3}T^{2/3}
    }
    \\
    &+
    \frac{
        \rho^{2/3}b^{2/3}
    }{
        (1-\rho)^{2/3}T^{2/3}
    }.
\end{aligned}
\]
Знак \(\lesssim\) скрывает численные константы, гладкость и начальный
функциональный разрыв.

Первое слагаемое совпадает с ведущим стохастическим членом
параллельного SGD. Второе возникает из-за отсутствия точной глобальной
синхронизации, третье — из-за различия локальных функций. Сетевые
члены убывают как \(T^{-2/3}\), то есть асимптотически быстрее, чем
\(T^{-1/2}\). Поэтому при достаточно большом числе итераций ведущим
становится
\[
    O\left(\frac{\sigma}{\sqrt{NT}}\right),
\]
и возникает линейное ускорение по числу агентов.

Для дальнейшего анализа используются порядки \(T^{-1/2}\), \(T^{-2/3}\) и зависимости от \(\rho\), \(1-\rho\), \(\sigma\) и \(b\); точные константы не влияют на качественное сравнение режимов.

\section{Переходный период}

Переходный период определяется моментом, когда сетевые члены становятся
не больше ведущего стохастического слагаемого. При однородных данных,
то есть \(b=0\), сравнение первых двух членов даёт порядок
\[
    T_{\mathrm{tr}}
    =
    O\left(
        \frac{\rho^4N^3}{(1-\rho)^2}
    \right),
\]
если опустить зависимость от \(\sigma\) и других констант.

При неоднородных данных порядок ухудшается до
\[
    T_{\mathrm{tr}}
    =
    O\left(
        \frac{\rho^4N^3}{(1-\rho)^4}
    \right).
\]
Следовательно, линейное ускорение является асимптотическим свойством:
на начальном участке сетевое рассогласование может доминировать.
Переходный период особенно велик при \(\rho\), близком к единице.

Коммуникационная стоимость одной итерации зависит от максимальной
степени графа. Если она равна \(d_{\max}\), то один агент передаёт
порядка
\[
    O(d_{\max}d)
\]
чисел. Поэтому требуется одновременно ограничивать степень графа и
обеспечивать малое значение \(\rho\).

Полный граф быстро смешивает информацию, но имеет степень порядка
\(N\). Кольцо имеет постоянную степень, однако
\[
    1-\rho=\Theta(N^{-2}),
\]
из-за чего его переходный период быстро растёт.

\section{Статический экспоненциальный граф}

Пронумеруем вершины числами
\[
    0,\ldots,N-1.
\]
Из вершины \(i\) проводятся дуги к вершинам
\[
    i+2^r\pmod N,
    \qquad
    r=0,\ldots,\lceil\log_2N\rceil-1.
\]
Таким образом, агент соединяется с вершинами на циклических
расстояниях
\[
    1,2,4,8,\ldots.
\]
Исходящая степень имеет порядок
\[
    O(\log N).
\]

Пусть используется \(q\) ненулевых смещений. Один из вариантов
матрицы смешивания задаётся формулой
\[
    w_{ij}
    =
    \begin{cases}
        \dfrac1{q+1},
        &
        j=i,
        \\[2mm]
        \dfrac1{q+1},
        &
        j=i+2^r\pmod N
        \text{ для некоторого }r,
        \\[2mm]
        0,
        &
        \text{иначе}.
    \end{cases}
\]
Благодаря циклической симметрии матрица является строчно- и
столбцово-стохастической.

Экспоненциальный граф распространяет информацию существенно быстрее
пути или кольца, сохраняя степень порядка \(\log N\). В лекции он
выделяется как удачный компромисс между стоимостью одной итерации и
переходным периодом, но оптимальность среди всех возможных топологий
не утверждается.

\section{Отклоняющийся агент}

Рассмотрим локальные функции
\[
    f_i(x)=\frac12(x-i)^2,
    \qquad
    i=1,\ldots,5.
\]
Глобальный минимум равен
\[
    x^\star=3.
\]
Пусть агенты \(1,\ldots,4\) выполняют протокол, а агент \(5\)
игнорирует согласование и минимизирует только
\[
    f_5(x)=\frac12(x-5)^2.
\]
Его обновление имеет вид
\[
    x_5^{k+1}
    =
    x_5^k-\eta(x_5^k-5),
\]
поэтому он стремится к точке \(5\), продолжая передавать соседям своё
состояние. Обычное усреднение способно сместить честных агентов от
социального решения \(3\).

Это более узкая модель, чем полный византийский противник: агент
отклоняется от протокола, но не посылает различным соседям
произвольные и разные сообщения.

\section{Отсечённое смешивание}

Для \(z\in\mathbb R^d\) и \(\tau>0\) определим
\[
    \operatorname{clip}(z,\tau)
    =
    \min
    \left\{
        1,
        \frac{\tau}{\lVert z\rVert}
    \right\}z
\]
при \(z\ne0\), а
\[
    \operatorname{clip}(0,\tau)=0.
\]
Тогда
\[
    \lVert\operatorname{clip}(z,\tau)\rVert
    \le\tau.
\]
При \(\lVert z\rVert\le\tau\) вектор не изменяется; при большей норме
сохраняется его направление, но ограничивается величина.

Вместо обычного gossip агент вычисляет
\[
    z_i^k
    =
    x_i^k
    +
    \sum_{j\in\mathcal N_i}
    w_{ij}
    \operatorname{clip}
    \left(
        x_j^k-x_i^k,
        \tau
    \right)
\]
и затем выполняет локальный шаг
\[
    x_i^{k+1}
    =
    z_i^k-\eta\nabla f_i(z_i^k).
\]
Вклад одного соседа в смешанное состояние ограничен величиной,
пропорциональной \(\tau\), поэтому удалённая локальная модель не может
иметь неограниченного влияния.

При большом \(\tau\) алгоритм близок к обычному усреднению и слабо
ограничивает отклоняющегося агента. Малое \(\tau\) повышает
устойчивость, но замедляет согласование честных агентов и может внести
смещение. Численный пример лекции показывает улучшение поведения
честных агентов при уменьшении порога, но сам по себе не является
общей теоремой.

Для строгой защиты от византийских противников необходимы
дополнительные предположения о числе противников, связности графа,
локальных функциях, пороге отсечения и модели передаваемых сообщений.
Одно только отсечение таких гарантий не даёт.

\lecture{42}{Введение в федеративное обучение}

\sourcevideo
    {https://www.youtube.com/playlist?list=PLOzRYVm0a65fhKDS8cYIC7YCuVGJ4FOJx}
    {Week 11: Lecture 42: Introduction to Federated Learning}

Федеративное обучение предназначено для обучения общей модели по
данным, распределённым между клиентами. Исходные данные остаются на
устройствах или в организациях, а сервер получает только результаты
локальных вычислений. Клиентами могут быть мобильные телефоны,
устройства интернета вещей, больницы, банки и другие организации с
закрытыми наборами данных.

\section{Постановка задачи}

Пусть имеется \(K\) клиентов. Клиент \(k\) хранит набор
\(\mathcal D_k\) мощности \(n_k=\lvert\mathcal D_k\rvert\), а общее
число объектов равно
\[
    n=\sum_{k=1}^{K}n_k.
\]
Для функции потерь \(\ell(w;\zeta)\) локальная и глобальная функции
задаются формулами
\[
    F_k(w)
    =
    \frac1{n_k}
    \sum_{\zeta\in\mathcal D_k}
    \ell(w;\zeta),
    \qquad
    F(w)
    =
    \sum_{k=1}^{K}\frac{n_k}{n}F_k(w).
\]
Задача обучения имеет вид
\[
    \min_{w\in\mathbb R^d}F(w).
\]
Вес \(n_k/n\) учитывает объём локальных данных.

Типичный пример --- предсказание следующего слова на мобильной
клавиатуре. История ввода создаётся пользователем, может содержать
чувствительную информацию и существенно различается между
устройствами. Централизованный сбор таких данных увеличивает
требования к хранению, передаче и защите информации.

\section{Коммуникационный раунд}

Обычный сервер параметров на каждом шаге рассылает модель, получает
локальные градиенты и сразу обновляет параметры. При миллионах
клиентов эта схема требует высокой пропускной способности и
одновременной доступности большого числа устройств. Кроме того,
клиенты различаются по скорости вычислений, памяти, заряду батареи и
качеству соединения, а передаваемые обновления могут раскрывать
свойства локальных данных.

В федеративной схеме сервер на раунде \(t\) выбирает подмножество
\[
    S_t\subseteq\{1,\ldots,K\}.
\]
Если в раунде участвует доля \(C\in(0,1]\) клиентов, обычно берут
\[
    m=\max\{1,\lfloor CK\rfloor\}.
\]
Сервер отправляет выбранным клиентам текущую модель \(w^t\), клиенты
обучают её на своих данных и возвращают локальные модели или их
изменения. Наборы \(\mathcal D_k\) устройство не покидают.

Частичное участие необходимо, поскольку устройство может быть
отключено, занято пользователем, недостаточно заряжено или связано с
сервером медленным каналом. В практических системах локальное обучение
часто запускают только при выполнении заданных условий, например при
подключении телефона к питанию и беспроводной сети.

\section{Локальное обучение и агрегирование}

Получив \(w^t\), клиент \(k\) полагает
\[
    w_k^{t,0}=w^t
\]
и выполняет стохастические градиентные шаги
\[
    w_k^{t,s+1}
    =
    w_k^{t,s}
    -
    \eta\,
    g_k\bigl(w_k^{t,s};\mathcal B_k^{t,s}\bigr),
\]
где \(\eta>0\) --- локальный шаг, а
\(\mathcal B_k^{t,s}\subseteq\mathcal D_k\) --- мини-пакет. После
\(\tau_k\) шагов клиент отправляет серверу \(w_k^{t,\tau_k}\) или
приращение
\[
    \Delta_k^t=w_k^{t,\tau_k}-w^t.
\]

Если размер мини-пакета равен \(B\), а клиент выполняет \(E\)
локальных эпох, то
\[
    \tau_k
    =
    E\left\lceil\frac{n_k}{B}\right\rceil.
\]
При делимости \(n_k\) на \(B\) получаем \(\tau_k=En_k/B\). Поэтому
даже при одинаковых \(E\) и \(B\) клиенты с различными объёмами данных
выполняют разное число обновлений.

После получения локальных моделей сервер вычисляет
\[
    w^{t+1}
    =
    \sum_{k\in S_t}
    \frac{n_k}{\sum_{j\in S_t}n_j}
    w_k^{t,\tau_k}.
\]
Эквивалентная запись через локальные приращения имеет вид
\[
    w^{t+1}
    =
    w^t
    +
    \sum_{k\in S_t}
    \frac{n_k}{\sum_{j\in S_t}n_j}
    \Delta_k^t.
\]
Такое усреднение присваивает больший вес клиентам с большим числом
объектов. Подробная схема FedAvg рассматривается в следующей лекции.

Несколько локальных шагов уменьшают число коммуникационных раундов,
но усиливают расхождение локальных моделей. В сервере параметров
клиент обычно передаёт градиент после одного шага, тогда как в
федеративной схеме \(\tau_k\) может быть значительно больше единицы.
Чем дольше клиент оптимизирует собственную функцию \(F_k\), тем сильнее
его траектория отклоняется от направления глобальной функции \(F\).

\section{Статистическая и системная неоднородность}

В федеративных задачах локальные наборы обычно не имеют одинакового
распределения. Пользователи вводят разные слова, используют разные
языки и обладают различными привычками, поэтому функции \(F_i\) и
\(F_j\) могут существенно отличаться. При большом числе локальных
шагов это различие создаёт дрейф локальной модели к минимизатору
\(F_k\), а не глобальной функции.

Системная неоднородность связана с различием процессоров, объёмов
памяти, скоростей соединения, уровней заряда и времени локального
обучения. Часть выбранных клиентов может не завершить вычисления к
моменту агрегации; такие участники образуют проблему отстающих
клиентов. Таким образом, параметры \(C\), \(B\), \(E\), \(\eta\) и
\(\tau_k\) одновременно определяют вычислительную нагрузку,
коммуникационную стоимость и величину локального дрейфа.

\section{Конфиденциальность}

Локальное хранение \(\mathcal D_k\) уменьшает необходимость передачи
исходных данных, но не обеспечивает приватность автоматически. По
градиенту или обновлению модели иногда можно восстановить свойства
локальной выборки. Поэтому федеративную схему дополняют защищённой
агрегацией, шифрованием, дифференциальной приватностью, ограничением
нормы и зашумлением обновлений. При защищённой агрегации сервер видит
сумму клиентских обновлений, но не каждое обновление отдельно. Эти
механизмы являются дополнительными компонентами, а не следствием
самого факта локального обучения.

\section{Основные архитектуры}

В cross-device постановке участвует очень большое число мобильных и
IoT-устройств. В одном раунде доступна лишь малая доля клиентов, связь
нестабильна, вычислительные возможности различаются, а локальные
наборы обычно малы и неоднородны.

Cross-silo постановка объединяет меньшее число организаций, например
больниц, банков или исследовательских центров. Такие клиенты чаще
имеют стабильное соединение и более мощную инфраструктуру, однако их
данные также остаются локальными и могут существенно различаться.
Обе постановки используют серверную агрегацию общей модели и требуют
отдельных механизмов защиты информации.

Федеративное обучение не следует отождествлять с децентрализованным
SGD. В децентрализованном методе центрального сервера нет, и агенты
обмениваются моделями с соседями по графу. В стандартной федеративной
схеме сервер выбирает клиентов, рассылает модель, собирает локальные
обновления и выполняет агрегацию. Её определяющими признаками являются
локальное хранение данных, частичное участие и несколько локальных
шагов между коммуникациями.

\lecture{43}{FedAvg}

\sourcevideo
    {https://www.youtube.com/playlist?list=PLOzRYVm0a65fhKDS8cYIC7YCuVGJ4FOJx}
    {Week 11: Lecture 43: FedAvg Algorithm}

FedAvg реализует основной цикл федеративного обучения: сервер
рассылает текущую модель, выбранные клиенты выполняют несколько
локальных шагов, после чего сервер усредняет полученные модели. В
отличие от обычного распределённого градиентного спуска, между двумя
коммуникационными раундами здесь выполняется несколько локальных
обновлений.

\section{Локальная и глобальная задачи}

Пусть имеется \(K\) клиентов. Клиент \(k\) хранит набор
\[
    \mathcal D_k
    =
    \{\zeta_{k,1},\ldots,\zeta_{k,n_k}\},
    \qquad
    n_k=\lvert\mathcal D_k\rvert,
\]
и локальную функцию
\[
    F_k(x)
    =
    \frac1{n_k}
    \sum_{j=1}^{n_k}
    \ell(x;\zeta_{k,j}),
    \qquad
    x\in\mathbb R^d.
\]
Если
\[
    n=\sum_{k=1}^{K}n_k,
    \qquad
    p_k=\frac{n_k}{n},
\]
то глобальная функция равна
\[
    F(x)
    =
    \sum_{k=1}^{K}p_kF_k(x),
    \qquad
    \sum_{k=1}^{K}p_k=1,
\]
а задача обучения имеет вид
\[
    \min_{x\in\mathbb R^d}F(x).
\]
Вес \(p_k\) учитывает объём данных клиента.

Алгоритм задаётся числом коммуникационных раундов \(T\), долей
участвующих клиентов \(C\in(0,1]\), размером локального мини-пакета
\(B\), числом локальных эпох \(E\) и шагом \(\eta>0\). В одном
раунде сервер выбирает
\[
    m=\max\{1,\lfloor CK\rfloor\}
\]
клиентов. Если клиент \(k\) за одну эпоху обрабатывает весь набор
\(\mathcal D_k\), то число его локальных шагов равно
\[
    \tau_k
    =
    E\left\lceil\frac{n_k}{B}\right\rceil.
\]
При одинаковых \(E\) и \(B\) клиенты с различными объёмами данных
могут выполнить разное число обновлений.

\section{Один раунд FedAvg}

На раунде \(t\) сервер выбирает множество
\[
    S_t\subseteq\{1,\ldots,K\},
    \qquad
    \lvert S_t\rvert=m,
\]
и отправляет каждому \(k\in S_t\) модель \(x^t\). Клиент
инициализирует
\[
    x_k^{t,0}=x^t
\]
и для \(s=0,\ldots,\tau_k-1\) выбирает мини-пакет
\(\mathcal B_k^{t,s}\subseteq\mathcal D_k\), вычисляет
\[
    g_k^{t,s}
    =
    \frac1{\lvert\mathcal B_k^{t,s}\rvert}
    \sum_{\zeta\in\mathcal B_k^{t,s}}
    \nabla_x\ell(x_k^{t,s};\zeta)
\]
и выполняет шаг
\[
    x_k^{t,s+1}
    =
    x_k^{t,s}-\eta g_k^{t,s}.
\]
После \(\tau_k\) шагов сервер получает модель
\(x_k^{t,\tau_k}\) или эквивалентное приращение
\[
    \Delta_k^t
    =
    x_k^{t,\tau_k}-x^t.
\]

Положим
\[
    n_{S_t}=\sum_{j\in S_t}n_j.
\]
Федеративное усреднение задаётся формулой
\[
    x^{t+1}
    =
    \sum_{k\in S_t}
    \frac{n_k}{n_{S_t}}x_k^{t,\tau_k}.
\]
Через приращения то же обновление записывается как
\[
    x^{t+1}
    =
    x^t
    +
    \sum_{k\in S_t}
    \frac{n_k}{n_{S_t}}\Delta_k^t.
\]
Если участвуют все клиенты, то \(n_{S_t}=n\), и серверные веса
совпадают с весами глобальной функции.

\section{Связь с градиентным спуском и локальный дрейф}

Предположим, что участвуют все клиенты и каждый выполняет один точный
градиентный шаг. Тогда
\[
    x_k^{t,1}
    =
    x^t-\eta\nabla F_k(x^t),
\]
и после усреднения
\[
\begin{aligned}
    x^{t+1}
    &=
    \sum_{k=1}^{K}p_k
    \bigl(x^t-\eta\nabla F_k(x^t)\bigr)
    \\
    &=
    x^t-\eta\nabla F(x^t).
\end{aligned}
\]
В этом частном случае FedAvg точно воспроизводит централизованный
градиентный шаг.

При \(\tau_k>1\) последующие градиенты вычисляются уже в различных
точках
\[
    x_k^{t,0},x_k^{t,1},\ldots,x_k^{t,\tau_k-1},
\]
поэтому серверное усреднение в общем случае не совпадает с одним
шагом по \(F\). Если локальные функции различаются, направления
\(\nabla F_i\) и \(\nabla F_j\) также различаются, и локальные модели
дрейфуют друг от друга. Увеличение числа локальных шагов сокращает
число коммуникаций, но усиливает этот дрейф; при неоднородных данных
он может заметно ухудшить направление глобального обновления.

\section{Эксперимент с MNIST}

Рассмотрим набор MNIST из \(60000\) изображений,
распределённых между \(100\) клиентами по \(600\) объектов. При
однородном разбиении изображения случайно перемешиваются, поэтому
локальные наборы приблизительно воспроизводят общее распределение
десяти классов.

Для неоднородного разбиения изображения сначала сортируются по меткам,
затем делятся на \(200\) блоков по \(300\) объектов, и каждый клиент
получает два блока. Поэтому его локальный набор обычно содержит только
две цифры. В экспериментах с полносвязной и свёрточной сетями такое
разбиение требовало большего числа коммуникационных раундов, а часть
настроек не достигала заданной точности. Причина состоит в том, что
локальный градиент строится по малому числу классов и может существенно
отклоняться от глобального направления.

\section{Параметры и стоимость}

Увеличение доли участия \(C\) даёт серверу более представительную
выборку клиентов и обычно уменьшает число требуемых раундов. При этом
число передаваемых моделей в одном раунде растёт примерно как
\(m\approx CK\), поэтому коммуникационная стоимость раунда
увеличивается.

Увеличение числа эпох \(E\) или уменьшение размера мини-пакета \(B\)
повышает число локальных шагов
\[
    \tau_k
    =
    E\left\lceil\frac{n_k}{B}\right\rceil.
\]
Это может сократить число обращений к серверу, но увеличивает локальную
вычислительную работу и потенциальный дрейф моделей. В таблицах
лекции обозначение \(B=\infty\) соответствует вычислению градиента по
всему локальному набору.

Следовательно, число коммуникационных раундов само по себе не
определяет эффективность FedAvg. Необходимо учитывать размер модели,
число участвующих клиентов, объём локальных вычислений, доступность
устройств и степень неоднородности данных.

\section{Ограничения базовой схемы}

Базовый FedAvg предполагает, что выбранные клиенты завершают локальное
обучение и возвращают корректные обновления. Он непосредственно не
устраняет задержки медленных клиентов, разрывы связи, злонамеренные
обновления, утечку информации из локальных моделей и сильный дрейф
при неоднородных данных. Для этих случаев требуются дополнительные
механизмы отбора клиентов, робастной и защищённой агрегации,
приватности и коррекции локального обучения.

Главный компромисс FedAvg состоит в том, что несколько сравнительно
дешёвых локальных шагов выполняются между двумя дорогими обменами с
сервером. Поэтому параметры \(C\), \(E\), \(B\) и \(\eta\) должны
выбираться совместно, а не независимо.

\lecture{44}{Сходимость FedAvg}

\sourcevideo
    {https://www.youtube.com/playlist?list=PLOzRYVm0a65fhKDS8cYIC7YCuVGJ4FOJx}
    {Week 11: Lecture 44: Convergence Analysis of FL}

Рассматривается задача
\[
    \min_{x\in\mathbb R^d}F(x),
    \qquad
    F(x)=\sum_{i=1}^{K}p_iF_i(x),
\]
где
\[
    p_i\ge0,
    \qquad
    \sum_{i=1}^{K}p_i=1.
\]
В FedAvg на каждом коммуникационном раунде выбирается часть клиентов,
каждый из них выполняет несколько локальных стохастических шагов, а
сервер усредняет полученные модели. Сходимость определяется совместным
действием стохастического шума, неполного участия клиентов и дрейфа
локальных траекторий.

\section{Предпосылки и динамика FedAvg}

Предполагается, что каждая локальная функция имеет
\(L\)-липшицев градиент:
\[
    \lVert\nabla F_i(x)-\nabla F_i(y)\rVert
    \le
    L\lVert x-y\rVert.
\]
Эквивалентная верхняя оценка имеет вид
\[
    F_i(y)
    \le
    F_i(x)
    +\langle\nabla F_i(x),y-x\rangle
    +\frac L2\lVert y-x\rVert^2.
\]

Клиент использует стохастический градиент \(g_i(x;\xi)\), для которого
\[
    \mathbb E_\xi[g_i(x;\xi)]
    =
    \nabla F_i(x)
\]
и
\[
    \mathbb E_\xi
    \left[
        \lVert
            g_i(x;\xi)-\nabla F_i(x)
        \rVert^2
    \right]
    \le
    \sigma^2.
\]
Следовательно,
\[
    \mathbb E_\xi\lVert g_i(x;\xi)\rVert^2
    \le
    \lVert\nabla F_i(x)\rVert^2+\sigma^2.
\]

Неоднородность локальных функций описывается условием
\[
    \sum_{i=1}^{K}
    p_i\lVert\nabla F_i(x)\rVert^2
    \le
    \beta^2\lVert\nabla F(x)\rVert^2
    +\kappa^2,
\]
где \(\beta^2\ge1\) и \(\kappa^2\ge0\). Параметр \(\beta\)
характеризует относительное различие направлений, а \(\kappa\) ---
остаточное различие локальных градиентов около стационарной точки
глобальной функции. В однородном идеализированном случае можно взять
\[
    \beta=1,
    \qquad
    \kappa=0.
\]

На раунде \(t\) выбранный клиент начинает из глобальной модели
\[
    x_i^{t,0}=x^t
\]
и выполняет \(\tau\) шагов
\[
    x_i^{t,s+1}
    =
    x_i^{t,s}
    -
    \eta g_i(x_i^{t,s};\xi_i^{t,s}),
    \qquad
    s=0,\ldots,\tau-1.
\]
Затем сервер агрегирует конечные локальные модели:
\[
    x^{t+1}
    =
    \sum_{i\in S_t}q_i^t x_i^{t,\tau},
    \qquad
    \sum_{i\in S_t}q_i^t=1.
\]

\section{Структура оценки сходимости}

При перечисленных предпосылках типичная невыпуклая оценка имеет
схематический вид
\[
\begin{aligned}
    \frac1T
    \sum_{t=0}^{T-1}
    \mathbb E
    \lVert\nabla F(x^t)\rVert^2
    \lesssim{}&
    \frac{F(x^0)-F_{\inf}}{\eta\tau T}
    \\
    &+
    \text{стохастическая ошибка}
    \\
    &+
    \text{ошибка локального дрейфа}.
\end{aligned}
\]
Первое слагаемое отвечает за начальную ошибку оптимизации.
Стохастическая часть зависит от дисперсии, размера мини-пакета и числа
участвующих клиентов. Дрейф возникает при \(\tau>1\), поскольку
градиенты вычисляются уже не в общей точке \(x^t\), а в различных
локальных точках \(x_i^{t,s}\).

Для анализа достаточно структуры оценки и её зависимости от параметров задачи; точные константы далее не используются.

Увеличение числа участников и размера мини-пакета обычно уменьшает
дисперсию агрегированного обновления. Увеличение числа локальных шагов
сокращает коммуникации, но усиливает накопленное отличие локальных
траекторий от глобального градиентного шага. Параметры \(\beta\) и
\(\kappa\) входят именно в оценку дрейфа и ухудшают её при росте
неоднородности.

\section{Участие клиентов и локальная работа}

Пусть в одном раунде участвует доля \(C\in(0,1]\) клиентов, так что
\[
    m\approx CK.
\]
При большем \(m\) агрегированное направление обычно имеет меньшую
дисперсию, поэтому может потребоваться меньше раундов. Одновременно
увеличиваются объём передачи, нагрузка на сервер и вероятность
ожидания медленного клиента.

Если клиент \(i\) имеет \(n_i\) объектов, использует мини-пакет
размера \(B\) и выполняет \(E\) локальных эпох, то число локальных
обновлений равно
\[
    \tau_i
    =
    E
    \left\lceil
        \frac{n_i}{B}
    \right\rceil.
\]
При фиксированном \(E\) увеличение \(B\) уменьшает число
последовательных шагов, но повышает объём обработки за один шаг и
требования к памяти. Поэтому реальное время зависит от оборудования и
параллелизма, а не только от числа итераций.

Увеличение \(E\) переносит больше работы внутрь коммуникационного
раунда. При умеренном значении это уменьшает число обменов с сервером.
При слишком большом \(E\) локальная модель начинает двигаться главным
образом к минимуму \(F_i\), а не к минимуму глобальной функции.

Время одного раунда складывается из рассылки модели, локального
обучения, обратной передачи и серверной агрегации. Параметры
\(\beta\) и \(\kappa\) непосредственно не задают длительность
локального шага, но влияют на число раундов, необходимое для достижения
заданной точности.

\section{Происхождение локального дрейфа}

Разворачивая локальную рекурсию, получаем
\[
    x_i^{t,\tau}
    =
    x^t
    -
    \eta
    \sum_{s=0}^{\tau-1}
    g_i(x_i^{t,s};\xi_i^{t,s}).
\]
После серверного усреднения
\[
    x^{t+1}
    =
    x^t
    -
    \eta
    \sum_{i\in S_t}q_i^t
    \sum_{s=0}^{\tau-1}
    g_i(x_i^{t,s};\xi_i^{t,s}).
\]
При \(\tau=1\) все локальные градиенты вычисляются в общей точке
\(x^t\). При \(\tau>1\) точки \(x_i^{t,s}\) начинают различаться,
поэтому усредняется сумма градиентов, вычисленных вдоль разных
траекторий.

Если локальные функции близки, то
\[
    \nabla F_i(x)
    \approx
    \nabla F(x),
\]
и несколько локальных шагов остаются хорошим приближением глобального
обновления. При сильной неоднородности направления
\(\nabla F_i(x)\) могут заметно различаться, а накопленный локальный
шаг становится смещённым относительно \(\nabla F(x)\). Поэтому
ошибка обычно растёт с \(\tau\), шагом \(\eta\) и параметрами
неоднородности.

\section{Точный пример с двумя клиентами}

Рассмотрим
\[
    F_1(x)=(x-1)^2,
    \qquad
    F_2(x)=2(x-5)^2
\]
и веса
\[
    p_1=p_2=\frac12.
\]
Тогда
\[
    F(x)
    =
    \frac12(x-1)^2+(x-5)^2
\]
и
\[
    F'(x)
    =
    (x-1)+2(x-5)
    =
    3x-11.
\]
Следовательно,
\[
    x^\star=\frac{11}{3}.
\]

Пусть оба клиента начинают из \(x^t\) и выполняют по одному точному
градиентному шагу:
\[
    x_1^{t,1}
    =
    x^t-2\eta(x^t-1),
\]
\[
    x_2^{t,1}
    =
    x^t-4\eta(x^t-5).
\]
После равномерного усреднения
\[
\begin{aligned}
    x^{t+1}
    &=
    \frac12x_1^{t,1}
    +
    \frac12x_2^{t,1}
    \\
    &=
    x^t
    -
    \eta\bigl[(x^t-1)+2(x^t-5)\bigr]
    \\
    &=
    x^t-\eta F'(x^t).
\end{aligned}
\]
При одном локальном шаге FedAvg здесь точно совпадает с градиентным
спуском по глобальной функции.

Если же каждый клиент полностью минимизирует собственную функцию, то
\[
    x_1^\star=1,
    \qquad
    x_2^\star=5,
\]
и сервер получает
\[
    \frac{x_1^\star+x_2^\star}{2}=3.
\]
Однако
\[
    3\ne\frac{11}{3}.
\]
Таким образом, если \(x_i^\star\in\operatorname*{arg\,min}F_i\),
то в общем случае
\[
    \sum_i p_i x_i^\star
    \ne
    x^\star,
    \qquad
    x^\star
    \in
    \operatorname*{arg\,min}_x
    \sum_i p_iF_i(x).
\]
Пример точно показывает источник смещения при чрезмерной локальной
оптимизации.

\section{Компромисс и границы анализа}

Увеличение \(m\) улучшает оценку глобального направления, но делает
каждый раунд дороже. Увеличение \(E\) и \(\tau\) уменьшает частоту
коммуникаций, но усиливает дрейф. Уменьшение \(B\) увеличивает число
локальных шагов и стохастический шум. Поэтому параметры должны
выбираться совместно с учётом стоимости связи, мощности клиентов,
неоднородности данных и требуемой точности.

Рассмотренная схема предполагает гладкость локальных функций,
несмещённые стохастические градиенты, ограниченную дисперсию и условие
ограниченной несхожести. Она не описывает вычислительную
неоднородность устройств, выпадение клиентов, асинхронность, задержки,
защищённую агрегацию и ограничения конфиденциальности. Эти эффекты
изменяют как время одного раунда, так и саму динамику алгоритма.

\lecture{45}{Неоднородность клиентов}

\sourcevideo
    {https://www.youtube.com/playlist?list=PLOzRYVm0a65fhKDS8cYIC7YCuVGJ4FOJx}
    {Week 12: Lecture 45}

В федеративном обучении выбранные клиенты начинают локальную работу
из одной серверной модели, но могут выполнить различный объём
вычислений. Причинами служат неодинаковые размеры локальных выборок,
скорости устройств, временные бюджеты, шаги и оптимизаторы. В
результате сервер усредняет обновления, полученные по различным
локальным траекториям.

\section{Локальная работа в FedAvg}

Пусть на коммуникационном раунде \(t\) сервер хранит модель
\(x^t\in\mathbb R^d\) и выбирает множество клиентов \(S_t\). Клиент
\(i\in S_t\) инициализирует
\[
    x_i^{t,0}=x^t
\]
и выполняет \(\tau_i\) локальных стохастических шагов
\[
    x_i^{t,j+1}
    =
    x_i^{t,j}
    -
    \eta_i g_i
    \left(
        x_i^{t,j};\xi_i^{t,j}
    \right),
    \qquad
    j=0,\ldots,\tau_i-1.
\]
После этого сервер агрегирует полученные модели:
\[
    x^{t+1}
    =
    \sum_{i\in S_t}p_i^t x_i^{t,\tau_i},
    \qquad
    \sum_{i\in S_t}p_i^t=1.
\]
При взвешивании по числу объектов обычно берут
\[
    p_i^t
    =
    \frac{n_i}{\sum_{r\in S_t}n_r},
\]
где \(n_i\) --- размер локальной выборки.

Если клиент использует мини-пакет размера \(B_i\) и выполняет
\(E_i\) эпох, то число локальных обновлений равно
\[
    \tau_i
    =
    E_i
    \left\lceil\frac{n_i}{B_i}\right\rceil.
\]
При делимости \(n_i\) на \(B_i\) получаем
\(\tau_i=E_in_i/B_i\). Таким образом, одинаковое число эпох не
означает одинакового числа шагов.

\section{Виды неоднородности}

Статистическая неоднородность означает различие локальных
распределений и функций:
\[
    \mathcal D_i\ne\mathcal D_j,
    \qquad
    F_i\ne F_j.
\]
Вычислительная неоднородность относится к различиям в самом процессе
локального обучения: значениях \(\tau_i\), \(\eta_i\), времени одного
шага и типе оптимизатора. Эти эффекты независимы: одинаковые
устройства могут хранить выборки разного размера, а одинаковые по
распределению данные могут обрабатываться с разной скоростью.

При общих \(E\) и \(B\)
\[
    n_i>n_j
    \quad\Longrightarrow\quad
    \tau_i\ge\tau_j.
\]
Локальное изменение модели имеет вид
\[
    \Delta_i^t
    :=
    x_i^{t,\tau_i}-x^t
    =
    -\eta_i
    \sum_{j=0}^{\tau_i-1}
    g_i
    \left(
        x_i^{t,j};\xi_i^{t,j}
    \right).
\]
Поэтому размер выборки влияет на глобальное обновление дважды: через
вес \(p_i^t\) и через длину локальной траектории \(\tau_i\). Первый
эффект отражает статистический вес клиента, второй изменяет само
направление и масштаб передаваемого обновления.

\section{Скорость устройств и отстающие клиенты}

Пусть \(c_i>0\) --- время одного локального шага клиента \(i\). При
фиксированном числе шагов \(\tau\) его локальное время примерно равно
\[
    T_i^{\mathrm{loc}}=\tau c_i.
\]
В синхронной схеме сервер начинает следующий раунд только после
получения требуемых ответов, поэтому
\[
    T_t^{\mathrm{round}}
    \approx
    \max_{i\in S_t}T_i^{\mathrm{loc}}
    +
    T_t^{\mathrm{comm}}.
\]
Клиенты с большим \(c_i\) становятся отстающими и определяют
продолжительность раунда.

Фиксация общего числа шагов выравнивает алгоритмический объём
локальной работы, но не время завершения. Кроме того, равное
\(\tau\) не означает одинакового числа просмотренных эпох, если
\(n_i\) различаются.

\section{Фиксированный временной бюджет}

Другой вариант состоит в выделении каждому клиенту одинакового
времени \(H>0\). Тогда
\[
    \tau_i
    =
    \left\lfloor\frac{H}{c_i}\right\rfloor,
\]
и более быстрый клиент выполняет больше шагов. Продолжительность
локальной стадии становится контролируемой, но сервер получает
обновления, построенные по траекториям различной длины.

Тем самым при разных скоростях устройств нельзя одновременно
обеспечить одинаковое число локальных шагов и одинаковое время
завершения. Выбор между этими режимами определяется тем, важнее ли
синхронизировать объём работы или ограничить задержку раунда.

\section{Различия шагов и локальных оптимизаторов}

Даже при одинаковом \(\tau_i\) масштаб \(\Delta_i^t\) зависит от
шага \(\eta_i\). Слишком большой шаг усиливает локальный дрейф и
может нарушить устойчивость, а слишком малый делает вклад клиента
незначительным.

Клиенты могут также использовать различные оптимизаторы. Например,
локальный метод с импульсом имеет вид
\[
    v_i^{t,j+1}
    =
    \beta_i v_i^{t,j}
    +
    g_i
    \left(
        x_i^{t,j};\xi_i^{t,j}
    \right),
\]
\[
    x_i^{t,j+1}
    =
    x_i^{t,j}
    -
    \eta_i v_i^{t,j+1}.
\]
Различия в \(\eta_i\), \(\beta_i\), состоянии импульса и типе
оптимизатора означают, что одно и то же число шагов не задаёт один и
тот же эффективный объём локального изменения. Сервер, получающий
только конечную модель, не всегда может восстановить эту информацию.

\section{Нормализация обновлений}

Обычная агрегация через локальные приращения записывается как
\[
    x^{t+1}
    =
    x^t
    +
    \sum_{i\in S_t}p_i^t\Delta_i^t.
\]
При различных \(\tau_i\) она смешивает обновления разного масштаба.
Естественная коррекция состоит в нормализации
\[
    \widehat\Delta_i^t
    =
    \frac{\Delta_i^t}{a_i^t},
\]
где \(a_i^t\) описывает эффективную накопленную длину локального
шага. Для обычного SGD с постоянным шагом простейшим ориентиром служит
\(a_i^t=\tau_i\), но при импульсе или меняющихся шагах коэффициент
должен учитывать всю локальную рекурсию. Точная нормализация приводит
к проблеме несогласованности целевой функции и к алгоритму FedNova,
которые рассматриваются в следующих лекциях.

Реальные клиенты различаются также по памяти, доступной мощности,
заряду и качеству соединения. В этой лекции эти факторы служат
мотивацией, но явно не входят в математическую модель. Основные
моделируемые величины --- \(n_i\), \(B_i\), \(E_i\), \(\tau_i\),
\(c_i\), \(\eta_i\) и параметры локального оптимизатора.

\lecture{46}{Несогласованность цели}

\sourcevideo
    {https://www.youtube.com/playlist?list=PLOzRYVm0a65fhKDS8cYIC7YCuVGJ4FOJx}
    {Week 12: Lecture 46: Objective Inconsistency Problem}

В федеративном обучении клиенты могут выполнять различное число
локальных шагов. Тогда обычное усреднение изменяет не только скорость
алгоритма: уже в первом порядке оно перевзвешивает локальные функции
и может направить метод к минимуму другой целевой функции.

\section{Глобальная задача и локальные обновления}

Пусть требуется минимизировать
\[
    F(x)=\sum_{i=1}^{K}p_iF_i(x),
    \qquad
    p_i\ge0,
    \qquad
    \sum_{i=1}^{K}p_i=1.
\]
Для эмпирического риска обычно используют
\[
    p_i=\frac{n_i}{\sum_{r=1}^{K}n_r},
\]
где \(n_i\) --- число объектов у клиента \(i\).

На раунде \(t\) все выбранные клиенты начинают из общей модели
\[
    x_i^{t,0}=x^t
\]
и выполняют \(\tau_i\) локальных градиентных шагов
\[
    x_i^{t,j+1}
    =
    x_i^{t,j}-\eta\nabla F_i(x_i^{t,j}),
    \qquad
    j=0,\ldots,\tau_i-1.
\]
Обычный FedAvg формирует
\[
    x^{t+1}
    =
    \sum_{i=1}^{K}p_i x_i^{t,\tau_i}
    =
    x^t+
    \sum_{i=1}^{K}p_i\Delta_i^t,
\]
где
\[
    \Delta_i^t=x_i^{t,\tau_i}-x^t.
\]
Если значения \(\tau_i\) различаются, масштабы и направления
локальных приращений также различаются.

\section{Точный пример с двумя клиентами}

Рассмотрим
\[
    F_1(x)=(x-1)^2,
    \qquad
    F_2(x)=2(x-5)^2,
    \qquad
    p_1=p_2=\frac12.
\]
Глобальная функция равна
\[
    F(x)
    =
    \frac12(x-1)^2+(x-5)^2,
\]
поэтому
\[
    F'(x)=3x-11,
    \qquad
    x^\star=\frac{11}{3}.
\]
Локальные минимизаторы равны \(1\) и \(5\), а их обычное среднее
равно \(3\). Следовательно, усреднение локальных решений не
эквивалентно минимизации средней функции.

Для первого клиента
\[
    x_1^{t,j+1}-1
    =
    (1-2\eta)(x_1^{t,j}-1),
\]
откуда
\[
    x_1^{t,\tau_1}
    =
    1+(1-2\eta)^{\tau_1}(x^t-1).
\]
Для второго клиента аналогично
\[
    x_2^{t,\tau_2}
    =
    5+(1-4\eta)^{\tau_2}(x^t-5).
\]
Локальные итерации устойчивы при \(0<\eta<1\) для первой функции и
при \(0<\eta<1/2\) для второй, поэтому далее достаточно считать
\[
    0<\eta<\frac12.
\]

\section{Серверная динамика и её предел}

Обозначим
\[
    a_1=(1-2\eta)^{\tau_1},
    \qquad
    a_2=(1-4\eta)^{\tau_2}.
\]
Тогда серверное усреднение задаёт аффинную рекурсию
\[
    x^{t+1}
    =
    3+
    \frac{a_1}{2}(x^t-1)
    +
    \frac{a_2}{2}(x^t-5).
\]
Коэффициент при \(x^t\) равен
\[
    q=\frac{a_1+a_2}{2}.
\]
При \(0<\eta<1/2\) выполнено
\[
    |a_1|<1,
    \qquad
    |a_2|<1,
\]
а значит, \(|q|<1\). Поэтому последовательность сходится к
единственной неподвижной точке
\[
    \widetilde x
    =
    \frac{6-a_1-5a_2}{2-a_1-a_2}.
\]
В общем случае \(\widetilde x\ne x^\star\), причём смещение зависит
не только от функций, но и от \(\eta\), \(\tau_1\) и \(\tau_2\).

При \(\eta\to0\) и фиксированных \(\tau_1,\tau_2\)
\[
    a_1=1-2\eta\tau_1+O(\eta^2),
    \qquad
    a_2=1-4\eta\tau_2+O(\eta^2),
\]
откуда
\[
    \widetilde x
    \longrightarrow
    \frac{\tau_1+10\tau_2}{\tau_1+2\tau_2}.
\]
Это значение совпадает с \(11/3\) при
\(\tau_1=\tau_2\), но обычно отличается от глобального минимума.

\section{Эффективная целевая функция}

Пусть градиенты \(F_i\) липшицевы, шаг \(\eta\) мал, а числа
локальных итераций фиксированы. Тогда
\[
    x_i^{t,\tau_i}
    =
    x^t
    -
    \eta\tau_i\nabla F_i(x^t)
    +
    O(\eta^2\tau_i^2).
\]
После усреднения
\[
\begin{aligned}
    x^{t+1}
    ={}&
    x^t
    -
    \eta
    \sum_{i=1}^{K}
    p_i\tau_i\nabla F_i(x^t)
    \\
    &+
    O\left(
        \eta^2
        \sum_{i=1}^{K}p_i\tau_i^2
    \right).
\end{aligned}
\]
Положим
\[
    s=\sum_{r=1}^{K}p_r\tau_r,
    \qquad
    \widetilde p_i=\frac{p_i\tau_i}{s}.
\]
Тогда главный член является градиентным шагом размера \(\eta s\) по
функции
\[
    \widetilde F(x)
    =
    \sum_{i=1}^{K}\widetilde p_iF_i(x).
\]
Отличие \(\widetilde F\) от требуемой функции \(F\) называется
несогласованностью целевой функции. При весах \(p_i=n_i/n\)
эффективный вес клиента пропорционален \(n_i\tau_i\), а не только
числу его объектов.

Эта интерпретация является локальной по шагу \(\eta\). При конечном
шаге и нелинейных функциях градиенты вычисляются в промежуточных
точках \(x_i^{t,j}\), поэтому обновление FedAvg не обязано быть точным
градиентным шагом по одной фиксированной функции. Если \(\tau_i\)
меняются от раунда к раунду, меняются и эффективные веса.

\section{Когда смещение исчезает}

Если все клиенты выполняют одинаковое число шагов,
\[
    \tau_i=\tau,
\]
то
\[
    \widetilde p_i
    =
    \frac{p_i\tau}{\sum_r p_r\tau}
    =p_i.
\]
Следовательно, в первом порядке эффективная и исходная целевые
функции совпадают. Однако равное число шагов не выравнивает время
работы: синхронный сервер по-прежнему ждёт самых медленных клиентов.

Если все локальные функции имеют одно и то же множество минимизаторов,
перевзвешивание может не изменить предельное решение, хотя изменяет
траекторию и скорость. В частности, при
\[
    F_1=\cdots=F_K
\]
любые положительные веса задают одну и ту же функцию с точностью до
множителя. Наиболее заметное смещение возникает при сочетании
неодинаковых \(\tau_i\) и статистически неоднородных локальных задач.

\section{Коррекция локальной работы}

В простейшем случае постоянного локального шага можно изменить веса
серверной агрегации. Если
\[
    q_i
    =
    \frac{p_i/\tau_i}{\sum_{r=1}^{K}p_r/\tau_r},
\]
то
\[
    q_i\tau_i
    =
    \frac{p_i}{\sum_r p_r/\tau_r}.
\]
Поэтому главный член обновления имеет требуемое направление
\[
    -\sum_{i=1}^{K}p_i\nabla F_i(x^t)
\]
с точностью до общего множителя, который можно включить в серверный
шаг.

Одного деления на \(\tau_i\) недостаточно, если клиент использует
меняющиеся шаги, импульс или другой локальный оптимизатор. Тогда
приращение представляется как взвешенная сумма локальных градиентов,
а нормирующий коэффициент должен учитывать всю локальную рекурсию.
Эта идея приводит к нормализованной агрегации FedNova, изучаемой в
следующей лекции.

\lecture{47}{FedNova и справедливость}

\sourcevideo
    {https://www.youtube.com/playlist?list=PLOzRYVm0a65fhKDS8cYIC7YCuVGJ4FOJx}
    {Week 12: Lecture 47: General Update Rule}

Рассматривается федеративная задача
\[
    \min_{x\in\mathbb R^d}F(x),
    \qquad
    F(x)=\sum_{i=1}^{K}p_iF_i(x),
\]
где
\[
    p_i\ge0,
    \qquad
    \sum_{i=1}^{K}p_i=1.
\]
На раунде \(t\) сервер рассылает модель \(x^t\), клиент \(i\)
выполняет \(\tau_i\) локальных шагов, а затем сервер агрегирует
полученные обновления. Общее правило позволяет отделить локальное
направление от масштаба обновления и увидеть, почему FedAvg
перевзвешивает клиентов с разным числом локальных шагов.

\section{Локальная траектория и изменение модели}

Клиент начинает из общей точки
\[
    x_i^{t,0}=x^t
\]
и выполняет рекурсию
\[
    x_i^{t,j+1}
    =
    x_i^{t,j}-\eta g_i^{t,j},
    \qquad
    j=0,\ldots,\tau_i-1,
\]
где \(g_i^{t,j}\) --- локальный стохастический градиент. После
\(\tau_i\) шагов
\[
    x_i^{t,\tau_i}
    =
    x^t
    -
    \eta
    \sum_{j=0}^{\tau_i-1}g_i^{t,j}.
\]
Определим локальное изменение модели
\[
    \Delta_i^t
    =
    x_i^{t,\tau_i}-x^t.
\]
Тогда
\[
    \Delta_i^t
    =
    -\eta
    \sum_{j=0}^{\tau_i-1}g_i^{t,j}.
\]

Обычное федеративное усреднение имеет вид
\[
    x^{t+1}
    =
    \sum_{i=1}^{K}p_i x_i^{t,\tau_i}
    =
    x^t+
    \sum_{i=1}^{K}p_i\Delta_i^t,
\]
то есть
\[
    x^{t+1}
    =
    x^t
    -
    \eta
    \sum_{i=1}^{K}p_i
    \sum_{j=0}^{\tau_i-1}g_i^{t,j}.
\]
Полное изменение клиента растёт не только из-за направления
градиентов, но и из-за числа выполненных локальных шагов.

\section{Нормализованное направление и общее правило}

Пусть локальный метод задаёт неотрицательные коэффициенты
\[
    a_{i,0},\ldots,a_{i,\tau_i-1}
\]
и вектор
\[
    a_i
    =
    (a_{i,0},\ldots,a_{i,\tau_i-1})^{\mathsf T}.
\]
Нормализованное локальное направление определяется формулой
\[
    d_i^t
    =
    \frac1{\lVert a_i\rVert_1}
    \sum_{j=0}^{\tau_i-1}
    a_{i,j}g_i^{t,j},
    \qquad
    \lVert a_i\rVert_1
    =
    \sum_{j=0}^{\tau_i-1}a_{i,j}.
\]
Коэффициенты \(a_{i,j}\) позволяют описывать переменные локальные
шаги и накопление импульса.

Широкий класс федеративных методов записывается как
\[
    x^{t+1}
    =
    x^t
    -
    \eta\tau_{\mathrm{eff}}
    \sum_{i=1}^{K}w_i d_i^t,
\]
где
\[
    w_i\ge0,
    \qquad
    \sum_{i=1}^{K}w_i=1.
\]
Здесь \(d_i^t\) задаёт направление клиента, \(w_i\) --- его
относительный серверный вес, а \(\tau_{\mathrm{eff}}\) --- общий
масштаб раунда. Различные алгоритмы получаются разным выбором
\(a_i\), \(w_i\) и \(\tau_{\mathrm{eff}}\).

Если каждый клиент выполняет один шаг, то
\[
    a_i=(1),
    \qquad
    d_i^t=g_i^{t,0},
    \qquad
    w_i=p_i,
    \qquad
    \tau_{\mathrm{eff}}=1,
\]
и общее правило превращается в обычное агрегирование одного
локального градиента:
\[
    x^{t+1}
    =
    x^t-
    \eta\sum_{i=1}^{K}p_i g_i^{t,0}.
\]

\section{FedAvg и несогласованность цели}

Для FedAvg все внутренние коэффициенты одинаковы:
\[
    a_i=\mathbf1_{\tau_i}.
\]
Тогда
\[
    \lVert a_i\rVert_1=\tau_i,
    \qquad
    d_i^t
    =
    \frac1{\tau_i}
    \sum_{j=0}^{\tau_i-1}g_i^{t,j}.
\]
Это средний локальный градиент клиента. Если положить
\[
    \tau_{\mathrm{eff}}
    =
    \sum_{r=1}^{K}p_r\tau_r
\]
и
\[
    w_i
    =
    \frac{p_i\tau_i}
    {\displaystyle\sum_{r=1}^{K}p_r\tau_r},
\]
то
\[
\begin{aligned}
    \tau_{\mathrm{eff}}
    \sum_{i=1}^{K}w_i d_i^t
    &=
    \sum_{i=1}^{K}
    p_i\tau_i
    \frac1{\tau_i}
    \sum_{j=0}^{\tau_i-1}g_i^{t,j}
    \\
    &=
    \sum_{i=1}^{K}p_i
    \sum_{j=0}^{\tau_i-1}g_i^{t,j}.
\end{aligned}
\]
Следовательно, общее правило точно воспроизводит FedAvg.

Однако исходная целевая функция использует веса \(p_i\), тогда как
эффективные веса в представлении FedAvg равны
\[
    w_i
    =
    \frac{p_i\tau_i}
    {\sum_r p_r\tau_r}.
\]
Поэтому при различных \(\tau_i\) клиент с большим числом локальных
шагов получает больший относительный вес. Это несогласованность между исходной целью и направлением, агрегируемым сервером. При равных
локальных длинах
\[
    \tau_i=\tau
\]
получаем \(w_i=p_i\), и проблема исчезает.

\section{Нормализация FedNova}

FedNova отделяет направление локального обновления от числа
выполненных шагов. Клиент нормализует накопленное изменение:
\[
    \widehat\Delta_i^t
    =
    \frac{\Delta_i^t}{\tau_i}.
\]
Так как
\[
    \Delta_i^t
    =
    -\eta
    \sum_{j=0}^{\tau_i-1}g_i^{t,j},
\]
то
\[
    \widehat\Delta_i^t
    =
    -\eta
    \frac1{\tau_i}
    \sum_{j=0}^{\tau_i-1}g_i^{t,j}.
\]
Клиент передаёт среднее изменение на один локальный шаг, а не полное
накопленное изменение.

Сервер сохраняет исходные веса задачи
\[
    w_i=p_i
\]
и использует масштаб
\[
    \tau_{\mathrm{eff}}
    =
    \sum_{i=1}^{K}p_i\tau_i.
\]
Обновление принимает вид
\[
    x^{t+1}
    =
    x^t
    +
    \tau_{\mathrm{eff}}
    \sum_{i=1}^{K}p_i
    \frac{\Delta_i^t}{\tau_i},
\]
или, через локальные градиенты,
\[
    x^{t+1}
    =
    x^t
    -
    \eta\tau_{\mathrm{eff}}
    \sum_{i=1}^{K}p_i
    \left(
        \frac1{\tau_i}
        \sum_{j=0}^{\tau_i-1}g_i^{t,j}
    \right).
\]
Теперь \(\tau_i\) влияет на общий масштаб раунда, но не изменяет
относительный вес клиента \(p_i\). FedAvg агрегирует
\(\Delta_i^t\), а FedNova --- \(\Delta_i^t/\tau_i\). При
одинаковых \(\tau_i\) оба метода совпадают.

Нормализация устраняет только перевзвешивание по числу локальных
шагов. Она сама по себе не устраняет статистическую неоднородность,
стохастический шум, задержки и различия локальных оптимизаторов.

\section{Меры справедливости}

Минимизация среднего риска может давать хорошее глобальное качество,
но плохой результат для отдельных клиентов. Поэтому дополнительно
рассматривают распределение локальных значений
\[
    F_1(x),\ldots,F_K(x).
\]
Положим
\[
    \overline F(x)
    =
    \frac1K
    \sum_{i=1}^{K}F_i(x).
\]
Простейшая мера разброса --- дисперсия локальных потерь:
\[
    V_F(x)
    =
    \frac1K
    \sum_{i=1}^{K}
    \bigl(F_i(x)-\overline F(x)\bigr)^2.
\]
Малое значение \(V_F\) означает близость клиентских потерь, но не
гарантирует, что сами потери малы. Поэтому дисперсию следует
рассматривать вместе со средней потерей.

Если \(F_i(x)\ge0\) и \(\sum_iF_i(x)>0\), определим
\[
    q_i(x)
    =
    \frac{F_i(x)}
    {\displaystyle\sum_{r=1}^{K}F_r(x)}.
\]
Энтропийная мера равна
\[
    H(x)
    =
    -\sum_{i=1}^{K}q_i(x)\log q_i(x).
\]
При равных значениях \(q_i=1/K\) получаем \(H(x)=\log K\), а при
концентрации значений у небольшого числа клиентов энтропия
уменьшается. Эта мера оценивает равномерность, но не различает
равномерно хорошие и равномерно плохие результаты.

Для неотрицательных показателей качества \(u_i(x)\) индекс Джейна
определяется формулой
\[
    J(x)
    =
    \frac{
        \left(\displaystyle\sum_{i=1}^{K}u_i(x)\right)^2
    }{
        K\displaystyle\sum_{i=1}^{K}u_i(x)^2
    }.
\]
При ненулевых показателях
\[
    \frac1K\le J(x)\le1,
\]
причём \(J(x)=1\) тогда и только тогда, когда все показатели равны.
Обычно \(u_i\) выбирают как точность или положительную полезность;
равенство потерь само по себе ещё не означает высокого качества.

\section{Средняя цель и равномерность}

Стандартная федеративная задача минимизирует
\[
    \sum_{i=1}^{K}p_iF_i(x),
\]
тогда как справедливая постановка дополнительно стремится уменьшить
различия между клиентами. Эти цели могут конфликтовать: уменьшение
средней потери не обязано улучшать худшего клиента, а строгая
равномерность может ухудшить среднее качество.

FedNova решает более узкую задачу: возвращает серверные веса к
коэффициентам \(p_i\), заданным исходной целевой функцией, при
различных \(\tau_i\). Это исправляет несогласованность, вызванную
вычислительной неоднородностью, но не превращает средневзвешенную цель
в самостоятельный критерий справедливости. Дисперсия, энтропия и
индекс Джейна используются как дополнительные показатели и должны
интерпретироваться вместе с общим качеством модели.


\end{document}